\documentclass[11pt,a4paper]{article}

\usepackage[T1]{fontenc}
\usepackage{lmodern}
\usepackage[a4paper,margin=20mm]{geometry}
\usepackage{amsmath,amssymb,amsthm,mathtools}
\usepackage{xcolor}
\usepackage{fancyhdr}
\usepackage{booktabs}
\usepackage{tabularx}
\usepackage{multicol}
\usepackage{enumerate}
\usepackage{hyperref}

\definecolor{primary}{RGB}{24, 62, 114}
\definecolor{accent}{RGB}{150, 38, 42}
\definecolor{muted}{RGB}{90, 96, 106}

\hypersetup{
  colorlinks=true,
  linkcolor=primary,
  urlcolor=primary,
  pdftitle={NBHM Ph.D. Question Papers and Answer Keys (2005--2022)},
  pdfauthor={National Board for Higher Mathematics}
}

\setlength{\parskip}{5pt plus 2pt minus 1pt}
\setlength{\parindent}{0pt}
\setlength{\headheight}{14pt}

\pagestyle{fancy}
\fancyhf{}
\fancyhead[L]{\small\color{muted}NBHM Ph.D. Question Papers \& Answer Keys (2005--2022)}
\fancyhead[R]{\small\color{muted}\nouppercase{\leftmark}}
\fancyfoot[C]{\small\thepage}
\renewcommand{\headrulewidth}{0.4pt}

\newcommand{\N}{\mathbb{N}}
\newcommand{\Z}{\mathbb{Z}}
\newcommand{\Q}{\mathbb{Q}}
\newcommand{\R}{\mathbb{R}}
\newcommand{\C}{\mathbb{C}}
\newcommand{\F}{\mathbb{F}}
\newcommand{\M}{\mathbb{M}}
\newcommand{\T}{\mathbb{T}}
\newcommand{\tr}{\operatorname{tr}}
\newcommand{\diag}{\operatorname{diag}}
\newcommand{\rank}{\operatorname{rank}}
\newcommand{\divg}{\operatorname{div}}
\newcommand{\curl}{\operatorname{curl}}
\newcommand{\RePart}{\operatorname{Re}}
\newcommand{\ImPart}{\operatorname{Im}}
\newcommand{\sgn}{\operatorname{sgn}}
\newcommand{\ansbox}[1]{\fbox{\strut\,#1\,}}

\newcommand{\examheader}[4]{%
  \clearpage
  \section{#1}
  \begin{center}
    {\large\bfseries NATIONAL BOARD FOR HIGHER MATHEMATICS}\\[3pt]
    {\bfseries #2}\\[2pt]
    {#3}\\[2pt]
    {\small #4}
  \end{center}
  \vspace{4pt}
}

\newcommand{\examsec}[1]{%
  \subsection*{#1}
  \vspace{-2pt}
}

\newcommand{\qitem}[1]{\par\medskip\noindent\textbf{#1}\hspace{0.5em}}

\begin{document}

\begin{titlepage}
  \centering
  \vspace*{3cm}
  {\Huge\bfseries\color{primary} NBHM Ph.D. Screening Tests}\\[0.6cm]
  {\LARGE\bfseries\color{accent} 2005 -- 2022}\\[0.8cm]
  {\Large Question Papers \& Answer Keys}\\[1.5cm]
  {\large National Board for Higher Mathematics (NBHM)\\
  Research Awards / Research Scholarships / Doctoral Scholarship Scheme}\\[2cm]
  \rule{0.6\textwidth}{0.6pt}\\[0.6cm]
  {\normalsize\color{muted} Complete Typeset Compilation in \LaTeX}
  \vfill
\end{titlepage}

\tableofcontents
\clearpage
% ==================== NBHM 2005 ====================
\examheader{NBHM 2005}{Research Awards Screening Test}{June 25, 2005}{Time Allowed: 90 Minutes \quad|\quad Maximum Marks: 40}

\subsection*{Instructions to Candidates}
\begin{itemize}
  \item Do not forget to write your name and roll number on the cover page. In the box marked `Institution', fill in the name of the institution where you are working towards a Ph.D.\ degree. In case you have not yet joined any institution for research, write \emph{Not Applicable}.
  \item Please ensure that your answer booklet contains 13 numbered (and printed) pages. The back of each printed page is blank and can be used for rough work.
  \item There are \textbf{five} sections, containing \textbf{ten} questions each, entitled Algebra, Analysis, Topology, Applied Mathematics and Miscellaneous. Answer as many questions as possible. The assessment of the paper will be based on the best \textbf{four} sections. Each question carries one point and the maximum marks to be scored is \textbf{forty}.
  \item Answer each question, as directed, in the space provided at the end of it. Answers are to be given in the form of a word (or words, if required), a numerical value (or values) or a simple mathematical expression. \textbf{Do not write sentences.}
  \item In certain questions (Qns.\ 1.7 to 1.10, 2.7 to 2.10) you are required to pick out the qualifying statement(s) from multiple choices. None of the statements, or more than one statement may qualify. Write \textbf{none} if none of the statements qualify, or list the labels of \emph{all} the qualifying statements (amongst (a), (b), (c) and (d)).
  \item Points will be awarded in the above questions and in Questions 3.5 to 3.10 only if \textbf{all} the correct choices are made. There will be no partial credit.
  \item $\N$ denotes the set of natural numbers, $\Z$ -- the integers, $\Q$ -- the rationals, $\R$ -- the reals and $\C$ -- the field of complex numbers. $\R^n$ denotes the $n$-dimensional Euclidean space, which is assumed to be endowed with its `usual' topology. The symbol $]a, b[$ will stand for the open interval $\{x \in \R \mid a < x < b\}$ while $[a, b]$ will stand for the corresponding closed interval; $[a, b[$ and $]a, b]$ will stand for the corresponding left-closed-right-open and left-open-right-closed intervals respectively. The symbol $I$ will denote the identity matrix of appropriate order.
\end{itemize}

\examsec{Section 1: Algebra}

\qitem{1.1} Find the value of $a \in \Z$ such that $2 + \sqrt{3}$ is a root of the polynomial
\[
  x^3 - 5x^2 + ax - 1.
\]

\qitem{1.2} Let $V = \R^5$ and $W = \R^7$. Let $T : V \to W$ be a linear map. If $\mathcal{N}(T)$ denotes the null space of $T$ and $\mathcal{R}(T)$ denotes its range, then
\[
  \dim \mathcal{N}(T) + \dim \mathcal{R}(T) = ?
\]

\qitem{1.3} Let $A$ be a $3 \times 3$ matrix whose eigenvalues are $-1, 1, 2$. Find $\alpha, \beta$ and $\gamma$ such that
\[
  A^{-1} = \alpha A^2 + \beta A + \gamma I.
\]

\qitem{1.4} What is the number of groups of order $6$ (upto isomorphism)?

\qitem{1.5} Let $G$ be a cyclic group of order $10$. For $a \in G$, let $\langle a \rangle$ denote the subgroup generated by $a$. How many elements are there in the set
\[
  \{a \in G \mid \langle a \rangle = G\}?
\]

\qitem{1.6} Let $\alpha = 2^{\frac{1}{3}}$ and $\beta = 5^{\frac{1}{4}}$. Let $L$ be the field obtained by adjoining $\alpha$ and $\beta$ to $\Q$. What is the degree of the extension $[L : \Q]$?

\qitem{1.7} Pick out the matrices which are diagonalizable over $\C$:
\begin{enumerate}[(a)]
  \item Any $n \times n$ unitary matrix with complex entries.
  \item Any $n \times n$ hermitian matrix with complex entries.
  \item Any $n \times n$ strictly upper triangular matrix with complex entries.
  \item Any $n \times n$ matrix with complex entries whose eigenvalues are real.
\end{enumerate}

\qitem{1.8} Pick out the units in $\Z[\sqrt{3}]$.\\
(a) $-7 + 4\sqrt{3}$ \qquad (b) $5 + 3\sqrt{3}$ \qquad (c) $2 - \sqrt{3}$ \qquad (d) $-3 - 2\sqrt{3}$.

\qitem{1.9} Pick out the integral domains from the following list of rings:
\begin{enumerate}[(a)]
  \item $\{a + b\sqrt{5} \mid a, b \in \Q\}$.
  \item The ring of continuous functions from $[0, 1]$ into $\R$.
  \item The ring of complex analytic functions on the disc $\{z \in \C \mid |z| < 1\}$.
  \item The polynomial ring $\Z[x]$.
\end{enumerate}

\qitem{1.10} Pick out the abelian groups from the following list:
\begin{enumerate}[(a)]
  \item Any group of order $4$.
  \item Any group of order $36$.
  \item Any group of order $47$.
  \item Any group of order $49$.
\end{enumerate}

\examsec{Section 2: Analysis}

\qitem{2.1} Let $f : [0, 1] \to \R$ be a continuous function. Then
\[
  \lim_{n \to \infty} \frac{1}{n} \sum_{k=0}^{n} f\left(\frac{k}{n}\right) = ?
\]

\qitem{2.2} What is the radius of convergence of the following series?
\[
  \sum_{n=0}^{\infty} \frac{z^n}{n!}.
\]

\qitem{2.3} Let $k \in [0, \infty[$ be a real number. Define
\[
  f_k(t) = \begin{cases}
    t^k \sin\frac{1}{t}, & t \neq 0, \\
    0, & t = 0.
  \end{cases}
\]
Let $A = \{k \in [0, \infty[ \mid f_k \text{ is differentiable}\}$. Then $A = ?$

\qitem{2.4} What is the least value of $K > 0$ such that
\[
  |\sin^2 x - \sin^2 y| \le K |x - y|
\]
for all real numbers $x$ and $y$?

\qitem{2.5} Let $\Gamma$ be the circle in the complex plane with centre at $z = 1$ and of radius unity. Evaluate:
\[
  \int_{\Gamma} \frac{z\,dz}{(z - 1)^4}.
\]

\qitem{2.6} If the plane $\R^2$ is provided with the Lebesgue measure, what is the measure of the set
\[
  A = \{(x, y) \in \R^2 \mid x^2 + y^2 = 1\}?
\]

\qitem{2.7} Pick out the sequences which are uniformly convergent:
\begin{enumerate}[(a)]
  \item $f_n(x) = \sin^n x$ on $[0, \pi/2[$.
  \item $f_n(x) = \frac{x^n}{n} + 1$ on $[0, 1[$.
  \item $f_n(x) = \frac{1}{1 + (x - n)^2}$ on $]-\infty, 0[$.
  \item $f_n(x) = \frac{1}{1 + (x - n)^2}$ on $]0, +\infty[$.
\end{enumerate}

\qitem{2.8} Pick out the functions which are Riemann integrable on the interval $[0, 1]$:
\begin{enumerate}[(a)]
  \item $f(x) = \begin{cases} 1, & \text{if } x \text{ is rational} \\ 0, & \text{if } x \text{ is irrational.} \end{cases}$
  \item $f(x) = \begin{cases} 1, & \text{if } x \in \{\alpha_1, \alpha_2, \ldots, \alpha_n\} \\ 0, & \text{otherwise} \end{cases}$\\
  where $\alpha_1, \ldots, \alpha_n$ are fixed, but arbitrarily chosen numbers in $[0, 1]$.
  \item $f(x) = \begin{cases} 0, & \text{if } x \text{ is irrational or if } x = 0 \\ \sin q\pi, & \text{if } x = p/q,\ p \text{ and } q \text{ positive and coprime integers.} \end{cases}$
\end{enumerate}

\qitem{2.9} Pick out the functions from the following list which are analytic in $\C$:
\begin{enumerate}[(a)]
  \item $f(z) = |z|^2$
  \item $f(z) = \overline{z}$
  \item $f(z) = \RePart(z)$
\end{enumerate}

\qitem{2.10} Pick out the statements which are true:
\begin{enumerate}[(a)]
  \item $|\sin z| \le 1$ for all $z \in \C$.
  \item $\sin^2 z + \cos^2 z = 1$ for all $z \in \C$.
  \item $\sin z = (e^{iz} - e^{-iz})/2$ for all $z \in \C$.
\end{enumerate}

\examsec{Section 3: Topology}

\qitem{3.1} Let $f : ]0, 1[ \to \R$ be continuous. It can be extended to a continuous function $\widetilde{f} : [0, 1] \to \R$ if, and only if, it is \underline{\hspace{3cm}}.

\qitem{3.2} Consider the disjoint closed sets in $\R^2$ given by
\[
  A = \{(x, y) \in \R^2 \mid y = 0\} \quad \text{and} \quad B = \{(x, y) \in \R^2 \mid xy = 1\}.
\]
What is the distance $d(A, B)$ between them?

\medskip
\noindent\emph{In Questions 3.3 and 3.4 below, write `0' if the set $A$ is empty, the exact number of elements in it if the set is finite, and `infinite' if the set is infinite.}

\qitem{3.3} Let $f : [0, 1] \to [0, 1]$ be such that $|f(x) - f(y)| \le \frac{1}{2}|x - y|$ for all $x, y \in [0, 1]$. Let $A = \{x \in [0, 1] \mid f(x) = x\}$. The number of elements in $A$ is \underline{\hspace{2cm}}.

\qitem{3.4} Let $f : [0, 1] \to [0, 1]$ be continuous and such that $f(0) = f(1)$. Let
\[
  A = \{(t, s) \in [0, 1] \times [0, 1] \mid t \neq s \text{ and } f(t) = f(s)\}.
\]
The number of elements in $A$ is \underline{\hspace{2cm}}.

\medskip
\noindent\emph{In Questions 3.5 to 3.10 below, mark a \textbf{tick} over the topological properties true for the set $A$ and \textbf{strike out} the properties that do not hold. Your answer will be treated as correct only if \textbf{all} the choices are correctly made.}

\qitem{3.5} Identify the space of all $n \times n$ matrices (with real entries) with $\R^{n^2}$. Let $A$ be the set of all invertible matrices.\\
Choices: open, closed, connected, dense.

\qitem{3.6} $A = f(B) \subset X$ where $B = \{(x, y) \in \R^2 \mid 1 \le x^2 + y^2 \le 2\}$, $X$ is an arbitrary topological space and $f : \R^2 \to X$ is an arbitrary continuous map.\\
Choices: open, closed, compact, connected.

\qitem{3.7} $A = X \setminus \{x_o\}$ where $X$ is an arbitrary Hausdorff topological space and $x_o \in X$.\\
Choices: open, closed, connected, dense.

\qitem{3.8} $A = f(B) \subset \R$ where $B$ is a closed interval contained in $]0, \infty[$ and $f(t) = \log t$.\\
Choices: open, closed, connected, compact.

\qitem{3.9} $A = \{(x, y) \in \R^2 \mid y = mx\} \setminus \{(0, 0)\} \subset \R^2$.\\
Choices: open, closed, connected, nowhere dense.

\qitem{3.10} $A$ is the closure in $\mathcal{C}[0, 1]$ of the set $B$ where
\[
  B = \{f \in \mathcal{C}^1[0, 1] \mid |f(x)| \le 1 \text{ and } |f'(x)| \le 1 \text{ for all } x \in [0, 1]\}.
\]
Choices: closed, compact, connected, dense.

\examsec{Section 4: Applied Mathematics}

\qitem{4.1} Let $S = \{(x, y, z) \in \R^3 \mid x^2 + y^2 + z^2 = 1\}$. Let $\mathbf{v} = (v_1, v_2, v_3)$ be a solenoidal vector field on $\R^3$. Evaluate:
\[
  \int_S [x(x + v_1(x, y, z)) + y(y + v_2(x, y, z)) + z(z + v_3(x, y, z))]\,dS.
\]

\qitem{4.2} Let $\mathbf{v} = (v_1, v_2, v_3)$ be a vector field on $\R^3$ where $v_1 = \sqrt{1 + x^2 + y^2}$, $v_2 = \sqrt{1 + z^2}$ and $v_3 = \sqrt{1 + x^2 y^2 z^2}$. Evaluate $\divg(\mathbf{curl}\,\mathbf{v})$.

\qitem{4.3} What is the smallest value of $\lambda \in \R$ such that the boundary value problem:
\[
  u''(x) + \lambda u(x) = 0 \quad \text{in } ]0, 1[ \quad \text{and} \quad u(0) = u(1) = 0
\]
has a non-trivial solution (\emph{i.e.}\ $u \not\equiv 0$)?

\qitem{4.4} Let $\mathbf{u}(t) = (u_1(t), u_2(t))$ be the unique solution of the problem:
\begin{align*}
  \frac{d\mathbf{u}}{dt}(t) &= A\mathbf{u}(t), \quad t > 0 \\
  \mathbf{u}(0) &= \mathbf{u}_o
\end{align*}
where $\mathbf{u}_o = (1, 1)$ and $A$ is a symmetric $2 \times 2$ matrix such that $\tr(A) < 0$ and $\det(A) > 0$. Evaluate:
\[
  \lim_{t \to \infty} u_1(t).
\]

\qitem{4.5} Simpson's rule gives the exact value of $\int_0^1 p(t)\,dt$ for every polynomial of degree less than or equal to \underline{\hspace{2cm}}.

\qitem{4.6} Consider the linear programming problem: Maximize $z = 5x + 7y$ such that
\begin{align*}
  x - y &\le 1 \\
  2x + y &\ge 2 \\
  x + 2y &\le 4 \\
  x \ge 0,&\quad y \ge 0.
\end{align*}
What is the optimal value of $z$?

\qitem{4.7} According to the classification of second order linear partial differential operators, the operator
\[
  \frac{\partial^2 u}{\partial x^2} - 4\frac{\partial^2 u}{\partial x \partial y} + 5\frac{\partial^2 u}{\partial y^2}
\]
is of \underline{\hspace{2cm}} type.

\qitem{4.8} Evaluate:
\[
  \iint_{\R^2} e^{-(x+2y)^2 - (x+y)^2}\,dxdy.
\]

\qitem{4.9} A necessary and sufficient condition that the boundary value problem:
\begin{align*}
  \frac{\partial^2 u}{\partial x^2} + \frac{\partial^2 u}{\partial y^2} &= f(x, y) \quad \text{in } \Omega \\
  \frac{\partial u}{\partial n} &= 0 \quad \text{on } \partial\Omega
\end{align*}
(where $\Omega \subset \R^2$ is a bounded domain with boundary $\partial\Omega$ and $\frac{\partial u}{\partial n}$ denotes the outer normal derivative of the function $u$) has a solution is \underline{\hspace{2cm}}.

\qitem{4.10} The radius $r$ and height $h$ of a right circular cylinder of fixed volume $V$ and least total surface area are connected by the relation \underline{\hspace{2cm}}.

\examsec{Section 5: Miscellaneous}

\qitem{5.1} What is the maximum number of pieces that a pizza can be cut into by $7$ knife strokes?

\qitem{5.2} Let $n$ be a fixed positive integer. Let $C_r$ denote the number of ways of choosing $r$ objects from a collection of $n$ objects. Evaluate:
\[
  C_1 + 2\cdot C_2 + \dots + n\cdot C_n.
\]

\qitem{5.3} Inside a square of side $2$ units, five points are marked at random. What is the probability that there are at least two points such that the distance between them is at most $\sqrt{2}$ units?

\qitem{5.4} What is the area of the triangle in the complex plane formed by the points representing $1, \omega$ and $\omega^2$, where $\omega$ is a complex cube root of unity?

\qitem{5.5} What is the number of points of intersection, in $\R^2$, of the two plane curves $(1 + x^2 + y^2)(x^2 + y^2 - 4) = 0$ and $y = 7x$?

\qitem{5.6} What geometric figure is formed by the locus of a point which moves so that the sum of four times its distance from the $x$-axis and nine times its distance from the $y$-axis is equal to $10$?

\qitem{5.7} A real number is \emph{algebraic} if it is the root of a polynomial with integer coefficients. Define $A : [0, 1] \to \R$ by
\[
  A(x) = \begin{cases} 1 & \text{if } x \text{ is algebraic} \\ 0 & \text{otherwise.} \end{cases}
\]
Evaluate: $\int_0^1 A(x)\,dx$.

\qitem{5.8} In the rectangle $[0, \pi/2] \times [0, 1] \subset \R^2$, a point $(x, y)$ is chosen at random. What is the probability that $y \le \sin x$?

\qitem{5.9} If $p$ is a prime greater than, or equal to, $11$, then, either $p^3 - 1$ or $p^3 + 1$ is divisible by $14$. True or False?

\qitem{5.10} Evaluate:
\[
  \sum_{n=1}^{\infty} \frac{n^2 - n + 1}{n!}.
\]

\examsec{Answer Key --- NBHM 2005}
\begin{multicols}{2}
\textbf{Section 1: Algebra}\\
\textbf{1.1} $5$\\
\textbf{1.2} $5$\\
\textbf{1.3} $\alpha = -1/2,\ \beta = 1,\ \gamma = 1/2$\\
\textbf{1.4} $2$\\
\textbf{1.5} $4$\\
\textbf{1.6} $12$\\
\textbf{1.7} (a), (b)\\
\textbf{1.8} (a) (c)\\
\textbf{1.9} (a) (c) (d)\\
\textbf{1.10} (a) (c) (d)

\medskip
\textbf{Section 2: Analysis}\\
\textbf{2.1} $\int_0^1 f(t)\,dt$\\
\textbf{2.2} $\infty$\\
\textbf{2.3} $]1, \infty[$\\
\textbf{2.4} $1$\\
\textbf{2.5} $0$\\
\textbf{2.6} $0$\\
\textbf{2.7} (b) (c)\\
\textbf{2.8} (b) (c)\\
\textbf{2.9} None\\
\textbf{2.10} (b)

\medskip
\textbf{Section 3: Topology}\\
\textbf{3.1} uniformly continuous\\
\textbf{3.2} $0$\\
\textbf{3.3} $1$\\
\textbf{3.4} infinite\\
\textbf{3.5} open, dense\\
\textbf{3.6} compact, connected\\
\textbf{3.7} open\\
\textbf{3.8} closed, connected\\
\textbf{3.9} nowhere dense\\
\textbf{3.10} closed, compact, connected

\medskip
\textbf{Section 4: Applied Mathematics}\\
\textbf{4.1} $4\pi$\\
\textbf{4.2} $0$\\
\textbf{4.3} $\pi^2$\\
\textbf{4.4} $0$\\
\textbf{4.5} $3$\\
\textbf{4.6} $17$\\
\textbf{4.7} elliptic\\
\textbf{4.8} $\pi$\\
\textbf{4.9} $\int_{\Omega} f(x, y)\,dxdy = 0$\\
\textbf{4.10} $h = 2r$

\medskip
\textbf{Section 5: Miscellaneous}\\
\textbf{5.1} $29$\\
\textbf{5.2} $n 2^{n-1}$\\
\textbf{5.3} $1$\\
\textbf{5.4} $3\sqrt{3}/4$\\
\textbf{5.5} $2$\\
\textbf{5.6} parallelogram\\
\textbf{5.7} $0$\\
\textbf{5.8} $2/\pi$\\
\textbf{5.9} True\\
\textbf{5.10} $2e - 1$

\medskip
\noindent\emph{Note: Correct answers to Qns.\ 2.1, 2.3 and 4.9 in any other equivalent notation can, obviously, be accepted!!!}
\end{multicols}
% ==================== NBHM 2006 ====================
\examheader{NBHM 2006}{Research Awards Screening Test}{February 25, 2006}{Time Allowed: 90 Minutes \quad|\quad Maximum Marks: 40}

\subsection*{Instructions to Candidates}
\begin{itemize}
  \item Do not forget to write your name and roll number on the cover page. In the box marked `Institution', fill in the name of the institution where you are working towards a Ph.D.\ degree. In case you have not yet joined any institution for research, write \emph{Not Applicable}.
  \item Please ensure that your answer booklet contains 16 numbered (and printed) pages. The back of each printed page is blank and can be used for rough work.
  \item There are \textbf{five} sections, containing \textbf{ten} questions each, entitled Algebra, Analysis, Topology, Applied Mathematics and Miscellaneous. Answer as many questions as possible. The assessment of the paper will be based on the best \textbf{four} sections. Each question carries one point and the maximum marks to be scored is \textbf{forty}.
  \item Answer each question, as directed, in the space provided at the end of it. Answers are to be given in the form of a word (or words, if required), a numerical value (or values) or a simple mathematical expression. \textbf{Do not write sentences.}
  \item In certain questions you are required to pick out the qualifying statement(s) from multiple choices. None of the statements, or more than one statement may qualify. Write \textbf{none} if none of the statements qualify, or list the labels of \emph{all} the qualifying statements (amongst (a), (b), (c) and (d)).
  \item Points will be awarded in the above questions only if \textbf{all} the correct choices are made. There will be no partial credit.
  \item $\N$ denotes the set of natural numbers, $\Z$ -- the integers, $\Q$ -- the rationals, $\R$ -- the reals and $\C$ -- the field of complex numbers. $\R^n$ denotes the $n$-dimensional Euclidean space, which is assumed to be endowed with its `usual' topology. The symbol $]a, b[$ will stand for the open interval $\{x \in \R \mid a < x < b\}$ while $[a, b]$ will stand for the corresponding closed interval; $[a, b[$ and $]a, b]$ will stand for the corresponding left-closed-right-open and left-open-right-closed intervals respectively. The symbol $I$ will denote the identity matrix of appropriate order.
\end{itemize}

\examsec{Section 1: Algebra}

\qitem{1.1} Let $f : (\Q, +) \to (\Q, +)$ be a non-zero homomorphism. Pick out the true statements:
\begin{enumerate}[a.]
  \item $f$ is always one-one.
  \item $f$ is always onto.
  \item $f$ is always a bijection.
  \item $f$ need be neither one-one nor onto.
\end{enumerate}

\qitem{1.2} Consider the element
\[
  \alpha = \begin{pmatrix} 1 & 2 & 3 & 4 & 5 \\ 2 & 1 & 4 & 5 & 3 \end{pmatrix}
\]
of the symmetric group $S_5$ on five elements. Pick out the true statements:
\begin{enumerate}[a.]
  \item The order of $\alpha$ is $5$.
  \item $\alpha$ is conjugate to $\begin{pmatrix} 4 & 5 & 2 & 3 & 1 \\ 5 & 4 & 3 & 1 & 2 \end{pmatrix}$.
  \item $\alpha$ is the product of two cycles.
  \item $\alpha$ commutes with all elements of $S_5$.
\end{enumerate}

\qitem{1.3} Let $G$ be a group of order $60$. Pick out the true statements:
\begin{enumerate}[a.]
  \item $G$ is abelian.
  \item $G$ has a subgroup of order $30$.
  \item $G$ has subgroups of order $2, 3$ and $5$.
  \item $G$ has subgroups of order $6, 10$ and $15$.
\end{enumerate}

\qitem{1.4} Consider the polynomial ring $R[x]$ where $R = \Z/12\Z$ and write the elements of $R$ as $\{0, 1, \dots, 11\}$. Write down all the distinct roots of the polynomial $f(x) = x^2 + 7x$ of $R[x]$.

\qitem{1.5} Let $R$ be the polynomial ring $\Z_2[x]$ and write the elements of $\Z_2$ as $\{0, 1\}$. Let $(f(x))$ denote the ideal generated by the element $f(x) \in R$. If $f(x) = x^2 + x + 1$, then the quotient ring $R/(f(x))$ is
\begin{enumerate}[a.]
  \item a ring but not an integral domain.
  \item an integral domain but not a field.
  \item a finite field of order $4$.
  \item an infinite field.
\end{enumerate}

\qitem{1.6} Consider the set of all linear transformations $T : \R^3 \to \R^4$ over $\R$. What is the dimension of this set, considered as a vector space over $\R$ with point-wise operations?

\qitem{1.7} Consider the matrix $A = \begin{bmatrix} 1 & 1 & 0 \\ 0 & 2 & 3 \\ 0 & 0 & 3 \end{bmatrix}$. Write down a matrix $P$ such that $P^{-1}AP$ is a diagonal matrix.

\qitem{1.8} Let $A$ be an orthogonal $3 \times 3$ matrix with real entries. Pick out the true statements:
\begin{enumerate}[a.]
  \item The determinant of $A$ is a rational number.
  \item $d(Ax, Ay) = d(x, y)$ for any two vectors $x$ and $y \in \R^3$, where $d(u, v)$ denotes the usual Euclidean distance between vectors $u$ and $v \in \R^3$.
  \item All the entries of $A$ are positive.
  \item All the eigenvalues of $A$ are real.
\end{enumerate}

\qitem{1.9} Pick out the correct statements from the following list:
\begin{enumerate}[a.]
  \item A homomorphic image of a UFD (unique factorization domain) is again a UFD.
  \item The element $2 \in \Z[\sqrt{-5}]$ is irreducible in $\Z[\sqrt{-5}]$.
  \item Units of the ring $\Z[\sqrt{-5}]$ are the units of $\Z$.
  \item The element $2$ is a prime element in $\Z[\sqrt{-5}]$.
\end{enumerate}

\qitem{1.10} Let $p$ and $q$ be two distinct primes. Pick the correct statements from the following:
\begin{enumerate}[a.]
  \item $\Q(\sqrt{p})$ is isomorphic to $\Q(\sqrt{q})$ as fields.
  \item $\Q(\sqrt{p})$ is isomorphic to $\Q(\sqrt{-q})$ as vector spaces over $\Q$.
  \item $[\Q(\sqrt{p}, \sqrt{q}) : \Q] = 4$.
  \item $\Q(\sqrt{p}, \sqrt{q}) = \Q(\sqrt{p} + \sqrt{q})$.
\end{enumerate}

\examsec{Section 2: Analysis}

\qitem{2.1} Let $f$ be a real valued function on $\R$. Consider the functions
\[
  w_j(x) = \sup\left\{|f(u) - f(v)| : u, v \in \left[x - \frac{1}{j}, x + \frac{1}{j}\right]\right\},
\]
where $j$ is a positive integer and $x \in \R$. Define next,
\[
  A_{j,n} = \left\{x \in \R : w_j(x) < \frac{1}{n}\right\},\quad n = 1, 2, \dots
\]
and
\[
  A_n = \bigcup_{j=1}^{\infty} A_{j,n},\quad n = 1, 2, \dots
\]
Now let
\[
  C = \{x \in \R : f \text{ is continuous at } x\}.
\]
Express $C$ in terms of the sets $A_n$.

\qitem{2.2} Let $f$ be a continuous real valued function on $\R$ and $n$, a positive integer. Find
\[
  \frac{d}{dx} \int_0^x (2x - t)^n f(t)\,dt.
\]

\qitem{2.3} For each $n \ge 1$, let $f_n$ be a monotonic increasing real valued function on $[0, 1]$ such that the sequence of functions $\{f_n\}$ converges pointwise to the function $f \equiv 0$. Pick out the true statements from the following:
\begin{enumerate}[a.]
  \item $f_n$ converges to $f$ uniformly.
  \item If the functions $f_n$ are also non-negative, then $f_n$ must be continuous for sufficiently large $n$.
\end{enumerate}

\qitem{2.4} Let $\Q$ denote the set of all rational numbers in the open interval $]0, 1[$. Let $\lambda(U)$ denote the Lebesgue measure of a subset $U$ of $]0, 1[$. Pick out the correct statements from the following:
\begin{enumerate}[a.]
  \item $\lambda(U) = 1$ for every open set $U \subset ]0, 1[$ which contains $\Q$.
  \item Given any $\varepsilon > 0$, there exists an open set $U \subset ]0, 1[$ containing $\Q$ such that $\lambda(U) < \varepsilon$.
\end{enumerate}

\qitem{2.5} A real valued function on an interval $[a, b]$ is said to be a function of bounded variation if there exists $M > 0$, such that for any finite set of points $a = a_0 < a_1 < a_2 < \dots < a_n = b$, we have $\sum_{i=0}^{n-1} |f(a_i) - f(a_{i+1})| < M$. Which of the following statements are necessarily true?
\begin{enumerate}[a.]
  \item Any continuous function on $[0, 1]$ is of bounded variation.
  \item If $f : \R \to \R$ is continuously differentiable, then its restriction to the interval $[-n, n]$ is of bounded variation on that interval, for any positive integer $n$.
  \item Any monotone function on $[0, 1]$ is of bounded variation.
\end{enumerate}

\qitem{2.6} Let $f$ be a differentiable function of one variable and let $g$ be the function of two variables given by $g(x, y) = f(ax + by)$, where $a, b$ are fixed nonzero numbers. Write down a partial differential equation satisfied by the function $g$.

\qitem{2.7} The curve $x^3 - y^3 = 1$ is asymptotic to the line $x = y$. Find the point on the curve farthest from the line $x = y$.

\qitem{2.8} Let $k$ be a fixed positive integer. Find $R_k$, the radius of convergence of the power series $\sum \left(\frac{n+1}{n}\right)^{n^2} z^{kn}$.

\qitem{2.9} Let $\gamma$ be a closed and continuously differentiable path in the upper half plane
\[
  \{z \in \C : z = x + iy,\ x, y \in \R,\ y > 0\}
\]
not passing through the point $i$. Describe the set of all possible values of the integral
\[
  \frac{1}{2\pi i} \int_{\gamma} \frac{2i}{z^2 + 1}\,dz.
\]

\qitem{2.10} Let $f$ be a function of three (real) variables having continuous partial derivatives. For each direction vector $h = (h_1, h_2, h_3)$ such that $h_1^2 + h_2^2 + h_3^2 = 1$, let $D_h f(x, y, z)$ be the directional derivative of $f$ along $h$ at $(x, y, z)$. For a point $(x_0, y_0, z_0)$ where the partial derivative $\frac{\partial}{\partial x} f(x_0, y_0, z_0)$ is not zero, maximize $D_h f(x_0, y_0, z_0)$ (as a function of $h$).

\examsec{Section 3: Topology}

\qitem{3.1} Let $f$ be the function on $\R$ defined by $f(t) = \frac{p+\sqrt{2}}{q+\sqrt{2}} - \frac{p}{q}$ if $t = \frac{p}{q}$ with $p, q \in \Z$ and $p$ and $q$ coprime to each other, and $f(t) = 0$ if $t$ is irrational. Answer the following:
\begin{enumerate}[i)]
  \item At which irrational numbers $t$ is $f$ continuous?
  \item At which rational numbers $t$ is $f$ continuous?
\end{enumerate}

\qitem{3.2} Let $f$ and $g$ be two continuous functions on $\R$. For any $a \in \R$ we define $J_a(f, g)$ to be the function given by $J_a(f, g)(t) = f(t)$ for all $t \le a$ and $J_a(f, g)(t) = g(t)$ if $t > a$. For what values of $a$ is $J_a(f, g)$ a continuous function?

\qitem{3.3} Let $A$ and $B$ be two finite subsets of $\R$. Describe a necessary and sufficient condition for the spaces $\R \setminus A$ and $\R \setminus B$ to be homeomorphic.

\qitem{3.4} Let $f : \R^2 \to \R$ be a continuous function. Let $D$ be the closed unit disc in $\R^2$. Is $f(D)$ necessarily an interval in $\R$? If it is an interval, which of the forms $]a, b[$, $[a, b[$, $]a, b]$ and $[a, b]$, with $a, b \in \R$ can it have?

\qitem{3.5} For $v \in \R^2$ and $r > 0$ let $D(v, r)$ denote the closed disc with centre at $v$ and radius $r$. Let $v = (5, 0) \in \R^2$. For $\alpha > 0$ let $X_\alpha$ be the subset
\[
  X_\alpha = D(-v, 3) \cup D(v, 3) \cup \{(x, \alpha x) : x \in \R\}.
\]
Determine the condition on $\alpha$ for $X_\alpha$ to be connected; when it is not connected how many connected components does $X_\alpha$ have?

\qitem{3.6} Which two of the following spaces are homeomorphic to each other?
\begin{enumerate}[i)]
  \item $X_1 = \{(x, y) \in \R^2 : xy = 0\}$;
  \item $X_2 = \{(x, y) \in \R^2 : x + y \ge 0 \text{ and } xy = 0\}$;
  \item $X_3 = \{(x, y) \in \R^2 : xy = 1\}$;
  \item $X_4 = \{(x, y) \in \R^2 : x + y \ge 0 \text{ and } xy = 1\}$.
\end{enumerate}

\qitem{3.7} Which of the following spaces are compact?
\begin{enumerate}[i)]
  \item $X_1 = \{(x, y) \in \R^2 : |x| + |y| < 10^{-100}\}$;
  \item $X_2 = \{(x, y) \in \R^2 : |x| + |y| \le 10^{100}\}$;
  \item $X_3 = \{(x, y) \in \R^2 : 1 \le x^2 + y^2 \le 2\}$;
  \item $X_4 = \{(x, y) \in \R^2 : x^2 + y^2 = 1 \text{ and } xy \neq 0\}$.
\end{enumerate}

\qitem{3.8} Which of the following spaces are locally compact?
\begin{enumerate}[i)]
  \item $X_1 = \{(x, y) \in \R^2 : x, y \text{ odd integers}\}$;
  \item $X_2 = \{(x, y) \in \R^2 : x^2 + 103xy + 7y^2 > 5\}$;
  \item $X_3 = \{(x, y) \in \R^2 : 0 \le x < 1,\ 0 < y \le 1\}$;
  \item $X_4 = \{(x, y) \in \R^2 : x, y \text{ irrational}\}$.
\end{enumerate}

\qitem{3.9} Which of the following metric spaces $(X_i, d_i)$, $1 \le i \le 4$, are complete?
\begin{enumerate}[i)]
  \item $X_1 = ]0, \pi/2[ \subset \R$, $d_1$ defined by $d_1(x, y) = |\tan x - \tan y|$ for all $x, y \in X_1$.
  \item $X_2 = [0, 1] \subset \R$, $d_2$ defined by $d_2(x, y) = \frac{|x-y|}{1+|x-y|}$ for all $x, y \in X_2$.
  \item $X_3 = \Q$, and $d_3$ defined by $d_3(x, y) = 1$ for all $x, y \in X_3$, $x \neq y$.
  \item $X_4 = \R$, $d_4$ defined by $d_4(x, y) = |e^x - e^y|$ for all $x, y \in X_4$.
\end{enumerate}

\qitem{3.10} On which of the following spaces is every continuous (real-valued) function bounded?
\begin{enumerate}[i)]
  \item $X_1 = ]0, 1[$;
  \item $X_2 = [0, 1]$;
  \item $X_3 = [0, 1[$;
  \item $X_4 = \{t \in [0, 1] : t \text{ irrational}\}$.
\end{enumerate}

\examsec{Section 4: Applied Mathematics}

\qitem{4.1} Let $\Gamma(s)$ stand for the usual Gamma function. Given that $\Gamma(1/2) = \sqrt{\pi}$, evaluate $\Gamma(5/2)$.

\qitem{4.2} Let $S = \{(x, y, z) \in \R^3 : x^2 + y^2 + z^2 = 1,\ z > 0\}$. Let $C = \{(x, y) \in \R^2 : x^2 + y^2 = 1\}$. Let $\tau$ be the unit tangent vector to $C$ in the $xy$-plane pointing left as we move clockwise along $C$. Let $\varphi(x, y, z) = x^2 + y^3 + z^4$. Evaluate:
\[
  \int_C \nabla \varphi \cdot \tau\,ds.
\]

\qitem{4.3} Let $a > 0$ and let $S = \{(x, y, z) \in \R^3 : x^2 + y^2 + z^2 = a^2\}$. Evaluate:
\[
  \iint_S (x^4 + y^4 + z^4)\,dS.
\]

\qitem{4.4} Let $f(x) = x^2 - 5$ for $x \in \R$. Let $x_0 = 1$. If $\{x_n\}$ denotes the sequence of iterates defined by the Newton-Raphson method to approximate a solution of $f(x) = 0$, find $x_1$.

\qitem{4.5} Let $A$ be a $2 \times 2$ matrix with real entries. Consider the linear system of ordinary differential equations given in vector notation as:
\[
  \frac{d\mathbf{x}}{dt}(t) = A\mathbf{x}(t), \qquad \text{where } \mathbf{x}(t) = \begin{pmatrix} u(t) \\ v(t) \end{pmatrix}.
\]
Pick out the cases from the following when we have $\lim_{t \to \infty} u(t) = 0$ and $\lim_{t \to \infty} v(t) = 0$:
\begin{enumerate}[a.]
  \item $A = \begin{pmatrix} 1 & 2 \\ 0 & 3 \end{pmatrix}$.
  \item $A = \begin{pmatrix} -1 & 2 \\ 0 & -3 \end{pmatrix}$.
  \item $A = \begin{pmatrix} 1 & -6 \\ 1 & -4 \end{pmatrix}$.
  \item $A = \begin{pmatrix} -1 & -6 \\ 1 & 4 \end{pmatrix}$.
\end{enumerate}

\qitem{4.6} Let $\Delta = \partial^2/\partial x^2 + \partial^2/\partial y^2$ denote the Laplace operator. Let $\Omega = \{(x, y) \in \R^2 : x^2 + y^2 < 1\}$. Let $\partial\Omega$ denote the boundary of the domain $\Omega$. Consider the following boundary value problem:
\begin{align*}
  \Delta u &= c \quad \text{in } \Omega \\
  \frac{\partial u}{\partial \nu} &= 1 \quad \text{on } \partial\Omega
\end{align*}
where $c$ is a real constant and $\partial u/\partial \nu$ denotes the outward normal derivative of $u$ on $\partial\Omega$. For what values of $c$ does the above problem admit a solution?

\qitem{4.7} Consider the Tricomi equation:
\[
  \frac{\partial^2 u}{\partial y^2} - y \frac{\partial^2 u}{\partial x^2} = 0.
\]
Describe the region in the $xy$-plane where this equation is elliptic.

\qitem{4.8} Evaluate:
\[
  \iint_{\R^2} e^{-(3x+2y)^2 - (4x+y)^2}\,dxdy.
\]

\qitem{4.9} Let $J_p$ denote the Bessel function of the first kind, of order $p$ and let $\{P_n\}$ denote the sequence of Legendre polynomials defined on the interval $[-1, 1]$. Pick out the true statements from the following:
\begin{enumerate}[a.]
  \item $\frac{d}{dx} J_0(x) = -J_1(x)$.
  \item Between any two positive zeroes of $J_0$, there exists a zero of $J_1$.
  \item $P_{n+1}(x)$ can be written as a linear combination of $P_n(x)$ and $P_{n-1}(x)$.
  \item $P_{n+1}(x)$ can be written as a linear combination of $xP_n(x)$ and $P_{n-1}(x)$.
\end{enumerate}

\qitem{4.10} Consider the linear programming problem: Maximize $z = 2x_1 + 3x_2 + x_3$ such that
\begin{align*}
  4x_1 + 3x_2 + x_3 &= 6 \\
  x_1 + 2x_2 + 5x_3 &\ge 4 \\
  x_1, x_2, x_3 &\ge 0.
\end{align*}
Write down the objective function of the dual problem.

\examsec{Section 5: Miscellaneous}

\qitem{5.1} A unimodular matrix is a matrix with integer entries and having determinant $1$ or $-1$. If $m$ and $n$ are positive integers, write down a necessary and sufficient condition such that there exists a unimodular matrix of order $2$ whose first row is the vector $(m, n)$.

\qitem{5.2} For any integer $n$ define $k(n) = \frac{n^7}{7} + \frac{n^3}{3} + \frac{11n}{21} + 1$ and
\[
  f(n) = \begin{cases} 0 & \text{if } k(n) \text{ is an integer,} \\ \frac{1}{n^2} & \text{if } k(n) \text{ is not an integer.} \end{cases}
\]
Find $\sum_{n=-\infty}^{\infty} f(n)$.

\qitem{5.3} Let $n \ge 2$. Evaluate:
\[
  \sum_{k=2}^{n} \frac{n!}{(n-k)!(k-2)!}.
\]

\qitem{5.4} A fair coin is tossed ten times. What is the probability that we can observe a string of eight heads, in succession, at some time?

\qitem{5.5} Evaluate the product $\prod_{n=2}^{\infty} \left(1 + \frac{1}{n^2} + \frac{1}{n^4} + \frac{1}{n^6} + \dots\right)$.

\qitem{5.6} Find all solutions of the equation
\[
  (x^2 + y^2 + z^2 - 1)^2 + (x + y + z - 3)^2 = 0.
\]

\qitem{5.7} For any real number $x$, let $f(x)$ denote the distance of $x$ from the nearest integer. Let $I(k) = [k\pi, k\pi + 1]$. Find $f(I(k))$ for all integers $k$.

\qitem{5.8} Let $K$ be a finite field. Can you always find a non-constant polynomial over $K$ which has no root in $K$? If yes, give one such polynomial.

\qitem{5.9} Evaluate:
\[
  \sum_{k=1}^{\infty} \frac{k^2}{k!}.
\]

\qitem{5.10} Pick out the countable sets from the following:
\begin{enumerate}[a.]
  \item $\{\alpha \in \R : \alpha \text{ is a root of a polynomial with integer coefficients}\}$.
  \item The complement in $\R$ of the set described in statement (a) above.
  \item The set of all points in the plane whose coordinates are rational.
  \item Any subset of $\R$ whose Lebesgue measure is zero.
\end{enumerate}

\examsec{Answer Key --- NBHM 2006}
\begin{multicols}{2}
\textbf{Section 1: Algebra}\\
\textbf{1.1} a, b, c\\
\textbf{1.2} b, c\\
\textbf{1.3} c\\
\textbf{1.4} $0, 5, 8, 9$\\
\textbf{1.5} c\\
\textbf{1.6} $12$\\
\textbf{1.7} Any matrix of the form ($a, b, c$ all non-zero):
$\begin{bmatrix} a & b & c/2 \\ 0 & b & c \\ 0 & 0 & c/3 \end{bmatrix}$\\
\textbf{1.8} a, b\\
\textbf{1.9} b, c\\
\textbf{1.10} b, c, d

\medskip
\textbf{Section 2: Analysis}\\
\textbf{2.1} $C = \bigcap_{n=1}^{\infty} A_n$\\
\textbf{2.2} $2n\int_0^x (2x - t)^{n-1} f(t)\,dt + x^n f(x)$\\
\textbf{2.3} a\\
\textbf{2.4} b\\
\textbf{2.5} b, c\\
\textbf{2.6} $b\frac{\partial g}{\partial x} = a\frac{\partial g}{\partial y}$\\
\textbf{2.7} $\left(\frac{1}{2^{1/3}}, -\frac{1}{2^{1/3}}\right)$\\
\textbf{2.8} $e^{-\frac{1}{k}}$\\
\textbf{2.9} All integers\\
\textbf{2.10} $\left[\left(\frac{\partial f}{\partial x}(x_0, y_0, z_0)\right)^2 + \left(\frac{\partial f}{\partial y}(x_0, y_0, z_0)\right)^2 + \left(\frac{\partial f}{\partial z}(x_0, y_0, z_0)\right)^2\right]^{\frac{1}{2}}$

\medskip
\textbf{Section 3: Topology}\\
\textbf{3.1} (i) continuous at all irrationals, (ii) continuous only at $t = 1$\\
\textbf{3.2} $f(a) = g(a)$\\
\textbf{3.3} $A$ and $B$ have the same cardinality\\
\textbf{3.4} (i) $f(D)$ is necessarily an interval; (ii) $[a, b]$\\
\textbf{3.5} (i) $X_\alpha$ is connected if and only if $\alpha \le \frac{3}{4}$. (ii) When not connected, it has $3$ components\\
\textbf{3.6} $X_2$ and $X_4$ are homeomorphic\\
\textbf{3.7} Compact sets are $X_2$ and $X_3$\\
\textbf{3.8} Locally compact sets are $X_1, X_2$ and $X_3$\\
\textbf{3.9} Complete metric spaces are $X_1, X_2$ and $X_3$\\
\textbf{3.10} $X_2$

\medskip
\textbf{Section 4: Applied Mathematics}\\
\textbf{4.1} $\frac{3}{4}\sqrt{\pi}$\\
\textbf{4.2} $0$\\
\textbf{4.3} $\frac{12\pi}{5}a^6$\\
\textbf{4.4} $3$\\
\textbf{4.5} b, c\\
\textbf{4.6} $c = 2$\\
\textbf{4.7} elliptic in the region $\{(x, y) \in \R^2 : y < 0\}$\\
\textbf{4.8} $\pi/5$\\
\textbf{4.9} a, b, d\\
\textbf{4.10} A linear functional in $3$ variables with coefficients $6, -6$ and $-4$; example: $6w_1 - 6w_2 - 4w_3$

\medskip
\textbf{Section 5: Miscellaneous}\\
\textbf{5.1} $m$ and $n$ are coprime\\
\textbf{5.2} $0$\\
\textbf{5.3} $n(n-1)2^{n-2}$\\
\textbf{5.4} $2^{-7}$\\
\textbf{5.5} $2$\\
\textbf{5.6} there is no solution\\
\textbf{5.7} $[0, 1/2]$ for each $k$\\
\textbf{5.8} Yes; if $K = \{a_1, \dots, a_n\}$, then take $(x - a_1)\dots(x - a_n) + 1$, for example.\\
\textbf{5.9} $2e$\\
\textbf{5.10} a, c
\end{multicols}
% ==================== NBHM 2007 ====================
\examheader{NBHM 2007}{Research Scholarships Screening Test}{January 27, 2007}{Time Allowed: Two Hours \quad|\quad Maximum Marks: 40}

\subsection*{Instructions to Candidates}
\begin{itemize}
  \item Please ensure that this booklet contains 12 numbered (and printed) pages. The back of each printed page is blank and can be used for rough work.
  \item There are \textbf{five} sections, containing \textbf{ten} questions each, entitled Algebra, Analysis, Topology, Applied Mathematics and Miscellaneous. Answer as many questions as possible. The assessment of the paper will be based on the best \textbf{four} sections. Each question carries one point and the maximum possible score is \textbf{forty}.
  \item Answer each question, as directed, in the space provided in the \textbf{answer booklet}, which is being supplied separately. This question paper is meant to be retained by you and so \textbf{do not} answer questions on it.
  \item In certain questions you are required to pick out the qualifying statement(s) from multiple choices. None of the statements, or more than one statement may qualify. Write \textbf{none} if none of the statements qualify, or list the labels of \emph{all} the qualifying statements (amongst (a), (b), and (c)).
  \item Points will be awarded in the above questions only if \textbf{all} the correct choices are made. \textbf{There will be no partial credit.}
  \item $\N$ denotes the set of natural numbers, $\Z$ -- the integers, $\Q$ -- the rationals, $\R$ -- the reals and $\C$ -- the field of complex numbers. $\R^n$ denotes the $n$-dimensional Euclidean space, which is assumed to be endowed with its `usual' topology. The symbol $]a, b[$ will stand for the open interval $\{x \in \R \mid a < x < b\}$ while $[a, b]$ will stand for the corresponding closed interval; $[a, b[$ and $]a, b]$ will stand for the corresponding left-closed-right-open and left-open-right-closed intervals respectively. The symbol $I$ will denote the identity matrix of appropriate order. The space of continuous real valued functions on an interval $[a, b]$ is denoted by $\mathcal{C}[a, b]$ and is endowed with its usual `sup' norm. The space of continuously differentiable real valued functions on $[a, b]$ is denoted by $\mathcal{C}^1[a, b]$ and its usual norm is the maximum of the sup-norms of the function and its derivative.
\end{itemize}

\examsec{Section 1: ALGEBRA}

\qitem{1.1} Let $G$ be a group of order $n$. Which of the following conditions imply that $G$ is abelian?
\begin{enumerate}[a.]
  \item $n = 15$.
  \item $n = 21$.
  \item $n = 36$.
\end{enumerate}

\qitem{1.2} Which of the following subgroups are necessarily normal subgroups?
\begin{enumerate}[a.]
  \item The kernel of a group homomorphism.
  \item The center of a group.
  \item The subgroup consisting of all matrices with positive determinant in the group of all invertible $n \times n$ matrices with real entries (under matrix multiplication).
\end{enumerate}

\qitem{1.3} List all the units in the ring of Gaussian integers.

\qitem{1.4} List all possible values occurring as $\deg f$ (degree of $f$) where $f$ is an irreducible polynomial in $\R[x]$.

\qitem{1.5} Write down an irreducible polynomial of degree $3$ over the field $\F_3$ of three elements.

\qitem{1.6} Let $A = (a_{ij})$ be an $n \times n$ matrix with real entries. Let $A_{ij}$ be the cofactor of the entry $a_{ij}$ of $A$. Let $\widetilde{A} = (A_{ij})$ be the matrix of cofactors. What is the rank of $\widetilde{A}$ under the following conditions:
\begin{enumerate}[(a)]
  \item the rank of $A$ is $n$?
  \item the rank of $A$ is $\le n - 2$?
\end{enumerate}

\qitem{1.7} Let $A$ be an $n \times n$ matrix with complex entries which is not a diagonal matrix. Pick out the cases when $A$ is diagonalizable.
\begin{enumerate}[a.]
  \item $A$ is idempotent.
  \item $A$ is nilpotent.
  \item $A$ is unitary.
\end{enumerate}

\qitem{1.8} For $n \ge 2$, let $\mathcal{M}(n)$ denote the ring of all $n \times n$ matrices with real entries. Which of the following statements are true?
\begin{enumerate}[a.]
  \item If $A \in \mathcal{M}(2)$ is nilpotent and non-zero, then there exists a matrix $B \in \mathcal{M}(2)$ such that $B^2 = A$.
  \item If $A \in \mathcal{M}(n)$, $n \ge 2$, is symmetric and positive definite, then there exists a symmetric matrix $B \in \mathcal{M}(n)$ such that $B^2 = A$.
  \item If $A \in \mathcal{M}(n)$, $n \ge 2$, is symmetric, then there exists a symmetric matrix $B \in \mathcal{M}(n)$ such that $B^3 = A$.
\end{enumerate}

\qitem{1.9} Which of the following matrices are non-singular?
\begin{enumerate}[a.]
  \item $I + A$ where $A \neq 0$ is a skew-symmetric real $n \times n$ matrix, $n \ge 2$.
  \item Every skew-symmetric non-zero real $5 \times 5$ matrix.
  \item Every skew-symmetric non-zero real $2 \times 2$ matrix.
\end{enumerate}

\qitem{1.10} Let $V$ be a real finite-dimensional vector space and $f$ and $g$ non-zero linear functionals on $V$. Assume that $\ker(f) \subset \ker(g)$. Pick out the true statements.
\begin{enumerate}[a.]
  \item $\ker(f) = \ker(g)$.
  \item $\ker(g)/\ker(f) \cong \R^k$ for some $k$ such that $1 \le k < n$.
  \item There exists a constant $c \neq 0$ such that $g = cf$.
\end{enumerate}

\examsec{Section 2: ANALYSIS}

\qitem{2.1} In each of the following cases, state whether the series is absolutely convergent, conditionally convergent (\emph{i.e.}\ convergent but not absolutely convergent) or divergent.
\begin{enumerate}[a.]
  \item $\displaystyle\sum_{n=1}^{\infty} (-1)^n \frac{1}{2n+3}$.
  \item $\displaystyle\sum_{n=1}^{\infty} (-1)^n \frac{n}{n+2}$.
  \item $\displaystyle\sum_{n=1}^{\infty} (-1)^n \frac{n\log n}{e^n}$.
\end{enumerate}

\qitem{2.2} Determine the interval of convergence of the series:
\[
  \sum_{n=1}^{\infty} (-1)^{n-1} \frac{(x-1)^n}{n}.
\]

\qitem{2.3} What is the cardinality of the following set?
\[
  A = \{f \in \mathcal{C}^1[0, 1] : f(0) = 0,\ f(1) = 1,\ |f'(t)| \le 1 \text{ for all } t \in [0, 1]\}.
\]

\qitem{2.4} Evaluate:
\[
  \lim_{n \to \infty} \frac{1}{n} \sum_{k=1}^{\left[\frac{n}{2}\right]} \cos\left(\frac{k\pi}{n}\right)
\]
where $\left[\frac{n}{2}\right]$ denotes the largest integer not exceeding $\frac{n}{2}$.

\qitem{2.5} Which of the following improper integrals are convergent?
\begin{enumerate}[a.]
  \item $\displaystyle\int_1^{\infty} \frac{dx}{\sqrt{x^3 + 2x + 2}}$.
  \item $\displaystyle\int_0^5 \frac{dx}{x^2 - 5x + 6}$.
  \item $\displaystyle\int_0^5 \frac{dx}{\sqrt[3]{7x + 2x^4}}$.
\end{enumerate}

\qitem{2.6} Which of the following series converge uniformly?
\begin{enumerate}[a.]
  \item $\displaystyle\sum_{n=1}^{\infty} (a_n \cos nt + b_n \sin nt)$ over the interval $[-\pi, \pi]$ where $\sum_n |a_n| < \infty$ and $\sum_n |b_n| < \infty$.
  \item $\displaystyle\sum_{n=0}^{\infty} e^{-nx} \cos nx$ over the interval $]0, \infty[$.
  \item $\displaystyle x^2 + \frac{x^2}{1+x^2} + \frac{x^2}{(1+x^2)^2} + \cdots$ over the interval $[-1, 1]$.
\end{enumerate}

\qitem{2.7} Let $f \in \mathcal{C}[0, \pi]$. Determine the cases where the given condition implies that $f \equiv 0$.
\begin{enumerate}[a.]
  \item $\displaystyle\int_0^{\pi} x^n f(x)\,dx = 0$ for all integers $n \ge 0$.
  \item $\displaystyle\int_0^{\pi} f(x) \cos nx\,dx = 0$ for all integers $n \ge 0$.
  \item $\displaystyle\int_0^{\pi} f(x) \sin nx\,dx = 0$ for all integers $n \ge 1$.
\end{enumerate}

\qitem{2.8} Let $C$ be the circle defined by $|z| = 3$ in the complex plane, described in the anticlockwise direction. Evaluate:
\[
  \int_C \frac{2z^2 - z - 2}{z - 2}\,dz.
\]

\qitem{2.9} Pick out the true statements:
\begin{enumerate}[a.]
  \item Let $f$ and $g$ be analytic in the disc $|z| < 2$ and let $f = g$ on the interval $[-1, 1]$. Then $f \equiv g$.
  \item If $f$ is a non-constant polynomial with complex coefficients, then it can be factorized into (not necessarily distinct) linear factors.
  \item There exists a non-constant analytic function in the disc $|z| < 1$ which assumes only real values.
\end{enumerate}

\qitem{2.10} Let $\Omega \subset \C$ be an open and connected set and let $f : \Omega \to \C$ be an analytic function. Pick out the true statements:
\begin{enumerate}[a.]
  \item $f$ is bounded if $\Omega$ is bounded.
  \item $f$ is bounded only if $\Omega$ is bounded.
  \item $f$ is bounded if, and only if, $\Omega$ is bounded.
\end{enumerate}

\examsec{Section 3: TOPOLOGY}

\qitem{3.1} In each of the following, $f$ is assumed to be continuous. Pick out the cases when $f$ \emph{cannot} be onto.
\begin{enumerate}[a.]
  \item $f : [-1, 1] \to \R$.
  \item $f : [-1, 1] \to \Q \cap [-1, 1]$.
  \item $f : \R \to [-1, 1]$.
\end{enumerate}

\qitem{3.2} Consider the set of all $n \times n$ matrices with real entries identified with $\R^{n^2}$, endowed with its usual topology. Pick out the true statements.
\begin{enumerate}[a.]
  \item The subset of all invertible matrices is connected.
  \item The subset of all invertible matrices is dense.
  \item The subset of all orthogonal matrices is compact.
\end{enumerate}

\qitem{3.3} Pick out the functions that are uniformly continuous on the given domain.
\begin{enumerate}[a.]
  \item $f(x) = \frac{1}{x}$ on the interval $]0, 1[$.
  \item $f(x) = x^2$ on $\R$.
  \item $f(x) = \sin^2 x$ on $\R$.
\end{enumerate}

\qitem{3.4} Let $(X, d)$ be a metric space and let $A$ and $B$ be subsets of $X$. Define
\[
  d(A, B) = \inf\{d(a, b) : a \in A,\ b \in B\}.
\]
Pick out the true statements.
\begin{enumerate}[a.]
  \item If $A$ and $B$ are disjoint, then $d(A, B) > 0$.
  \item If $A$ and $B$ are closed and disjoint, then $d(A, B) > 0$.
  \item If $A$ and $B$ are compact and disjoint, then $d(A, B) > 0$.
\end{enumerate}

\qitem{3.5} Pick out the sets that are homeomorphic to the set $\{(x, y) \in \R^2 : xy = 1\}$.
\begin{enumerate}[a.]
  \item $\{(x, y) \in \R^2 : xy = 0\}$.
  \item $\{(x, y) \in \R^2 : x^2 - y^2 = 1\}$.
  \item $\{(x, y) \in \R^2 : x^2 + y^2 = 1\}$.
\end{enumerate}

\qitem{3.6} Let $(X_i, d_i)$, $i = 1, 2, 3$, be the metric spaces where $X_1 = X_2 = X_3 = \mathcal{C}[0, 1]$ and
\begin{align*}
  d_1(f, g) &= \sup_{t \in [0, 1]} |f(x) - g(x)| \\
  d_2(f, g) &= \int_0^1 |f(x) - g(x)|\,dx \\
  d_3(f, g) &= \left(\int_0^1 |f(x) - g(x)|^2\,dx\right)^{\frac{1}{2}}.
\end{align*}
Let $id$ be the identity map of $\mathcal{C}[0, 1]$ onto itself. Pick out the true statements.
\begin{enumerate}[a.]
  \item $id : X_1 \to X_2$ is continuous.
  \item $id : X_2 \to X_1$ is continuous.
  \item $id : X_3 \to X_2$ is continuous.
\end{enumerate}

\qitem{3.7} Pick out the compact sets.
\begin{enumerate}[a.]
  \item $\{(z_1, z_2) \in \C \times \C : z_1^2 + z_2^2 = 1\}$.
  \item The unit sphere in $\ell_2$, the space of all square summable real sequences, with its usual metric
  \[
    d(\{x_i\}, \{y_i\}) = \left(\sum_{i=1}^{\infty} |x_i - y_i|^2\right)^{\frac{1}{2}}.
  \]
  \item The closure of the unit ball of $\mathcal{C}^1[0, 1]$ in $\mathcal{C}[0, 1]$.
\end{enumerate}

\qitem{3.8} Let $f : S^1 \to \R$ be any continuous map, where $S^1$ is the unit circle in the plane. Let
\[
  A = \{(x, y) \in S^1 \times S^1 : x \neq y,\ f(x) = f(y)\}.
\]
Is $A$ non-empty? If the answer is `yes', is it finite, countable or uncountable?

\qitem{3.9} Let $f : S^1 \to \R$ be any continuous map, where $S^1$ is the unit circle in the plane. Let
\[
  A = \{(x, y) \in S^1 \times S^1 : x = -y,\ f(x) = f(y)\}.
\]
Is $A$ non-empty?

\qitem{3.10} Let $f \in \mathcal{C}^1[-1, 1]$ such that $|f(t)| \le 1$ and $|f'(t)| \le \frac{1}{2}$ for all $t \in [-1, 1]$. Let
\[
  A = \{t \in [-1, 1] : f(t) = t\}.
\]
Is $A$ non-empty? If the answer is `yes', what is its cardinality?

\examsec{Section 4: APPLIED MATHEMATICS}

\qitem{4.1} Let $u$ be a smooth function defined on the ball centered at the origin and of radius $a > 0$ in $\R^3$. Assume that
\[
  \frac{\partial^2 u}{\partial x^2} + \frac{\partial^2 u}{\partial y^2} + \frac{\partial^2 u}{\partial z^2} = 1
\]
throughout the ball. Compute:
\[
  \int_S \frac{\partial u}{\partial n}\,dS
\]
where $S$ is the sphere with centre at the origin and radius $a$ and $\frac{\partial u}{\partial n}$ denotes the outer normal derivative of $u$ on $S$.

\qitem{4.2} Consider a homogeneous fluid moving with velocity $u$ in space. Write down the equation which expresses the principle of conservation of mass.

\qitem{4.3} Let $C$ be the equatorial circle on the unit sphere in $\R^3$ and let $\tau$ be the unit tangent vector to $C$ taken in the anticlockwise sense. Compute:
\[
  \int_C \mathbf{F} \cdot \tau\,ds
\]
where $\mathbf{F}(x, y, z) = x\mathbf{i} - y\mathbf{j} - z\mathbf{k}$.

\qitem{4.4} Determine the value of the least possible positive number $\lambda$ such that the following problem has a non-trivial solution:
\begin{align*}
  u''(x) + \lambda u(x) &= 0, \quad 0 < x < 1 \\
  u'(0) = u'(1) &= 0.
\end{align*}

\qitem{4.5} A pendulum of mass $m$ and length $\ell$ is pulled to an angle $\alpha$ from the vertical and released from rest. Write down the differential equation satisfied by the angle $\theta(t)$ made by the pendulum with the vertical at time $t$, using the principle of conservation of energy. (If $s$ is the arc length measured from the vertical position, then the velocity $v$ is given by $v = \frac{ds}{dt}$.)

\qitem{4.6} Find d'Alembert's solution to the problem:
\begin{align*}
  \frac{\partial^2 u}{\partial t^2} &= \frac{\partial^2 u}{\partial x^2}, \quad x \in \R,\ t > 0 \\
  u(x, 0) &= x^2, \quad x \in \R \\
  \frac{\partial u}{\partial t}(x, 0) &= 0, \quad x \in \R.
\end{align*}

\qitem{4.7} Solve: Minimize $z = 2x_1 + 3x_2$, such that
\begin{align*}
  x_1 + x_2 &\le 4 \\
  3x_1 + x_2 &\ge 4 \\
  x_1 + 5x_2 &\ge 4
\end{align*}
and such that $0 \le x_1 \le 3$, and $0 \le x_2 \le 3$.

\qitem{4.8} Consider the iterative scheme $x_{n+1} = B x_n + c$ for $n \ge 0$, where $B$ is a real $N \times N$ matrix and $c \in \R^N$. The scheme is said to be convergent if the sequence $\{x_n\}$ of iterates converges for \emph{every} choice of initial vector $x_0$. Pick out the true statements.
\begin{enumerate}[a.]
  \item The scheme is convergent if, and only if, the spectral radius of $B$ is $< 1$.
  \item The scheme is convergent if, and only if, for some matrix norm $\|\cdot\|$, we have $\|B\| < 1$.
  \item The scheme is convergent if, and only if, $B$ has an eigenvalue $\lambda$ such that $0 < \lambda < 1$.
\end{enumerate}

\qitem{4.9} Write down the Laplace transform $L[f](p)$ of the function $f(x) = \sin ax$, where $a > 0$.

\qitem{4.10} What is the necessary and sufficient condition for the following problem to admit a solution?
\begin{align*}
  -\Delta u &= f \quad \text{in } \Omega \\
  \frac{\partial u}{\partial n} &= g \quad \text{on } \partial\Omega
\end{align*}
where $\Omega \subset \R^n$ is a bounded domain with boundary $\partial\Omega$, $\Delta$ is the Laplace operator, $f$ and $g$ are given smooth functions and $\frac{\partial u}{\partial n}$ denotes the outer normal derivative of $u$.

\examsec{Section 5: MISCELLANEOUS}

\qitem{5.1} Find the area of the polygon whose vertices are the $n$-th roots of unity in the complex plane, when $n \ge 3$.

\qitem{5.2} Define $p_n(t) = \cos(n \cos^{-1} t)$ for $t \in [-1, 1]$. Express $p_4(t)$ as a polynomial in $t$.

\qitem{5.3} What is the probability that a point $(x, y)$, chosen at random in the rectangle $[-1, 1] \times [0, 1]$ is such that $y > x^2$?

\qitem{5.4} An urn contains four white balls and two red balls. A ball is drawn at random and is replaced in the urn each time. What is the probability that after two successive draws, both balls drawn are white?

\qitem{5.5} Let $ABC$ be a triangle in the plane such that $BC$ is perpendicular to $AC$. Let $a, b, c$ be the lengths of $BC$, $AC$ and $AB$ respectively. Suppose that $a, b, c$ are integers and have no common divisor other than $1$. Which of the following statements are necessarily true?
\begin{enumerate}[a.]
  \item Either $a$ or $b$ is an even integer.
  \item The area of the triangle $ABC$ is an even integer.
  \item Either $a$ or $b$ is divisible by $3$.
\end{enumerate}

\qitem{5.6} What are the last two digits in the usual decimal representation of $3^{400}$?

\qitem{5.7} Find the number of integers less than $3600$ and prime to it.

\qitem{5.8} Let $n$ be a positive integer. Give an example of a sequence of $n$ consecutive composite numbers.

\qitem{5.9} For a point $P = (x, y)$ in the plane, write $f(P) = ax + by$, where $a$ and $b$ are given real numbers. Let $f(A) = f(B) = 10$. Let $C$ be a point not on the line joining $A$ and $B$ and let $C'$ be the reflection of $C$ with respect to this line. If $f(C) = 15$, find $f(C')$.

\qitem{5.10} Let $V$ be a four dimensional vector space over the field $\F_3$ of three elements. Find the number of distinct one-dimensional subspaces of $V$.

\examsec{Answer Key --- NBHM 2007}
\begin{multicols}{2}
\textbf{Section 1: Algebra}\\
\textbf{1.1} a.\\
\textbf{1.2} all.\\
\textbf{1.3} $\pm 1, \pm i$.\\
\textbf{1.4} $1, 2$.\\
\textbf{1.5} e.g.\ $x^3 + 2x + 1$ (any polynomial of degree $3$, for which $0, 1$ and $2$ are not roots $\pmod{3}$).\\
\textbf{1.6} (a) $n$; (b) $0$.\\
\textbf{1.7} a, c.\\
\textbf{1.8} b, c.\\
\textbf{1.9} a, c.\\
\textbf{1.10} a, c.

\medskip
\textbf{Section 2: Analysis}\\
\textbf{2.1} (a) conditionally convergent; (b) divergent; (c) absolutely convergent.\\
\textbf{2.2} $]0, 2]$.\\
\textbf{2.3} $1$.\\
\textbf{2.4} $1/\pi$.\\
\textbf{2.5} a, c.\\
\textbf{2.6} a.\\
\textbf{2.7} all.\\
\textbf{2.8} $8\pi i$.\\
\textbf{2.9} a, b.\\
\textbf{2.10} none.

\medskip
\textbf{Section 3: Topology}\\
\textbf{3.1} a, b.\\
\textbf{3.2} b, c.\\
\textbf{3.3} c.\\
\textbf{3.4} c.\\
\textbf{3.5} b.\\
\textbf{3.6} a, c.\\
\textbf{3.7} c.\\
\textbf{3.8} Yes; uncountable.\\
\textbf{3.9} Yes.\\
\textbf{3.10} Yes; $1$.

\medskip
\textbf{Section 4: Applied Mathematics}\\
\textbf{4.1} $\frac{4}{3}\pi a^3$.\\
\textbf{4.2} $\divg u = 0$.\\
\textbf{4.3} $0$.\\
\textbf{4.4} $\pi^2$.\\
\textbf{4.5} $\frac{1}{2}\ell^2 \left(\frac{d\theta}{dt}\right)^2 = g\ell(\cos\theta - \cos\alpha)$\\
\textbf{4.6} $u(x, t) = x^2 + t^2$.\\
\textbf{4.7} $\min z = 4$ at the point $(8/7, 4/7)$. (Either data can be accepted as full answer).\\
\textbf{4.8} a, b.\\
\textbf{4.9} $L[f](p) = a/(a^2 + p^2)$.\\
\textbf{4.10} $\int_{\Omega} f\,dx + \int_{\partial\Omega} g\,dS = 0$.

\medskip
\textbf{Section 5: Miscellaneous}\\
\textbf{5.1} $\frac{n}{2}\sin\frac{2\pi}{n}$.\\
\textbf{5.2} $8t^4 - 8t^2 + 1$.\\
\textbf{5.3} $2/3$.\\
\textbf{5.4} $4/9$.\\
\textbf{5.5} all.\\
\textbf{5.6} $01$.\\
\textbf{5.7} $960$.\\
\textbf{5.8} Example: $(n+1)! + 2, \dots, (n+1)! + (n+1)$.\\
\textbf{5.9} $5$.\\
\textbf{5.10} $40$.
\end{multicols}
% ==================== NBHM 2008 ====================
\examheader{NBHM 2008}{Research Scholarships Screening Test}{Saturday, February 2, 2008}{Time Allowed: Two Hours \quad|\quad Maximum Marks: 40}

\subsection*{Instructions to Candidates}
\begin{itemize}
  \item Please ensure that this booklet contains 9 numbered (and printed) pages. The back of each printed page is blank and can be used for rough work.
  \item There are \textbf{five} sections, containing \textbf{ten} questions each, entitled Algebra, Analysis, Topology, Applied Mathematics and Miscellaneous. Answer as many questions as possible. The assessment of the paper will be based on the best \textbf{four} sections. Each question carries one point and the maximum possible score is \textbf{forty}.
  \item Answer each question, as directed, in the space provided in the \textbf{answer booklet}, which is being supplied separately. This question paper is meant to be retained by you and so \textbf{do not} answer questions on it.
  \item In certain questions you are required to pick out the qualifying statement(s) from multiple choices. None of the statements, or more than one statement may qualify. Write \textbf{none} if none of the statements qualify, or list the labels of \emph{all} the qualifying statements (amongst (a), (b), and (c)).
  \item Points will be awarded in the above questions only if \textbf{all} the correct choices are made. \textbf{There will be no partial credit.}
  \item $\N$ denotes the set of natural numbers, $\Z$ -- the integers, $\Q$ -- the rationals, $\R$ -- the reals and $\C$ -- the field of complex numbers. $\R^n$ denotes the $n$-dimensional Euclidean space, which is assumed to be endowed with its `usual' topology. The symbol $\Z_n$ will denote the ring of integers modulo $n$. The symbol $]a, b[$ will stand for the open interval $\{x \in \R \mid a < x < b\}$ while $[a, b]$ will stand for the corresponding closed interval; $[a, b[$ and $]a, b]$ will stand for the corresponding left-closed-right-open and left-open-right-closed intervals respectively. The symbol $I$ will denote the identity matrix of appropriate order. The space of continuous real valued functions on an interval $[a, b]$ is denoted by $\mathcal{C}[a, b]$ and is endowed with its usual `sup' norm. The space of continuously differentiable real valued functions on $[a, b]$ is denoted by $\mathcal{C}^1[a, b]$ and its usual norm is the maximum of the sup-norms of the function and its derivative.
\end{itemize}

\examsec{Section 1: Algebra}

\qitem{1.1} Let $S_7$ denote the group of permutations of $7$ symbols. Find the order of the permutation:
\[
  \begin{pmatrix} 1 & 2 & 3 & 4 & 5 & 6 & 7 \\ 6 & 4 & 5 & 7 & 3 & 1 & 2 \end{pmatrix}.
\]

\qitem{1.2} Write down the number of mutually nonisomorphic abelian groups of order $19^5$.

\qitem{1.3} For two ideals $\mathcal{I}$ and $\mathcal{J}$ in a commutative ring $\mathcal{R}$, define $\mathcal{I} : \mathcal{J} = \{a \in \mathcal{R} : a\mathcal{J} \subset \mathcal{I}\}$. In the ring $\Z$ of all integers, if $\mathcal{I} = 12\Z$ and $\mathcal{J} = 8\Z$, find $\mathcal{I} : \mathcal{J}$.

\qitem{1.4} Let $\mathcal{P}$ be a prime ideal in a commutative ring $\mathcal{R}$ and let $S = \mathcal{R} \setminus \mathcal{P}$, \emph{i.e.}\ the complement of $\mathcal{P}$ in $\mathcal{R}$. Pick out the true statements:
\begin{enumerate}[(a)]
  \item $S$ is closed under addition.
  \item $S$ is closed under multiplication.
  \item $S$ is closed under addition and multiplication.
\end{enumerate}

\qitem{1.5} Let $p$ be a prime and consider the field $\Z_p$. List the primes for which the following system of linear equations DOES NOT have a solution in $\Z_p$:
\begin{align*}
  5x + 3y &= 4 \\
  3x + 6y &= 1.
\end{align*}

\qitem{1.6} Let $A$ be a $227 \times 227$ matrix with entries in $\Z_{227}$, such that all its eigenvalues are distinct. Write down its trace.

\qitem{1.7} Let $B$ be a nilpotent $n \times n$ matrix with complex entries. Set $A = B - I$. Write down the determinant of $A$.

\qitem{1.8} Let $\mathbf{x}$ and $\mathbf{y}$ be two non-zero $n \times 1$ vectors. If $\mathbf{y}^T$ denotes the transpose of $\mathbf{y}$, what are the eigenvalues of the $n \times n$ matrix $\mathbf{x}\mathbf{y}^T$?

\qitem{1.9} Let $A$ be a real symmetric $n \times n$ matrix whose only eigenvalues are $0$ and $1$. Let the dimension of the null space of $A - I$ be $m$. Pick out the true statements:
\begin{enumerate}[(a)]
  \item The characteristic polynomial of $A$ is $(\lambda - 1)^m \lambda^{m-n}$.
  \item $A^k = A^{k+1}$ for all positive integers $k$.
  \item The rank of $A$ is $m$.
\end{enumerate}

\qitem{1.10} What is the dimension of the space of all $n \times n$ matrices with real entries which are such that the sum of the entries in the first row and the sum of the diagonal entries are both zero?

\examsec{Section 2: Analysis}

\qitem{2.1} Let $f(x) = \frac{1}{1+x^2}$. Consider its Taylor expansion about a point $a \in \R$, given by $f(x) = \sum_{n=0}^{\infty} a_n(x - a)^n$. What is the radius of convergence of this series?

\qitem{2.2} Consider the series
\[
  \sum_{n=1}^{\infty} (-1)^n \frac{x^2 + n}{n^2}.
\]
Pick out the true statements:
\begin{enumerate}[(a)]
  \item The series converges for all real values of $x$.
  \item The series converges uniformly on $\R$.
  \item The series does not converge absolutely for any real value of $x$.
\end{enumerate}

\qitem{2.3} Let $f : \R \to \R$ be a continuous function. Which of the following imply that it is uniformly continuous?
\begin{enumerate}[(a)]
  \item $f$ is $2\pi$-periodic.
  \item $f$ is differentiable and its derivative is bounded on $\R$.
  \item $f$ is absolutely continuous.
\end{enumerate}

\qitem{2.4} Let $f : [-1, 1] \to \R$ be continuous. Assume that $\int_{-1}^{1} f(t)\,dt = 1$. Evaluate:
\[
  \lim_{n \to \infty} \int_{-1}^{1} f(t)\cos^2 nt\,dt.
\]

\qitem{2.5} Let $f$ be a continuously differentiable $2\pi$-periodic real valued function on the real line. Let $a_n = \int_{-\pi}^{\pi} f(t)\cos nt\,dt$ where $n$ is a non-negative integer. Pick out the true statements:
\begin{enumerate}[(a)]
  \item The derivative of $f$ is also a $2\pi$-periodic function.
  \item $|a_n| \le C\frac{1}{n}$ for all $n$, where $C > 0$ is a constant independent of $n$.
  \item $a_n \to 0$, as $n \to \infty$.
\end{enumerate}

\qitem{2.6} Let $f_n$ and $f$ be continuous functions on an interval $[a, b]$ and assume that $f_n \to f$ uniformly on $[a, b]$. Pick out the true statements:
\begin{enumerate}[(a)]
  \item If $f_n$ are all Riemann integrable, then $f$ is Riemann integrable.
  \item If $f_n$ are all continuously differentiable, then $f$ is continuously differentiable.
  \item If $x_n \to x$ in $[a, b]$, then $f_n(x_n) \to f(x)$.
\end{enumerate}

\qitem{2.7} Let $f : \R \to \R$ be a continuously differentiable function such that $f'(0) = 0$. Define, for $x$ and $y \in \R$,
\[
  g(x, y) = f\left(\sqrt{x^2 + y^2}\right).
\]
Pick out the true statements:
\begin{enumerate}[(a)]
  \item $g$ is a differentiable function on $\R^2$.
  \item $g$ is a differentiable function on $\R^2$ if, and only if, $f(0) = 0$.
  \item $g$ is differentiable only on $\R^2 \setminus \{(0, 0)\}$.
\end{enumerate}

\qitem{2.8} Find the square roots of the complex number $1 + 2i$.

\qitem{2.9} Evaluate:
\[
  \int_{|z|=1} \frac{4 + z}{(2 - z)z}\,dz
\]
the circle $\{|z| = 1\}$ being described in the anticlockwise direction.

\qitem{2.10} Pick out the true statements:
\begin{enumerate}[(a)]
  \item There exists an analytic function $f$ on $\C$ such that $f(2i) = 0$, $f(0) = 2i$ and $|f(z)| \le 2$ for all $z \in \C$.
  \item There exists an analytic function $f$ in the open unit disc $\{z \in \C : |z| < 1\}$ such that $f\left(\frac{1}{2}\right) = 1$ and $f\left(\frac{1}{2^n}\right) = 0$ for all integers $n \ge 2$.
  \item There exists an analytic function $f$ on $\C$ whose real part is given by $u(x, y) = x^2 + y^2$, where $z = x + iy$.
\end{enumerate}

\examsec{Section 3: Topology}

\qitem{3.1} Which of the following define a metric on $\R$?
\begin{enumerate}[(a)]
  \item $d_1(x, y) = \dfrac{\big|\,|x| - |y|\,\big|}{1 + |x|\cdot|y|}$.
  \item $d_2(x, y) = \sqrt{|x - y|}$.
  \item $d_3(x, y) = |f(x) - f(y)|$, where $f : \R \to \R$ is a strictly monotonically increasing function.
\end{enumerate}

\qitem{3.2} Let $X$ and $Y$ be topological spaces and let $f : X \to Y$ be a continuous bijection. Under which of the following conditions will $f$ be a homeomorphism?
\begin{enumerate}[(a)]
  \item $X$ and $Y$ are complete metric spaces.
  \item $X$ and $Y$ are Banach spaces and $f$ is linear.
  \item $X$ is a compact topological space and $Y$ is Hausdorff.
\end{enumerate}

\qitem{3.3} Let $(X, d)$ be a compact metric space and let $\{f_\alpha : \alpha \in A\}$ be a uniformly bounded and equicontinuous family of functions on $X$. Define
\[
  f(x) = \sup_{\alpha \in A} f_\alpha(x).
\]
Pick out the true statements:
\begin{enumerate}[(a)]
  \item For any $t \in \R$, the set $\{x \in X : f(x) < t\}$ is an open set in $X$.
  \item The function $f$ is continuous.
  \item Every sequence $\{f_{\alpha_n}\}$ contained in the above family admits a uniformly convergent subsequence.
\end{enumerate}

\qitem{3.4} Let $D = \{x \in \R^2 : |x| \le 1\}$ where $|x|$ is the usual euclidean norm of the vector $x$. Let $f : D \to X$ be a continuous function into a topological space $X$. Pick out the cases below when $f$ will NEVER be onto.
\begin{enumerate}[(a)]
  \item $X = [-1, 1]$.
  \item $X = [-1, 1] \setminus \{0\}$.
  \item $X = ]-1, 1[$.
\end{enumerate}

\qitem{3.5} Let $\M_n(\R)$ denote the set of all $n \times n$ matrices with real entries, considered as the space $\R^{n^2}$. Which of the following subsets are compact?
\begin{enumerate}[(a)]
  \item The set of all invertible matrices.
  \item The set of all orthogonal matrices.
  \item The set of all matrices whose trace is zero.
\end{enumerate}

\qitem{3.6} With the notations as in the preceding question, which of the following sets are connected?
\begin{enumerate}[(a)]
  \item The set of all invertible matrices.
  \item The set of all orthogonal matrices.
  \item The set of all matrices whose trace is zero.
\end{enumerate}

\qitem{3.7} Let $f : \R \to \R$ be a continuous function. Define
\[
  G = \{(x, f(x)) : x \in \R\} \subset \R^2.
\]
Pick out the true statements:
\begin{enumerate}[(a)]
  \item $G$ is closed in $\R^2$.
  \item $G$ is open in $\R^2$.
  \item $G$ is connected in $\R^2$.
\end{enumerate}

\qitem{3.8} Let $X$ be a topological space and let $A \subset X$. Let $\partial A$ denote the boundary of $A$, \emph{i.e.}\ the set of points in the closure of $A$ which are not in the interior of $A$. A closed set is \emph{nowhere dense} if its interior is the empty set. Pick out the true statements:
\begin{enumerate}[(a)]
  \item If $A$ is open, then $\partial A$ is nowhere dense.
  \item If $A$ is closed, then $\partial A$ is nowhere dense.
  \item If $A$ is any subset, then $\partial A$ is always nowhere dense.
\end{enumerate}

\qitem{3.9} Let $V$ be a complete normed linear space and let $B$ be a basis for $V$ as a vector space. Pick out the true statements:
\begin{enumerate}[(a)]
  \item $B$ can be a finite set.
  \item $B$ can be a countably infinite set.
  \item If $B$ is infinite, then it must be an uncountable set.
\end{enumerate}

\qitem{3.10} Let $V_1 = \mathcal{C}[0, 1]$ with the metric $d_1(f, g) = \max_{t \in [0, 1]} |f(t) - g(t)|$. Let $V_2 = \mathcal{C}[0, 1]$ with the metric $d_2(f, g) = \int_0^1 |f(t) - g(t)|\,dt$. Let $id$ denote the identity map on $\mathcal{C}[0, 1]$. Pick out the true statements:
\begin{enumerate}[(a)]
  \item $id : V_1 \to V_2$ is continuous.
  \item $id : V_2 \to V_1$ is continuous.
  \item $id : V_1 \to V_2$ is a homeomorphism.
\end{enumerate}

\examsec{Section 4: Applied Mathematics}

\qitem{4.1} Write down the Laplace transform of the function $f(x) = \cos 2x$.

\qitem{4.2} Let $B$ denote the unit ball in $\R^N$, $N \ge 2$. Let $\alpha_N$ be its ($N$-dimensional) volume and let $\beta_N$ be its ($(N-1)$-dimensional) surface measure. Apply Gauss' divergence theorem to the vector field $\mathbf{v}(\mathbf{x}) = \mathbf{x}$ and derive the relation connecting $\alpha_N$ and $\beta_N$.

\qitem{4.3} Assume that the rate at which a body cools is proportional to the difference in temperature between the body and its surroundings. A body is heated to $110^\circ\mathrm{C}$ and is placed in air at $10^\circ\mathrm{C}$. After one hour, its temperature is $60^\circ\mathrm{C}$. At what time will its temperature reach $30^\circ\mathrm{C}$?

\qitem{4.4} A body of mass $m$ falls under gravity and is retarded by a force proportional to its velocity. Write down the differential equation satisfied by the velocity $v(t)$ at time $t$.

\qitem{4.5} Solve the above equation given that the velocity is zero at time $t = 0$.

\qitem{4.6} Write down the critical points of the nonlinear system of differential equations:
\begin{align*}
  \frac{dx}{dt} &= y(x^2 + 1) \\
  \frac{dy}{dt} &= 2xy^2.
\end{align*}

\qitem{4.7} Classify the following differential operator as elliptic, hyperbolic or parabolic:
\[
  \mathcal{L}(u) = 2\frac{\partial^2 u}{\partial x^2} - 2\frac{\partial^2 u}{\partial x \partial y} + 2\frac{\partial^2 u}{\partial y^2}.
\]

\qitem{4.8} Let $B$ be the unit ball in the plane and let $u$ be a solution of the boundary value problem:
\begin{align*}
  \Delta u &= C \quad \text{in } B \\
  \frac{\partial u}{\partial n} &= 1 \quad \text{on } \partial B
\end{align*}
where $\Delta$ denotes the Laplace operator, $\partial B$ denotes the boundary of $B$ and $\frac{\partial u}{\partial n}$ denotes the outer normal derivative on the boundary. Evaluate $C$, given that it is a constant.

\qitem{4.9} Write down the dual of the linear programming problem:
\[
  \text{Max.}: 2x + 3y
\]
such that
\begin{align*}
  x + 2y &= 3 \\
  2x + y &\ge 4 \\
  x + y &\le 5 \\
  x \ge 0,&\quad y \ge 0.
\end{align*}

\qitem{4.10} Write down the Newton-Raphson iteration scheme to find the square root of $a > 0$, by solving the equation $x^2 = a$.

\examsec{Section 5: Miscellaneous}

\qitem{5.1} Differentiate with respect to $t$:
\[
  f(t) = \int_{-t^2}^{a} e^{-x^2}\,dx,\quad (a > 0).
\]

\qitem{5.2} Sum the series:
\[
  1 + \frac{1 + 3}{2!} + \frac{1 + 3 + 3^2}{3!} + \cdots
\]

\qitem{5.3} Sum the series:
\[
  1 + \frac{1}{3}\cdot\frac{1}{4} + \frac{1}{5}\cdot\frac{1}{4^2} + \frac{1}{7}\cdot\frac{1}{4^3} + \cdots
\]

\qitem{5.4} Let $A$ be the point $(0, 1)$ and $B$ the point $(2, 2)$ in the plane. Consider all paths made up of the two line segments $AC$ and $CB$ as the point $C$ varies on the $x$-axis. Find the coordinates of $C$ for which the corresponding path has the shortest length.

\qitem{5.5} Find the area of the pentagon formed in the plane with the fifth roots of unity as its vertices.

\qitem{5.6} Let $A, B, C, D$ and $E$ be five points marked, in clockwise order, on the unit circle in the plane (with centre at origin). Let $\alpha$ and $\beta$ be real numbers and set $f(P) = \alpha x + \beta y$ where $P$ is a point whose coordinates are $(x, y)$. Assume that $f(A) = 10$, $f(B) = 5$, $f(C) = 4$ and $f(D) = 10$. Which of the following are impossible?\\
(a) $f(E) = 2$ \qquad (b) $f(E) = 4$ \qquad (c) $f(E) = 5$

\qitem{5.7} Let $r$ identical red balls and $b$ identical black balls be arranged in a row. Write down the number of arrangements for which the last ball is black.

\qitem{5.8} It is known that a family has two children.
\begin{enumerate}[(a)]
  \item If it is known that one of the children is a girl, what is the probability that the other child is also a girl?
  \item If it is known that the elder child is a girl, what is the probability that the younger child is also a girl?
\end{enumerate}

\qitem{5.9} A real number is called \emph{algebraic} if it occurs as the root of a polynomial with integer coefficients. Otherwise it is said to be a \emph{transcendental} number. Consider the interval $[0, 1]$ considered as a probability space when it is provided with the Lebesgue measure. What is the probability that a number chosen at random in $[0, 1]$ is transcendental?

\qitem{5.10} List all primes $p \le 13$ such that $p$ divides $n^{13} - n$ for every integer $n$.

\examsec{Answer Key --- NBHM 2008}
\begin{multicols}{2}
\textbf{Section 1: Algebra}\\
\textbf{1.1} $6$\\
\textbf{1.2} $7$\\
\textbf{1.3} $3\Z$\\
\textbf{1.4} b\\
\textbf{1.5} $3$\\
\textbf{1.6} $0$\\
\textbf{1.7} $(-1)^n$\\
\textbf{1.8} $0, \mathbf{x}^T\mathbf{y}$\\
\textbf{1.9} a, b, c\\
\textbf{1.10} $n^2 - 2$

\medskip
\textbf{Section 2: Analysis}\\
\textbf{2.1} $\sqrt{a^2 + 1}$\\
\textbf{2.2} a, c\\
\textbf{2.3} a, b, c\\
\textbf{2.4} $\frac{1}{2}$\\
\textbf{2.5} a, b, c\\
\textbf{2.6} a, c\\
\textbf{2.7} a\\
\textbf{2.8} $\pm\left(\sqrt{\frac{\sqrt{5}+1}{2}} + i\sqrt{\frac{\sqrt{5}-1}{2}}\right)$\\
\textbf{2.9} $4\pi i$\\
\textbf{2.10} None

\medskip
\textbf{Section 3: Topology}\\
\textbf{3.1} b, c\\
\textbf{3.2} b, c\\
\textbf{3.3} a, b, c\\
\textbf{3.4} b, c\\
\textbf{3.5} b\\
\textbf{3.6} c\\
\textbf{3.7} a, c\\
\textbf{3.8} a, b\\
\textbf{3.9} a, c\\
\textbf{3.10} a

\medskip
\textbf{Section 4: Applied Mathematics}\\
\textbf{4.1} $\frac{s}{s^2+4}$\\
\textbf{4.2} $\beta_N = N\alpha_N$\\
\textbf{4.3} $\frac{\log 5}{\log 2}$ hours\\
\textbf{4.4} $\frac{dv}{dt} = g - cv^2$ \emph{(as printed in key)}\\
\textbf{4.5} $v(t) = \sqrt{\frac{g}{c}}\frac{1 - e^{-2\sqrt{gc}\,t}}{1 + e^{-2\sqrt{gc}\,t}}$\\
\textbf{4.6} $\{(x, 0) : x \in \R\}$\\
\textbf{4.7} Elliptic\\
\textbf{4.8} $2$\\
\textbf{4.9} Min.: $3w_1 - 3w_2 - 4w_3 + 5w_4$ such that $w_1 - w_2 - 2w_3 + w_4 \ge 2$, $2w_1 - 2w_2 - w_3 + w_4 \ge 3$, $w_i \ge 0$, $1 \le i \le 4$.\\
\textbf{4.10} $x_{n+1} = \frac{x_n^2 + a}{2x_n}$

\medskip
\textbf{Section 5: Miscellaneous}\\
\textbf{5.1} $2te^{-t^4}$\\
\textbf{5.2} $\frac{e^3 - e}{2}$\\
\textbf{5.3} $\log 3$\\
\textbf{5.4} $\left(\frac{2}{3}, 0\right)$\\
\textbf{5.5} $\frac{5}{2}\sin\frac{2\pi}{5}$\\
\textbf{5.6} a, b, c\\
\textbf{5.7} $\binom{r+b-1}{r}$\\
\textbf{5.8} (a) $\frac{1}{3}$; (b) $\frac{1}{2}$\\
\textbf{5.9} $1$\\
\textbf{5.10} $2, 3, 5, 7, 13$
\end{multicols}
% ==================== NBHM 2009 ====================
\examheader{NBHM 2009}{Research Scholarships Screening Test}{Saturday, January 24, 2009}{Time Allowed: 150 Minutes \quad|\quad Maximum Marks: 40}

\subsection*{Instructions to Candidates}
\begin{itemize}
  \item Please ensure that this booklet contains 12 numbered (and printed) pages. The back of each printed page is blank and can be used for rough work.
  \item There are \textbf{five} sections, containing \textbf{ten} questions each, entitled Algebra, Analysis, Topology, Applied Mathematics and Miscellaneous. Answer as many questions as possible. The assessment of the paper will be based on the best \textbf{four} sections. Each question carries one point and the maximum possible score is \textbf{forty}.
  \item Answer each question, as directed, in the space provided in the \textbf{answer booklet}, which is being supplied separately. This question paper is meant to be retained by you and so \textbf{do not} answer questions on it.
  \item In certain questions you are required to pick out the qualifying statement(s) from multiple choices. None of the statements, or one or more than one statement may qualify. Write \textbf{none} if none of the statements qualify, or list the labels of \emph{all} the qualifying statements (amongst (a), (b), and (c)).
  \item Points will be awarded in the above questions only if \textbf{all} the correct choices are made. \textbf{There will be no partial credit.}
  \item $\N$ denotes the set of natural numbers, $\Z$ -- the integers, $\Q$ -- the rationals, $\R$ -- the reals and $\C$ -- the field of complex numbers. $\R^n$ denotes the $n$-dimensional Euclidean space, which is assumed to be endowed with its `usual' topology. The symbol $\Z_n$ will denote the ring of integers modulo $n$. The symbol $]a, b[$ will stand for the open interval $\{x \in \R \mid a < x < b\}$ while $[a, b]$ will stand for the corresponding closed interval; $[a, b[$ and $]a, b]$ will stand for the corresponding left-closed-right-open and left-open-right-closed intervals respectively. The symbol $I$ will denote the identity matrix of appropriate order. The space of continuous real valued functions on an interval $[a, b]$ is denoted by $\mathcal{C}[a, b]$ and is endowed with its usual `sup' norm.
\end{itemize}

\examsec{Section 1: Algebra}

\qitem{1.1} Pick out the cases where the given subgroup $H$ is a normal subgroup of the group $G$.
\begin{enumerate}[(a)]
  \item $G$ is the group of all $2 \times 2$ invertible upper triangular matrices with real entries, under matrix multiplication, and $H$ is the subgroup of all such matrices $(a_{ij})$ such that $a_{11} = 1$.
  \item $G$ is the group of all $2 \times 2$ invertible upper triangular matrices with real entries, under matrix multiplication, and $H$ is the subgroup of all such matrices $(a_{ij})$ such that $a_{11} = a_{22}$.
  \item $G$ is the group of all $n \times n$ invertible matrices with real entries, under matrix multiplication, and $H$ is the subgroup of such matrices with positive determinant.
\end{enumerate}

\qitem{1.2} Let $GL(n, \R)$ denote the group of all invertible $n \times n$ matrices with real entries, under matrix multiplication, and let $SL(n, \R)$ denote the subgroup of such matrices whose determinant is equal to unity. Identify the quotient group $GL(n, \R)/SL(n, \R)$.

\qitem{1.3} Let $S_n$ denote the symmetric group of permutations of $n$ symbols. Does $S_7$ contain an element of order $10$? If `yes', write down an example of such an element.

\qitem{1.4} What is the largest possible order of an element in $S_7$?

\qitem{1.5} Write down all the units in the ring $\Z_8$ of all integers modulo $8$.

\qitem{1.6} Pick out the cases where the given ideal is a maximal ideal.
\begin{enumerate}[(a)]
  \item The ideal $15\Z$ in $\Z$.
  \item The ideal $\mathcal{I} = \{f : f(0) = 0\}$ in the ring $\mathcal{C}[0, 1]$ of all continuous real valued functions on the interval $[0, 1]$.
  \item The ideal generated by $x^3 + x + 1$ in the ring of polynomials $\F_3[x]$, where $\F_3$ is the field of three elements.
\end{enumerate}

\qitem{1.7} Let $A$ be a $2 \times 2$ matrix with complex entries which is non-zero and non-diagonal. Pick out the cases when $A$ is diagonalizable.
\begin{enumerate}[(a)]
  \item $A^2 = I$.
  \item $A^2 = 0$.
  \item All eigenvalues of $A$ are equal to $2$.
\end{enumerate}

\qitem{1.8} Let $\mathbf{x}$ and $\mathbf{y} \in \R^n$ be two non-zero (column) vectors. Let $\mathbf{y}^T$ denote the transpose of $\mathbf{y}$. Let $A = \mathbf{x}\mathbf{y}^T$, \emph{i.e.}\ $A = (a_{ij})$ where $a_{ij} = x_i y_j$. What is the rank of $A$?

\qitem{1.9} Let $\mathbf{x}$ be a non-zero (column) vector in $\R^n$. What is the necessary and sufficient condition for the matrix $A = I - 2\mathbf{x}\mathbf{x}^T$ to be orthogonal?

\qitem{1.10} Let $A$ be an $n \times n$ matrix with complex entries. Pick out the true statements.
\begin{enumerate}[(a)]
  \item $A$ is always similar to an upper-triangular matrix.
  \item $A$ is always similar to a diagonal matrix.
  \item $A$ is similar to a block diagonal matrix, with each diagonal block of size strictly less than $n$, provided $A$ has at least $2$ distinct eigenvalues.
\end{enumerate}

\examsec{Section 2: Analysis}

\qitem{2.1} Evaluate:
\[
  \lim_{n \to \infty} n\sin(2\pi e n!).
\]

\qitem{2.2} Evaluate:
\[
  \lim_{n \to \infty} \sum_{k=1}^{n} \frac{1}{\sqrt{4n^2 - k^2}}.
\]

\qitem{2.3} Pick out the convergent series:
\begin{enumerate}[(a)]
  \item $\displaystyle\sum_{n=1}^{\infty} \frac{(n+1)^n}{n^{n + \frac{3}{2}}}$.
  \item $\displaystyle 1 + \frac{1}{2^2} + \frac{2^2}{3^3} + \frac{3^3}{4^4} + \cdots$
  \item $\displaystyle\sum_{n=1}^{\infty} \sqrt{\frac{1 + 4^n}{1 + 5^n}}$.
\end{enumerate}

\qitem{2.4} Which of the following functions are continuous?
\begin{enumerate}[(a)]
  \item $\displaystyle f(x) = x^2 + \frac{x^2}{1+x^2} + \frac{x^2}{(1+x^2)^2} + \cdots,\quad x \in \R$.
  \item $\displaystyle f(x) = \sum_{n=1}^{\infty} (-1)^n \frac{\cos nx}{n^{\frac{3}{2}}},\quad x \in [-\pi, \pi]$.
  \item $\displaystyle f(x) = \sum_{n=1}^{\infty} n^2 x^n,\quad x \in \left[-\frac{1}{2}, \frac{1}{2}\right]$.
\end{enumerate}

\qitem{2.5} Which of the following functions are uniformly continuous?
\begin{enumerate}[(a)]
  \item $f(x) = \frac{1}{x}$ in $(0, 1)$.
  \item $f(x) = x^2$ in $\R$.
  \item $f(x) = \sin^2 x$ in $\R$.
\end{enumerate}

\qitem{2.6} Pick out the sequences $\{f_n\}$ which are uniformly convergent.
\begin{enumerate}[(a)]
  \item $f_n(x) = nxe^{-nx}$ on $(0, \infty)$.
  \item $f_n(x) = x^n$ on $[0, 1]$.
  \item $f_n(x) = \dfrac{\sin nx}{\sqrt{n}}$ on $\R$.
\end{enumerate}

\qitem{2.7} Which of the following functions are Riemann integrable on the interval $[0, 1]$?
\begin{enumerate}[(a)]
  \item $f(x) = \lim_{n \to \infty} \cos^{2n}(24\pi x)$.
  \item $f(x) = \begin{cases} 0, & \text{if } x \text{ is rational} \\ 1, & \text{if } x \text{ is irrational.} \end{cases}$
  \item $f(x) = \begin{cases} \cos x, & \text{if } 0 \le x \le \frac{1}{2} \\ \sin x, & \text{if } \frac{1}{2} < x \le 1. \end{cases}$
\end{enumerate}

\qitem{2.8} Let $z = x + iy$ be a complex number, where $x$ and $y \in \R$, and let $f(z) = u(x, y) + iv(x, y)$, where $u$ and $v$ are real valued, be an analytic function on $\C$. Express $f'(z)$ in terms of the partial derivatives of $u$ and $v$.

\qitem{2.9} Let $z \in \C$ be as in the previous question. Find the image of the set $S = \{z : x > 0,\ 0 < y < 2\}$ under the transformation $f(z) = iz + 1$.

\qitem{2.10} Find the residue at $z = 0$ for the function
\[
  f(z) = \frac{1 + 2z}{z^2 + z^3}.
\]

\examsec{Section 3: Topology}

\qitem{3.1} Let $X$ be a metric space and let $f : X \to \R$ be a continuous function. Pick out the true statements.
\begin{enumerate}[(a)]
  \item $f$ always maps Cauchy sequences into Cauchy sequences.
  \item If $X$ is compact, then $f$ always maps Cauchy sequences into Cauchy sequences.
  \item If $X = \R^n$, then $f$ always maps Cauchy sequences into Cauchy sequences.
\end{enumerate}

\qitem{3.2} Let $B$ be the closed ball in $\R^2$ with centre at the origin and radius unity. Pick out the true statements.
\begin{enumerate}[(a)]
  \item There exists a continuous function $f : B \to \R$ which is one-one.
  \item There exists a continuous function $f : B \to \R$ which is onto.
  \item There exists a continuous function $f : B \to \R$ which is one-one and onto.
\end{enumerate}

\qitem{3.3} Let $A$ and $B$ be subsets of $\R$. Define $C = \{a + b : a \in A,\ b \in B\}$. Pick out the true statements.
\begin{enumerate}[(a)]
  \item $C$ is closed if $A$ and $B$ are closed.
  \item $C$ is closed if $A$ is closed and $B$ is compact.
  \item $C$ is compact if $A$ is closed and $B$ is compact.
\end{enumerate}

\qitem{3.4} Which of the following subsets of $\R^2$ are compact?
\begin{enumerate}[(a)]
  \item $\{(x, y) : xy = 1\}$
  \item $\{(x, y) : x^{\frac{2}{3}} + y^{\frac{2}{3}} = 1\}$
  \item $\{(x, y) : x^2 + y^2 < 1\}$
\end{enumerate}

\qitem{3.5} Which of the following sets in $\R^2$ are connected?
\begin{enumerate}[(a)]
  \item $\{(x, y) : x^2 y^2 = 1\}$
  \item $\{(x, y) : x^2 + y^2 = 1\}$
  \item $\{(x, y) : 1 < x^2 + y^2 < 2\}$
\end{enumerate}

\qitem{3.6} Let $\mathcal{P}$ denote the set of all polynomials in the real variable $x$ which varies over the interval $[0, 1]$. What is the closure of $\mathcal{P}$ in $\mathcal{C}[0, 1]$ (with its usual sup-norm topology)?

\qitem{3.7} Let $\{f_n\}$ be a sequence of functions which are continuous over $[0, 1]$ and continuously differentiable in $]0, 1[$. Assume that $|f_n(x)| \le 1$ and that $|f_n'(x)| \le 1$ for all $x \in ]0, 1[$ and for each positive integer $n$. Pick out the true statements.
\begin{enumerate}[(a)]
  \item $f_n$ is uniformly continuous for each $n$.
  \item $\{f_n\}$ is a convergent sequence in $\mathcal{C}[0, 1]$.
  \item $\{f_n\}$ contains a subsequence which converges in $\mathcal{C}[0, 1]$.
\end{enumerate}

\qitem{3.8} Pick out the true statements.
\begin{enumerate}[(a)]
  \item Let $f : [0, 2] \to [0, 1]$ be a continuous function. Then, there always exists $x \in [0, 1]$ such that $f(x) = x$.
  \item Let $f : [0, 1] \to [0, 1]$ be a continuous function which is continuously differentiable in $]0, 1[$ and such that $|f'(x)| \le \frac{1}{2}$ for all $x \in ]0, 1[$. Then, there exists a unique $x \in [0, 1]$ such that $f(x) = x$.
  \item Let $S = \{\mathbf{p} = (x, y) \in \R^2 : x^2 + y^2 = 1\}$. Let $f : S \to S$ be a continuous function. Then, there always exists $\mathbf{p} \in S$ such that $f(\mathbf{p}) = \mathbf{p}$.
\end{enumerate}

\qitem{3.9} Let $(X, d)$ be a metric space. Let $A$ and $B$ be subsets of $X$. Define $d(A, B) = \inf\{d(a, b) : a \in A,\ b \in B\}$. For $x \in X$, define $d(x, A) = \inf\{d(x, a) : a \in A\}$. Pick out the true statements.
\begin{enumerate}[(a)]
  \item The function $x \mapsto d(x, A)$ is uniformly continuous.
  \item $d(x, A) = 0$ if, and only if, $x \in A$.
  \item $d(A, B) = 0$ implies that $A \cap B \neq \emptyset$.
\end{enumerate}

\qitem{3.10} Let $B = \{(x, y) \in \R^2 : x^2 + y^2 \le 1\}$ and $D = \{(x, y) \in \R^2 : x^2 + y^2 < 1\}$. Pick out the true statements.
\begin{enumerate}[(a)]
  \item Given a continuous function $g : B \to \R$, there always exists a continuous function $f : \R^2 \to \R$ such that $f = g$ on $B$.
  \item Given a continuous function $g : D \to \R$, there always exists a continuous function $f : \R^2 \to \R$ such that $f = g$ on $D$.
  \item There exists a continuous function $f : \R^2 \to \R$ such that $f \equiv 1$ on the set $\{(x, y) \in \R^2 : x^2 + y^2 = \frac{3}{2}\}$ and $f \equiv 0$ on the set $B \cup \{(x, y) \in \R^2 : x^2 + y^2 \ge 2\}$.
\end{enumerate}

\examsec{Section 4: Applied Mathematics}

\qitem{4.1} A spherical ball of volatile material evaporates (\emph{i.e.}\ its volume decreases) at a rate proportional to its surface area. If the initial radius is $r_0$ and at time $t = 1$, the radius is $r_0/2$, find the time at which the ball disappears completely.

\qitem{4.2} A body of mass $m$ falling from rest under gravity experiences air resistance proportional to the square of its velocity. Write down the initial value problem for the vertical displacement $x$ of the body.

\qitem{4.3} A body falling from rest under gravity travels a distance $y$ and has a velocity $v$ at time $t$. Write down the relationship between $v$ and $y$.

\qitem{4.4} Let $\mathbf{x} = (x_1, \dots, x_n) \in \R^n$ and let $\mathbf{v}(\mathbf{x}) = \mathbf{x}$. Apply Gauss' divergence theorem to $\mathbf{v}$ over the unit ball
\[
  B = \left\{\mathbf{x} \in \R^n : \sum_{i=1}^{n} x_i^2 \le 1\right\}
\]
and deduce the relationship between $\omega_n$, the ($n$-dimensional) volume of $B$, and $\sigma_n$, the ($(n-1)$-dimensional) surface measure of $B$.

\qitem{4.5} Write down the general solution of the linear system:
\begin{align*}
  \frac{dx}{dt} &= x + y \\
  \frac{dy}{dt} &= 4x - 2y.
\end{align*}

\qitem{4.6} What is the smallest positive value of $\lambda$ such that the problem:
\begin{align*}
  u'' + \lambda u &= 0 \quad \text{in } ]0, 1[ \\
  u(0) = u(1) &\quad \text{and} \quad u'(0) = u'(1)
\end{align*}
has a solution such that $u \not\equiv 0$?

\qitem{4.7} Let $f : \R \to \R$ be a continuous function. Write down the solution of the problem:
\begin{align*}
  \frac{\partial^2 u}{\partial t^2} - \frac{\partial^2 u}{\partial x^2} &= 0, \quad t > 0, \\
  u(x, 0) &= f(x), \quad x \in \R, \\
  \frac{\partial u}{\partial t}(x, 0) &= 0, \quad x \in \R.
\end{align*}

\qitem{4.8} Use duality to find the optimal value of the cost function in the following linear programming problem:
\[
  \max\ x + y + z \quad \text{such that } 3x + 2y + 2z = 1,\quad x \ge 0,\ y \ge 0,\ z \ge 0.
\]

\qitem{4.9} The value of $\sqrt{10}$ is computed by solving the equation $x^2 = 10$ using the Newton-Raphson method. Starting from some value $x_0 > 0$, write down the iteration scheme.

\qitem{4.10} Write down the Laplace transform $F(s)$ of the function $f(x) = x^3$.

\examsec{Section 5: Miscellaneous}

\qitem{5.1} Let $P = (0, 1)$ and $Q = (4, 1)$ be points on the plane. Let $A$ be a point which moves on the $x$-axis between the points $(0, 0)$ and $(4, 0)$. Let $B$ be a point which moves on the line $y = 2$ between the points $(0, 2)$ and $(4, 2)$. Consider all possible paths consisting of the line segments $PA$, $AB$ and $BQ$. What is the shortest possible length of such a path?

\qitem{5.2} A convex polygon has its interior angles in arithmetic progression, the least angle being $120^\circ$ and the common difference being $5^\circ$. Find the number of sides of the polygon.

\qitem{5.3} Let $a, b$ and $c$ be the lengths of the sides of an arbitrary triangle. Define
\[
  x = \frac{ab + bc + ca}{a^2 + b^2 + c^2}.
\]
Pick out the true statements.\\
(a) $\frac{1}{2} \le x \le 2$. \qquad (b) $\frac{1}{2} \le x \le 1$. \qquad (c) $\frac{1}{2} < x \le 1$.

\qitem{5.4} What is the maximum number of pieces that can be obtained from a pizza by making $7$ cuts with a knife?

\qitem{5.5} In arithmetic base $3$, a number is expressed as $210100$. Find its square root and express it in base $3$.

\qitem{5.6} Evaluate:
\[
  \left(\frac{-1 + i\sqrt{3}}{\sqrt{2} + i\sqrt{2}}\right)^{20}.
\]

\qitem{5.7} Let $n$ be a fixed positive integer and let $C_r$ stand for the usual binomial coefficients \emph{i.e.}, the number of ways of choosing $r$ objects from $n$ objects. Evaluate:
\[
  C_1 + 2C_2 + \dots + nC_n.
\]

\qitem{5.8} Let $x, y$ and $z$ be real numbers such that $x^2 + y^2 + z^2 = 1$. Find the maximum and minimum values of $2x + 3y + z$.

\qitem{5.9} Let $x_0 = 0$. For $n \ge 0$, define $x_{n+1} = x_n^2 + \frac{1}{4}$. Pick out the true statements:
\begin{enumerate}[(a)]
  \item The sequence $\{x_n\}$ is bounded.
  \item The sequence $\{x_n\}$ is monotonic.
  \item The sequence $\{x_n\}$ is convergent.
\end{enumerate}

\qitem{5.10} Seven tickets are numbered consecutively from $1$ to $7$. Two of them are selected in order without replacement. Let $A$ denote the event that the numbers on the two tickets add up to $9$. Let $B$ be the event that the numbers on the two tickets differ by $3$. If each draw has equal probability $\frac{1}{42}$ (the draw $(1, 7)$ being considered as distinct from the draw $(7, 1)$, for example) find the probability $P(B \mid A)$.

\examsec{Answer Key --- NBHM 2009}
\begin{multicols}{2}
\textbf{Section 1: Algebra}\\
\textbf{1.1} a, b, c\\
\textbf{1.2} Multiplicative group of non-zero reals.\\
\textbf{1.3} Yes. Eg: $(12345)(67)$\\
\textbf{1.4} $12$\\
\textbf{1.5} $1, 3, 5, 7$\\
\textbf{1.6} b\\
\textbf{1.7} a\\
\textbf{1.8} $1$\\
\textbf{1.9} $\mathbf{x}^T\mathbf{x} = 1$\\
\textbf{1.10} a, c

\medskip
\textbf{Section 2: Analysis}\\
\textbf{2.1} $2\pi$\\
\textbf{2.2} $\frac{\pi}{6}$\\
\textbf{2.3} a, c\\
\textbf{2.4} b, c\\
\textbf{2.5} c\\
\textbf{2.6} c\\
\textbf{2.7} a, c\\
\textbf{2.8} $f'(z) = u_x(x, y) + iv_x(x, y)$\\
\textbf{2.9} $f(S) = \{u + iv : |u| < 1,\ v > 0\}$\\
\textbf{2.10} $1$

\medskip
\textbf{Section 3: Topology}\\
\textbf{3.1} b, c\\
\textbf{3.2} none\\
\textbf{3.3} b\\
\textbf{3.4} b\\
\textbf{3.5} b, c\\
\textbf{3.6} $\mathcal{C}[0, 1]$\\
\textbf{3.7} a, c\\
\textbf{3.8} a, b\\
\textbf{3.9} a\\
\textbf{3.10} a, c

\medskip
\textbf{Section 4: Applied Mathematics}\\
\textbf{4.1} $t = 2$\\
\textbf{4.2} $mx'' = mg - k(x')^2,\ x(0) = x'(0) = 0$\\
\textbf{4.3} $v^2 = 2gy$\\
\textbf{4.4} $\sigma_n = n\omega_n$\\
\textbf{4.5} $x(t) = c_1 e^{-3t} + c_2 e^{2t},\ y(t) = -4c_1 e^{-3t} + c_2 e^{2t}$\\
\textbf{4.6} $4\pi^2$\\
\textbf{4.7} $u(x, t) = \frac{1}{2}[f(x + t) + f(x - t)]$\\
\textbf{4.8} $\frac{1}{2}$\\
\textbf{4.9} $x_{n+1} = \frac{1}{2}\left(x_n + \frac{10}{x_n}\right)$\\
\textbf{4.10} $F(s) = \frac{6}{s^4}$

\medskip
\textbf{Section 5: Miscellaneous}\\
\textbf{5.1} $4\sqrt{2}$\\
\textbf{5.2} $9$\\
\textbf{5.3} a, b, c\\
\textbf{5.4} $29$\\
\textbf{5.5} $220$\\
\textbf{5.6} $\frac{1}{2} + i\frac{\sqrt{3}}{2}$\\
\textbf{5.7} $n2^{n-1}$\\
\textbf{5.8} Maximum $= \sqrt{14}$; Minimum $= -\sqrt{14}$\\
\textbf{5.9} a, b, c\\
\textbf{5.10} $\frac{1}{3}$

\medskip
\noindent\emph{Note: 1.\ Accept any correct example for Qn.\ 1.3. 2.\ Accept any correct equivalent form of the answers for Qns.\ 2.8, 4.2, 4.3, 4.5 and 4.9. 3.\ If (c) is marked in Qn.\ 5.3, accept even if others are not marked.}
\end{multicols}
% ==================== NBHM 2010 ====================
\examheader{NBHM 2010}{Research Scholarships Screening Test}{Saturday, January 23, 2010}{Time Allowed: 150 Minutes \quad|\quad Maximum Marks: 40}

\subsection*{Instructions to Candidates}
\begin{itemize}
  \item Please ensure that this booklet contains 11 numbered (and printed) pages. The back of each printed page is blank and can be used for rough work.
  \item There are \textbf{five} sections, containing \textbf{ten} questions each, entitled Algebra, Analysis, Topology, Applied Mathematics and Miscellaneous. Answer as many questions as possible. The assessment of the paper will be based on the best \textbf{four} sections. Each question carries one point and the maximum possible score is \textbf{forty}.
  \item Answer each question, as directed, in the space provided in the \textbf{answer booklet}, which is being supplied separately. This question paper is meant to be retained by you and so \textbf{do not} answer questions on it.
  \item In certain questions you are required to pick out the qualifying statement(s) from multiple choices. None of the statements, or one or more than one statement may qualify. Write \textbf{none} if none of the statements qualify, or list the labels of \emph{all} the qualifying statements (amongst (a), (b), and (c)).
  \item Points will be awarded in the above questions only if \textbf{all} the correct choices are made. \textbf{There will be no partial credit.}
  \item $\N$ denotes the set of natural numbers, $\Z$ -- the integers, $\Q$ -- the rationals, $\R$ -- the reals and $\C$ -- the field of complex numbers. $\R^n$ denotes the $n$-dimensional Euclidean space, which is assumed to be endowed with its `usual' topology. The symbol $\Z_n$ will denote the ring of integers modulo $n$. The symbol $]a, b[$ will stand for the open interval $\{x \in \R \mid a < x < b\}$ while $[a, b]$ will stand for the corresponding closed interval; $[a, b[$ and $]a, b]$ will stand for the corresponding left-closed-right-open and left-open-right-closed intervals respectively. The symbol $I$ will denote the identity matrix of appropriate order. The space of continuous real valued functions on an interval $[a, b]$ is denoted by $\mathcal{C}[a, b]$ and is endowed with its usual `sup' norm.
\end{itemize}

\examsec{Section 1: Algebra}

\qitem{1.1} Solve the equation
\[
  x^4 - 2x^3 + 4x^2 + 6x - 21 = 0
\]
given that two of its roots are equal in magnitude but opposite in sign.

\qitem{1.2} Let $G$ be a group. A subgroup $H$ of $G$ is called \emph{characteristic} if $\varphi(H) \subset H$ for all automorphisms $\varphi$ of $G$. Pick out the true statement(s):
\begin{enumerate}[(a)]
  \item Every characteristic subgroup is normal.
  \item Every normal subgroup is characteristic.
  \item If $N$ is a normal subgroup of a group $G$, and $M$ is a characteristic subgroup of $N$, then $M$ is a normal subgroup of $G$.
\end{enumerate}

\qitem{1.3} Let $G$ be a group and let $H$ and $K$ be subgroups of $G$. The \emph{commutator subgroup} $(H, K)$ is defined as the smallest subgroup containing all elements of the form $hkh^{-1}k^{-1}$, where $h \in H$ and $k \in K$. Pick out the true statement(s):
\begin{enumerate}[(a)]
  \item If $H$ and $K$ are normal subgroups, then $(H, K)$ is a normal subgroup.
  \item If $H$ and $K$ are characteristic subgroups, then $(H, K)$ is a characteristic subgroup.
  \item $(G, G)$ is normal in $G$ and $G/(G, G)$ is abelian.
\end{enumerate}

\qitem{1.4} Write the following permutation as a product of disjoint cycles:
\[
  \begin{pmatrix} 1 & 2 & 3 & 4 & 5 & 6 \\ 6 & 5 & 4 & 3 & 1 & 2 \end{pmatrix}.
\]

\qitem{1.5} Pick out the true statement(s):
\begin{enumerate}[(a)]
  \item The set of all $2 \times 2$ matrices with rational entries (with the usual operations of matrix addition and matrix multiplication) is a ring which has no non-trivial ideals.
  \item Let $R = \mathcal{C}[0, 1]$ be considered as a ring with the usual operations of pointwise addition and pointwise multiplication. Let $\mathcal{I} = \{f : [0, 1] \to \R \mid f(1/2) = 0\}$. Then $\mathcal{I}$ is a maximal ideal.
  \item Let $R$ be a commutative ring and let $\mathcal{P}$ be a prime ideal of $R$. Then $R/\mathcal{P}$ is an integral domain.
\end{enumerate}

\qitem{1.6} What is the degree of the following numbers over $\Q$?\\
(a) $\sqrt{2} + \sqrt{3}$ \qquad (b) $\sqrt{2}\sqrt{3}$

\qitem{1.7} Let $V$ be the real vector space of all polynomials of degree $\le 3$ with real coefficients. Define the linear transformation
\[
  T(\alpha_0 + \alpha_1 x + \alpha_2 x^2 + \alpha_3 x^3) = \alpha_0 + \alpha_1(x+1) + \alpha_2(x+1)^2 + \alpha_3(x+1)^3.
\]
Write down the matrix of $T$ with respect to the basis $\{1, x, x^2, x^3\}$ of $V$.

\qitem{1.8} Let $A$ be an $n \times n$ upper triangular matrix with complex entries. Pick out the true statement(s):
\begin{enumerate}[(a)]
  \item If $A \neq 0$, and if $a_{ii} = 0$, for all $1 \le i \le n$, then $A^n = 0$.
  \item If $A \neq I$ and if $a_{ii} = 1$ for all $1 \le i \le n$, then $A$ is not diagonalizable.
  \item If $A \neq 0$, then $A$ is invertible.
\end{enumerate}

\qitem{1.9} Pick out the true statement(s):
\begin{enumerate}[(a)]
  \item There exist $n \times n$ matrices $A$ and $B$ with real entries such that $(I - (AB - BA))^n = 0$.
  \item If $A$ is a symmetric and positive definite $n \times n$ matrix, then $(\tr(A))^n \ge n^n \det(A)$, where `$\tr$' denotes the trace and `$\det$' denotes the determinant of a matrix.
  \item Let $A$ be a $5 \times 5$ skew-symmetric matrix with real entries. Then $A$ is singular.
\end{enumerate}

\qitem{1.10} Let $A$ be a $5 \times 5$ matrix whose characteristic polynomial is given by $(\lambda - 2)^3(\lambda + 2)^2$. If $A$ is diagonalizable, find $\alpha$ and $\beta$ such that
\[
  A^{-1} = \alpha A + \beta I.
\]

\examsec{Section 2: Analysis}

\qitem{2.1} Let $\{a_n\}$ be a sequence of positive real numbers such that $\lim_{n \to \infty} \frac{a_{n+1}}{a_n} = r < 1$. Can we evaluate $\lim_{n \to \infty} a_n$? If `yes', write down that limit.

\qitem{2.2} Test the following series for convergence:
\begin{enumerate}[(a)]
  \item $\displaystyle\sum_{n=1}^{\infty} \frac{(n+1)^n}{n^{n+\frac{5}{4}}}$.
  \item $\displaystyle\sum_{n=1}^{\infty} \frac{1}{\sqrt{n}}\tan\left(\frac{1}{n}\right)$.
\end{enumerate}

\qitem{2.3} Consider the polynomial $p(x) = a_0 + a_1 x + a_2 x^2 + \dots + a_n x^n$, with real coefficients. Pick out the case(s) which ensure that the polynomial $p(\cdot)$ has a root in the interval $[0, 1]$.
\begin{enumerate}[(a)]
  \item $a_0 < 0$ and $a_0 + a_1 + \dots + a_n > 0$.
  \item $\displaystyle a_0 + \frac{a_1}{2} + \dots + \frac{a_n}{n+1} = 0$.
  \item $\displaystyle \frac{a_0}{1\cdot 2} + \frac{a_1}{2\cdot 3} + \dots + \frac{a_n}{(n+1)(n+2)} = 0$.
\end{enumerate}

\qitem{2.4} Pick out the true statement(s):
\begin{enumerate}[(a)]
  \item The function $f(x) = \dfrac{\sin(x^2)}{\sin^2 x}$ is uniformly continuous on the interval $]0, 1[$.
  \item A continuous function $f : \R \to \R$ is uniformly continuous if it maps Cauchy sequences into Cauchy sequences.
  \item If a continuous function $f : \R \to \R$ is uniformly continuous, then it maps Cauchy sequences into Cauchy sequences.
\end{enumerate}

\qitem{2.5} Test the following for uniform convergence:
\begin{enumerate}[(a)]
  \item The sequence of functions $\left\{\dfrac{x^n}{1+x^n}\right\}$ over the interval $[0, 2]$.
  \item The series $\displaystyle\sum_{n=1}^{\infty} \frac{\sin nx}{n^2 + 1}$ over $\R$.
  \item The sequence of functions $\{n^2 x^2 e^{-nx}\}$ over the interval $]0, \infty[$.
\end{enumerate}

\qitem{2.6} Evaluate:
\[
  \lim_{n \to \infty} \left\{\left(1 + \frac{1}{n}\right)\left(1 + \frac{2}{n}\right)\cdots\left(1 + \frac{n}{n}\right)\right\}^{\frac{1}{n}}.
\]

\qitem{2.7} Let $f : [-\pi, \pi] \to \R$ be continuous. Pick out the case(s) which imply that $f \equiv 0$.
\begin{enumerate}[(a)]
  \item $\displaystyle\int_{-\pi}^{\pi} x^n f(x)\,dx = 0$, for all $n \ge 0$.
  \item $\displaystyle\int_{-\pi}^{\pi} f(x)\cos nx\,dx = 0$, for all $n \ge 0$.
  \item $\displaystyle\int_{-\pi}^{\pi} f(x)\sin nx\,dx = 0$, for all $n \ge 1$.
\end{enumerate}

\qitem{2.8} Evaluate:
\[
  \int_{\Gamma} \frac{dz}{(z^2 + 4)^2}
\]
where $\Gamma = \{z \in \C \mid |z - i| = 2\}$, described in the anticlockwise (\emph{i.e.}\ positive) direction.

\qitem{2.9} Find the residue at $z = 1$ of the function:
\[
  f(z) = \frac{5z - 2}{z(z - 1)}.
\]

\qitem{2.10} Let $f : \C \to \C$ be analytic. Which of the following conditions imply that $f$ is a constant function?
\begin{enumerate}[(a)]
  \item $\RePart f(z) > 0$ for all $z \in \C$.
  \item $|f(z)| \in \Z$ for all $z \in \C$.
  \item $f(z) = i$ when $z = \left(1 + \frac{k}{n}\right) + i$ for every positive integer $k$.
\end{enumerate}

\examsec{Section 3: Topology}

\qitem{3.1} Let $S^1$ denote the unit circle in the plane $\R^2$. Pick out the true statement(s):
\begin{enumerate}[(a)]
  \item There exists $f : S^1 \to \R$ which is continuous and one-one.
  \item For every continuous function $f : S^1 \to \R$, there exist uncountably many pairs of distinct points $x$ and $y$ in $S^1$ such that $f(x) = f(y)$.
  \item There exists $f : S^1 \to \R$ which is continuous and one-one and onto.
\end{enumerate}

\qitem{3.2} Which of the following metric spaces are separable?
\begin{enumerate}[(a)]
  \item $\mathcal{C}[0, 1]$ with its usual `sup-norm' topology.
  \item The space $\ell^\infty$ of all bounded real sequences with the metric $d(x, y) = \sup_n |x_n - y_n|$, where $x = (x_n)$ and $y = (y_n)$.
  \item The space $\ell^2$ of all square summable real sequences with the metric $d(x, y) = \left(\sum_{n=1}^{\infty} |x_n - y_n|^2\right)^{\frac{1}{2}}$, where $x = (x_n)$ and $y = (y_n)$.
\end{enumerate}

\qitem{3.3} Which of the following sets are nowhere dense?
\begin{enumerate}[(a)]
  \item The Cantor set in $[0, 1]$.
  \item The $xy$-plane in $\R^3$.
  \item Any countable set in $\R$.
\end{enumerate}

\qitem{3.4} Pick out the true statement(s).
\begin{enumerate}[(a)]
  \item If $f : ]-1, 1[ \to \R$ is bounded and continuous, it is uniformly continuous.
  \item If $f : S^1 \to \R$ is continuous, it is uniformly continuous.
  \item If $(X, d)$ is a metric space and $A \subset X$, then the function $f(x) = d(x, A)$ defined by $d(x, A) = \inf_{y \in A} d(x, y)$ is uniformly continuous.
\end{enumerate}

\qitem{3.5} Which of the following maps define a homeomorphism?
\begin{enumerate}[(a)]
  \item $f : \R \to ]0, \infty[$, where $f(x) = e^x$.
  \item $f : [0, 1] \to S^1$, where $f(t) = (\cos 2\pi t, \sin 2\pi t)$.
  \item Any map $f : X \to Y$ which is continuous, one-one and onto, if $X$ is compact and $Y$ is Hausdorff.
\end{enumerate}

\qitem{3.6} Consider the set of all $n \times n$ matrices with real entries as the space $\R^{n^2}$. Which of the following sets are compact?
\begin{enumerate}[(a)]
  \item The set of all orthogonal matrices.
  \item The set of all matrices with determinant equal to unity.
  \item The set of all invertible matrices.
\end{enumerate}

\qitem{3.7} In the set of all $n \times n$ matrices with real entries, considered as the space $\R^{n^2}$, which of the following sets are connected?
\begin{enumerate}[(a)]
  \item The set of all orthogonal matrices.
  \item The set of all matrices with trace equal to unity.
  \item The set of all symmetric and positive definite matrices.
\end{enumerate}

\qitem{3.8} Let $X$ be an arbitrary topological space. Pick out the true statement(s):
\begin{enumerate}[(a)]
  \item If $X$ is compact, then every sequence in $X$ has a convergent subsequence.
  \item If every sequence in $X$ has a convergent subsequence, then $X$ is compact.
  \item $X$ is compact if, and only if, every sequence in $X$ has a convergent subsequence.
\end{enumerate}

\qitem{3.9} Which of the following metric spaces are complete?
\begin{enumerate}[(a)]
  \item The space $\mathcal{C}^1[0, 1]$ of continuously differentiable real-valued functions on $[0, 1]$ with the metric $d(f, g) = \max_{t \in [0, 1]} |f(t) - g(t)|$.
  \item The space of all polynomials in a single variable with real coefficients, with the same metric as above.
  \item The space $\mathcal{C}[0, 1]$ with the metric $d(f, g) = \int_0^1 |f(t) - g(t)|\,dt$.
\end{enumerate}

\qitem{3.10} Classify the following alphabets into homeomorphism classes:
\[
  \mathbf{N,\ B,\ H,\ M}
\]

\examsec{Section 4: Applied Mathematics}

\qitem{4.1} A body, falling under gravity, experiences a resisting force of air proportional to the square of the velocity of the body. Write down the differential equation governing the motion satisfied by the distance $y(t)$ travelled by the body in time $t$.

\qitem{4.2} Reduce the following differential equation to a linear system of first order equations:
\[
  \frac{d^2 x}{dt^2} + P(t)\frac{dx}{dt} + Q(t)x = 0.
\]

\qitem{4.3} The volume of the unit ball in $\R^N$ is given by $\omega_N = \dfrac{\pi^{\frac{N}{2}}}{\Gamma\left(\frac{N}{2} + 1\right)}$, where $\Gamma(\cdot)$ denotes the usual gamma function. Write down the explicit value of $\omega_5$.

\qitem{4.4} Consider the differential equation $(1 + x)y' = py$, where $p$ is a constant. Assume that the equation has a power series solution $y = \sum_{n=0}^{\infty} a_n x^n$. Write down the recurrence relation for the coefficients $a_n$.

\qitem{4.5} In the above problem, if $y(0) = 1$, use the above series to find a closed form solution to the equation.

\qitem{4.6} Classify the following partial differential operators as elliptic, parabolic or hyperbolic:\\
(a) $5u_{xx} + 6u_{xy} + 2u_{yy}$. \qquad (b) $2u_{xx} + 6u_{xy} + 2u_{yy}$.

\qitem{4.7} Let $f$ and $g$ be two smooth scalar valued functions. Compute $\divg(\nabla f \times \nabla g)$.

\qitem{4.8} Let $S$ denote the sphere centred at the origin and of radius $a > 0$ in $\R^3$. Write down the coordinates of the unit outward normal to $S$ at the point $(x, y, z) \in S$.

\qitem{4.9} Use Gauss' divergence theorem to evaluate
\[
  \iint_S (x^4 + y^4 + z^4)\,dS
\]
where $S$ is the sphere mentioned in the preceding problem.

\qitem{4.10} Consider the domain $[0, 1] \times [0, T]$. Let $h > 0$ and $k > 0$. Let $x_n = nh$ and $t_m = mk$ for positive integers $m$ and $n$. Let $u_n^m = u(x_n, t_m)$. Write down the partial differential equation for which the following discretization defines a numerical scheme:
\[
  \frac{u_n^{m+1} - u_n^m}{k} = \frac{u_{n+1}^m - 2u_n^m + u_{n-1}^m}{h^2}.
\]

\examsec{Section 5: Miscellaneous}

\qitem{5.1} Let $n$ be a fixed positive integer and let $C_k$ denote the usual binomial coefficient ${}^n C_k$, the number of ways of choosing $k$ objects from $n$ objects. Evaluate:
\[
  C_0 + \frac{C_1}{2} + \frac{C_2}{3} + \dots + \frac{C_n}{n+1}.
\]

\qitem{5.2} Find the number of ways $2n$ persons can be seated at $2$ round tables, with $n$ persons at each table.

\qitem{5.3} Let a point $(x, y)$ be chosen at random in the square $[0, 1] \times [0, 1]$. Find the probability that $y \ge x^2$.

\qitem{5.4} Pick out the true statement(s):
\begin{enumerate}[(a)]
  \item If $n$ is an odd positive integer, then $8$ divides $n^2 - 1$.
  \item If $n$ and $m$ are odd positive integers, then $n^2 + m^2$ is not a perfect square.
  \item For every positive integer $n$, $\dfrac{n^5}{5} + \dfrac{n^3}{3} + \dfrac{7n}{15}$ is an integer.
\end{enumerate}

\qitem{5.5} Consider a circle of unit radius centered at $O$ in the plane. Let $AB$ be a chord which makes an angle $\theta$ with the tangent to the circle at $A$. Find the area of the triangle $OAB$.

\qitem{5.6} Evaluate: $\dfrac{1}{2\cdot 3} + \dfrac{1}{4\cdot 5} + \dfrac{1}{6\cdot 7} + \cdots$

\qitem{5.7} Evaluate: $1 + \dfrac{3}{4} + \dfrac{3\cdot 5}{4\cdot 8} + \dfrac{3\cdot 5\cdot 7}{4\cdot 8\cdot 12} + \cdots$

\qitem{5.8} Find the sum to $n$ terms as well as the sum to infinity of the series:
\[
  \frac{1}{3} + \frac{1}{4}\cdot\frac{1}{2!} + \frac{1}{5}\cdot\frac{1}{3!} + \cdots
\]

\qitem{5.9} If $a, b$ and $c$ are all distinct real numbers, find the condition that the following determinant vanishes:
\[
  \begin{vmatrix}
    a & a^2 & 1 + a^3 \\
    b & b^2 & 1 + b^3 \\
    c & c^2 & 1 + c^3
  \end{vmatrix}.
\]

\qitem{5.10} Assume that the line segment $[0, 2]$ in the $x$-axis of the plane acts as a mirror. A light ray from the point $(0, 1)$ gets reflected off this mirror and reaches the point $(2, 2)$. Find the point of incidence on the mirror.

\examsec{Answer Key --- NBHM 2010}
\begin{multicols}{2}
\textbf{Section 1: Algebra}\\
\textbf{1.1} $1 \pm i\sqrt{6},\ \pm\sqrt{3}$\\
\textbf{1.2} a, c\\
\textbf{1.3} a, b, c\\
\textbf{1.4} $(1625)(34)$\\
\textbf{1.5} a, b, c\\
\textbf{1.6} (a) $4$, (b) $2$\\
\textbf{1.7} $\begin{bmatrix} 1 & 1 & 1 & 1 \\ 0 & 1 & 2 & 3 \\ 0 & 0 & 1 & 3 \\ 0 & 0 & 0 & 1 \end{bmatrix}$\\
\textbf{1.8} a, b\\
\textbf{1.9} b, c\\
\textbf{1.10} $\alpha = 1/4,\ \beta = 0$

\medskip
\textbf{Section 2: Analysis}\\
\textbf{2.1} yes, $0$\\
\textbf{2.2} (a) convergent, (b) convergent\\
\textbf{2.3} a, b, c\\
\textbf{2.4} a, c\\
\textbf{2.5} (a) not uniformly convergent; (b) uniformly convergent; (c) not uniformly convergent\\
\textbf{2.6} $4/e$\\
\textbf{2.7} a, b, c\\
\textbf{2.8} $\pi/16$\\
\textbf{2.9} $3$\\
\textbf{2.10} a, b, c

\medskip
\textbf{Section 3: Topology}\\
\textbf{3.1} b\\
\textbf{3.2} a, c\\
\textbf{3.3} a, b\\
\textbf{3.4} b, c\\
\textbf{3.5} a, c\\
\textbf{3.6} a\\
\textbf{3.7} b, c\\
\textbf{3.8} none\\
\textbf{3.9} none\\
\textbf{3.10} $\{N, M\},\ \{B\},\ \{H\}$

\medskip
\textbf{Section 4: Applied Mathematics}\\
\textbf{4.1} $my'' = mg - c(y')^2$\\
\textbf{4.2} $x' = y;\ y' = -Py - Qx$\\
\textbf{4.3} $8\pi^2/15$\\
\textbf{4.4} $(n+1)a_{n+1} = (p - n)a_n,\ n \ge 0$\\
\textbf{4.5} $y = (1 + x)^p$\\
\textbf{4.6} (a) Elliptic, (b) Hyperbolic\\
\textbf{4.7} $0$\\
\textbf{4.8} $(x/a, y/a, z/a)$\\
\textbf{4.9} $12\pi a^6/5$\\
\textbf{4.10} $u_t = u_{xx}$

\medskip
\textbf{Section 5: Miscellaneous}\\
\textbf{5.1} $(2^{n+1} - 1)/(n + 1)$\\
\textbf{5.2} $(2n)!/n^2$\\
\textbf{5.3} $2/3$\\
\textbf{5.4} a, b, c\\
\textbf{5.5} $\sin\theta\cos\theta$\\
\textbf{5.6} $1 - \log 2$\\
\textbf{5.7} $2\sqrt{2}$\\
\textbf{5.8} $\frac{1}{2} - \frac{1}{(n+2)!};\ \frac{1}{2}$\\
\textbf{5.9} $abc + 1 = 0$\\
\textbf{5.10} $(2/3, 0)$

\medskip
\noindent\emph{Note: Accept any correct equivalent form of the answers.}
\end{multicols}
% ==================== NBHM 2011 ====================
\examheader{NBHM 2011}{Research Scholarships Screening Test}{Saturday, January 22, 2011}{Time Allowed: 150 Minutes \quad|\quad Maximum Marks: 40}

\subsection*{Instructions to Candidates}
\begin{itemize}
  \item Please ensure that this booklet contains 11 numbered (and printed) pages. The back of each printed page is blank and can be used for rough work.
  \item There are \textbf{five} sections, containing \textbf{ten} questions each, entitled Algebra, Analysis, Topology, Applied Mathematics and Miscellaneous. Answer as many questions as possible. The assessment of the paper will be based on the best \textbf{four} sections. Each question carries one point and the maximum possible score is \textbf{forty}.
  \item Answer each question, as directed, in the space provided in the \textbf{answer booklet}, which is being supplied separately. This question paper is meant to be retained by you and so \textbf{do not} answer questions on it.
  \item In certain questions you are required to pick out the qualifying statement(s) from multiple choices. None of the statements, or one or more than one statement may qualify. Write \textbf{none} if none of the statements qualify, or list the labels of \emph{all} the qualifying statements (amongst (a), (b), and (c)).
  \item Points will be awarded in the above questions only if \textbf{all} the correct choices are made. \textbf{There will be no partial credit.}
  \item $\N$ denotes the set of natural numbers, $\Z$ -- the integers, $\Q$ -- the rationals, $\R$ -- the reals and $\C$ -- the field of complex numbers. $\R^n$ denotes the $n$-dimensional Euclidean space, which is assumed to be endowed with its `usual' topology. The symbol $]a, b[$ will stand for the open interval $\{x \in \R \mid a < x < b\}$ while $[a, b]$ will stand for the corresponding closed interval; $[a, b[$ and $]a, b]$ will stand for the corresponding left-closed-right-open and left-open-right-closed intervals respectively. The symbol $I$ will denote the identity matrix of appropriate order. The space of continuous real valued functions on an interval $[a, b]$ is denoted by $\mathcal{C}[a, b]$ and is endowed with its usual `sup' norm.
\end{itemize}

\examsec{Section 1: Algebra}

\qitem{1.1} Solve: $x^4 - 3x^3 + 4x^2 - 3x + 1 = 0$.

\qitem{1.2} Pick out the true statements:
\begin{enumerate}[a.]
  \item Let $H$ and $K$ be subgroups of a group $G$. For $g \in G$, define the double coset $HgK = \{hgk \mid h \in H,\ k \in K\}$. Then, if $H$ is normal, we have $HgH = gH$ for all $g \in G$.
  \item Let $GL(n; \C)$ be the group of all $n \times n$ invertible matrices with complex entries. The set of all $n \times n$ invertible upper triangular matrices is a normal subgroup.
  \item Let $\M(n; \R)$ denote the set of all $n \times n$ matrices with real entries (identified with $\R^{n^2}$ and endowed with its usual topology) and let $GL(n; \R)$ denote the group of invertible matrices. Let $G$ be a subgroup of $GL(n; \R)$. Define
  \[
    H = \left\{A \in G \;\middle|\; \begin{array}{l}\text{there exists } \varphi : [0, 1] \to G \text{ continuous,}\\ \text{such that } \varphi(0) = A,\ \varphi(1) = I\end{array}\right\}.
  \]
  Then, $H$ is a normal subgroup of $G$.
\end{enumerate}

\qitem{1.3} How many (non-isomorphic) groups of order $15$ are there?

\qitem{1.4} Pick out the true statements:
\begin{enumerate}[a.]
  \item Let $R$ be a commutative ring with identity. Let $M$ be an ideal such that every element of $R$ not in $M$ is a unit. Then $R/M$ is a field.
  \item Let $R$ be as above and let $M$ be an ideal such that $R/M$ is an integral domain. Then $M$ is a prime ideal.
  \item Let $R = \mathcal{C}[0, 1]$ be the ring of real-valued continuous functions on $[0, 1]$ with respect to pointwise addition and pointwise multiplication. Let $M = \{f \in R \mid f(0) = f(1) = 0\}$. Then $M$ is a maximal ideal.
\end{enumerate}

\qitem{1.5} Write down all the possible values for the degree of an irreducible polynomial in $\R[x]$.

\qitem{1.6} Let $V$ be the real vector space of all polynomials in $\R[x]$ with degree less than, or equal to $4$. Consider the linear transformation which maps $p \in V$ to its derivative $p'$. If the matrix of this transformation with respect to the basis $\{1, x, x^2, x^3, x^4\}$ is $A$, write down the matrix $A^3$.

\qitem{1.7} Let $\T(n; \R) \subset \M(n; \R)$ denote the set of all matrices whose trace is zero. Write down a basis for $\T(2; \R)$.

\qitem{1.8} What is the quotient space $\M(n; \R)/\T(n; \R)$ isomorphic to?

\qitem{1.9} Construct a $2 \times 2$ matrix $A (\neq I)$ with real entries such that $A^3 = I$.

\qitem{1.10} If $A \in \M(n; \R)$, let ${}^t A$ denote its transpose. A matrix $S \in \M(n; \R)$ is said to be \emph{skew-symmetric} if ${}^t S = -S$. Pick out the true statements:
\begin{enumerate}[a.]
  \item If $S \in \M(n; \R)$ is skew-symmetric and non-singular, then $n$ is even.
  \item Let $G = \{T \in GL(n; \R) \mid {}^t T S T = S,\ \text{for all skew-symmetric } S \in \M(n; \R)\}$. Then $G$ is a subgroup of $GL(n; \R)$.
  \item Let $I_n$ and $O_n$ denote the $n \times n$ identity and null matrices respectively. Let $S$ be the $2n \times 2n$ matrix given in block form by
  \[
    \begin{bmatrix} O_n & I_n \\ -I_n & O_n \end{bmatrix}.
  \]
  If $X$ is a $2n \times 2n$ matrix such that ${}^t X S + SX = 0$, then the trace of $X$ is zero.
\end{enumerate}

\examsec{Section 2: Analysis}

\qitem{2.1} Let $\{a_n\}$ be a sequence of positive terms. Pick out the cases which imply that $\sum a_n$ is convergent.
\begin{enumerate}[a.]
  \item $\displaystyle\lim_{n \to \infty} n^{\frac{3}{2}} a_n = \frac{3}{2}$.
  \item $\displaystyle\sum n^2 a_n^2 < \infty$.
  \item $\displaystyle\frac{a_{n+1}}{a_n} < \left(\frac{n}{n+1}\right)^2$, for all $n$.
\end{enumerate}

\qitem{2.2} Evaluate:
\[
  \lim_{n \to \infty} \left\{\frac{1}{1+n^3} + \frac{4}{8+n^3} + \dots + \frac{n^2}{n^3+n^3}\right\}.
\]

\qitem{2.3} Find the points in $\R$ where the following function is differentiable:
\[
  f(x) = \begin{cases} \tan^{-1} x, & \text{if } |x| \le 1 \\ \frac{\pi}{4}\sgn(x) + \frac{|x|-1}{2}, & \text{if } |x| > 1, \end{cases}
\]
where $\sgn(x)$ equals $+1$ if $x > 0$, $-1$ if $x < 0$ and is equal to zero if $x = 0$ and $\tan^{-1}(x)$ takes its values in the range $]-\pi/2, \pi/2[$ for real numbers $x$.

\qitem{2.4} Pick out the true statements:
\begin{enumerate}[a.]
  \item If $P$ is a polynomial in one variable with real coefficients which has all its roots real, then its derivative $P'$ has all its roots real as well.
  \item The equation $\cos(\sin x) = x$ has exactly one solution in the interval $\left[0, \frac{\pi}{2}\right]$.
  \item $\cos x > 1 - \frac{x^2}{2}$ for all $x > 0$.
\end{enumerate}

\qitem{2.5} Let $f, f_n : [0, 1] \to \R$ be continuous functions. Complete the following sentence such that both statements (a) and (b) below are true: ``Let $f_n \to f$ \underline{\hspace{2cm}}''
\begin{enumerate}[a.]
  \item $\displaystyle\lim_{n \to \infty} \int_0^1 f_n(x)\,dx = \int_0^1 f(x)\,dx$.
  \item $\displaystyle\lim_{n \to \infty}\lim_{x \to 0} f_n(x) = \lim_{x \to 0}\lim_{n \to \infty} f_n(x)$.
\end{enumerate}

\qitem{2.6} Let $f : ]0, 1[ \to \R$ be continuous. Pick out the statements which imply that $f$ is uniformly continuous.
\begin{enumerate}[a.]
  \item $|f(x) - f(y)| \le \sqrt{|x - y|}$, for all $x, y \in ]0, 1[$.
  \item $f(1/n) \to 1/2$ and $f(1/n^2) \to 1/4$.
  \item $f(x) = x^{\frac{1}{2}} \sin\dfrac{1}{x^3}$.
\end{enumerate}

\qitem{2.7} Evaluate: $\displaystyle\iint_{[0,1]\times[0,1]} \max\{x, y\}\,dxdy$.

\qitem{2.8} Let $f : \C \to \C$ be analytic. Pick out the cases when $f$ is \textbf{not} necessarily a constant.
\begin{enumerate}[a.]
  \item $\ImPart(f'(z)) > 0$ for all $z \in \C$.
  \item $f(n) = 3$ for all $n \in \Z$.
  \item $f'(0) = 0$ and $|f'(z)| \le 3$ for all $z \in \C$.
\end{enumerate}

\qitem{2.9} Let $f : \C \to \C$ be analytic. Write $z = x + iy$ and $f = u + iv$, where $u$ and $v$ are real valued functions of $x$ and $y$. Pick out the true statements.
\begin{enumerate}[a.]
  \item $f'(z) = \dfrac{\partial u}{\partial x} + i\dfrac{\partial v}{\partial x}$.
  \item $\dfrac{\partial^2 u}{\partial x^2} + \dfrac{\partial^2 u}{\partial y^2} = 0$.
  \item $f''(0) = \dfrac{1}{2\pi i}\displaystyle\int_{|z|=1} \frac{f(z)}{z^3}\,dz$.
\end{enumerate}

\qitem{2.10} Find the square roots of $1 + i\sqrt{3}$.

\examsec{Section 3: Topology}

\qitem{3.1} Which of the following define a metric?
\begin{enumerate}[a.]
  \item $d((x, y), (x', y')) = \min\{|x - x'|, |y - y'|\}$ on $\R^2$.
  \item $d((x, y), (x', y')) = |x| + |y| + |x'| + |y'|$ on $\R^2$.
  \item $D((x, y), (x', y')) = d(x, x') + d(y, y')$ on $X \times X$, where $(X, d)$ is a metric space.
\end{enumerate}

\qitem{3.2} Let $(X, d)$ be a metric space and let $A \subset X$. For $x \in X$ define $d(x, A) = \inf\{d(x, y) \mid y \in A\}$. Pick out the true statements:
\begin{enumerate}[a.]
  \item $x \mapsto d(x, A)$ is a uniformly continuous function.
  \item If $\partial A = \{x \in X \mid d(x, A) = 0\} \cap \{x \in X \mid d(x, X \setminus A) = 0\}$, then $\partial A$ is closed for any $A \subset X$.
  \item Let $A$ and $B$ be subsets of $X$ and define $d(A, B) = \inf\{d(a, B) \mid a \in A\}$. Then $d(A, B) = d(B, A)$.
\end{enumerate}

\qitem{3.3} Let $X$ be a topological space and for $A \subset X$, denote by $\overline{A}$ and $A^\circ$, the closure and interior of $A$ respectively. Pick out the true statements.
\begin{enumerate}[a.]
  \item $\overline{A \cup B} = \overline{A} \cup \overline{B}$.
  \item $\overline{A \cap B} = \overline{A} \cap \overline{B}$.
  \item Consider $\R$ as the $x$-axis in $\R^2$. Then $\R^\circ = \emptyset$.
\end{enumerate}

\qitem{3.4} Pick out the true statements.
\begin{enumerate}[a.]
  \item Let $\{X_i\}_{i \in \mathcal{I}}$ be topological spaces. Then, the product topology is the smallest topology on $X = \prod_{i \in \mathcal{I}} X_i$ such that each of the canonical projections $p_i : X \to X_i$ is continuous.
  \item Let $X$ be a topological space and $W \subset X$. Then, the induced subspace topology on $W$ is the smallest topology such that $id|_W : W \to X$, where $id$ is the identity map, is continuous.
  \item Let $X = \R^n$ with the usual topology. This is the smallest topology such that all linear functionals on $X$ are continuous.
\end{enumerate}

\qitem{3.5} Which of the following subsets are dense in the given spaces?
\begin{enumerate}[a.]
  \item The set of trigonometric polynomials in the space of continuous functions on $[-\pi, \pi]$ which are $2\pi$-periodic (with the sup-norm topology).
  \item The subset of $\mathcal{C}^\infty$ functions with compact support in $\R$ in the space of bounded real-valued continuous functions on $\R$ (with the sup-norm topology).
  \item $GL(n; \R)$ in $\M(n; \R)$ (with its usual topology after identification with $\R^{n^2}$).
\end{enumerate}

\qitem{3.6} Pick out the compact sets.
\begin{enumerate}[a.]
  \item $\{(x, y) \mid x^2 - y^2 = 1\} \subset \R^2$.
  \item $\{\operatorname{Tr}(A) \mid A \in \M(n; \R),\ A \text{ orthogonal}\} \subset \R$, where $\operatorname{Tr}(A)$ denotes the trace of the matrix $A$.
  \item The set of all matrices in $\M(n; \R)$ all of whose eigenvalues satisfy the condition $|\lambda| \le 2$.
\end{enumerate}

\qitem{3.7} Pick out the connected sets.
\begin{enumerate}[a.]
  \item $\{(x, y) \mid xy = 1\} \subset \R^2$.
  \item The set of all upper triangular matrices in $\M(n; \R)$.
  \item The set of all invertible diagonal matrices in $\M(n; \R)$.
\end{enumerate}

\qitem{3.8} Pick out the true statements.
\begin{enumerate}[a.]
  \item Let $f : \Z \to \Z^2$ be a bijection. There exists a continuous function from $\R$ to $\R^2$ which extends $f$.
  \item Let $D$ denote the closed unit disc in $\R^2$. There exists a continuous mapping $f : D \setminus \{(0, 0)\} \to \{x \in \R \mid |x| \le 1\}$ which is onto.
  \item Let $D$ denote the closed unit disc in $\R^2$. There exists a continuous mapping $f : D \setminus \{(0, 0)\} \to \{x \in \R \mid |x| > 1\}$ which is onto.
\end{enumerate}

\qitem{3.9} A Hausdorff topological space is said to be normal if given any two disjoint closed sets $A$ and $B$, there exist disjoint open sets $U$ and $V$ such that $A \subset U$ and $B \subset V$. Pick out the true statements.
\begin{enumerate}[a.]
  \item Every metric space is normal.
  \item If $X$ is a normal space with at least two distinct points, then there exist non-constant real-valued continuous functions on $X$.
  \item If $X$ is normal and $Y \subset X$ is closed, then $Y$ is normal for the induced topology.
\end{enumerate}

\qitem{3.10} Which of the following pairs of sets are homeomorphic?
\begin{enumerate}[a.]
  \item $A = \{(x, y) \mid x^2 + y^2 - 2x + 4y - 5 = 0\}$ and $B = \{(x, y) \mid 5x^2 + 3y^2 = 1\}$.
  \item $A = \{(x, y) \mid x^2 + y^2 - 2x + 4y - 5 = 0\}$ and $B = \{(x, y) \mid 5x^2 - 3y^2 = 1\}$.
  \item $A = \{(x, y) \mid x^2 + y^2 - 2x + 4y - 5 \le 0\}$ and $B = \{(x, y) \mid 5x^2 + 3y^2 \ge 1\}$.
\end{enumerate}

\examsec{Section 4: Applied Mathematics}

\qitem{4.1} Simpson's rule is used to approximate the integral $\int_0^1 f(x)\,dx$. If $f$ is a polynomial, what is the maximum possible degree it can have so that Simpson's rule gives the exact value of this integral?

\qitem{4.2} A right circular cylinder of fixed volume has maximum total surface area. What is the relationship between its height $h$ and radius $r$?

\qitem{4.3} In the equations governing the flow of an incompressible fluid of uniform density, if $\mathbf{u}$ is the velocity vector and $p$ is the pressure, write down the equation which expresses the law of conservation of mass.

\qitem{4.4} A particle of mass $M$ is attached to a fixed wall by a spring. The spring exerts no force when the particle is at its equilibrium position at $x = 0$ and exerts a restoring force proportional to the displacement when it is displaced to a distance $x$. In addition, there is a damping force due to the medium in which the displacement takes place, which is a force opposing the motion and is proportional to the velocity of the particle. If the particle is pulled to a position $x_0$ at time $t = 0$ and is released without any velocity, write down the initial value problem governing the motion of the particle.

\qitem{4.5} Solve the following linear programming problem:
\begin{align*}
  \max z &= 5x + 7y \\
  x - y &\le 1 \\
  2x + y &\ge 2 \\
  x + 2y &\le 4 \\
  x, y &\ge 0.
\end{align*}

\qitem{4.6} Write down the dual of the above problem.

\qitem{4.7} Find the general solution of the system:
\begin{align*}
  \frac{dx}{dt} &= x + y \\
  \frac{dy}{dt} &= 4x - 2y.
\end{align*}

\qitem{4.8} Let $\mathbf{a}, \mathbf{b}$ and $\mathbf{c}$ be vectors in $\R^3$. Express $\mathbf{a} \times (\mathbf{b} \times \mathbf{c})$ as a linear combination of $\mathbf{b}$ and $\mathbf{c}$.

\qitem{4.9} Solve: $\dfrac{\partial^2 u}{\partial x^2} = 6xy;\quad u(0, y) = y;\quad \dfrac{\partial u}{\partial x}(1, y) = 0$.

\qitem{4.10} Let $\Omega$ be a smooth plane domain of unit area. Let $u(x, y) = 3x^2 + y^2$. If $\frac{\partial u}{\partial n}$ denotes its outer normal derivative on $\partial\Omega$, the boundary of $\Omega$, compute $\int_{\partial\Omega} \frac{\partial u}{\partial n}\,ds$.

\examsec{Section 5: Miscellaneous}

\qitem{5.1} Let $V$ be a real vector space of real-valued functions on a given set. Assume that constant functions are in $V$ and that if $f \in V$, then $f^2 \in V$ and that $|f| \in V$. Pick out the true statements.
\begin{enumerate}[a.]
  \item If $f, g \in V$, then $fg \in V$.
  \item If $f, g \in V$, then $\max\{f, g\} \in V$.
  \item If $f \in V$ and $p$ is any polynomial in one variable, with real coefficients, then $p(f) \in V$.
\end{enumerate}

\qitem{5.2} A fair coin is tossed $10$ times, the tosses being independent of each other. Find the probability that the results of the third, fourth and fifth tosses are identical.

\qitem{5.3} Determine if the following collections are countable or uncountable.
\begin{enumerate}[a.]
  \item The collection of all finite subsets of $\N$.
  \item The collection of all infinite sequences of positive integers.
  \item The collection of all roots of all polynomials in one variable, with integer coefficients.
\end{enumerate}

\qitem{5.4} Find the maximum value of $x + 2y + 3z$ subject to the constraint $x^2 + y^2 + z^2 = 1$.

\qitem{5.5} Let $A_n$ be the $n \times n$ matrix whose $(i, j)$-th entry is given by $2\delta_{ij} - \delta_{i+1,j} - \delta_{i,j+1}$, where $\delta_{ij}$ equals $1$ if $i = j$ and zero otherwise. Compute the determinant of $A_n$.

\qitem{5.6} How many real roots does the following equation have? $3^x + 4^x = 5^x$.

\qitem{5.7} Let $N > 1$ be a positive integer. Let $\phi(N)$ denote the number of positive integers less than $N$ and prime to it (unity being included in this count). Express the sum of all the integers less than $N$ and prime to it in terms of $\phi(N)$.

\qitem{5.8} Pick out the true statements.
\begin{enumerate}[a.]
  \item The sum of $r$ consecutive positive integers is divisible by $r$.
  \item The product of $r$ consecutive positive integers is divisible by $r!$.
  \item For each positive integer $r$, there exist $r$ consecutive positive integers which are all composite.
\end{enumerate}

\qitem{5.9} Let $n$ be a fixed positive integer and let $0 \le k \le n$. We denote by $C_k$, the number of ways of choosing $k$ objects from $n$ distinct objects. Sum to $n$ terms:
\[
  3C_1 + 7C_2 + 11C_3 + \cdots
\]

\qitem{5.10} Find the sum of the following infinite series:
\[
  \frac{1}{5} - \frac{1\cdot 4}{5\cdot 10} + \frac{1\cdot 4\cdot 7}{5\cdot 10\cdot 15} - \cdots
\]

\examsec{Answer Key --- NBHM 2011}
\begin{multicols}{2}
\textbf{Section 1: Algebra}\\
\textbf{1.1} $1, 1, \frac{1 \pm i\sqrt{3}}{2}$\\
\textbf{1.2} a, c\\
\textbf{1.3} one\\
\textbf{1.4} a, b\\
\textbf{1.5} $1, 2$\\
\textbf{1.6} $\begin{bmatrix} 0 & 0 & 0 & 6 & 0 \\ 0 & 0 & 0 & 0 & 24 \\ 0 & 0 & 0 & 0 & 0 \\ 0 & 0 & 0 & 0 & 0 \\ 0 & 0 & 0 & 0 & 0 \end{bmatrix}$\\
\textbf{1.7} Any three linearly independent $2 \times 2$ matrices with trace zero. Example: $\begin{bmatrix} 0 & 1 \\ 0 & 0 \end{bmatrix},\ \begin{bmatrix} 0 & 0 \\ 1 & 0 \end{bmatrix},\ \begin{bmatrix} 1 & 0 \\ 0 & -1 \end{bmatrix}$\\
\textbf{1.8} $\R$\\
\textbf{1.9} Any $2 \times 2$ matrix with trace $-1$ and determinant $1$.\\
\textbf{1.10} a, b, c

\medskip
\textbf{Section 2: Analysis}\\
\textbf{2.1} a, b, c\\
\textbf{2.2} $\frac{1}{3}\log 2$\\
\textbf{2.3} $\R \setminus \{-1\}$\\
\textbf{2.4} a, b, c\\
\textbf{2.5} uniformly on $[0, 1]^*$\\
\textbf{2.6} a, c\\
\textbf{2.7} $\frac{2}{3}$\\
\textbf{2.8} a, b\\
\textbf{2.9} a, b\\
\textbf{2.10} $\pm\frac{\sqrt{3}+i}{\sqrt{2}}$

\medskip
\textbf{Section 3: Topology}\\
\textbf{3.1} c\\
\textbf{3.2} a, b, c\\
\textbf{3.3} a, c\\
\textbf{3.4} a, b, c\\
\textbf{3.5} a, c\\
\textbf{3.6} b\\
\textbf{3.7} b\\
\textbf{3.8} a, b\\
\textbf{3.9} a, b, c\\
\textbf{3.10} a

\medskip
\textbf{Section 4: Applied Mathematics}\\
\textbf{4.1} $3$\\
\textbf{4.2} $h = 2r$\\
\textbf{4.3} $\divg(\mathbf{u}) = 0$\\
\textbf{4.4} $Mx'' + cx' + kx = 0;\ x(0) = x_0;\ x'(0) = 0$, where $k$ and $c$ are positive constants\\
\textbf{4.5} $\max z = 17;\ x = 2;\ y = 1$\\
\textbf{4.6} $\min f = u - 2v + 4w$ such that $u - 2v + w \ge 5;\ -u - v + 2w \ge 7;\ u, v, w \ge 0$\\
\textbf{4.7} $x(t) = c_1 e^{2t} + c_2 e^{-3t};\ y(t) = c_1 e^{2t} - 4c_2 e^{-3t}$\\
\textbf{4.8} $(\mathbf{a}\cdot\mathbf{c})\mathbf{b} - (\mathbf{a}\cdot\mathbf{b})\mathbf{c}$\\
\textbf{4.9} $u(x, y) = y(x^3 - 3x + 1)$\\
\textbf{4.10} $8$

\medskip
\textbf{Section 5: Miscellaneous}\\
\textbf{5.1} a, b, c\\
\textbf{5.2} $\frac{1}{4}$\\
\textbf{5.3} a.\ countable; b.\ uncountable; c.\ countable\\
\textbf{5.4} $\sqrt{14}$\\
\textbf{5.5} $n + 1$\\
\textbf{5.6} one\\
\textbf{5.7} $\frac{N}{2}\phi(N)$\\
\textbf{5.8} b, c\\
\textbf{5.9} $4n 2^{n-1} - 2^n + 1$\\
\textbf{5.10} $1 - \frac{1}{2}5^{\frac{1}{3}}$

\medskip
\noindent\emph{Note: Accept any correct equivalent form of the answers. $^*$Qn.\ 2.5: Accept even if the answer is just `uniformly'.}
\end{multicols}
% ==================== NBHM 2012 ====================
\examheader{NBHM 2012}{Research Scholarships Screening Test}{Saturday, January 28, 2012}{Time Allowed: 150 Minutes \quad|\quad Maximum Marks: 40}

\subsection*{Instructions to Candidates \& Notation}
\begin{itemize}
  \item Please ensure that this booklet contains 12 numbered (and printed) pages. The back of each printed page is blank and can be used for rough work.
  \item There are \textbf{five} sections, containing \textbf{ten} questions each, entitled Algebra, Analysis, Topology, Applied Mathematics and Miscellaneous. Answer as many questions as possible. The assessment of the paper will be based on the best \textbf{four} sections. Each question carries one point and the maximum possible score is \textbf{forty}.
  \item Answer each question, as directed, in the space provided in the \textbf{answer booklet}, which is being supplied separately. This question paper is meant to be retained by you and so \textbf{do not} answer questions on it.
  \item In certain questions you are required to pick out the qualifying statement(s) from multiple choices. None of the statements, or one or more than one statement may qualify. Write \textbf{none} if none of the statements qualify, or list the labels of \emph{all} the qualifying statements (amongst (a), (b), and (c)).
  \item Points will be awarded for questions involving more than one answer only if \textbf{all} the correct answers are given. \textbf{There will be no partial credit.} Calculators are \textbf{not} allowed.
  \item \textbf{Notation:} $\N, \Z, \Q, \R, \C$ denote standard number sets; $\R^n, \C^n$ have their usual topology; $\M_n(\R), \M_n(\C)$ denote $n \times n$ matrices identified with $\R^{n^2}, \C^{n^2}$; $\Z_n$ is the ring of integers modulo $n$; $]a, b[, [a, b], [a, b[, ]a, b]$ denote intervals; $\mathcal{C}[a, b], \mathcal{C}^1[a, b], \mathcal{C}^\infty$ denote continuous, $\mathcal{C}^1$, and smooth functions; $f', f''$ denote derivatives; $I$ is the identity matrix; $\det(A), \tr(A)$ denote determinant and trace; $GL_n(\R), GL_n(\C)$ denote invertible $n \times n$ matrices.
\end{itemize}

\examsec{Section 1: Algebra}

\qitem{1.1} Which of the following are subgroups of $GL_3(\C)$?
\begin{enumerate}[a.]
  \item $H = \{A \in \M_3(\C) \mid \det(A) = 2^l,\ l \in \Z\}$.
  \item $H = \left\{\begin{bmatrix} 1 & \alpha & \beta \\ 0 & 1 & \gamma \\ 0 & 0 & 1 \end{bmatrix} \;\middle|\; \alpha, \beta, \gamma \in \C\right\}$.
  \item $H = \left\{\begin{bmatrix} 1 & 0 & \alpha \\ 0 & 1 & 0 \\ 0 & 0 & 1 \end{bmatrix} \;\middle|\; \alpha \in \C\right\}$.
\end{enumerate}

\qitem{1.2} Let $S_7$ denote the symmetric group of all permutations of the symbols $\{1, 2, 3, 4, 5, 6, 7\}$. Pick out the true statements:
\begin{enumerate}[a.]
  \item $S_7$ has an element of order $10$;
  \item $S_7$ has an element of order $15$;
  \item the order of any element of $S_7$ is at most $12$.
\end{enumerate}

\qitem{1.3} Let $\mathcal{C}(\R)$ denote the ring of all continuous real-valued functions on $\R$, with the operations of pointwise addition and pointwise multiplication. Which of the following form an ideal in this ring?
\begin{enumerate}[a.]
  \item The set of all $\mathcal{C}^\infty$ functions with compact support.
  \item The set of all continuous functions with compact support.
  \item The set of all continuous functions which vanish at infinity, \emph{i.e.}\ functions $f$ such that $\lim_{|x| \to \infty} f(x) = 0$.
\end{enumerate}

\qitem{1.4} Find the number of non-zero elements in the field $\Z_p$, where $p$ is an odd prime number, which are squares, \emph{i.e.}\ of the form $m^2$, $m \in \Z_p$, $m \neq 0$.

\qitem{1.5} Find the inverse in $\Z_5$ of the following matrix:
\[
  \begin{bmatrix} 1 & 2 & 0 \\ 0 & 2 & 4 \\ 0 & 0 & 3 \end{bmatrix}.
\]

\qitem{1.6} Let $\mathcal{P}_3$ denote the (real) vector space of all polynomials (in one variable), with real coefficients and of degree less than, or equal to, $3$, equipped with the standard basis $\{1, x, x^2, x^3\}$. Write down the matrix (with respect to this basis) of the linear transformation
\[
  L(p) = p'' - 2p' + p,\quad p \in \mathcal{P}_3.
\]

\qitem{1.7} Find the unique polynomial $p \in \mathcal{P}_3$ (see Question 1.6 above) such that $p'' - 2p' + p = x^3$.

\qitem{1.8} Let $A = (a_{ij}) \in \M_n(\R)$, $n \ge 3$. Let $B = (b_{ij})$ be the matrix of its cofactors, \emph{i.e.}\ $b_{ij}$ is the cofactor of the entry $a_{ij}$ in $A$. What is the rank of $B$ when
\begin{enumerate}[a.]
  \item the rank of $A$ is $n$?
  \item the rank of $A$ is less than, or equal to, $n - 2$?
\end{enumerate}

\qitem{1.9} Let $A \in \M_3(\R)$ which is \textbf{not} a diagonal matrix. Pick out the cases when $A$ is diagonalizable over $\R$:
\begin{enumerate}[a.]
  \item when $A^2 = A$;
  \item when $(A - 3I)^2 = 0$;
  \item when $A^2 + I = 0$.
\end{enumerate}

\qitem{1.10} Let $A \in \M_3(\R)$ which is \textbf{not} a diagonal matrix. Let $p$ be a polynomial (in one variable), with real coefficients and of degree $3$ such that $p(A) = 0$. Pick out the true statements:
\begin{enumerate}[a.]
  \item $p = cp_A$ where $c \in \R$ and $p_A$ is the characteristic polynomial of $A$;
  \item if $p$ has a complex root (\emph{i.e.}\ a root with non-zero imaginary part), then $p = cp_A$, with $c$ and $p_A$ as above;
  \item if $p$ has a complex root, then $A$ is diagonalizable over $\C$.
\end{enumerate}

\examsec{Section 2: Analysis}

\qitem{2.1} Which of the following statements are true?
\begin{enumerate}[a.]
  \item Let $\{a_{mn}\}$, $m, n \in \N$, be an arbitrary double sequence of real numbers. Then $\sum_{m=1}^{\infty}\sum_{n=1}^{\infty} a_{mn}^3 = \sum_{n=1}^{\infty}\sum_{m=1}^{\infty} a_{mn}^3$.
  \item Let $\{a_{mn}\}$, $m, n \in \N$, be an arbitrary double sequence of real numbers. Then $\sum_{m=1}^{\infty}\sum_{n=1}^{\infty} a_{mn}^2 = \sum_{n=1}^{\infty}\sum_{m=1}^{\infty} a_{mn}^2$.
  \item Let $\{a_{mn}\}$, $m, n \in \N$, be a double sequence of real numbers such that $|a_{mn}| \le \sqrt{m/n}$ for all $m, n \in \N$. Then $\sum_{m=1}^{\infty}\sum_{n=1}^{\infty} \frac{a_{mn}}{m^2 n} = \sum_{n=1}^{\infty}\sum_{m=1}^{\infty} \frac{a_{mn}}{m^2 n}$.
\end{enumerate}

\qitem{2.2} Let $f \in \mathcal{C}[-1, 1]$. Evaluate: $\displaystyle\lim_{h \to 0} \frac{1}{h}\int_{-h}^{h} f(t)\,dt$.

\qitem{2.3} Let $f \in \mathcal{C}^1[-1, 1]$. Evaluate: $\displaystyle\lim_{n \to \infty} \frac{1}{n}\sum_{k=1}^{n} f'\left(\frac{k}{3n}\right)$.

\qitem{2.4} Let $f \in \mathcal{C}[-\pi, \pi]$. Evaluate:
\begin{enumerate}[a.]
  \item $\displaystyle\lim_{n \to \infty} \int_{-\pi}^{\pi} f(t)\cos nt\,dt$;
  \item $\displaystyle\lim_{n \to \infty} \int_{-\pi}^{\pi} f(t)\cos^2 nt\,dt$;
\end{enumerate}

\qitem{2.5} In each of the following cases, examine whether the given sequence (or series) of functions converges uniformly over the given domain:
\begin{enumerate}[a.]
  \item $f_n(x) = \dfrac{nx}{1+nx}$, $x \in ]0, \infty[$;
  \item $\displaystyle\sum_{n=1}^{\infty} \frac{n\sin nx}{e^n}$, $x \in [0, \pi]$;
  \item $f_n(x) = \dfrac{x^n}{1+x^n}$, $x \in [0, 2]$.
\end{enumerate}

\qitem{2.6} Compute $F'(x)$ where $F(x) = \displaystyle\int_{-x}^{x} \frac{1 - e^{-xy}}{y}\,dy,\ x > 0$.

\qitem{2.7} Let $a > 0$ and let $k \in \N$. Evaluate: $\displaystyle\lim_{n \to \infty} a^{-nk} \prod_{j=1}^{k} \left(a + \frac{j}{n}\right)^n$.

\qitem{2.8} Write down the power series expansion of the function $f(z) = 1/z^2$ about the point $z = 2$.

\qitem{2.9} Let $C$ be the circle $|z + 2| = 3$ described in the anti-clockwise (\emph{i.e.}\ positive) sense in the complex plane. Evaluate: $\displaystyle\int_C \frac{dz}{z^3(z+4)}$.

\qitem{2.10} Which of the following statements are true?
\begin{enumerate}[a.]
  \item There exists an entire function $f : \C \to \C$ which takes only real values and is such that $f(0) = 0$ and $f(1) = 1$.
  \item There exists an entire function $f : \C \to \C$ such that $f\left(n + \frac{1}{n}\right) = 0$ for all $n \in \N$.
  \item There exists an entire function $f : \C \to \C$ which is onto and which is such that $f(1/n) = 0$ for all $n \in \N$.
\end{enumerate}

\examsec{Section 3: Topology}

\qitem{3.1} Let $A$ and $B$ be subsets of $\R^n$. Define $A + B = \{x + y \mid x \in A,\ y \in B\}$. Pick out the true statements:
\begin{enumerate}[a.]
  \item if $A$ and $B$ are closed sets, then $A + B$ is a closed set;
  \item if $A$ is an open set and if $B$ is a closed set, then $A + B$ is an open set;
  \item if $A$ and $B$ are compact sets, then so is $A + B$.
\end{enumerate}

\qitem{3.2} Let $X$ and $Y$ be metric spaces and let $f : X \to Y$ be a mapping. Pick out the true statements:
\begin{enumerate}[a.]
  \item if $f$ is uniformly continuous, then the image of every Cauchy sequence in $X$ is a Cauchy sequence in $Y$;
  \item if $X$ is complete and if $f$ is continuous, then the image of every Cauchy sequence in $X$ is a Cauchy sequence in $Y$;
  \item if $Y$ is complete and if $f$ is continuous, then the image of every Cauchy sequence in $X$ is a Cauchy sequence in $Y$;
\end{enumerate}

\qitem{3.3} Which of the following statements are true?
\begin{enumerate}[a.]
  \item If $A$ is a dense subset of a topological space $X$, then $X \setminus A$ is nowhere dense in $X$.
  \item If $A$ is a nowhere dense subset of a topological space $X$, then $X \setminus A$ is dense in $X$.
  \item The set $\R$, identified with the $x$-axis in $\R^2$, is nowhere dense in $\R^2$.
\end{enumerate}

\qitem{3.4} Which of the following metric spaces are separable?
\begin{enumerate}[a.]
  \item The space $\mathcal{C}[0, 1]$, with the usual `sup-norm' metric.
  \item The space $\ell_1$ of all absolutely convergent real sequences, with the metric $d_1(\{a_i\}, \{b_i\}) = \sum_{i=1}^{\infty} |a_i - b_i|$.
  \item The space $\ell_\infty$ of all bounded real sequences, with the metric $d_\infty(\{a_i\}, \{b_i\}) = \sup_{1 \le i < \infty} |a_i - b_i|$.
\end{enumerate}

\qitem{3.5} Which of the following sets are compact?
\begin{enumerate}[a.]
  \item The closed unit ball centred at $0$ and of radius $1$ of $\ell_1$ (see Question 3.4(b) above).
  \item The set of all unitary matrices in $\M_2(\C)$.
  \item The set of all matrices in $\M_2(\C)$ with determinant equal to unity.
\end{enumerate}

\qitem{3.6} Which of the following sets are connected?
\begin{enumerate}[a.]
  \item The set $\{(x, y) \in \R^2 \mid xy = 1\}$ in $\R^2$.
  \item The set of all symmetric matrices in $\M_n(\R)$.
  \item The set of all orthogonal matrices in $\M_n(\R)$.
\end{enumerate}

\qitem{3.7} Which of the following metric spaces are complete?
\begin{enumerate}[a.]
  \item The space of all continuous real-valued functions on $\R$ with compact support, with the usual `sup-norm' metric.
  \item The space $\mathcal{C}[0, 1]$ with the metric $d_1(f, g) = \int_0^1 |f(t) - g(t)|\,dt$.
  \item The space $\mathcal{C}^1[0, 1]$ with the metric $d(f, g) = \max_{t \in [0, 1]} |f(t) - g(t)|$.
\end{enumerate}

\qitem{3.8} Let $X_j = \mathcal{C}[0, 1]$ with the metric $d_1$ (see Question 3.7(b) above) when $j = 1$, the metric $d_2(f, g) = \left(\int_0^1 |f(t) - g(t)|^2\,dt\right)^{\frac{1}{2}}$ when $j = 2$ and the usual `sup-norm' metric when $j = 3$. Let $\mathrm{id} : \mathcal{C}[0, 1] \to \mathcal{C}[0, 1]$ be the identity map. Pick out the true statements:
\begin{enumerate}[a.]
  \item $\mathrm{id} : X_2 \to X_1$ is continuous;
  \item $\mathrm{id} : X_1 \to X_3$ is continuous;
  \item $\mathrm{id} : X_3 \to X_2$ is continuous.
\end{enumerate}

\qitem{3.9} Which of the following statements are true?
\begin{enumerate}[a.]
  \item Consider the subspace $S^1 = \{(x, y) \in \R^2 \mid x^2 + y^2 = 1\}$ of $\R^2$. Then, there exists a continuous function $f : S^1 \to \R$ which is onto.
  \item There exists a continuous function $f : S^1 \to \R$ which is one-one.
  \item Let $X = \{A = (a_{ij}) \in \M_2(\R) \mid \tr(A) = 0 \text{ and } |a_{ij}| \le 2 \text{ for all } 1 \le i, j \le 2\}$. Let $Y = \{\det(A) \mid A \in X\} \subset \R$. Then, there exist $\alpha < 0$ and $\beta > 0$ such that $Y = [\alpha, \beta]$.
\end{enumerate}

\qitem{3.10} Let $S^1 \subset \R^2$ be as in Question 3.9(a) above. Let $D = \{(x, y) \in \R^2 \mid x^2 + y^2 \le 1\}$ and $E = \{(x, y) \in \R^2 \mid 2x^2 + 3y^2 \le 1\}$ be also considered as subspaces of $\R^2$. Which of the following statements are true?
\begin{enumerate}[a.]
  \item If $f : D \to S^1$ is a continuous mapping, then there exists $x \in S^1$ such that $f(x) = x$.
  \item If $f : S^1 \to S^1$ is a continuous mapping, then there exists $x \in S^1$ such that $f(x) = x$.
  \item If $f : E \to E$ is a continuous mapping, then there exists $x \in E$ such that $f(x) = x$.
\end{enumerate}

\examsec{Section 4: Applied Mathematics}

\qitem{4.1} Find the family of orthogonal trajectories of the family of curves $y = cx^2$.

\qitem{4.2} Let $f \in \mathcal{C}[0, 1]$ be given. Consider the problem: find a curve $u$ such that $u(0) = u(1) = 0$ which minimizes the functional
\[
  J(v) = \frac{1}{2}\int_0^1 (v'(x))^2\,dx - \int_0^1 f(x)v(x)\,dx
\]
over all admissible curves $v$. Write down the boundary value problem (Euler-Lagrange equation) satisfied by the solution $u$.

\qitem{4.3} Let $\omega \in \R$ be a constant. Solve:
\[
  \frac{dx(t)}{dt} = x(t) - \omega y(t),\quad \frac{dy(t)}{dt} = \omega x(t) + y(t),\quad t > 0,\qquad x(0) = y(0) = 1.
\]

\qitem{4.4} Let $b \in \R$ be a constant. Solve:
\[
  \frac{\partial u}{\partial t}(x, t) + b\frac{\partial u}{\partial x}(x, t) = 0,\quad x \in \R,\ t > 0,\qquad u(x, 0) = x^2.
\]

\qitem{4.5} Let $v$ be a smooth harmonic function on $\R^n$. If $r^2 = \sum_{i=1}^{n} |x_i|^2$, where $x = (x_1, \dots, x_n) \in \R^n$, and if $v$ is a radial function, \emph{i.e.}\ $v(x) = v(r)$, write down the ordinary differential equation satisfied by $v$.

\qitem{4.6} Let $a$ and $b$ be positive constants. Let $y$ satisfy $y''(t) + ay'(t) + by(t) = 0$, $t > 0$, $y(0) = 1$ and $y'(0) = 1$. Write down the Laplace transform of $y$.

\qitem{4.7} Write down the Newton-Raphson iteration scheme to find $1/\sqrt{a}$, where $a > 0$, by solving the equation $x^{-2} - a = 0$.

\qitem{4.8} Let $B \in \M_n(\R)$ and let $b \in \R^n$ be a given fixed vector. Consider the iteration scheme $x_{n+1} = Bx_n + b$, $x_0$ given. Pick out the true statements:
\begin{enumerate}[a.]
  \item the scheme is always convergent for any initial vector $x_0$.
  \item if the scheme is always convergent for any initial vector $x_0$, then $I - B$ is invertible.
  \item if the scheme is always convergent for any initial vector $x_0$, then every eigenvalue $\lambda$ of $B$ satisfies $|\lambda| < 1$.
\end{enumerate}

\qitem{4.9} Let
\[
  A = \begin{bmatrix} 1 & -2 & 3 & -2 \\ 1 & 1 & 0 & 3 \\ -1 & 1 & 1 & -1 \\ 0 & -3 & 1 & 1 \end{bmatrix}.
\]
Pick out the smallest disc in the complex plane containing all the eigenvalues of $A$ from amongst the following:\\
a.\ $|z - 1| \le 7$; \qquad b.\ $|z - 1| \le 6$; \qquad c.\ $|z - 1| \le 4$.

\qitem{4.10} Solve: maximize $z = 7x + 5y$ such that $x \ge 0$, $y \ge 0$ and $x + 2y \le 6$, $4x + 3y \le 12$.

\examsec{Section 5: Miscellaneous}

\qitem{5.1} Which of the following sets are countable?
\begin{enumerate}[a.]
  \item The set of all sequences of non-negative integers.
  \item The set of all sequences of non-negative integers with only a finite number of non-zero terms.
  \item The set of all roots of all monic polynomials in one variable with rational coefficients.
\end{enumerate}

\qitem{5.2} A magic square of order $N$ is an $N \times N$ matrix with positive integral entries such that the elements of every row, every column and the two diagonals all add up to the same number. If a magic square is filled with numbers in arithmetic progression starting with $a \in \N$ and common difference $d \in \N$, what is the value of this common sum?

\qitem{5.3} A committee consists of $n$ members and a group photograph is to be taken by seating them in a row. If two particular members do not get along with each other, in how many ways can the committee members be seated so that these two are never adjacent to each other?

\qitem{5.4} Let $n \ge 2$ and let $D_n$ be the number of permutations of $\{1, 2, \dots, n\}$ which leave no symbol fixed. (For example: $D_2 = 1$). Write down an expression for $D_n$ in terms of $D_k$, $2 \le k \le n - 1$.

\qitem{5.5} Five letters are addressed to five different persons and the corresponding envelopes are prepared. The letters are put into the envelopes at random. What is the probability that no letter is in its proper envelope?

\qitem{5.6} Which of the following statements are true?
\begin{enumerate}[a.]
  \item The $9$-th power of any positive integer is of the form $19m$ or $19m \pm 1$.
  \item For any positive integer $n$, the number $n^{13} - n$ is divisible by $2730$.
  \item The number $18! + 1$ is divisible by $437$.
\end{enumerate}

\qitem{5.7} Let $a_i$, $1 \le i \le n$ be non-negative real numbers. Let $S$ denote their sum. Pick out the true statements:
\begin{enumerate}[a.]
  \item $\prod_{k=1}^{n} (1 + a_k) \ge 1 + S$;
  \item $\prod_{k=1}^{n} (1 + a_k) \le 1 + \frac{S}{1!} + \frac{S^2}{2!} + \dots + \frac{S^n}{n!}$;
  \item $\prod_{k=1}^{n} (1 + a_k) \le \frac{1}{1 - S}$, if $S < 1$.
\end{enumerate}

\qitem{5.8} Consider the Fibonacci sequence $\{a_n\}$ defined by $a_0 = 0$, $a_1 = 1$, $a_{n+1} = a_n + a_{n-1}$, $n \ge 1$. Write down its enumerating function, \emph{i.e.}\ the function with the formal power series expansion $\sum_{n=0}^{\infty} a_n x^n$.

\qitem{5.9} Find the lengths of the semi-axes of the ellipse whose equation is given by $5x^2 - 8xy + 5y^2 = 1$.

\qitem{5.10} Let $x_0 = a$ and $x_1 = b$. Define $x_{n+1} = \left(1 - \frac{1}{2n}\right)x_n + \frac{1}{2n}x_{n-1}$, $n \ge 1$. Find $\lim_{n \to \infty} x_n$.

\examsec{Answer Key --- NBHM 2012}
\begin{multicols}{2}
\textbf{Section 1: Algebra}\\
\textbf{1.1} a, b, c\\
\textbf{1.2} a, c\\
\textbf{1.3} b\\
\textbf{1.4} $(p - 1)/2$\\
\textbf{1.5} $\begin{bmatrix} 1 & 4 & 3 \\ 0 & 3 & 1 \\ 0 & 0 & 2 \end{bmatrix}$\\
\textbf{1.6} $\begin{bmatrix} 1 & -2 & 2 & 0 \\ 0 & 1 & -4 & 6 \\ 0 & 0 & 1 & -6 \\ 0 & 0 & 0 & 1 \end{bmatrix}$\\
\textbf{1.7} $x^3 + 6x^2 + 18x + 24$\\
\textbf{1.8} (a) $n$; (b) $0$\\
\textbf{1.9} a\\
\textbf{1.10} b, c

\medskip
\textbf{Section 2: Analysis}\\
\textbf{2.1} b, c\\
\textbf{2.2} $2f(0)$\\
\textbf{2.3} $3[f(1/3) - f(0)]$\\
\textbf{2.4} (a) $0$; (b) $\frac{1}{2}\int_{-\pi}^{\pi} f(t)\,dt$\\
\textbf{2.5} (a) Not uniformly convergent; (b) uniformly convergent; (c) not uniformly convergent\\
\textbf{2.6} $\frac{2}{x}\left(e^{x^2} - e^{-x^2}\right)$\\
\textbf{2.7} $e^{\frac{k(k+1)}{2a}}$\\
\textbf{2.8} $\frac{1}{4} + \frac{1}{4}\sum_{n=1}^{\infty} (-1)^n (n+1)\left(\frac{z-2}{2}\right)^n$\\
\textbf{2.9} $0$\\
\textbf{2.10} b

\medskip
\textbf{Section 3: Topology}\\
\textbf{3.1} b, c\\
\textbf{3.2} a, b\\
\textbf{3.3} b, c\\
\textbf{3.4} a, b\\
\textbf{3.5} b\\
\textbf{3.6} b\\
\textbf{3.7} none\\
\textbf{3.8} a, c\\
\textbf{3.9} c\\
\textbf{3.10} a, c

\medskip
\textbf{Section 4: Applied Mathematics}\\
\textbf{4.1} $x^2 + 2y^2 = c^2$\\
\textbf{4.2} $-u'' = f$ on $]0, 1[$; $u(0) = u(1) = 0$\\
\textbf{4.3} $x(t) = e^t(\cos\omega t - \sin\omega t),\ y(t) = e^t(\cos\omega t + \sin\omega t)$\\
\textbf{4.4} $u(x, t) = (x - bt)^2$\\
\textbf{4.5} $v''(r) + \frac{n-1}{r}v'(r) = 0$\\
\textbf{4.6} $L[y](s) = \frac{1 + a + s}{s^2 + as + b}$\\
\textbf{4.7} $x_{n+1} = \frac{1}{2}(3x_n - ax_n^3)$\\
\textbf{4.8} b, c\\
\textbf{4.9} b\\
\textbf{4.10} $\max z = 21$ at $x = 3;\ y = 0$

\medskip
\textbf{Section 5: Miscellaneous}\\
\textbf{5.1} b, c\\
\textbf{5.2} $\frac{N}{2}(2a + (N^2 - 1)d)$\\
\textbf{5.3} $n! - 2(n-1)!$\\
\textbf{5.4} $D_n = n! - 1 - \sum_{k=1}^{n-2} \binom{n}{k} D_{n-k}$\\
\textbf{5.5} $44/120 = 11/30$\\
\textbf{5.6} a, b, c\\
\textbf{5.7} a, b, c\\
\textbf{5.8} $\frac{x}{1 - x - x^2}$\\
\textbf{5.9} semi-major axis $= 1$; semi-minor axis $= 1/3$\\
\textbf{5.10} $a + (b - a)e^{-\frac{1}{2}}$

\medskip
\noindent\emph{Note: Accept any correct equivalent form of the answers.}
\end{multicols}
% ==================== NBHM 2013 ====================
\examheader{NBHM 2013}{Research Scholarships Screening Test}{Saturday, January 19, 2013}{Time Allowed: 150 Minutes \quad|\quad Maximum Marks: 40}

\subsection*{Instructions to Candidates \& Notation}
\begin{itemize}
  \item Booklet contains 11 numbered pages; 5 sections of 10 questions each (Algebra, Analysis, Topology, Applied Mathematics, Miscellaneous); assessment based on best 4 sections (max 40 marks); no partial credit; calculators not allowed.
  \item \textbf{Notation:} Standard $\N, \Z, \Q, \R, \C, \R^n, \C^n, \M_n(\R), \M_n(\C)$; open/closed intervals $]a, b[, [a, b]$; $\mathcal{C}[a, b], \mathcal{C}^1[a, b]$; $f', f''$; transpose $x^T, A^T$; identity $I$; $\det(A), \tr(A)$; $GL_n(\R), GL_n(\C)$; all logarithms are to the base $e$.
\end{itemize}

\examsec{Section 1: Algebra}

\qitem{1.1} Find the number of elements of order two in the symmetric group $S_4$ of all permutations of the four symbols $\{1, 2, 3, 4\}$.

\qitem{1.2} Let $G$ be the group of all invertible $2 \times 2$ upper triangular matrices (under matrix multiplication). Pick out the normal subgroups of $G$ from the following:
\begin{enumerate}[a.]
  \item $H = \{A \in G : a_{12} = 0\}$;
  \item $H = \{A \in G : a_{11} = 1\}$;
  \item $H = \{A \in G : a_{11} = a_{22}\}$, where $A = \begin{bmatrix} a_{11} & a_{12} \\ 0 & a_{22} \end{bmatrix}$.
\end{enumerate}

\qitem{1.3} Let $G = GL_n(\R)$ and let $H$ be the (normal) subgroup of all matrices with positive determinant. Identify the quotient group $G/H$.

\qitem{1.4} Which of the following rings are integral domains?
\begin{enumerate}[a.]
  \item $\R[x]$, the ring of all polynomials in one variable with real coefficients.
  \item $\M_n(\R)$.
  \item The ring of complex analytic functions defined on the unit disc of the complex plane (with pointwise addition and multiplication as the ring operations).
\end{enumerate}

\qitem{1.5} Find the condition on the real numbers $a, b$ and $c$ such that the following system of equations has a solution:
\begin{align*}
  2x + y + 3z &= a \\
  x + z &= b \\
  y + z &= c.
\end{align*}

\qitem{1.6} Let $\mathcal{P}_n$ denote the vector space of all polynomials in one variable with real coefficients and of degree less than, or equal to, $n$, equipped with the standard basis $\{1, x, x^2, \dots, x^n\}$. Define $T : \mathcal{P}_2 \to \mathcal{P}_3$ by
\[
  T(p)(x) = \int_0^x p(t)\,dt + p'(x) + p(2).
\]
Write down the matrix of this transformation with respect to the standard bases of $\mathcal{P}_2$ and $\mathcal{P}_3$.

\qitem{1.7} Determine the dimension of the kernel of the linear transformation $T$ defined in Question 1.6 above.

\qitem{1.8} A symmetric matrix in $\M_n(\R)$ is said to be \emph{non-negative definite} if $x^T A x \ge 0$ for all (column) vectors $x \in \R^n$. Which of the following statements are true?
\begin{enumerate}[a.]
  \item If a real symmetric $n \times n$ matrix is non-negative definite, then all of its eigenvalues are non-negative.
  \item If a real symmetric $n \times n$ matrix has all its eigenvalues non-negative, then it is non-negative definite.
  \item If $A \in \M_n(\R)$, then $AA^T$ is non-negative definite.
\end{enumerate}

\qitem{1.9} Only one of the following matrices is non-negative definite. Find it.\\
a.\ $\begin{bmatrix} 5 & -3 \\ -3 & 5 \end{bmatrix}$ \qquad
b.\ $\begin{bmatrix} 1 & -3 \\ -3 & 5 \end{bmatrix}$ \qquad
c.\ $\begin{bmatrix} 1 & 3 \\ 3 & 5 \end{bmatrix}$.

\qitem{1.10} Let $B$ be the real symmetric non-negative definite $2 \times 2$ matrix such that $B^2 = A$ where $A$ is the non-negative definite matrix in Question 1.9 above. Write down the characteristic polynomial of $B$.

\examsec{Section 2: Analysis}

\qitem{2.1} Evaluate: $\displaystyle\lim_{n \to \infty} \sin\left(\left(2n\pi + \frac{1}{2n\pi}\right)\sin\left(2n\pi + \frac{1}{2n\pi}\right)\right)$.

\qitem{2.2} Evaluate: $\displaystyle\lim_{n \to \infty} \frac{1}{n}\left[(n+1)(n+2)\cdots(n+n)\right]^{\frac{1}{n}}$.

\qitem{2.3} Which of the following series are convergent?
\begin{enumerate}[a.]
  \item $\displaystyle\sum_{n=1}^{\infty} \frac{\frac{1}{2} + (-1)^n}{n}$.
  \item $\displaystyle\sum_{n=1}^{\infty} (-1)^n(\sqrt{n+1} - \sqrt{n})$.
  \item $\displaystyle\sum_{n=1}^{\infty} \frac{\sin(n^{\frac{3}{2}})}{n^{\frac{3}{2}}}$.
\end{enumerate}

\qitem{2.4} Which of the following functions are uniformly continuous?
\begin{enumerate}[a.]
  \item $f(x) = x\sin\frac{1}{x}$ on $]0, 1[$.
  \item $f(x) = \sin^2 x$ on $]0, \infty[$.
  \item $f(x) = \sin(x\sin x)$ on $]0, \infty[$.
\end{enumerate}

\qitem{2.5} Find the points where the following function is differentiable:
\[
  f(x) = \begin{cases} \tan^{-1} x, & \text{if } |x| \le 1, \\ \frac{\pi x}{4|x|} + \frac{|x|-1}{2}, & \text{if } |x| > 1. \end{cases}
\]

\qitem{2.6} Which of the following sequences/series of functions are uniformly convergent on $[0, 1]$?
\begin{enumerate}[a.]
  \item $f_n(x) = (\cos(\pi n! x))^{2n}$.
  \item $\displaystyle\sum_{m=1}^{\infty} \frac{\cos(m^6 x)}{m^3}$.
  \item $f_n(x) = n^2 x(1 - x^2)^n$.
\end{enumerate}

\qitem{2.7} Let $f \in \mathcal{C}^1[0, 1]$. For a partition $(\mathcal{P}) : 0 = x_0 < x_1 < x_2 < \dots < x_n = 1$, define $S(\mathcal{P}) = \sum_{i=1}^{n} |f(x_i) - f(x_{i-1})|$. Compute the supremum of $S(\mathcal{P})$ taken over all possible partitions $\mathcal{P}$.

\qitem{2.8} Write down the Taylor series expansion about the origin in the region $\{|x| < 1\}$ for the function
\[
  f(x) = x\tan^{-1}(x) - \frac{1}{2}\log(1 + x^2).
\]

\qitem{2.9} Write down all possible values of $i^{-2i}$.

\qitem{2.10} What is the image of the set $\{z \in \C : z = x + iy,\ x \ge 0,\ y \ge 0\}$ under the mapping $z \mapsto z^2$?

\examsec{Section 3: Topology}

\qitem{3.1} Let $(X, d)$ be a metric space. For subsets $A$ and $B$ of $X$, define $d(A, B) = \inf\{d(a, b) : a \in A,\ b \in B\}$. Which of the following statements are true?
\begin{enumerate}[a.]
  \item If $\overline{A} \cap \overline{B} = \emptyset$, then $d(A, B) > 0$.
  \item If $d(A, B) > 0$, then there exist open sets $U$ and $V$ such that $A \subset U$, $B \subset V$, $U \cap V = \emptyset$.
  \item $d(A, B) = 0$ if, and only if, there exists a sequence of points $\{x_n\}$ in $A$ converging to a point in $B$.
\end{enumerate}

\qitem{3.2} Let $X$ be a set and let $(Y, \tau)$ be a topological space. Let $g : X \to Y$ be a given map. Define $\tau' = \{U \subset X : U = g^{-1}(V) \text{ for some } V \in \tau\}$. Which of the following statements are true?
\begin{enumerate}[a.]
  \item $\tau'$ defines a topology on $X$.
  \item $\tau'$ defines a topology on $X$ only if $g$ is onto.
  \item Let $g$ be onto. Define the equivalence relation $x \sim y$ if, and only if, $g(x) = g(y)$. Then the quotient space of $X$ with respect to this relation, with the topology inherited from $\tau'$, is homeomorphic to $(Y, \tau)$.
\end{enumerate}

\qitem{3.3} Find pairs of homeomorphic sets from the following:\\
$A = \{(x, y) \in \R^2 : xy = 0\}$; \quad $B = \{(x, y) \in \R^2 : x + y \ge 0,\ xy = 0\}$;\\
$C = \{(x, y) \in \R^2 : xy = 1\}$; \quad $D = \{(x, y) \in \R^2 : x + y \ge 0,\ xy = 1\}$.

\qitem{3.4} Let $(X, \tau)$ be a topological space. A map $f : X \to \R$ is said to be \emph{lower semi-continuous} if for every $\alpha \in \R$, the set $f^{-1}(]-\infty, \alpha])$ is closed in $X$. It is said to be \emph{upper semi-continuous} if, for every $\alpha \in \R$, the set $f^{-1}([\alpha, \infty[)$ is closed in $X$. Which of the following statements are true?
\begin{enumerate}[a.]
  \item If $\{f_n\}$ is a sequence of lower semi-continuous real valued functions on $X$, then $f = \sup_n f_n$ is also lower semi-continuous.
  \item Every continuous real valued function on $X$ is lower semi-continuous.
  \item If a real valued function is both upper and lower semi-continuous, then it is continuous.
\end{enumerate}

\qitem{3.5} Let $S = \{A \in \M_n(\R) : \tr(A) = 0\}$. Which of the following statements are true?
\begin{enumerate}[a.]
  \item $S$ is nowhere dense in $\M_n(\R)$.
  \item $S$ is connected in $\M_n(\R)$.
  \item $S$ is compact in $\M_n(\R)$.
\end{enumerate}

\qitem{3.6} Let $S$ be the set of all symmetric non-negative definite matrices (see Question 1.8) in $\M_n(\R)$. Which of the following statements are true?
\begin{enumerate}[a.]
  \item $S$ is closed in $\M_n(\R)$.
  \item $S$ is connected in $\M_n(\R)$.
  \item $S$ is compact in $\M_n(\R)$.
\end{enumerate}

\qitem{3.7} Which of the following sets are compact in $\M_n(\R)$?
\begin{enumerate}[a.]
  \item The set of all upper triangular matrices all of whose eigenvalues satisfy $|\lambda| \le 2$.
  \item The set of all real symmetric matrices all of whose eigenvalues satisfy $|\lambda| \le 2$.
  \item The set of all diagonalizable matrices all of whose eigenvalues satisfy $|\lambda| \le 2$.
\end{enumerate}

\qitem{3.8} Let $X$ be the set of all real sequences. Consider the subset
\[
  S = \left\{x = (x_n) \in X : \begin{array}{l}x_n \in \Q \text{ for all } n,\\ x_n = 0, \text{ except for a finite number of } n\end{array}\right\}.
\]
Which of the following statements are true?
\begin{enumerate}[a.]
  \item $S$ is dense in $\ell_1$, the space of absolutely summable sequences, provided with the metric $d_1(x, y) = \sum_{n=1}^{\infty} |x_n - y_n|$.
  \item $S$ is dense in $\ell_2$, the space of square summable sequences, provided with the metric $d_2(x, y) = \left(\sum_{n=1}^{\infty} |x_n - y_n|^2\right)^{\frac{1}{2}}$.
  \item $S$ is dense in $\ell_\infty$, the space of bounded sequences, provided with the metric $d_\infty(x, y) = \sup_n \{|x_n - y_n|\}$.
\end{enumerate}

\qitem{3.9} Which of the following statements are true?
\begin{enumerate}[a.]
  \item There exists a continuous function $f : \{(x, y) \in \R^2 : 2x^2 + 3y^2 = 1\} \to \R$ which is one-one.
  \item There exists a continuous function $f : ]-1, 1[ \to ]-1, 1]$ which is one-one and onto.
  \item There exists a continuous function $f : \{(x, y) \in \R^2 : y^2 = 4x\} \to \R$ which is one-one.
\end{enumerate}

\qitem{3.10} Which of the following statements are true?
\begin{enumerate}[a.]
  \item Let $f : ]0, \infty[ \to ]0, \infty[$ be such that $|f(x) - f(y)| \le \frac{1}{2}|x - y|$ for all $x$ and $y$. Then $f$ has a fixed point.
  \item Let $f : [-1, 1] \to [-1, 1]$ be continuous. Then $f$ has a fixed point.
  \item Let $f : \R \to \R$ be continuous and periodic with period $T > 0$. Then there exists a point $x_0 \in \R$ such that $f(x_0) = f\left(x_0 + \frac{T}{2}\right)$.
\end{enumerate}

\examsec{Section 4: Applied Mathematics}

\qitem{4.1} Find all the solutions $(\lambda, u)$, $u \not\equiv 0$, of the problem:
\[
  u'' + \lambda u = 0 \text{ in } ]0, 1[,\qquad u(0) = 0 = u'(1).
\]

\qitem{4.2} Find the constant $c$ such that the following problem has a solution:
\[
  -u'' = c \text{ in } ]a, b[,\qquad u'(a) = -1,\quad u'(b) = 1.
\]

\qitem{4.3} Evaluate: $\displaystyle\int_{-\infty}^{\infty}\int_{-\infty}^{\infty} e^{-(3x^2 + 2\sqrt{2}xy + 3y^2)}\,dxdy$.

\qitem{4.4} Find the stationary function $y = y(x)$ of the integral $\int_0^4 [xy' - (y')^2]\,dx$ satisfying the conditions $y(0) = 0$ and $y(4) = 3$.

\qitem{4.5} Let $L(y)$ denote the Laplace transform of a function $y = y(x)$. If $y$ and $y'$ are bounded, express $L(y'')$ in terms of $L(y)$, $y$ and $y'$.

\qitem{4.6} Find the singular points of the differential equation $x^3(x - 1)y'' - 2(x - 1)y' + 3xy = 0$ and state whether they are regular singular points or irregular singular points.

\qitem{4.7} Let $(\lambda_1, y_1)$ and $(\lambda_2, y_2)$ be two solutions of the problem
\[
  (p(x)y'(x))' + \lambda q(x)y(x) = 0 \text{ in } ]a, b[,\qquad y(a) = 0 = y(b)
\]
where $p$ and $q$ are positive and continuous functions on $]a, b[$. If $\lambda_1 \neq \lambda_2$, evaluate $\int_a^b q(x)y_1(x)y_2(x)\,dx$.

\qitem{4.8} Solve: $xy'' - y' = 3x^2$.

\qitem{4.9} Let $f \in \mathcal{C}[a, b]$. Write down Simpson's rule to approximate $\int_a^b f(x)\,dx$ using the points $x = a$, $x = (a + b)/2$ and $x = b$.

\qitem{4.10} What is the highest value of $n$ such that Simpson's rule (see Question 4.9 above) gives the exact value of the integral of $f$ on $[a, b]$ when $f$ is a polynomial of degree less than, or equal to, $n$?

\examsec{Section 5: Miscellaneous}

\qitem{5.1} Let $m > n$. In how many ways can we seat $m$ men and $n$ women in a row for a photograph if no two women are to be seated adjacent to each other?

\qitem{5.2} Let $n \in \N$ be fixed. For $r \le n$, let $C_r$ denote the usual binomial coefficient $\binom{n}{r}$ which gives the number of ways of choosing $r$ objects from a given set of $n$ objects. Evaluate:
\[
  C_0 + 4C_1 + 7C_2 + \dots + (3n + 1)C_n.
\]

\qitem{5.3} Let $A =$ the set of all sequences of real numbers, $B =$ the set of all sequences of positive real numbers, $C = \mathcal{C}[0, 1]$ and $D = \R$. Which of the following statements are true?
\begin{enumerate}[a.]
  \item All the four sets have the same cardinality.
  \item $A$ and $B$ have the same cardinality.
  \item $A, B$ and $D$ have the same cardinality, which is different from that of $C$.
\end{enumerate}

\qitem{5.4} For a positive integer $n$, define
\[
  \Lambda(n) = \begin{cases} \log p, & \text{if } n = p^r,\ p \text{ a prime and } r \in \N, \\ 0, & \text{otherwise.} \end{cases}
\]
Given a positive integer $N$, evaluate: $\displaystyle\sum_{d \mid N} \Lambda(d)$, where the sum ranges over all divisors $d$ of $N$.

\qitem{5.5} Let $a, b$ and $c$ be real numbers. Evaluate:
\[
  \begin{vmatrix}
    b^2 c^2 & bc & b + c \\
    c^2 a^2 & ca & c + a \\
    a^2 b^2 & ab & a + b
  \end{vmatrix}.
\]

\qitem{5.6} Write down the equation (with leading coefficient equal to unity) whose roots are the squares of the roots of the equation $x^3 - 6x^2 + 10x - 3 = 0$.

\qitem{5.7} Let $A = (0, 1)$ and $B = (1, 1)$ in the plane $\R^2$. Determine the length of the shortest path from $A$ to $B$ consisting of the line segments $AP$, $PQ$ and $QB$, where $P$ varies on the $x$-axis between the points $(0, 0)$ and $(1, 0)$ and $Q$ varies on the line $\{y = 3\}$ between the points $(0, 3)$ and $(1, 3)$.

\qitem{5.8} Let $x_0 = a$, $x_1 = b$. If $x_{n+2} = \frac{1}{3}(x_n + 2x_{n+1})$, $n \ge 0$, find $\lim_{n \to \infty} x_n$.

\qitem{5.9} Which of the following statements are true?
\begin{enumerate}[a.]
  \item If $a, b$ and $c$ are the sides of a triangle, then $\dfrac{ab + bc + ca}{a^2 + b^2 + c^2} \ge \dfrac{1}{2}$.
  \item If $a, b$ and $c$ are the sides of a triangle, then $\dfrac{ab + bc + ca}{a^2 + b^2 + c^2} \le 1$.
  \item Both statements above are true for all triples $(a, b, c)$ of strictly positive real numbers.
\end{enumerate}

\qitem{5.10} Let $f \in \mathcal{C}[a, b]$. Assume that $\min_{x \in [a, b]} f(x) = m > 0$ and let $M = \max_{x \in [a, b]} f(x)$. Which of the following inequalities are true?
\begin{enumerate}[a.]
  \item $\dfrac{1}{M}\displaystyle\int_a^b f(x)\,dx + m\int_a^b \frac{1}{f(x)}\,dx \ge 2\sqrt{\frac{m}{M}}(b - a)$.
  \item $\displaystyle\int_a^b f(x)\,dx \int_a^b \frac{1}{f(x)}\,dx \ge (b - a)^2$.
  \item $\displaystyle\int_a^b f(x)\,dx \int_a^b \frac{1}{f(x)}\,dx \le (b - a)^2$.
\end{enumerate}

\examsec{Answer Key --- NBHM 2013}
\begin{multicols}{2}
\textbf{Section 1: Algebra}\\
\textbf{1.1} $9$\\
\textbf{1.2} b, c\\
\textbf{1.3} The multiplicative group $\{-1, 1\}$\\
\textbf{1.4} a, c\\
\textbf{1.5} $a - 2b - c = 0$\\
\textbf{1.6} $\begin{bmatrix} 1 & 3 & 4 \\ 1 & 0 & 2 \\ 0 & \frac{1}{2} & 0 \\ 0 & 0 & \frac{1}{3} \end{bmatrix}$\\
\textbf{1.7} $0$\\
\textbf{1.8} a, b, c\\
\textbf{1.9} a\\
\textbf{1.10} $\lambda^2 - 3\sqrt{2}\lambda + 4$

\medskip
\textbf{Section 2: Analysis}\\
\textbf{2.1} $\sin 1$\\
\textbf{2.2} $\frac{4}{e}$\\
\textbf{2.3} b, c\\
\textbf{2.4} a, b\\
\textbf{2.5} $\R \setminus \{-1\}$\\
\textbf{2.6} b\\
\textbf{2.7} $\int_0^1 |f'(t)|\,dt$\\
\textbf{2.8} $\sum_{n=1}^{\infty} (-1)^{n-1}\frac{x^{2n}}{(2n-1)(2n)}$\\
\textbf{2.9} $e^{(4n+1)\pi},\ n \in \Z$\\
\textbf{2.10} $\{z = x + iy : y \ge 0\}$

\medskip
\textbf{Section 3: Topology}\\
\textbf{3.1} b\\
\textbf{3.2} a, c\\
\textbf{3.3} $B$ and $D$ are homeomorphic\\
\textbf{3.4} a, b, c\\
\textbf{3.5} a, b\\
\textbf{3.6} a, b\\
\textbf{3.7} b\\
\textbf{3.8} a, b\\
\textbf{3.9} c\\
\textbf{3.10} b, c

\medskip
\textbf{Section 4: Applied Mathematics}\\
\textbf{4.1} $\lambda = (2n+1)^2\frac{\pi^2}{4}$, $u = C\sin(2n+1)\frac{\pi}{2}x$, $n = 0, 1, 2, \dots$\\
\textbf{4.2} $c = \frac{-2}{b - a}$\\
\textbf{4.3} $\frac{\pi}{\sqrt{7}}$\\
\textbf{4.4} $y(x) = \frac{x^2 - x}{4}$\\
\textbf{4.5} $L(y'')(s) = s^2 L(y)(s) - sy(0) - y'(0)$\\
\textbf{4.6} $x = 0$, irregular singular point; $x = 1$, regular singular point\\
\textbf{4.7} $0$\\
\textbf{4.8} $y(x) = x^3 + C_1 x^2 + C_2$\\
\textbf{4.9} $\int_a^b f(x)\,dx \sim \frac{(b-a)}{6}[f(a) + 4f((a+b)/2) + f(b)]$\\
\textbf{4.10} $3$

\medskip
\textbf{Section 5: Miscellaneous}\\
\textbf{5.1} $\frac{m!(m+1)!}{(m-n+1)!}$\\
\textbf{5.2} $2^{n-1}(3n+2)$\\
\textbf{5.3} a, b\\
\textbf{5.4} $\log N$\\
\textbf{5.5} $0$\\
\textbf{5.6} $x^3 - 16x^2 + 64x - 9 = 0$\\
\textbf{5.7} $\sqrt{37}$\\
\textbf{5.8} $\frac{a + 3b}{4}$\\
\textbf{5.9} a, b\\
\textbf{5.10} a, b
\end{multicols}
% ==================== NBHM 2014 ====================
\examheader{NBHM 2014}{Research Scholarships Screening Test}{Saturday, January 25, 2014}{Time Allowed: 150 Minutes \quad|\quad Maximum Marks: 40}

\subsection*{Instructions to Candidates \& Notation}
\begin{itemize}
  \item Booklet contains 12 numbered pages; 5 sections of 10 questions each (Algebra, Analysis, Topology, Calculus \& Differential Equations, Miscellaneous); assessment based on best 4 sections (max 40 marks); no partial credit; calculators not allowed.
  \item \textbf{Notation:} $\N = \{1, 2, 3, \dots\}$; $\Z, \Q, \R, \C$; $\R^n, \C^n, \M_n(\R), \M_n(\C)$; $]a, b[, [a, b]$; $\mathcal{C}[a, b], \mathcal{C}^1[a, b]$; $f', f''$; transpose $x^T, A^T$ and adjoint $x^*, A^*$; identity $I$; $\det(A), \tr(A)$; $\ker(\varphi), \ker(A), \mathcal{R}(A)$; $GL_n(\R), GL_n(\C)$; $S_n$; all logarithms are to the base $e$.
\end{itemize}

\examsec{Section 1: Algebra}

\qitem{1.1} Let $G$ be a finite group of order $n \ge 2$. Which of the following statements are true?
\begin{enumerate}[a.]
  \item There always exists an injective homomorphism from $G$ into $S_n$.
  \item There always exists an injective homomorphism from $G$ into $S_m$ for some $m < n$.
  \item There always exists an injective homomorphism from $G$ into $GL_n(\R)$.
\end{enumerate}

\qitem{1.2} Let $\C^*$ denote the multiplicative group of non-zero complex numbers and let $P$ denote the subgroup of positive (real) numbers. Identify the quotient group $\C^*/P$.

\qitem{1.3} Given a finite group and a prime $p$ which divides its order, let $N(p)$ denote the number of $p$-Sylow subgroups of $G$. If $G$ is a group of order $21$, what are the possible values for $N(3)$ and $N(7)$?

\qitem{1.4} Let $V$ be the real vector space of all polynomials, in a single variable and with real coefficients, of degree at most $3$. Let $V^*$ be its dual space. Let $t_1 = 1, t_2 = 2, t_3 = 3, t_4 = 4$. Which of the following sets of functionals $\{f_i,\ 1 \le i \le 4\}$ form a basis for $V^*$?
\begin{enumerate}[a.]
  \item For $1 \le i \le 4$, and for all $p \in V$, $f_i(p) = p(t_i)$.
  \item For all $p \in V$, $f_i(p) = p(t_i)$ for $i = 1, 2$, $f_3(p) = p'(t_1)$ and $f_4(p) = p'(t_2)$.
  \item For all $p \in V$, $f_i(p) = p(t_i)$ for $1 \le i \le 3$ and $f_4(p) = \int_1^2 p'(t)\,dt$.
\end{enumerate}

\qitem{1.5} Let $V$ be a finite dimensional real vector space and let $f$ and $g$ be non-zero linear functionals on $V$. Assume that $\ker(f) \subset \ker(g)$. Which of the following statements are true?
\begin{enumerate}[a.]
  \item $\ker(f) = \ker(g)$.
  \item $f = \lambda g$ for some real number $\lambda \neq 0$.
  \item The linear map $A : V \to \R^2$ defined by $Ax = (f(x), g(x))$, for all $x \in V$, is onto.
\end{enumerate}

\qitem{1.6} Let $V$ be a finite dimensional real vector space and let $A : V \to V$ be a linear map such that $A^2 = A$. Assume that $A \neq 0$ and that $A \neq I$. Which of the following statements are true?
\begin{enumerate}[a.]
  \item $\ker(A) \neq \{0\}$.
  \item $V = \ker(A) \oplus \mathcal{R}(A)$.
  \item The map $I + A$ is invertible.
\end{enumerate}

\qitem{1.7} Let $A \in \M_2(\R)$ be a matrix which is not a diagonal matrix. Which of the following statements are true?
\begin{enumerate}[a.]
  \item If $\tr(A) = -1$ and $\det(A) = 1$, then $A^3 = I$.
  \item If $A^3 = I$, then $\tr(A) = -1$ and $\det(A) = 1$.
  \item If $A^3 = I$, then $A$ is diagonalizable over $\R$.
\end{enumerate}

\qitem{1.8} Let $x \in \R^n$ be a non-zero (column) vector. Define $A = xx^T \in \M_n(\R)$.
\begin{enumerate}[a.]
  \item What is the rank of $A$?
  \item What is the necessary and sufficient condition for $I - 2A$ to be an orthogonal matrix?
\end{enumerate}

\qitem{1.9} Let $A \in GL_n(\R)$ have integer entries. Let $b \in \R^n$ be a (column) vector, also with integer entries. Which of the following statements are true?
\begin{enumerate}[a.]
  \item If $Ax = b$, then the entries of $x$ are also integers.
  \item If $Ax = b$, then the entries of $x$ are rational.
  \item The matrix $A^{-1}$ has integer entries if, and only if, $\det(A) = \pm 1$.
\end{enumerate}

\qitem{1.10} In each of the following cases, describe the smallest subset of $\C$ which contains all the eigenvalues of every member of the set $S$.
\begin{enumerate}[a.]
  \item $S = \{A \in \M_n(\C) \mid A = BB^* \text{ for some } B \in \M_n(\C)\}$.
  \item $S = \{A \in \M_n(\C) \mid A = B + B^* \text{ for some } B \in \M_n(\C)\}$.
  \item $S = \{A \in \M_n(\C) \mid A + A^* = 0\}$.
\end{enumerate}

\examsec{Section 2: Analysis}

\qitem{2.1} Find the largest interval for which the following series is convergent at all points $x$ in it:
\[
  \sum_{n=1}^{\infty} \frac{2^n(3x - 1)^n}{n}.
\]

\qitem{2.2} Let $m$ and $k$ be fixed positive integers. Evaluate:
\[
  \lim_{n \to \infty} \left(\frac{(n+1)^m + (n+2)^m + \dots + (n+k)^m}{n^{m-1}} - kn\right).
\]

\qitem{2.3} Which of the following statements are true?
\begin{enumerate}[a.]
  \item If $f$ is twice continuously differentiable in $]a, b[$ and if for all $x \in ]a, b[$, $f''(x) + 2f'(x) + 3f(x) = 0$, then $f$ is infinitely differentiable in $]a, b[$.
  \item Let $f \in \mathcal{C}[a, b]$ be differentiable in $]a, b[$. If $f(a) = f(b) = 0$, then, for any real number $\alpha$, there exists $x \in ]a, b[$ such that $f'(x) + \alpha f(x) = 0$.
  \item The function defined below is not differentiable at $x = 0$:
  \[
    f(x) = \begin{cases} x^2 \left|\cos\frac{\pi}{x}\right|, & x \neq 0, \\ 0, & x = 0. \end{cases}
  \]
\end{enumerate}

\qitem{2.4} Let $f : \R \to \R$ be continuous. Which of the following statements are true?
\begin{enumerate}[a.]
  \item If $f$ is bounded, then $f$ is uniformly continuous.
  \item If $f$ is differentiable and if $f'$ is bounded, then $f$ is uniformly continuous.
  \item If $\lim_{|x| \to \infty} f(x) = 0$, then $f$ is uniformly continuous.
\end{enumerate}

\qitem{2.5} In which of the following cases, is the function $f$ of bounded variation on $[0, 1]$?
\begin{enumerate}[a.]
  \item The function $f : [0, 1] \to \R$ such that, for all $x, y \in [0, 1]$, $|f(x) - f(y)| \le 3|x - y|$.
  \item The function $f$ is monotonically decreasing on $[0, 1]$.
  \item If for some non-negative Riemann integrable function $g$ on $[0, 1]$, $f(x) = \int_0^x g(t)\,dt$ for all $x \in [0, 1]$.
\end{enumerate}

\qitem{2.6} Let $g_n(x) = n\left[f\left(x + \frac{1}{n}\right) - f(x)\right]$, where $f : \R \to \R$ is a continuous function. Which of the following statements are true?
\begin{enumerate}[a.]
  \item If $f(x) = x^3$, then $g_n \to f'$ uniformly on $\R$ as $n \to \infty$.
  \item If $f(x) = x^2$, then $g_n \to f'$ uniformly on $\R$ as $n \to \infty$.
  \item If $f$ is differentiable and if $f'$ is uniformly continuous on $\R$, then $g_n \to f'$ uniformly on $\R$ as $n \to \infty$.
\end{enumerate}

\qitem{2.7} Which of the following statements are true?
\begin{enumerate}[a.]
  \item The series $\displaystyle\sum_{n=1}^{\infty} \frac{x^2}{1 + n^2 x^2}$ does not converge uniformly on $\R$.
  \item The series in (a) above converges uniformly on $\R$.
  \item The sum of the series $\displaystyle\sum_{n=1}^{\infty} \frac{\sin nx^2}{1 + n^3}$ defines a continuously differentiable function on $\R$.
\end{enumerate}

\qitem{2.8} Find the sum of the series: $\displaystyle\sum_{n=1}^{\infty} \frac{(-1)^{n+1}}{n(n+1)}$.

\qitem{2.9} Let $\{f_n\}$ be a sequence of bounded real valued functions on $[0, 1]$ converging to $f$ at all points of this interval. Which of the following statements are true?
\begin{enumerate}[a.]
  \item If $f_n$ and $f$ are all continuous, then $\lim_{n \to \infty} \int_0^1 f_n(t)\,dt = \int_0^1 f(t)\,dt$.
  \item If $f_n \to f$ uniformly, as $n \to \infty$, on $[0, 1]$, then $\lim_{n \to \infty} \int_0^1 f_n(t)\,dt = \int_0^1 f(t)\,dt$.
  \item If $\int_0^1 |f_n(t) - f(t)|\,dt \to 0$ as $n \to \infty$, then $\lim_{n \to \infty} \int_0^1 f_n(t)\,dt = \int_0^1 f(t)\,dt$.
\end{enumerate}

\qitem{2.10} Let $f : [0, \pi] \to \R$ be a continuous function such that $f(0) = 0$. Which of the following statements are true?
\begin{enumerate}[a.]
  \item If $\int_0^{\pi} f(t)\cos nt\,dt = 0$ for all $n \in \{0\} \cup \N$, then $f \equiv 0$.
  \item If $\int_0^{\pi} f(t)\sin nt\,dt = 0$ for all $n \in \N$, then $f \equiv 0$.
  \item If $\int_0^{\pi} t^n f(t)\,dt = 0$ for all $n \in \{0\} \cup \N$, then $f \equiv 0$.
\end{enumerate}

\examsec{Section 3: Topology}

\qitem{3.1} Let $A$ and $B$ be subsets of $\R^n$. Define $A + B = \{a + b \mid a \in A,\ b \in B\}$. Consider the sets
\begin{align*}
  W &= \{(x, y) \in \R^2 \mid x > 0,\ y > 0\}, & X &= \{(x, y) \in \R^2 \mid x \in \R,\ y = 0\}, \\
  Y &= \{(x, y) \in \R^2 \mid xy = 1\}, & Z &= \{(x, y) \in \R^2 \mid |x| \le 1,\ |y| \le 1\}.
\end{align*}
Which of the following statements are true?\\
a.\ The set $W + X$ is open. \qquad b.\ The set $X + Y$ is closed. \qquad c.\ The set $Y + Z$ is closed.

\qitem{3.2} Let $X$ be a topological space and let $A$ be a subset of $X$. Which of the following statements are true?
\begin{enumerate}[a.]
  \item If $A$ is dense in $X$, then $A^\circ$ (the interior of $A$), is also dense in $X$.
  \item If $A$ is dense in $X$, then $X \setminus A$ is nowhere dense.
  \item If $A$ is nowhere dense, then $X \setminus A$ is dense.
\end{enumerate}

\qitem{3.3} Consider the space $X = \mathcal{C}[0, 1]$ with its usual `sup-norm' topology. Let $S = \left\{f \in X \;\middle|\; \int_0^1 f(t)\,dt \neq 0\right\}$. Which of the following statements are true?\\
a.\ The set $S$ is open. \qquad b.\ The set $S$ is dense in $X$. \qquad c.\ The set $S$ is connected.

\qitem{3.4} Consider the space $X = \mathcal{C}[0, 1]$ with its usual `sup-norm' topology. Let $S = \left\{f \in X \;\middle|\; \int_0^1 f(t)\,dt = 0\right\}$. Which of the following statements are true?\\
a.\ The set $S$ is closed. \qquad b.\ The set $S$ is connected. \qquad c.\ The set $S$ is compact.

\qitem{3.5} Let $(X, d)$ be a metric space. Which of the following statements are true?
\begin{enumerate}[a.]
  \item A sequence $\{x_n\}$ converges to $x$ in $X$ if, and only if, the sequence $\{y_n\}$ is a Cauchy sequence in $X$, where, for $k \ge 1$, $y_{2k-1} = x_k$ and $y_{2k} = x$.
  \item If $f : X \to X$ maps Cauchy sequences into Cauchy sequences, then $f$ is continuous.
  \item If $f : X \to X$ is continuous, then it maps Cauchy sequences into Cauchy sequences.
\end{enumerate}

\qitem{3.6} Which of the following spaces are separable?
\begin{enumerate}[a.]
  \item The space $\mathcal{C}[a, b]$ with its usual `sup-norm' topology.
  \item The space $\mathcal{C}[0, 1]$ with the metric defined by $d(f, g) = \int_0^1 |f(t) - g(t)|\,dt$.
  \item The space $\ell_\infty$ consisting of all bounded real sequences with the metric $d(\{x_n\}, \{y_n\}) = \sup_{n \in \N} |x_n - y_n|$.
\end{enumerate}

\qitem{3.7} Consider the space $\M_2(\R)$ with its usual topology. Which of the following sets are dense?
\begin{enumerate}[a.]
  \item The set of all invertible matrices.
  \item The set of all matrices with both eigenvalues real.
  \item The set of all matrices $A$ such that $\tr(A) = 0$.
\end{enumerate}

\qitem{3.8} Which of the following statements are true?
\begin{enumerate}[a.]
  \item If $f : \R \to \R$ is injective and continuous, then it is strictly monotonic.
  \item If $f \in \mathcal{C}[0, 2]$ is such that $f(0) = f(2)$, then there exist $x_1$ and $x_2$ in $[0, 2]$ such that $x_1 - x_2 = 1$ and $f(x_1) = f(x_2)$.
  \item Let $f$ and $g$ be continuous real valued functions on $\R$ such that for all $x \in \R$, we have $f(g(x)) = g(f(x))$. If there exists $x_0 \in \R$ such that $f(f(x_0)) = g(g(x_0))$, then there exists $x_1 \in \R$ such that $f(x_1) = g(x_1)$.
\end{enumerate}

\qitem{3.9} Which of the following statements are true?
\begin{enumerate}[a.]
  \item Let $V = \mathcal{C}_c(\R)$, the space of continuous functions on $\R$ with compact support endowed with the metric $d(f, g) = \left(\int_{-\infty}^{\infty} |f(t) - g(t)|^2\,dt\right)^{\frac{1}{2}}$. Let $f : \R \to \R$ be a continuous function which vanishes outside $[0, 1]$. Define $f_n(x) = f(x - n)$ for $n \in \N$. Then $\{f_n\}$ has a convergent subsequence in $V$.
  \item Let $\varphi, \psi$ be continuous functions on $[0, 1]$. Let $\{f_n\}$ be a sequence in $\mathcal{C}[0, 1]$ with its usual `sup-norm' topology such that, for all $n \in \N$, the functions $f_n$ are continuously differentiable and for all $x \in [0, 1]$, and for all $n \in \N$ we have $|f_n(x)| \le \varphi(x)$ and $|f_n'(x)| \le \psi(x)$. Then there exists a subsequence of $\{f_n\}$ which converges in $\mathcal{C}[0, 1]$.
  \item Let $\{A_n\}$ be a sequence of orthogonal matrices in $\M_2(\R)$. Then it has a convergent subsequence.
\end{enumerate}

\qitem{3.10} Which of the following pairs of sets are homeomorphic?
\begin{enumerate}[a.]
  \item The sets $\Q$ and $\Z$ with their usual topologies inherited from $\R$.
  \item The sets $]0, 1[$ and $]0, \infty[$ with their usual topologies inherited from $\R$.
  \item The sets $S^1 = \{z \in \C \mid z = e^{i\theta},\ 0 \le \theta < 2\pi\}$ and $A = \{z \in \C \mid z = re^{i\theta},\ 1 \le r \le 2,\ 0 \le \theta < 2\pi\}$ with their usual topologies inherited from $\C \cong \R^2$.
\end{enumerate}

\examsec{Section 4: Calculus and Differential Equations}

\qitem{4.1} Let $S = \{x = (x_1, x_2, \dots, x_n) \in \R^n \mid 0 \le x_1 \le x_2 \le \dots \le x_n \le 1\}$. Find the volume of the set $S$.

\qitem{4.2} Let $F : ]0, \infty[ \to \R$ be defined by $F(x) = \displaystyle\int_{-x}^{x} \frac{1 - e^{-xy}}{y}\,dy$. Compute $F'(x)$.

\qitem{4.3} Let $f(x, y) = x^2 + 5y^2 - 6x + 10y + 6$. Where are the maxima/minima of $f$ (if any) located?

\qitem{4.4} Evaluate: $\displaystyle\int_{-\infty}^{\infty}\int_{-\infty}^{\infty} e^{-(2x^2 + 2xy + 2y^2)}\,dxdy$.

\qitem{4.5} Write down the Taylor series expansion about the origin, up to the term involving $x^7$, for the function $f(x) = \frac{1}{2}\left[x\sqrt{1 - x^2} + \sin^{-1} x\right]$.

\qitem{4.6} Solve: $-\dfrac{d^2 u}{dr^2} - \dfrac{1}{r}\dfrac{du}{dr} = 1$, in $0 < r < 1$, $u'(0) = 0 = u(1)$.

\qitem{4.7} Which of the following two-point boundary value problems admit a unique solution?
\begin{enumerate}[a.]
  \item $-u''(x) = 2x$ in $0 < x < 1$ and $u(0) = u(1) = 0$.
  \item $-u''(x) = 2x$ in $0 < x < 1$ and $u(0) = u'(1) = 0$.
  \item $-u''(x) = 2x$ in $0 < x < 1$ and $u'(0) = u'(1) = 0$.
\end{enumerate}

\qitem{4.8} Which of the following statements are true?
\begin{enumerate}[a.]
  \item Let $\psi$ be a non-negative and continuously differentiable function on $]0, \infty[$ such that $\psi'(x) \le \psi(x)$ for all $x \in ]0, \infty[$. Then $\lim_{x \to \infty} \psi(x) = 0$.
  \item Let $\psi$ be a non-negative function continuous on $[0, \infty[$ and differentiable on $]0, \infty[$ such that $\psi(0) = 0$ and such that $\psi'(x) \le \psi(x)$ for all $x \in ]0, \infty[$. Then $\psi \equiv 0$.
  \item Let $\varphi$ be a non-negative and continuous function on $[0, \infty[$ and such that $\varphi(x) \le \int_0^x \varphi(t)\,dt$ for all $x \in [0, \infty[$. Then $\varphi \equiv 0$.
\end{enumerate}

\qitem{4.9} Write down the expression for the Laplace transform $F(s)$ of the function $f(x) = x^n$, where $n \in \N$.

\qitem{4.10} Amongst all smooth curves $y(x)$ passing through the points $(x_1, 0)$ and $(x_2, 0)$ in the plane, we wish to find that whose surface of revolution about the $x$-axis has the least surface area. Write down the functional that must be minimised to find this curve.

\examsec{Section 5: Miscellaneous}

\qitem{5.1} Let $A = (a_{ij}) \in \M_n(\R)$ be defined by $a_{ij} = \begin{cases} i, & \text{if } i + j = n + 1, \\ 0, & \text{otherwise.} \end{cases}$ Compute $\det(A)$.

\qitem{5.2} Let $n \in \N$ be fixed. For $0 \le k \le n$, let $C_k$ denote the usual binomial coefficient $\binom{n}{k}$ of choosing $k$ objects from a set of $n$ objects. Evaluate: $C_0^2 + C_1^2 + \dots + C_n^2$.

\qitem{5.3} Which of the following numbers are prime?\\
a.\ $179$. \qquad b.\ $197$. \qquad c.\ $199$.

\qitem{5.4} Given $f : \R \to \R$, define $f^2(x) = f(f(x))$. Which of the following statements are true?
\begin{enumerate}[a.]
  \item If $f$ is strictly monotonic, then $f^2$ is strictly increasing.
  \item If $f^2(x) = -x$ for all $x \in \R$, then $f$ is injective.
  \item There does not exist a continuous function $f : \R \to \R$ such that $f^2(x) = -x$ for all $x \in \R$.
\end{enumerate}

\qitem{5.5} Let $a$ be a fixed positive real number. Evaluate:
\[
  \max_{\substack{x_i \ge 0,\ 1 \le i \le n \\ \sum_{i=1}^{n} x_i = a}} x_1 x_2 \cdots x_n.
\]

\qitem{5.6} A real number is said to be \emph{algebraic} if it is the root of a non-zero polynomial of degree at least one with integer coefficients. Otherwise the number is said to be \emph{transcendental}. Which of the following statements are true?
\begin{enumerate}[a.]
  \item Algebraic numbers are dense in $\R$.
  \item Transcendental numbers are dense in $\R$.
  \item The number $\cos\left(\frac{\pi}{13}\right)$ is algebraic.
\end{enumerate}

\qitem{5.7} Let $f : ]a, b[ \to \R$ be a given function. Which of the following statements are true?
\begin{enumerate}[a.]
  \item If $f$ is convex in $]a, b[$, then the set $\Gamma = \{(x, y) \in \R^2 \mid x \in ]a, b[,\ y \ge f(x)\}$ is a convex set.
  \item If $f$ is convex in $]a, b[$, then the set $\Gamma = \{(x, y) \in \R^2 \mid x \in ]a, b[,\ y \le f(x)\}$ is a convex set.
  \item If $f$ is convex in $]a, b[$, then $|f|$ is also convex in $]a, b[$.
\end{enumerate}

\qitem{5.8} Two fair dice are rolled. What is the probability that the sum of the numbers on the top faces is $8$?

\qitem{5.9} Let $\{f_n\}_{n=1}^{\infty}$ and $f$ be real valued functions defined on $\R$. For $\varepsilon > 0$ and for $m \in \N$, define $E_m(\varepsilon) = \{x \in \R \mid |f_m(x) - f(x)| \ge \varepsilon\}$. Let $S = \{x \in \R \mid \text{the sequence } \{f_n(x)\} \text{ does not converge to } f(x)\}$. Express $S$ in terms of the sets $\{E_m(\varepsilon)\}_{m \in \N, \varepsilon > 0}$ (using the set theoretic operations of unions and intersections).

\qitem{5.10} Consider the Fibonacci sequence $\{a_n\}_{n=0}^{\infty}$ defined by $a_0 = a_1 = 1$ and $a_n = a_{n-1} + a_{n-2}$, $n \ge 2$. Let $F(z) = \sum_{n=0}^{\infty} a_n z^n$ be the generating function. Express $F$ in closed form as a function of $z$.

\examsec{Answer Key --- NBHM 2014}
\begin{multicols}{2}
\textbf{Section 1: Algebra}\\
\textbf{1.1} a, c\\
\textbf{1.2} $S^1$, the multiplicative group of complex numbers with modulus one\\
\textbf{1.3} $N(3) = 1$ or $7$; $N(7) = 1$\\
\textbf{1.4} a, b\\
\textbf{1.5} a, b\\
\textbf{1.6} a, b, c\\
\textbf{1.7} a, b\\
\textbf{1.8} a.\ rank of $A = 1$; b.\ $x^T x = 1$\\
\textbf{1.9} b, c\\
\textbf{1.10} a.\ $\{\lambda \in \C \mid \lambda \text{ real, } \ge 0\}$; b.\ $\{\lambda \in \C \mid \lambda \text{ real}\}$; c.\ $\{\lambda \in \C \mid \RePart(\lambda) = 0\}$

\medskip
\textbf{Section 2: Analysis}\\
\textbf{2.1} $\frac{1}{6} \le x < \frac{1}{2}$\\
\textbf{2.2} $\frac{k(k+1)}{2}m$\\
\textbf{2.3} a, b\\
\textbf{2.4} b, c\\
\textbf{2.5} a, b, c\\
\textbf{2.6} b, c\\
\textbf{2.7} b, c\\
\textbf{2.8} $2\log 2 - 1$\\
\textbf{2.9} b, c\\
\textbf{2.10} a, b, c

\medskip
\textbf{Section 3: Topology}\\
\textbf{3.1} a, c\\
\textbf{3.2} c\\
\textbf{3.3} a, b\\
\textbf{3.4} a, b\\
\textbf{3.5} a, b\\
\textbf{3.6} a, b\\
\textbf{3.7} a\\
\textbf{3.8} a, b, c\\
\textbf{3.9} b, c\\
\textbf{3.10} b

\medskip
\textbf{Section 4: Calculus \& Differential Equations}\\
\textbf{4.1} $\frac{1}{n!}$\\
\textbf{4.2} $\frac{2}{x}(e^{x^2} - e^{-x^2})$\\
\textbf{4.3} Minimum at $(3, -1)$.\\
\textbf{4.4} $\frac{\pi}{\sqrt{3}}$\\
\textbf{4.5} $x - \frac{1}{2}\frac{x^3}{3} - \frac{1\cdot 1}{2\cdot 4}\frac{x^5}{5} - \frac{1\cdot 1\cdot 3}{2\cdot 4\cdot 6}\frac{x^7}{7} - \cdots$\\
\textbf{4.6} $\frac{1 - r^2}{4}$\\
\textbf{4.7} a, b\\
\textbf{4.8} b, c\\
\textbf{4.9} $\frac{n!}{s^{n+1}}$\\
\textbf{4.10} $\int_{x_1}^{x_2} 2\pi y(x)\sqrt{1 + (y'(x))^2}\,dx$ (The constant $2\pi$ can be omitted.)

\medskip
\textbf{Section 5: Miscellaneous}\\
\textbf{5.1} $\det(A) = \begin{cases} (-1)^{\frac{n}{2}} n!, & \text{for } n \text{ even} \\ (-1)^{\frac{n-1}{2}} n!, & \text{for } n \text{ odd.} \end{cases}$\\
\textbf{5.2} $\binom{2n}{n}$\\
\textbf{5.3} a, b, c\\
\textbf{5.4} a, b, c\\
\textbf{5.5} $\left(\frac{a}{n}\right)^n$\\
\textbf{5.6} a, b, c\\
\textbf{5.7} a\\
\textbf{5.8} $\frac{5}{36}$\\
\textbf{5.9} $S = \bigcup_{\varepsilon > 0} \bigcap_{n=1}^{\infty} \bigcup_{m=n}^{\infty} E_m(\varepsilon)$\\
\textbf{5.10} $\frac{1}{1 - z - z^2}$

\medskip
\noindent\emph{Note: Accept any correct equivalent form of the answers.}
\end{multicols}
% ==================== NBHM 2015 ====================
\examheader{NBHM 2015}{Research Scholarships Screening Test}{Saturday, January 24, 2015}{Time Allowed: 150 Minutes \quad|\quad Maximum Marks: 40}

\subsection*{Instructions to Candidates \& Notation}
\begin{itemize}
  \item Booklet contains 10 numbered pages; 5 sections of 10 questions each (Algebra, Analysis, Topology, Calculus \& Differential Equations, Miscellaneous); assessment based on best 4 sections (max 40 marks); no partial credit; calculators not allowed.
  \item \textbf{Notation:} Standard $\N, \Z, \Q, \R, \C, \R^n, \C^n, \M_n(\R), \M_n(\C)$, binomial coefficient $\binom{n}{r}$, intervals $]a, b[, [a, b]$, $\mathcal{C}[a, b], \mathcal{C}^1[a, b]$, $f', f''$, $x^T, x^*, A^T, A^*$, $I$, $\det(A), \tr(A)$, logarithms to base $e$.
\end{itemize}

\examsec{Section 1: Algebra}

\qitem{1.1} Solve the following equation, given that its roots are in arithmetic progression:
\[
  x^3 - 9x^2 + 28x - 30 = 0.
\]

\qitem{1.2} Which of the following statements are true?
\begin{enumerate}[a.]
  \item Every group of order $51$ is cyclic.
  \item Every group of order $151$ is cyclic.
  \item Every group of order $505$ is cyclic.
\end{enumerate}

\qitem{1.3} Let $G$ be the multiplicative group of non-zero complex numbers. Consider the group homomorphism $\varphi : G \to G$ given by $\varphi(z) = z^4$.
\begin{enumerate}[a.]
  \item Identify $H$, the kernel of $\varphi$.
  \item Identify (up to isomorphism) the quotient space $G/H$.
\end{enumerate}

\qitem{1.4} How many elements of order $7$ are there in a group of order $28$?

\qitem{1.5} Which of the following equations can occur as the class equation of a group of order $10$?
\begin{enumerate}[a.]
  \item $10 = 1 + 1 + 1 + 2 + 5$
  \item $10 = 1 + 2 + 3 + 4$
  \item $10 = 1 + 1 + \dots + 1$ ($10$ times)
\end{enumerate}

\qitem{1.6} Let $W$ be the subspace of $\R^4$ spanned by the vectors $(1, 0, -1, 1)$ and $(2, 3, -1, 2)$. Write down a basis for $W^\perp$, the orthogonal complement of $W$ in $\R^4$ with respect to the usual euclidean inner-product.

\qitem{1.7} Let $B$ be a $5 \times 3$ matrix and let $C$ be a $3 \times 5$ matrix, both with real entries. Set $A = BC$. What are the possible values of the rank of $A$ when\\
a.\ both $B$ and $C$ have rank $3$? \qquad b.\ both $B$ and $C$ have rank $2$?

\qitem{1.8} Write down all the eigenvalues (along with their multiplicities) of the matrix $A = (a_{ij}) \in \M_n(\R)$ where $a_{ij} = 1$ for all $1 \le i, j \le n$.

\qitem{1.9} Let $V = \M_n(\C)$ be equipped with the inner-product $(A, B) = \tr(B^* A)$, $A, B \in V$. Let $M \in \M_n(\C)$. Define $T : V \to V$ by $T(A) = MA$. What is $T^*(A)$, where $T^*$ denotes the adjoint of the mapping $T$?

\qitem{1.10} Find a symmetric and positive definite matrix $B$ such that $B^2 = A$, where $A = \begin{bmatrix} 2 & -1 \\ -1 & 2 \end{bmatrix}$.

\examsec{Section 2: Analysis}

\qitem{2.1} Find the radius of convergence of the series: $\displaystyle\sum_{n=1}^{\infty} \frac{n^3 z^n}{3^n}$.

\qitem{2.2} Let $(X, d)$ be a metric space and let $\{x_n\}$ and $\{y_n\}$ be arbitrary sequences in $X$. Which of the following statements are true?
\begin{enumerate}[a.]
  \item If both $\{x_n\}$ and $\{y_n\}$ are Cauchy sequences, then the sequence of real numbers $\{d(x_n, y_n)\}$ is a Cauchy sequence.
  \item If $d(x_n, x_{n+1}) < \frac{1}{n+1}$, then the sequence $\{x_n\}$ is a Cauchy sequence.
  \item If $d(x_n, x_{n+1}) < \frac{1}{2^n}$, then the sequence $\{x_n\}$ is a Cauchy sequence.
\end{enumerate}

\qitem{2.3} Let $\{a_n\}$ be a sequence of real numbers such that the series $\sum_{n=1}^{\infty} |a_n|^2$ is convergent. Which of the following statements are true?
\begin{enumerate}[a.]
  \item The series $\sum_{n=1}^{\infty} \frac{a_n}{n}$ is convergent.
  \item The series $\sum_{n=1}^{\infty} |a_n|^p$ is convergent for all $2 < p < \infty$.
  \item The series $\sum_{n=1}^{\infty} |a_n|^p$ is convergent for all $1 < p < 2$.
\end{enumerate}

\qitem{2.4} Which of the following statements are true?
\begin{enumerate}[a.]
  \item If $f(x) = |x|^3$ for all $x \in \R$, then $f$ is twice differentiable on $\R$.
  \item If $f : \R \to \R$ is such that $|f(x) - f(y)| \le |x - y|^{\sqrt{2}}$ for all $x$ and $y$ in $\R$, then $f$ is a constant.
  \item If $f : \R \to \R$ is differentiable on $\R$ and if $|f'(t)| \le M$ for all $t \in \R$, then there exists $\varepsilon_0 > 0$ such that for all $0 < \varepsilon \le \varepsilon_0$, the function $g(x) = x + \varepsilon f(x)$ is injective.
\end{enumerate}

\qitem{2.5} Let $f : [0, 1] \to \R$ be continuous. Define $F(x) = \int_0^x f(t)\,dt$, $x \in [0, 1]$. Which of the following statements are true?
\begin{enumerate}[a.]
  \item The function $F$ is Lipschitz continuous on $[0, 1]$.
  \item The function $F$ is uniformly continuous on $[0, 1]$.
  \item The function $F$ is of bounded variation on $[0, 1]$.
\end{enumerate}

\qitem{2.6} Let $f_n(t) = t^n$ for $n \in \N$. Which of the following statements are true?
\begin{enumerate}[a.]
  \item The sequence $\{f_n\}$ converges uniformly on $\left[\frac{1}{4}, \frac{1}{2}\right]$.
  \item The sequence $\{f_n\}$ converges uniformly on $[0, 1]$.
  \item The sequence $\{f_n\}$ converges uniformly on $]0, 1[$.
\end{enumerate}

\qitem{2.7} Let $f \in \mathcal{C}^1[-\pi, \pi]$ be such that $f(-\pi) = f(\pi)$. Define $a_n = \int_{-\pi}^{\pi} f(t)\cos nt\,dt$, $n \in \N$. Which of the following statements are true?
\begin{enumerate}[a.]
  \item The sequence $\{a_n\}$ is bounded.
  \item The sequence $\{na_n\}$ converges to zero as $n \to \infty$.
  \item The series $\sum_{n=1}^{\infty} n^2 |a_n|^2$ is convergent.
\end{enumerate}

\qitem{2.8} Let $\Gamma$ be a simple closed curve in the complex plane (described in the positive sense) and let $z_0$ be a point in the interior of this curve. Evaluate:
\[
  \int_{\Gamma} \frac{z^3 + 2z}{(z - z_0)^3}\,dz.
\]

\qitem{2.9} Let $\Gamma$ stand for the unit circle $\{z = e^{i\theta} : -\pi \le \theta \le \pi\}$ in the complex plane. Let $k \in \R$ be a fixed constant.
\begin{enumerate}[a.]
  \item When $\Gamma$ is described in the positive sense, evaluate the integral $\displaystyle\int_{\Gamma} \frac{e^{kz}}{z}\,dz$.
  \item Hence, or otherwise, evaluate the integral $\displaystyle\int_0^{\pi} e^{k\cos\theta}\cos(k\sin\theta)\,d\theta$.
\end{enumerate}

\qitem{2.10} Which of the following statements are true?
\begin{enumerate}[a.]
  \item There exists a non-constant entire function which is bounded on the upper half-plane $H = \{z \in \C : \ImPart(z) > 0\}$.
  \item There exists a non-constant entire function which takes only real values on the imaginary axis.
  \item There exists a non-constant entire function which is bounded on the imaginary axis.
\end{enumerate}

\examsec{Section 3: Topology}

\qitem{3.1} Let $S^1 = \{(x, y) \in \R^2 : x^2 + y^2 = 1\}$ and let $S = \{(x, y) \in \R^2 : ax^2 + 2hxy + by^2 = 1\}$. In which of the following cases is $S$ homeomorphic to $S^1$?\\
a.\ $a = 2, h = 1, b = 2$ \qquad b.\ $a = 1, h = 2, b = 1$ \qquad c.\ $a = 1, h = 2, b = 4$

\qitem{3.2} Let $(X, d)$ be a compact metric space. Which of the following statements are true?
\begin{enumerate}[a.]
  \item $X$ is complete.
  \item $X$ is separable.
  \item If $f : X \to \R$ is a continuous mapping, then it maps Cauchy sequences into Cauchy sequences.
\end{enumerate}

\qitem{3.3} Let $J$ be any indexing set and let $(X_j, \tau_j)$ be topological spaces for each $j \in J$. Let $X = \prod_{j \in J} X_j$ be the product space with the corresponding product topology, $\tau$. Let $p_j : X \to X_j$, $j \in J$ be the coordinate projection. Which of the following statements are true?
\begin{enumerate}[a.]
  \item The product topology $\tau$ is the weakest (\emph{i.e.}\ smallest) topology on $X$ such that each coordinate projection $p_j$, $j \in J$ is continuous.
  \item For each $j \in J$, the mapping $p_j$ maps open sets in $X$ onto open sets in $X_j$.
  \item If $(X', \tau')$ is any topological space and if $f : X' \to X$ is a given mapping, then $f$ is continuous if, and only if, $p_j \circ f : X' \to X_j$ is continuous for each $j \in J$.
\end{enumerate}

\qitem{3.4} Let $X$ be the space of all polynomials in one variable, with real coefficients. If $p = a_0 + a_1 x + \dots + a_n x^n \in X$, define $\|p\| = |a_0| + |a_1| + \dots + |a_n|$, which gives the metric $d(p, q) = \|p - q\|$ on $X$. Which of the following statements are true?
\begin{enumerate}[a.]
  \item The metric space $X$ is complete.
  \item Define $T : X \to X$ by $Tp = a_0 + a_1 x + \frac{a_2}{2}x^2 + \dots + \frac{a_n}{n}x^n$, where $p$ is as described earlier. Then $T$ is continuous.
  \item The mapping $T$ defined above is bijective and is a homeomorphism.
\end{enumerate}

\qitem{3.5} State whether each of the following subsets of $\M_2(\R)$ are open, closed or neither open nor closed.
\begin{enumerate}[a.]
  \item The set of all matrices in $\M_2(\R)$ such that neither eigenvalue is real.
  \item The set of all matrices in $\M_2(\R)$ such that both eigenvalues are real.
\end{enumerate}

\qitem{3.6} Let $X = \R^2 \setminus \{(x, y) \in \R^2 : 3x + 5y + 1 = 0\}$. Which of the following points lie in the same connected component of $X$ as the origin?\\
a.\ $(-1, 2)$ \qquad b.\ $(2, -1)$ \qquad c.\ $(1, -2)$

\qitem{3.7} Which of the following sets in $\M_n(\R)$ are connected?
\begin{enumerate}[a.]
  \item $S = \{A \in \M_n(\R) : A^T A = AA^T = I\}$
  \item $S = \{A \in \M_n(\R) : \tr(A) = 1\}$
  \item $S = \{A \in \M_n(\R) : x^T A x \ge 0 \text{ for all } x \in \R^n\}$
\end{enumerate}

\qitem{3.8} Which of the following sequences $\{f_n\}$ in $\mathcal{C}[0, 1]$ must contain a uniformly convergent subsequence?
\begin{enumerate}[a.]
  \item When $|f_n(t)| \le 3$ for all $t \in [0, 1]$ and for all $n \in \N$.
  \item When $f_n \in \mathcal{C}^1[0, 1]$, $|f_n(t)| \le 3$ and $|f_n'(t)| \le 5$ for all $t \in [0, 1]$ and for all $n \in \N$.
  \item When $f_n \in \mathcal{C}^1[0, 1]$ and $\int_0^1 |f_n(t)|\,dt \le 1$ for all $n \in \N$.
\end{enumerate}

\qitem{3.9} Let $X = \{f \in \mathcal{C}[-5, 5] : f(-5) = f(5) = 0\}$. Which of the following statements are true?
\begin{enumerate}[a.]
  \item There exists $f \in X$ such that $f \equiv 2$ on $[-1, 0]$ and $f \equiv 3$ on $[1, 2] \cup [3, 4]$.
  \item For every $f \in X$, there exist distinct points $x_1$ and $x_2$ in $]-5, 5[$ such that $f(x_1) = f(x_2)$.
  \item For every $f \in X$, there exists $x \in ]-5, 5[$ such that $f(x) = x$.
\end{enumerate}

\qitem{3.10} Let $A \in \M_n(\C)$ and let $\rho(A) = \max\{|\lambda| : \lambda \text{ is an eigenvalue of } A\}$ denote its spectral radius. Which of the following subsets of $\M_n(\C)$ are compact?
\begin{enumerate}[a.]
  \item $S = \{A \in \M_n(\C) : \rho(A) \le 1\}$
  \item $S = \{A \in \M_n(\C) : A = A^* \text{ and } \rho(A) \le 1\}$
  \item $S = \{A \in \M_n(\C) : AA^* = A^* A = I\}$
\end{enumerate}

\examsec{Section 4: Calculus and Differential Equations}

\qitem{4.1} Find the outward unit normal to the curve $\Gamma$, given by the equation $x^2 + 4y^2 = 4$, at a point $P = (x, y)$ lying on it.

\qitem{4.2} Write down the equation of the tangent to the curve $\Gamma$ given in the preceding problem at the point $P = (\sqrt{3}, \frac{1}{2})$ lying on it.

\qitem{4.3} If $n = (n_1(x, y), n_2(x, y))$ is the outward unit normal at the point $P = (x, y)$ lying on the curve $\Gamma$ given in Problem 4.1, evaluate $\int_{\Gamma} (n_1(x, y)x + n_2(x, y)y)\,ds$.

\qitem{4.4} Let $\varphi : \R \to \R$ be a continuous function. Let $f(x, y, z) = \varphi(r)$, where $r = \sqrt{x^2 + y^2 + z^2}$. Let $B$ be the ball in $\R^3$ with centre at the origin and of radius $a > 0$. Express the integral $\int_B f(x, y, z)\,dxdydz$ as an integral with respect to $r$.

\qitem{4.5} Find the maximum area that a rectangle can have if its sides are parallel to the coordinate axes and if it is inscribed in the ellipse $\frac{x^2}{a^2} + \frac{y^2}{b^2} = 1$.

\qitem{4.6} Express the following iterated integral with the order of integration reversed:
\[
  \int_{-1}^{2}\int_{-x}^{2-x^2} f(x, y)\,dydx.
\]

\qitem{4.7} Find all the possible solutions $(\lambda, u)$, where $\lambda \in \R$ and $u \not\equiv 0$, to the boundary value problem:
\[
  u''(x) + \lambda u(x) = 0,\quad x \in ]0, 1[,\qquad u(0) = u(1) \text{ and } u'(0) = u'(1).
\]

\qitem{4.8} Using the change of the dependent variable $z = y^{-2}$, solve the differential equation: $xy' + y = x^4 y^3$.

\qitem{4.9} Find the general solution of the linear system: $x'(t) = 4x(t) - y(t)$, $y'(t) = 2x(t) + y(t)$.

\qitem{4.10} Find the extremal functions $y(x)$ of the integral: $\int_0^1 (y^2 - (y')^2)\,dx$.

\examsec{Section 5: Miscellaneous}

\qitem{5.1} Given that the equation $ax^2 + 2hxy + by^2 = 1$ represents an ellipse in the plane, what is its area?

\qitem{5.2} Find the area of the circle formed by the intersection of the sphere $x^2 + y^2 + z^2 - 2x - 4y - 6z - 2 = 0$ with the plane $x + 2y + 2z - 20 = 0$.

\qitem{5.3} For a positive integer $N$, let $\phi(N)$ denote the number of positive integers (including unity) which are less than $N$ and coprime to it. Which of the following statements are true?
\begin{enumerate}[a.]
  \item If $N \neq M$, then $\phi(NM) = \phi(N)\phi(M)$.
  \item If $N > 2$, then $\phi(N)$ is always even.
  \item If $p$ is a prime and if $N = p^k$, $k \in \N$, then $\phi(N) = N\left(1 - \frac{1}{p}\right)$.
\end{enumerate}

\qitem{5.4} Let $N$ be a fixed positive integer and let $S$ be the set of all positive integers (including unity) which are less than $N$ and coprime to it. What is the sum of all the elements of $S$?

\qitem{5.5} Which of the following statements are true?
\begin{enumerate}[a.]
  \item For every $r \in \N$, there exist $r$ consecutive composite numbers in $\N$.
  \item For every $r \in \N$, the product of $r$ consecutive numbers in $\N$ is always divisible by $r!$.
  \item If $p$ is a prime and if $r \in \N$ is such that $0 < r < p$, then $p$ divides $\binom{p}{r}$.
\end{enumerate}

\qitem{5.6} Let $n \in \N$. Which of the following statements are true?
\begin{enumerate}[a.]
  \item For every $n > 1$, $\left(\dfrac{1}{2}\dfrac{3}{4}\cdots\dfrac{2n-1}{2n}\right)^{\frac{1}{n}} > \dfrac{1}{2}$.
  \item For every $n \ge 1$, $\dfrac{1}{2}\dfrac{3}{4}\cdots\dfrac{2n-1}{2n} < \dfrac{1}{\sqrt{2n+1}}$.
  \item For every $n > 1$, $1\cdot 3\cdot 5\cdots(2n-1) < n^n$.
\end{enumerate}

\qitem{5.7} Let $u : \R^N \to \R$ be a given function. For $a \in \R$, define $a^+ = \max\{a, 0\}$. For a fixed $t \in \R$, set $v(x) = (u(x) - t)^+ + t$, $x \in \R^N$. Which of the following statements are true?
\begin{enumerate}[a.]
  \item $\{x \in \R^N : v(x) = t\} = \{x \in \R^N : u(x) = t\}$.
  \item $\{x \in \R^N : v(x) > t\} = \{x \in \R^N : u(x) > t\}$.
  \item $\{x \in \R^N : v(x) > \tau\} = \{x \in \R^N : u(x) > \tau\}$ for all $\tau \ge t$.
\end{enumerate}

\qitem{5.8} Find the number of divisors of $N = 2520$ (excluding unity and $N$).

\qitem{5.9} In how many ways can we rearrange the letters in the word \emph{INDIVISIBILITY} such that no two `I's are adjacent to each other? (Note: You are allowed to express the answer in terms of binomial coefficients, factorials etc., in which case you need not explicitly calculate this number.)

\qitem{5.10} BCCI has shortlisted $n$ cricketers for a forthcoming tour. It has to select a team of $r$ players \textbf{and} name the captain of the team. This can be done in two ways: (Australian method) First choose the team and then select the captain from amongst the team members. (British method) First choose the captain and then select the remaining members of the team. Write down the combinatorial identity which expresses the fact that both methods yield the same number of outcomes.

\examsec{Answer Key --- NBHM 2015}
\begin{multicols}{2}
\textbf{Section 1: Algebra}\\
\textbf{1.1} $3 - i,\ 3,\ 3 + i$\\
\textbf{1.2} a, b\\
\textbf{1.3} $H = \{\pm 1, \pm i\}$; $G/H \cong G$\\
\textbf{1.4} $6$\\
\textbf{1.5} c\\
\textbf{1.6} Any two linearly independent vectors satisfying $x - z + t = 0;\ 2x + 3y - z + 2t = 0$\\
\textbf{1.7} a.\ $3$; b.\ $1$ or $2$\\
\textbf{1.8} $\lambda = 0$ with multiplicity $n - 1$ and $\lambda = n$ with multiplicity $1$\\
\textbf{1.9} $T^*(A) = M^* A$\\
\textbf{1.10} $\begin{bmatrix} \frac{1+\sqrt{3}}{2} & \frac{1-\sqrt{3}}{2} \\ \frac{1-\sqrt{3}}{2} & \frac{1+\sqrt{3}}{2} \end{bmatrix}$

\medskip
\textbf{Section 2: Analysis}\\
\textbf{2.1} $3$\\
\textbf{2.2} a, c\\
\textbf{2.3} a, b\\
\textbf{2.4} a, b, c\\
\textbf{2.5} a, b, c\\
\textbf{2.6} a\\
\textbf{2.7} a, b, c\\
\textbf{2.8} $6\pi i z_0$\\
\textbf{2.9} a.\ $2\pi i$; b.\ $\pi$\\
\textbf{2.10} a, b, c

\medskip
\textbf{Section 3: Topology}\\
\textbf{3.1} a\\
\textbf{3.2} a, b, c\\
\textbf{3.3} a, b, c\\
\textbf{3.4} b\\
\textbf{3.5} a.\ open; b.\ closed\\
\textbf{3.6} a, b\\
\textbf{3.7} b, c\\
\textbf{3.8} b\\
\textbf{3.9} a, b, c\\
\textbf{3.10} b, c

\medskip
\textbf{Section 4: Calculus \& Differential Equations}\\
\textbf{4.1} $\left(\frac{x}{\sqrt{x^2+16y^2}}, \frac{4y}{\sqrt{x^2+16y^2}}\right)$\\
\textbf{4.2} $\sqrt{3}x + 2y - 4 = 0$\\
\textbf{4.3} $4\pi$\\
\textbf{4.4} $4\pi \int_0^a r^2 \varphi(r)\,dr$\\
\textbf{4.5} $2ab$\\
\textbf{4.6} $\int_{-2}^{1}\int_{-y}^{\sqrt{2-y}} f(x, y)\,dxdy + \int_{1}^{2}\int_{-\sqrt{2-y}}^{\sqrt{2-y}} f(x, y)\,dxdy$\\
\textbf{4.7} $\lambda = 0$ and $u =$ constant; $\lambda = 4n^2\pi^2$ and $u_n = A\cos 2n\pi x + B\sin 2n\pi x$, for $n \in \N$\\
\textbf{4.8} $\frac{1}{y^2} = cx^2 - x^4$\\
\textbf{4.9} $x(t) = Ae^{3t} + Be^{2t},\ y(t) = Ae^{3t} + 2Be^{2t}$\\
\textbf{4.10} $y = A\cos x + B\sin x$, or equivalently $y = c\sin(x - d)$

\medskip
\textbf{Section 5: Miscellaneous}\\
\textbf{5.1} $\frac{\pi}{\sqrt{ab - h^2}}$\\
\textbf{5.2} $7\pi$\\
\textbf{5.3} b, c\\
\textbf{5.4} $\frac{1}{2}N\phi(N)$\\
\textbf{5.5} a, b, c\\
\textbf{5.6} a, b, c\\
\textbf{5.7} b, c\\
\textbf{5.8} $46$\\
\textbf{5.9} $8!\binom{9}{6} = 8!\binom{9}{3} = 3,386,880$\\
\textbf{5.10} $r\binom{n}{r} = n\binom{n-1}{r-1}$

\medskip
\noindent\emph{Note: Accept any correct equivalent form of the answers.}
\end{multicols}
% ==================== NBHM 2016 ====================
\examheader{NBHM 2016}{Research Scholarships Screening Test}{Saturday, January 23, 2016}{Time Allowed: 150 Minutes \quad|\quad Maximum Marks: 40}

\subsection*{Instructions to Candidates \& Notation}
\begin{itemize}
  \item Booklet contains 10 numbered pages; 5 sections of 10 questions each (Algebra, Analysis, Topology, Calculus \& Differential Equations, Miscellaneous); assessment based on best 4 sections (max 40 marks); no partial credit; calculators not allowed.
  \item \textbf{Notation:} Standard $\N, \Z, \Q, \R, \C, \R^n, \C^n, \M_n(\R), \M_n(\C), \binom{n}{r}$, characteristic function $\chi_E$, intervals $]a, b[, [a, b]$, $\mathcal{C}[a, b]$, $d_1(f, g) = \int_X |f(x) - g(x)|\,dx$, $f', f'', x^T, x^*, A^T, A^*, I, \det(A), \tr(A), \ker(\varphi), \ker(A), GL_n(\R), GL_n(\C), S_n, \Z_n$, logarithms to base $e$.
\end{itemize}

\examsec{Section 1: Algebra}

\qitem{1.1} With the usual notations, compute $aba^{-1}$ in $S_5$ and express it as the product of disjoint cycles, where $a = (1\ 2\ 3)(4\ 5)$ and $b = (2\ 3)(1\ 4)$.

\qitem{1.2} Consider the following permutation:
\[
  \begin{pmatrix} 1 & 2 & 3 & 4 & 5 & 6 & 7 & 8 & 9 & 10 \\ 7 & 4 & 10 & 6 & 2 & 9 & 8 & 1 & 5 & 3 \end{pmatrix}.
\]
a.\ Is this an odd or an even permutation? \qquad b.\ What is its order in $S_{10}$?

\qitem{1.3} Which of the following statements are true?
\begin{enumerate}[a.]
  \item Let $G$ be a group of order $99$ and let $H$ be a subgroup of order $11$. Then $H$ is normal in $G$.
  \item Let $H$ be the subgroup of $S_3$ consisting of the two elements $\{e, a\}$ where $e$ is the identity and $a = (1\ 2)$. Then $H$ is normal in $S_3$.
  \item Let $G$ be a finite group and let $H$ be a subgroup of $G$. Define $W = \bigcap_{g \in G} gHg^{-1}$. Then $W$ is a normal subgroup of $G$.
\end{enumerate}

\qitem{1.4} Consider the ring $\mathcal{C}[0, 1]$ with the operations of pointwise addition and pointwise multiplication. Give an example of an ideal in this ring which is not a maximal ideal.

\qitem{1.5} Compute the (multiplicative) inverse of $4x + 3$ in the field $\Z_{11}[x]/(x^2 + 1)$.

\qitem{1.6} Let $A \in \M_5(\R)$. If $A = (a_{ij})$, let $A_{ij}$ denote the cofactor of the entry $a_{ij}$, $1 \le i, j \le 5$. Let $\widehat{A}$ denote the matrix whose $(ij)$-th entry is $A_{ij}$, $1 \le i, j \le 5$.\\
a.\ What is the rank of $\widehat{A}$ when the rank of $A$ is $5$? \qquad b.\ What is the rank of $\widehat{A}$ when the rank of $A$ is $3$?

\qitem{1.7} Write down the minimal polynomial of $A \in \M_n(\R)$, where $A = (a_{ij})$ and $a_{ij} = \begin{cases} 1 & \text{if } i + j = n + 1, \\ 0 & \text{otherwise.} \end{cases}$

\qitem{1.8} Let $V = \R^5$ be equipped with the usual euclidean inner-product. Which of the following statements are true?
\begin{enumerate}[a.]
  \item If $W$ and $Z$ are subspaces of $V$ such that both of them are of dimension $3$, then there exists $z \in Z$ such that $z \neq 0$ and $z \perp W$.
  \item There exists a non-zero linear map $T : V \to V$ such that $\ker(T) \cap W \neq \{0\}$ for every subspace $W$ of $V$ of dimension $4$.
  \item Let $W$ be a subspace of $V$ of dimension $3$. Let $T : V \to W$ be a linear map which is surjective and let $S : W \to V$ be a linear map which is injective. Then, there exists $x \in V$ such that $x \neq 0$ and such that $S \circ T(x) = 0$.
\end{enumerate}

\qitem{1.9} Which of the following statements are true?
\begin{enumerate}[a.]
  \item Let $A \in \M_3(\R)$ be such that $A^4 = I$, $A \neq \pm I$. Then $A^2 + I = 0$.
  \item Let $A \in \M_2(\R)$ be such that $A^3 = I$, $A \neq I$. Then $A^2 + A + I = 0$.
  \item Let $A \in \M_3(\R)$ be such that $A^3 = I$, $A \neq I$. Then $A^2 + A + I = 0$.
\end{enumerate}

\qitem{1.10} Find an orthogonal matrix $P$ and a diagonal matrix $D$, both in $\M_2(\R)$, such that $P^T A P = D$, where $A = \begin{bmatrix} 5 & -3 \\ -3 & 5 \end{bmatrix}$.

\examsec{Section 2: Analysis}

\qitem{2.1} Let $\{a_n\}$ be a sequence of real numbers such that $\displaystyle\lim_{n \to \infty} \left|a_n + 3\left(\frac{n-2}{n}\right)^n\right|^{\frac{1}{n}} = \frac{3}{5}$. Compute $\lim_{n \to \infty} a_n$.

\qitem{2.2} Let $f : [0, \infty[ \to [0, \infty[$ be a continuous function such that $\int_0^\infty f(t)\,dt < \infty$. Which of the following statements are true?
\begin{enumerate}[a.]
  \item The sequence $\{f(n)\}_{n \in \N}$ is bounded.
  \item $f(n) \to 0$ as $n \to \infty$.
  \item The series $\sum_{n=1}^{\infty} f(n)$ is convergent.
\end{enumerate}

\qitem{2.3} Let $\rho : \R \to \R$ be a continuous function such that $\rho(x) \ge 0$ for all $x \in \R$, $\rho(x) = 0$ if $|x| \ge 1$ and $\int_{-\infty}^{\infty} \rho(t)\,dt = 1$. Let $f : \R \to \R$ be a continuous function. Evaluate:
\[
  \lim_{\varepsilon \to 0} \frac{1}{\varepsilon}\int_{-\infty}^{\infty} \rho\left(\frac{x}{\varepsilon}\right)f(x)\,dx.
\]

\qitem{2.4} Let $I \subset \R$ be an interval. A real valued function $f$ defined on $I$ is said to have the \emph{intermediate value property} (IVP) if for every $a, b \in I$ such that $a < b$, the function $f$ assumes every value between $f(a)$ and $f(b)$ in the interval $(a, b)$. Which of the following statements are true?
\begin{enumerate}[a.]
  \item Define $f : [0, 1] \to \R$ by $f(x) = \begin{cases} \sin\frac{1}{x} & \text{if } 0 < x \le 1, \\ 0 & \text{if } x = 0. \end{cases}$ Then $f$ has IVP.
  \item If $f : \R \to \R$ is strictly increasing and has IVP, then $f$ is continuous.
  \item If $f : [a, b] \to \R$ is a differentiable function, then $f'$ has IVP.
\end{enumerate}

\qitem{2.5} Write down the Taylor expansion (about the origin) of the function $f(x) = \int_0^x \tan^{-1} t\,dt$.

\qitem{2.6} Use the preceding exercise to find the sum of the series: $1 - \frac{1}{2} - \frac{1}{3} + \frac{1}{4} + \frac{1}{5} - \frac{1}{6} - \frac{1}{7} + \cdots$

\qitem{2.7} Let $\{f_n\}$ be a sequence of continuous real valued functions defined on $\R$ converging uniformly on $\R$ to a function $f$. Which of the following statements are true?
\begin{enumerate}[a.]
  \item If each of the functions $f_n$ is bounded, then $f$ is also bounded.
  \item If each of the functions $f_n$ is uniformly continuous, then $f$ is also uniformly continuous.
  \item If each of the functions $f_n$ is integrable, then $\lim_{n \to \infty} \int_{-\infty}^{\infty} f_n(t)\,dt = \int_{-\infty}^{\infty} f(t)\,dt$.
\end{enumerate}

\qitem{2.8} Let $f : \R \to \R$ be a given function. Consider the following statements: A: The function $f$ is continuous almost everywhere. B: There exists a continuous function $g : \R \to \R$ such that $f = g$ almost everywhere. Which of the following implications are true?\\
a.\ A $\Rightarrow$ B. \qquad b.\ B $\Rightarrow$ A. \qquad c.\ A $\Leftrightarrow$ B.

\qitem{2.9} Give an example of an analytic function $f : \C \to \C$ such that $f(P) = H$, where $P = \{z \in \C \mid z = x + iy,\ x \ge 0,\ y \ge 0\}$ and $H = \{z \in \C \mid z = x + iy,\ y \ge 0\}$.

\qitem{2.10} Which of the following statements are true?
\begin{enumerate}[a.]
  \item There exists an analytic function $f : \C \to \C$ such that for every $z \in \C$, $z = x + iy$, $\RePart f(z) = e^x$.
  \item There exists an analytic function $f : \C \to \C$ such that $f$ is bounded on both the real and imaginary axes.
  \item There exists an analytic function $f : \C \to \C$ such that $f(0) = 1$ and for every $z \in \C$ such that $|z| \ge 1$, we have $|f(z)| \le e^{-|z|}$.
\end{enumerate}

\examsec{Section 3: Topology}

\qitem{3.1} Which of the following sequences $\{f_n\}$ are Cauchy?
\begin{enumerate}[a.]
  \item $f_n(x) = \begin{cases} 0 & \text{if } x \notin [n-1, n+1], \\ x - n + 1 & \text{if } x \in [n-1, n], \\ n + 1 - x & \text{if } x \in [n, n+1], \end{cases}$ in the space $X = \left\{f : \R \to \R \;\middle|\; f \text{ is continuous and } \int_{-\infty}^{\infty} |f(t)|\,dt < \infty\right\}$ equipped with the $d_1$ metric.
  \item $f_n(x) = \frac{x+n}{n}$ in the space $\mathcal{C}[0, 1]$ with the usual sup-norm metric.
  \item $f_n(x) = \frac{nx}{1+nx}$ in the space $\mathcal{C}[0, 1]$ equipped with the usual sup-norm metric.
\end{enumerate}

\qitem{3.2} Let $f_n(x) = \begin{cases} 1 - nx & \text{if } 0 \le x \le \frac{1}{n}, \\ 0 & \text{if } \frac{1}{n} \le x \le 1. \end{cases}$ Let $\mathcal{C}[0, 1]$ be equipped with the $d_1$ metric. Which of the following statements are true?\\
a.\ The sequence $\{f_n\}$ is Cauchy. \qquad b.\ The sequence $\{f_n\}$ is convergent. \qquad c.\ The sequence $\{f_n\}$ is not convergent.

\qitem{3.3} Which of the following normed linear spaces, all equipped with the sup-norm, are complete?
\begin{enumerate}[a.]
  \item The space of bounded uniformly continuous real valued functions defined on $\R$.
  \item The space of continuous real valued functions defined on $\R$ having compact support.
  \item The space of continuously differentiable real valued functions defined on $[0, 1]$.
\end{enumerate}

\qitem{3.4} Which of the following sets, $S$, are dense?
\begin{enumerate}[a.]
  \item $S = \bigcup_{m,n \in \Z} T_{m,n}$, in $\R^2$, where $T_{m,n}$ is the straight line passing through the origin and the point $(m, n)$.
  \item $S = GL_n(\R)$, in $\M_n(\R)$.
  \item $S = \{A \in \M_2(\R) \mid \text{both eigenvalues of } A \text{ are real}\}$, in $\M_2(\R)$.
\end{enumerate}

\qitem{3.5} Which of the following subsets of $\R^2$ are connected?
\begin{enumerate}[a.]
  \item $\R^2 \setminus \Q \times \Q$.
  \item $\left\{\left(x, \sin\frac{1}{x}\right) \in \R^2 \;\middle|\; 0 < x < \infty\right\} \cup \{(0, 0)\}$.
  \item $\{(x, y) \in \R^2 \mid xy = 1\} \cup \{(x, y) \in \R^2 \mid y = 0\}$.
\end{enumerate}

\qitem{3.6} Which of the following subsets are path-connected?
\begin{enumerate}[a.]
  \item $\{(x, y) \in \R^2 \mid x^2 + y^2 < 1\} \cup \{(x, y) \in \R^2 \mid y = 1\} \subset \R^2$.
  \item $\bigcup_{n=1}^{\infty} \{(x, y) \in \R^2 \mid x = ny\} \subset \R^2$.
  \item The set of all symmetric matrices all of whose eigenvalues are non-negative, in $\M_n(\R)$.
\end{enumerate}

\qitem{3.7} Which of the following statements are true?
\begin{enumerate}[a.]
  \item If $K \subset \M_n(\R)$ is a compact subset, then all the eigenvalues of all the elements of $K$ form a bounded set.
  \item Let $K \subset \M_n(\R)$ be defined by $K = \{A \in \M_n(\R) \mid A = A^T,\ \tr(A) = 1,\ x^T A x \ge 0 \text{ for all } x \in \R^n\}$. Then $K$ is compact.
  \item Let $K \subset \mathcal{C}[0, 1]$ (with the usual sup-norm metric) be defined by $K = \left\{f \in \mathcal{C}[0, 1] \;\middle|\; \int_0^1 f(t)\,dt = 1 \text{ and } f(x) \ge 0 \text{ for all } x \in [0, 1]\right\}$. Then $K$ is compact.
\end{enumerate}

\qitem{3.8} A function $f : \R \to \R$ is said to be \emph{lower semicontinuous} (lsc) if the set $f^{-1}(]-\infty, \alpha])$ is closed for every $\alpha \in \R$. Which of the following statements are true?
\begin{enumerate}[a.]
  \item If $E \subset \R$ is a closed set, then $f = \chi_E$ is lsc.
  \item If $E \subset \R$ is an open set, then $f = \chi_E$ is lsc.
  \item If $G = \{(x, y) \in \R^2 \mid y = f(x)\}$ is closed in $\R^2$, then $f$ is lsc.
\end{enumerate}

\qitem{3.9} Let $X$ be a non-empty compact Hausdorff space. Which of the following statements are true?
\begin{enumerate}[a.]
  \item If $X$ has at least $n$ distinct points, then the dimension of $\mathcal{C}(X)$, the space of continuous real valued functions defined on $X$, is at least $n$.
  \item If $A$ and $B$ are disjoint, non-empty and closed sets in $X$, there exists $f \in \mathcal{C}(X)$ such that $f(x) = -3$ for all $x \in A$ and $f(x) = 4$ for all $x \in B$.
  \item If $A \subset X$ is a closed and non-empty subset and if $g : A \to \R$ is a continuous function, then there exists $f \in \mathcal{C}(X)$ such that $f(x) = g(x)$ for all $x \in A$.
\end{enumerate}

\qitem{3.10} Which of the following subsets of $\R^2$ are homeomorphic to the set $\{(x, y) \in \R^2 \mid xy = 1\}$?
\begin{enumerate}[a.]
  \item $\{(x, y) \in \R^2 \mid xy - 2x - y + 2 = 0\}$.
  \item $\{(x, y) \in \R^2 \mid x^2 - 3x + 2 = 0\}$.
  \item $\{(x, y) \in \R^2 \mid 2x^2 - 2xy + 2y^2 = 1\}$.
\end{enumerate}

\examsec{Section 4: Calculus and Differential Equations}

\qitem{4.1} Evaluate: $\displaystyle\int_0^{\infty} x^4 e^{-x^2}\,dx$.

\qitem{4.2} Find the arc length of the curve in the plane, whose equation in polar coordinates is given by $r = a\cos\theta$, when $\theta$ varies over the interval $\left[-\frac{\pi}{2}, \frac{\pi}{2}\right]$.

\qitem{4.3} Let $S = [0, 1] \times [0, 1] \subset \R^2$. Evaluate: $\displaystyle\iint_S \max(x, y)\,dxdy$.

\qitem{4.4} Evaluate: $\displaystyle\int_{-\infty}^{\infty}\int_{-\infty}^{\infty} e^{-(5x^2 - 6xy + 5y^2)}\,dxdy$.

\qitem{4.5} Let $\mathbf{x} = (x, y) \in \R^2$. Let $\mathbf{n}(\mathbf{x})$ denote the unit outward normal to the ellipse $\gamma$ whose equation is given by $\frac{x^2}{4} + \frac{y^2}{9} = 1$ at the point $\mathbf{x}$ on it. Evaluate: $\displaystyle\int_\gamma \mathbf{x}\cdot\mathbf{n}(\mathbf{x})\,ds(\mathbf{x})$.

\qitem{4.6} Let $\omega > 0$ and let $(x_0, y_0) \in \R^2$. Solve: $\frac{dx}{dt}(t) = \omega y(t)$, $\frac{dy}{dt}(t) = -\omega x(t)$, $x(0) = x_0$, $y(0) = y_0$.

\qitem{4.7} Let $\omega > 0$. Compute the matrix $e^A$, where $A = \begin{bmatrix} 0 & \omega \\ -\omega & 0 \end{bmatrix}$.

\qitem{4.8} Write down the first order system of equations equivalent to the differential equation $\dfrac{d^3 y}{dx^3} = \dfrac{d^2 y}{dx^2} - x^2\left(\dfrac{dy}{dx}\right)^2$.

\qitem{4.9} Consider the system of differential equations: $x' = y(x^2 + 1)$, $y' = 2xy^2$.\\
a.\ Find the critical points of the system. \qquad b.\ Find all the solution paths of the system.

\qitem{4.10} Consider the boundary value problem: $-y''(x) = f(x)$ for $0 < x < 1$, $y'(0) = y'(1) = 0$. In which of the following cases does there exist a solution to this problem?\\
a.\ $f(x) = \cos\pi x$. \qquad b.\ $f(x) = x - \frac{1}{2}$. \qquad c.\ $f(x) = \sin\pi x$.

\examsec{Section 5: Miscellaneous}

\qitem{5.1} Write down the condition to be satisfied by the real numbers $a, b, c$ and $d$ in order that the sphere $x^2 + y^2 + z^2 = 1$ and the plane $ax + by + cz + d = 0$ have a non-empty intersection.

\qitem{5.2} In a triangle $ABC$, the base $AB = 6\text{ cms}$. The vertex $C$ varies such that the area is always equal to $12\text{ cm}^2$. Find the minimum value of the sum $CA + CB$.

\qitem{5.3} Find the maximum value the expression $2x + 3y + z$ takes as $(x, y, z)$ varies over the sphere $x^2 + y^2 + z^2 = 1$.

\qitem{5.4} Let $k, r$ and $n$ be positive integers such that $1 < k < r < n$. Find $\alpha_\ell$, $0 \le \ell \le k$ such that $\binom{n}{r} = \sum_{\ell=0}^{k} \alpha_\ell \binom{k}{\ell}$.

\qitem{5.5} Which of the following sets are countable?
\begin{enumerate}[a.]
  \item The set of all algebraic numbers.
  \item The set of all strictly increasing infinite sequences of positive integers.
  \item The set of all infinite sequences of integers which are in arithmetic progression.
\end{enumerate}

\qitem{5.6} Find all integer solutions of the following pair of congruences: $x \equiv 5 \pmod{8}$, $x \equiv 2 \pmod{7}$.

\qitem{5.7} Let $F : \R \to \R$ be defined by $F(s) = \begin{cases} 1 & \text{if } s \ge \frac{1}{2}, \\ 0 & \text{if } s < \frac{1}{2}. \end{cases}$ Evaluate: $\int_0^1 F(\sin\pi x)\,dx$.

\qitem{5.8} Let $\alpha = 1 + \frac{1}{9} + \frac{1}{25} + \frac{1}{49} + \cdots$, $\beta = 1 - \frac{1}{3} + \frac{1}{5} - \frac{1}{7} + \cdots$, $\gamma = 1 + \frac{1}{4} + \frac{1}{9} + \frac{1}{16} + \cdots$. Which of the following numbers are rational?\\
a.\ $\frac{\alpha}{\gamma}$. \qquad b.\ $\frac{\beta}{\gamma}$. \qquad c.\ $\frac{\beta^2}{\gamma}$.

\qitem{5.9} In how many ways can $7$ people be seated around a circular table such that two particular people are always seated next to each other?

\qitem{5.10} Find the sum of the following infinite series: $\dfrac{4}{20} + \dfrac{4\cdot 7}{20\cdot 30} + \dfrac{4\cdot 7\cdot 10}{20\cdot 30\cdot 40} + \cdots$

\examsec{Answer Key --- NBHM 2016}
\begin{multicols}{2}
\textbf{Section 1: Algebra}\\
\textbf{1.1} $(1\ 3)(2\ 5)$\\
\textbf{1.2} a.\ odd; b.\ $30$\\
\textbf{1.3} a, c\\
\textbf{1.4} Any example of the form $\mathcal{I} = \{f \mid f(x) = 0 \text{ for all } x \in S\}$ where $S \subset [0, 1]$ has at least two points.\\
\textbf{1.5} $6x + 1$\\
\textbf{1.6} a.\ $5$; b.\ $0$\\
\textbf{1.7} $\lambda^2 - 1$\\
\textbf{1.8} b, c\\
\textbf{1.9} b\\
\textbf{1.10} $P = \begin{bmatrix} \frac{1}{\sqrt{2}} & \frac{1}{\sqrt{2}} \\ \frac{1}{\sqrt{2}} & -\frac{1}{\sqrt{2}} \end{bmatrix};\ D = \begin{bmatrix} 2 & 0 \\ 0 & 8 \end{bmatrix}$

\medskip
\textbf{Section 2: Analysis}\\
\textbf{2.1} $-3e^{-2}$\\
\textbf{2.2} None\\
\textbf{2.3} $f(0)$\\
\textbf{2.4} a, b, c\\
\textbf{2.5} $\sum_{n=1}^{\infty} (-1)^{n-1}\frac{x^{2n}}{(2n-1)(2n)}$\\
\textbf{2.6} $\frac{\pi}{4} - \frac{1}{2}\log 2$\\
\textbf{2.7} a, b\\
\textbf{2.8} None\\
\textbf{2.9} Standard example: $f(z) = z^2$\\
\textbf{2.10} b

\medskip
\textbf{Section 3: Topology}\\
\textbf{3.1} b\\
\textbf{3.2} a, b\\
\textbf{3.3} a\\
\textbf{3.4} a, b\\
\textbf{3.5} a, b\\
\textbf{3.6} a, b, c\\
\textbf{3.7} a, b\\
\textbf{3.8} b, c\\
\textbf{3.9} a, b, c\\
\textbf{3.10} b

\medskip
\textbf{Section 4: Calculus \& Differential Equations}\\
\textbf{4.1} $\frac{3}{8}\sqrt{\pi}$\\
\textbf{4.2} $\pi a$\\
\textbf{4.3} $\frac{2}{3}$\\
\textbf{4.4} $\frac{\pi}{4}$\\
\textbf{4.5} $12\pi$\\
\textbf{4.6} $x(t) = x_0\cos\omega t + y_0\sin\omega t,\ y(t) = -x_0\sin\omega t + y_0\cos\omega t$\\
\textbf{4.7} $\begin{bmatrix} \cos\omega & \sin\omega \\ -\sin\omega & \cos\omega \end{bmatrix}$\\
\textbf{4.8} $y' = u;\ u' = v;\ v' = v - x^2 u^2$\\
\textbf{4.9} a.\ All points $(x, 0)$, $x \in \R$; b.\ $y = c(x^2 + 1)$\\
\textbf{4.10} a, b

\medskip
\textbf{Section 5: Miscellaneous}\\
\textbf{5.1} $d^2 \le a^2 + b^2 + c^2$\\
\textbf{5.2} $10\text{ cms}$\\
\textbf{5.3} $\sqrt{14}$\\
\textbf{5.4} $\alpha_\ell = \binom{n-k}{r-\ell},\ 0 \le \ell \le k$\\
\textbf{5.5} a, c\\
\textbf{5.6} $56k + 37,\ k \in \Z$\\
\textbf{5.7} $\frac{2}{3}$\\
\textbf{5.8} a, c\\
\textbf{5.9} $2 \times 5! = 240$\\
\textbf{5.10} $10\left(\frac{10}{7}\right)^{\frac{1}{3}} - 11$

\medskip
\noindent\emph{Note: Please accept any correct equivalent form of the answers.}
\end{multicols}
% ==================== NBHM 2017 ====================
\examheader{NBHM 2017}{Research Scholarships Screening Test}{Saturday, January 21, 2017}{Time Allowed: 150 Minutes \quad|\quad Maximum Marks: 40}

\subsection*{Instructions to Candidates \& Notation}
\begin{itemize}
  \item Booklet contains 11 numbered pages; 5 sections of 10 questions each (Algebra, Analysis, Topology, Calculus \& Differential Equations, Miscellaneous); assessment based on best 4 sections (max 40 marks); no partial credit; calculators not allowed.
  \item \textbf{Notation:} Standard $\N, \Z, \Q, \R, \C, \R^n, \C^n, \M_n(\R), \M_n(\C), \binom{n}{r}$, intervals $]a, b[, [a, b]$, $\mathcal{C}[a, b], \mathcal{C}_c(\R)$, $L^p(]a, b[), L^\infty(]a, b[)$, $f', f'', I, \det(A), \tr(A), GL_n(\R), GL_n(\C), SL_n(\R), S_n, \Z_n, \F_p$, logarithms to base $e$.
\end{itemize}

\examsec{Section 1: Algebra}

\qitem{1.1} Let $G$ be a group. Which of the following statements are true?
\begin{enumerate}[a.]
  \item Let $H$ and $K$ be subgroups of $G$ of orders $3$ and $5$ respectively. Then $H \cap K = \{e\}$, where $e$ is the identity element of $G$.
  \item If $G$ is an abelian group of odd order, then $\varphi(x) = x^2$ is an automorphism of $G$.
  \item If $G$ has exactly one element of order $2$, then this element belongs to the centre of $G$.
\end{enumerate}

\qitem{1.2} Let $n \in \N$, $n \ge 2$. Which of the following statements are true?
\begin{enumerate}[a.]
  \item Any finite group $G$ of order $n$ is isomorphic to a subgroup of $GL_n(\R)$.
  \item The group $\Z_n$ is isomorphic to a subgroup of $GL_2(\R)$.
  \item The group $\Z_{12}$ is isomorphic to a subgroup of $S_7$.
\end{enumerate}

\qitem{1.3} Which of the following statements are true?
\begin{enumerate}[a.]
  \item The matrices $\begin{bmatrix} 1 & 1 \\ 0 & 1 \end{bmatrix}$ and $\begin{bmatrix} 1 & 0 \\ 1 & 1 \end{bmatrix}$ are conjugate in $GL_2(\R)$.
  \item The matrices $\begin{bmatrix} 1 & 1 \\ 0 & 1 \end{bmatrix}$ and $\begin{bmatrix} 1 & 0 \\ 1 & 1 \end{bmatrix}$ are conjugate in $SL_2(\R)$.
  \item The matrices $\begin{bmatrix} 1 & 0 \\ 0 & 2 \end{bmatrix}$ and $\begin{bmatrix} 1 & 3 \\ 0 & 2 \end{bmatrix}$ are conjugate in $GL_2(\R)$.
\end{enumerate}

\qitem{1.4} Let $p$ be an odd prime. Find the number of non-zero squares in $\F_p$.

\qitem{1.5} Find a generator of $\F_7^\times$, the multiplicative group of non-zero elements of $\F_7$.

\qitem{1.6} The characteristic polynomial of a matrix $A \in \M_5(\R)$ is given by $x^5 + \alpha x^4 + \beta x^3$, where $\alpha$ and $\beta$ are non-zero real numbers. What are the possible values of the rank of $A$?

\qitem{1.7} Let $A \in \M_3(\R)$ be a symmetric matrix whose eigenvalues are $1, 1$ and $3$. Express $A^{-1}$ in the form $\alpha I + \beta A$, where $\alpha, \beta \in \R$.

\qitem{1.8} Let $A \in \M_n(\R)$, $n \ge 2$. Which of the following statements are true?\\
a.\ If $A^{2n} = 0$, then $A^n = 0$. \qquad b.\ If $A^2 = I$, then $A = \pm I$. \qquad c.\ If $A^{2n} = I$, then $A^n = \pm I$.

\qitem{1.9} Which of the following statements are true?
\begin{enumerate}[a.]
  \item There does not exist a non-diagonal matrix $A \in \M_2(\R)$ such that $A^3 = I$.
  \item There exists a non-diagonal matrix $A \in \M_2(\R)$ which is diagonalizable over $\R$ and which is such that $A^3 = I$.
  \item There exists a non-diagonal matrix $A \in \M_2(\R)$ such that $A^3 = I$ and such that $\tr(A) = -1$.
\end{enumerate}

\qitem{1.10} Let $n \ge 2$ and let $W$ be the subspace of $\M_n(\R)$ consisting of all matrices whose trace is zero. If $A = (a_{ij})$ and $B = (b_{ij})$, for $1 \le i, j \le n$, are elements in $\M_n(\R)$, define their inner-product by $(A, B) = \sum_{i,j=1}^{n} a_{ij} b_{ij}$. Identify the subspace $W^\perp$ of elements orthogonal to the subspace $W$.

\examsec{Section 2: Analysis}

\qitem{2.1} Let $\{x_n\}_{n=1}^{\infty}$ be a sequence of real numbers. Let $\alpha = \liminf_{n \to \infty} x_n$. Which of the following statements are true?
\begin{enumerate}[a.]
  \item For every $\varepsilon > 0$, there exists a subsequence $\{x_{n_k}\}$ such that $x_{n_k} \le \alpha + \varepsilon$ for all $k \in \N$.
  \item For every $\varepsilon > 0$, there exists a subsequence $\{x_{n_k}\}$ such that $x_{n_k} \le \alpha - \varepsilon$ for all $k \in \N$.
  \item There exists a subsequence $\{x_{n_k}\}$ such that $x_{n_k} \to \alpha$ as $k \to \infty$.
\end{enumerate}

\qitem{2.2} Let $\{a_n\}_{n=1}^{\infty}$ be a sequence of positive real numbers such that $\sum_{n=1}^{\infty} a_n$ is divergent. Which of the following series are convergent?\\
a.\ $\displaystyle\sum_{n=1}^{\infty} \frac{a_n}{1 + na_n}$. \qquad b.\ $\displaystyle\sum_{n=1}^{\infty} \frac{a_n}{1 + n^2 a_n}$. \qquad c.\ $\displaystyle\sum_{n=1}^{\infty} \frac{a_n}{1 + a_n}$.

\qitem{2.3} Let $\{a_n\}_{n=1}^{\infty}$ be a sequence of positive real numbers such that $\sum_{n=1}^{\infty} a_n$ is convergent. Which of the following series are convergent?\\
a.\ $\displaystyle\sum_{n=1}^{\infty} \frac{a_n}{1 + a_n}$. \qquad b.\ $\displaystyle\sum_{n=1}^{\infty} \frac{a_n^{\frac{1}{4}}}{n^{\frac{4}{5}}}$. \qquad c.\ $\displaystyle\sum_{n=1}^{\infty} na_n \sin\frac{1}{n}$.

\qitem{2.4} Let $f : \R \to \R$ be a given function. It is said to be \emph{lower semi-continuous} (respectively \emph{upper semi-continuous}) if the set $f^{-1}(]-\infty, \alpha])$ (respectively, the set $f^{-1}([\alpha, \infty[)$) is closed for every $\alpha \in \R$. Let $f$ and $g$ be two real valued functions defined on $\R$. Which of the following statements are true?
\begin{enumerate}[a.]
  \item If $f$ and $g$ are continuous, then $\max\{f, g\}$ is continuous.
  \item If $f$ and $g$ are lower semi-continuous, then $\max\{f, g\}$ is lower semi-continuous.
  \item If $f$ and $g$ are upper semi-continuous, then $\max\{f, g\}$ is upper semi-continuous.
\end{enumerate}

\qitem{2.5} Let $f : \R \to \R$ be a continuous function. Which of the following statements are true?
\begin{enumerate}[a.]
  \item If $f$ is continuously differentiable, then $f$ is uniformly continuous.
  \item If $f$ has compact support, then $f$ is uniformly continuous.
  \item If $\lim_{|x| \to \infty} |f(x)| = 0$, then $f$ is uniformly continuous.
\end{enumerate}

\qitem{2.6} Let $f : ]0, 2[ \to \R$ be defined by
\[
  f(x) = \begin{cases} x^2, & \text{if } x \in ]0, 2[ \cap \Q, \\ 2x - 1, & \text{if } x \in ]0, 2[ \setminus \Q. \end{cases}
\]
Check for the points of differentiability of $f$ and evaluate the derivative at those points.

\qitem{2.7} Let $\{f_n\}_{n=1}^{\infty}$ be a sequence of continuous real valued functions defined on $\R$ which converges pointwise to a continuous real valued function $f$. Which of the following statements are true?
\begin{enumerate}[a.]
  \item If $0 \le f_n \le f$ for all $n \in \N$, then $\lim_{n \to \infty} \int_{-\infty}^{\infty} f_n(t)\,dt = \int_{-\infty}^{\infty} f(t)\,dt$.
  \item If $|f_n(t)| \le |\sin t|$ for all $t \in \R$ and for all $n \in \N$, then $\lim_{n \to \infty} \int_{-\infty}^{\infty} f_n(t)\,dt = \int_{-\infty}^{\infty} f(t)\,dt$.
  \item If $|f_n(t)| \le e^t$ for all $t \in \R$ and for all $n \in \N$, then for all $a, b \in \R$, $a < b$, $\lim_{n \to \infty} \int_a^b f_n(t)\,dt = \int_a^b f(t)\,dt$.
\end{enumerate}

\qitem{2.8} Which of the following statements are true?
\begin{enumerate}[a.]
  \item The following series is uniformly convergent over $[-1, 1]$: $\displaystyle\sum_{n=0}^{\infty} \frac{x^2}{(1+x^2)^n}$.
  \item $\displaystyle\lim_{n \to \infty} \int_{\frac{\pi}{2}}^{\pi} \frac{\sin nx}{nx^5}\,dx = \pi$.
  \item Define, for $x \in \R$, $f(x) = \sum_{n=1}^{\infty} \frac{\sin nx^2}{1 + n^3}$. Then $f$ is a continuously differentiable function.
\end{enumerate}

\qitem{2.9} Write down the Laurent series expansion of the function $f(z) = \dfrac{-1}{(z - 1)(z - 2)}$ in the annulus $\{z \in \C \mid 1 < |z| < 2\}$.

\qitem{2.10} Which of the following statements are true?
\begin{enumerate}[a.]
  \item There exists a non-constant entire function which is bounded on the real and imaginary axes of $\C$.
  \item The ring of analytic functions on the open unit disc of $\C$ (with respect to the operations of pointwise addition and pointwise multiplication) is an integral domain.
  \item There exists an entire function $f$ such that $f(0) = 1$ and such that $|f(z)| \le \frac{1}{|z|}$ for all $|z| \ge 5$.
\end{enumerate}

\examsec{Section 3: Topology}

\qitem{3.1} Let $(X, d)$ be a metric space and let $\{x_n\}_{n=1}^{\infty}$ and $\{y_n\}_{n=1}^{\infty}$ be arbitrary Cauchy sequences in $X$. Which of the following statements are true?
\begin{enumerate}[a.]
  \item The sequence $\{d(x_n, y_n)\}$ converges as $n \to \infty$.
  \item The sequence $\{d(x_n, y_n)\}$ converges as $n \to \infty$ only if $X$ is complete.
  \item No conclusion can be drawn about the convergence of $\{d(x_n, y_n)\}$.
\end{enumerate}

\qitem{3.2} Which of the following statements are true?
\begin{enumerate}[a.]
  \item Let $X$ be a set equipped with two topologies $\tau_1$ and $\tau_2$. Assume that any given sequence in $X$ converges with respect to the topology $\tau_1$ if, and only if, it also converges with respect to the topology $\tau_2$. Then $\tau_1 = \tau_2$.
  \item Let $(X, \tau_1)$ and $(Y, \tau_2)$ be two topological spaces and let $f : X \to Y$ be a given map. Then $f$ is continuous if, and only if, given any sequence $\{x_n\}_{n=1}^{\infty}$ such that $x_n \to x$ in $X$, we have $f(x_n) \to f(x)$ in $Y$.
  \item Let $(X, \tau)$ be a compact topological space and let $\{x_n\}_{n=1}^{\infty}$ be a sequence in $X$. Then, it has a convergent subsequence.
\end{enumerate}

\qitem{3.3} Which of the following statements are true?
\begin{enumerate}[a.]
  \item Let $n \ge 2$. The subset of nilpotent matrices in $\M_n(\C)$ is closed in $\M_n(\C)$.
  \item Let $n \ge 2$. The set of all matrices in $\M_n(\C)$ which represent orthogonal projections is closed in $\M_n(\C)$.
  \item The set of all matrices in $\M_2(\R)$ such that both of their eigenvalues are purely imaginary, is closed in $\M_2(\R)$.
\end{enumerate}

\qitem{3.4} Which of the following sets are dense?
\begin{enumerate}[a.]
  \item The set of all numbers of the form $\frac{m}{2^n}$ where $0 \le m \le 2^n$ and $n \in \N$, in the space $[0, 1]$.
  \item The set of all polynomial functions in the space $L^1(]0, 1[)$.
  \item The linear span of the family $\{\sin nt\}_{n=1}^{\infty}$ in the space $L^2(]-\pi, \pi[)$.
\end{enumerate}

\qitem{3.5} Let $n \ge 2$. Which of the following subsets are nowhere dense in $\M_n(\R)$?\\
a.\ The set $GL_n(\R)$. \qquad b.\ The set of all matrices whose trace is zero. \qquad c.\ The set of all singular matrices.

\qitem{3.6} Which of the following topological spaces are separable?
\begin{enumerate}[a.]
  \item Any real Banach space which admits a Schauder basis $\{u_n\}_{n=1}^{\infty}$.
  \item The space $\mathcal{C}[0, 1]$.
  \item The space $L^p(]0, 1[)$, where $1 \le p \le \infty$.
\end{enumerate}

\qitem{3.7} Which of the following sets are connected?
\begin{enumerate}[a.]
  \item The set of all points in the plane with at least one coordinate irrational.
  \item An infinite set $X$ with the topology $\tau$ given by $\tau = \{X, \emptyset\} \cup \{A \subset X \mid X \setminus A \text{ is a finite set}\}$.
  \item The set $K = \left\{f \in \mathcal{C}[0, 1] \;\middle|\; \int_0^{\frac{1}{2}} f(t)\,dt - \int_{\frac{1}{2}}^{1} f(t)\,dt = 1\right\}$.
\end{enumerate}

\qitem{3.8} Which of the following statements are true?
\begin{enumerate}[a.]
  \item There exists a continuous bijection $f : [0, 1] \to [0, 1] \times [0, 1]$.
  \item There exists a continuous map $f : S^1 \to \R$ which is injective, where $S^1$ stands for the unit circle in the plane.
  \item There exists a continuous map $f : [0, 1] \to SL_2(\R)$ which is surjective.
\end{enumerate}

\qitem{3.9} Which of the following statements are true?
\begin{enumerate}[a.]
  \item Let $g \in \mathcal{C}[0, 1]$ be fixed. Then the set $A = \left\{f \in \mathcal{C}[0, 1] \;\middle|\; \int_0^1 f(t)g(t)\,dt = 0\right\}$ is closed in $\mathcal{C}[0, 1]$.
  \item Let $g \in \mathcal{C}_c(\R)$ be fixed. Then the set $A = \left\{f \in \mathcal{C}_c(\R) \;\middle|\; \int_{-\infty}^{\infty} f(t)g(t)\,dt = 0\right\}$ is closed in $\mathcal{C}_c(\R)$.
  \item Let $g \in L^2(\R)$ be fixed. Then the set $A = \left\{f \in L^2(\R) \;\middle|\; \int_{-\infty}^{\infty} f(t)g(t)\,dt = 0\right\}$ is closed in $L^2(\R)$.
\end{enumerate}

\qitem{3.10} Which of the following statements are true?
\begin{enumerate}[a.]
  \item Let $X$ be a compact topological space and let $\mathcal{F}$ be a family of real valued functions defined on $X$ with the following properties: (i) If $f, g \in \mathcal{F}$, then $fg \in \mathcal{F}$, where $(fg)(x) = f(x)g(x)$ for all $x \in X$. (ii) For every $x \in X$, there exists an open neighbourhood $U(x)$ of $x$ and a function $f \in \mathcal{F}$ such that the restriction of $f$ to $U(x)$ is identically zero. Then the function which is identically zero on all of $X$ belongs to $\mathcal{F}$.
  \item Let $X = \{f : [0, 1] \to [0, 1] \mid |f(t) - f(s)| \le |t - s| \text{ for all } s, t \in [0, 1]\}$. Define $d(f, g) = \max_{t \in [0, 1]} |f(t) - g(t)|$ for $f, g \in X$. Then $(X, d)$ is a compact metric space.
  \item Let $\{f_i\}_{i \in I}$ be a collection of functions in $\mathcal{C}[0, 1]$ such that given any finite subfamily of functions, its members vanish at some common point (which depends on that subfamily). Then there exists $x_0 \in [0, 1]$ such that $f_i(x_0) = 0$ for all $i \in I$.
\end{enumerate}

\examsec{Section 4: Calculus \& Differential Equations}

\qitem{4.1} Let $x > 1$. Define $F(x) = \displaystyle\int_{x^2}^{x^3} \tan(xy^2)\,dy$. Differentiate $F$ with respect to $x$.

\qitem{4.2} Evaluate: $\displaystyle\int_{-\infty}^{\infty} e^{-2x^2}\,dx$.

\qitem{4.3} Let $\mathbf{n}(x, y, z)$ denote the unit outer normal vector on the surface $S$ of the cylinder $x^2 + y^2 \le 4$, $0 \le z \le 3$. Compute $\int_S \mathbf{v}\cdot\mathbf{n}\,dS$ where $\mathbf{v}(x, y, z) = xz\mathbf{i} + 2yz\mathbf{j} + 3xy\mathbf{k}$.

\qitem{4.4} Evaluate the line integral $\int_C P\,dx + Q\,dy$, where $C$ is the circle centered at the origin and of radius $a > 0$ (described in the counter-clockwise sense) in the plane and $P(x, y) = \frac{-y}{x^2+y^2}$, $Q(x, y) = \frac{x}{x^2+y^2}$.

\qitem{4.5} Let $\Omega$ be a bounded open subset of $\R^3$ and let $\partial\Omega$ denote its boundary. Given sufficiently smooth real valued functions $u$ and $v$ on $\overline{\Omega}$, let $\frac{\partial u}{\partial n}$ and $\frac{\partial v}{\partial n}$ denote the outer normal derivatives of $u$ and $v$ respectively on $\partial\Omega$. Fill in the blank in the following identity:
\[
  \int_{\partial\Omega} \left(\frac{\partial u}{\partial n}v - \frac{\partial v}{\partial n}u\right)\,dS = \int_{\Omega} (\dots\dots\dots)\,dx\,dy\,dz.
\]

\qitem{4.6} Find the maximum value of $x^2 + xy$ subject to the condition $x^2 + y^2 \le 1$.

\qitem{4.7} Interchange the order of integration: $\displaystyle\int_{-1}^{2}\int_{-x}^{2-x^2} f(x, y)\,dy\,dx$.

\qitem{4.8} Find all the non-trivial solutions $(\lambda, u)$ (\emph{i.e.}\ $u \not\equiv 0$), of the boundary value problem: $-u''(x) = \lambda u(x)$, $0 < x < 1$, and $u(0) = u'(1) = 0$.

\qitem{4.9} Consider the initial value problem: $u'(t) = Au(t)$, $t > 0$, and $u(0) = u_0$, where $u_0$ is a given vector in $\R^2$ and $A = \begin{bmatrix} 1 & -2 \\ 1 & a \end{bmatrix}$. Find the range of values of $a$ such that $|u(t)| \to 0$ as $t \to \infty$.

\qitem{4.10} Let $u(x, t)$ be the solution of the wave equation: $\frac{\partial^2 u}{\partial t^2} = \frac{\partial^2 u}{\partial x^2}$, $x \in \R,\ t > 0$; $u(x, 0) = u_0(x)$, $x \in \R$; $u_t(x, 0) = 0$, $x \in \R$. Let $u_0(x)$ be the function defined by $u_0(x) = \begin{cases} 1, & \text{if } |x| < 2, \\ 0, & \text{otherwise.} \end{cases}$ Compute $u(x, 1)$ at all points $x \in \R$ where it is continuous.

\examsec{Section 5: Miscellaneous}

\qitem{5.1} Let $x \in \R$ and let $n \in \N$. Evaluate: $\displaystyle\sum_{k=0}^{n} \binom{n}{k}\sin\left(x + \frac{k\pi}{2}\right)$.

\qitem{5.2} Let $n \in \N$, $n \ge 2$. Let $x_1, \dots, x_n \in ]0, \pi[$. Set $x = (x_1 + \dots + x_n)/n$. Which of the following statements are true?\\
a.\ $\prod_{k=1}^{n} \sin x_k \ge \sin^n x$. \qquad b.\ $\prod_{k=1}^{n} \sin x_k \le \sin^n x$. \qquad c.\ Neither (a) nor (b) is necessarily true.

\qitem{5.3} Which of the following sets are convex?\\
a.\ $\{(x, y) \in \R^2 \mid xy \ge 1,\ x \ge 0,\ y \ge 0\}$. \qquad b.\ $\{(x, y) \in \R^2 \mid |x|^{\frac{1}{3}} + |y|^{\frac{1}{3}} \le 1\}$. \qquad c.\ $\{(x, y) \in \R^2 \mid y \ge x^2\}$.

\qitem{5.4} Find the area of the circle got by intersecting the sphere $x^2 + y^2 + z^2 = 1$ with the plane $x + y + z = 1$.

\qitem{5.5} Let $n \in \N$, $n \ge 3$. Find the area of the polygon with one vertex at $z = 1$ and whose other vertices are situated at the roots of the polynomial $1 + z + z^2 + \dots + z^{n-1}$ in the complex plane.

\qitem{5.6} Find the maximum value of $3x + 2y$ subject to the conditions: $2x + 3y \ge 6$, $y - x \le 2$, $0 \le x \le 3$, $y \ge 0$.

\qitem{5.7} A committee of six members is formed from a group of $7$ men and $4$ women. What is the probability that the committee contains\\
a.\ exactly two women? \qquad b.\ at least two women?

\qitem{5.8} Find the sum of the infinite series: $\dfrac{1}{2\cdot 3\cdot 4} + \dfrac{1}{4\cdot 5\cdot 6} + \dfrac{1}{6\cdot 7\cdot 8} + \cdots$

\qitem{5.9} Find the remainder when $8^{130}$ is divided by $13$.

\qitem{5.10} Let $a_i \in \R$, $1 \le i \le 4$. Evaluate:
\[
  \begin{vmatrix}
    1 & 1 & 1 & 1 \\
    a_1 & a_2 & a_3 & a_4 \\
    a_1^2 & a_2^2 & a_3^2 & a_4^2 \\
    a_1^3 & a_2^3 & a_3^3 & a_4^3
  \end{vmatrix}.
\]

\examsec{Answer Key --- NBHM 2017}
\begin{multicols}{2}
\textbf{Section 1: Algebra}\\
\textbf{1.1} a, b, c\\
\textbf{1.2} a, b, c\\
\textbf{1.3} a, c\\
\textbf{1.4} $\frac{p-1}{2}$\\
\textbf{1.5} $3$ or $5$\\
\textbf{1.6} $2, 3, 4$\\
\textbf{1.7} $A^{-1} = \frac{4}{3}I - \frac{1}{3}A$\\
\textbf{1.8} a\\
\textbf{1.9} c\\
\textbf{1.10} $W^\perp = \{aI \mid a \in \R\}$

\medskip
\textbf{Section 2: Analysis}\\
\textbf{2.1} a, c\\
\textbf{2.2} b\\
\textbf{2.3} a, b, c\\
\textbf{2.4} a, b, c\\
\textbf{2.5} b, c\\
\textbf{2.6} Only $x = 1$ and $f'(1) = 2$\\
\textbf{2.7} a, c\\
\textbf{2.8} c\\
\textbf{2.9} $\sum_{n=1}^{\infty} \frac{1}{z^n} + \sum_{n=0}^{\infty} \frac{z^n}{2^{n+1}}$\\
\textbf{2.10} a, b

\medskip
\textbf{Section 3: Topology}\\
\textbf{3.1} a\\
\textbf{3.2} None\\
\textbf{3.3} a, b\\
\textbf{3.4} a, b\\
\textbf{3.5} b, c\\
\textbf{3.6} a, b\\
\textbf{3.7} a, b, c\\
\textbf{3.8} None\\
\textbf{3.9} a, b, c\\
\textbf{3.10} a, b, c

\medskip
\textbf{Section 4: Calculus \& Differential Equations}\\
\textbf{4.1} $\int_{x^2}^{x^3} y^2\sec^2(xy^2)\,dy + 3x^2\tan x^7 - 2x\tan x^5$\\
\textbf{4.2} $\sqrt{\frac{\pi}{2}}$\\
\textbf{4.3} $54\pi$\\
\textbf{4.4} $2\pi$\\
\textbf{4.5} $\int_\Omega (v\Delta u - u\Delta v)\,dx\,dy\,dz$, where $\Delta$ is the Laplace operator.\\
\textbf{4.6} $\frac{1+\sqrt{2}}{2}$\\
\textbf{4.7} $\int_{-2}^{1}\int_{-y}^{\sqrt{2-y}} f(x, y)\,dx\,dy + \int_{1}^{2}\int_{-\sqrt{2-y}}^{\sqrt{2-y}} f(x, y)\,dx\,dy$\\
\textbf{4.8} $\lambda_n = \frac{(2n-1)^2\pi^2}{4}$, $u_n = C\sin\frac{(2n-1)\pi x}{2}$, $n \in \N$\\
\textbf{4.9} $-2 < a < -1$\\
\textbf{4.10} $u(x, 1) = \begin{cases} 1, & \text{if } |x| < 1, \\ \frac{1}{2}, & \text{if } 1 < |x| < 3, \\ 0, & \text{if } |x| > 3. \end{cases}$

\medskip
\textbf{Section 5: Miscellaneous}\\
\textbf{5.1} $2^{\frac{n}{2}}\sin\left(x + \frac{n\pi}{4}\right)$\\
\textbf{5.2} b\\
\textbf{5.3} a, c\\
\textbf{5.4} $\frac{2\pi}{3}$\\
\textbf{5.5} $\frac{n}{2}\sin\frac{2\pi}{n}$\\
\textbf{5.6} $19$\\
\textbf{5.7} a.\ $\frac{5}{11}$, b.\ $\frac{53}{66}$\\
\textbf{5.8} $\frac{3}{4} - \log 2$\\
\textbf{5.9} $12$\\
\textbf{5.10} $\prod_{1 \le i < j \le 4} (a_i - a_j)$

\medskip
\noindent\emph{Note: Please accept any correct equivalent form of the answers.}
\end{multicols}
% ==================== NBHM 2018 ====================
\examheader{NBHM 2018}{Research Scholarships Screening Test}{Saturday, January 20, 2018}{Time Allowed: 150 Minutes \quad|\quad Maximum Marks: 40}

\subsection*{Instructions to Candidates \& Notation}
\begin{itemize}
  \item Booklet contains 9 numbered pages; 5 sections of 10 questions each (Algebra, Analysis, Topology, Calculus \& Differential Equations, Miscellaneous); assessment based on best 4 sections (max 40 marks); no partial credit; calculators not allowed.
  \item \textbf{Notation:} Standard $\N, \Z, \Q, \R, \C, \R^n, \C^n, (\cdot,\cdot), \M_n(\R), \M_n(\C), \binom{n}{r}, ]a, b[, [a, b], \mathcal{C}[a, b], \mathcal{C}^1[a, b], \ell_2, f', f'', f''', x^T, x^*, A^T, A^*, I, \det(A), \tr(A), \ker(\varphi), \ker(A), GL_n(F), S_n, \F_p$, logarithms to base $e$.
\end{itemize}

\examsec{Section 1: Algebra}

\qitem{1.1} Find the number of elements conjugate to $(1\ 2\ 3\ 4\ 5\ 6\ 7)$ in $S_7$.

\qitem{1.2} What is the order of a $2$-Sylow subgroup in $GL_3(\F_5)$?

\qitem{1.3} Let $H$ be the subgroup generated by $(1\ 2)$ in $S_3$. Compute the normalizer, $N(H)$, of $H$.

\qitem{1.4} Let $G$ be a group. Which of the following statements are true?
\begin{enumerate}[a.]
  \item The normalizer of a subgroup of $G$ is a normal subgroup of $G$.
  \item The centre of $G$ is a normal subgroup of $G$.
  \item If $H$ is a normal subgroup of $G$ and is of order $2$, then $H$ is contained in the centre of $G$.
\end{enumerate}

\qitem{1.5} Which of the following are prime ideals in the ring $\mathcal{C}[0, 1]$?
\begin{enumerate}[a.]
  \item $J = \{f \in \mathcal{C}[0, 1] \mid f(x) = 0 \text{ for all } \frac{1}{3} \le x \le \frac{2}{3}\}$.
  \item $J = \{f \in \mathcal{C}[0, 1] \mid f(\frac{1}{3}) = f(\frac{2}{3}) = 0\}$.
  \item $J = \{f \in \mathcal{C}[0, 1] \mid f(\frac{1}{3}) = 0\}$.
\end{enumerate}

\qitem{1.6} Let $W = \left\{A \in \M_3(\R) \;\middle|\; A^T = -A \text{ and } \sum_{j=1}^{3} a_{1j} = 0\right\}$. Write down a basis for $W$.

\qitem{1.7} Let $A \in \M_5(\C)$ be such that $(A^2 - I)^2 = 0$. Assume that $A$ is not a diagonal matrix. Which of the following statements are true?\\
a.\ $A$ is diagonalizable. \qquad b.\ $A$ is not diagonalizable. \qquad c.\ No conclusion can be drawn about the diagonalizability of $A$.

\qitem{1.8} Which of the following statements are true?
\begin{enumerate}[a.]
  \item If $A \in \M_n(\R)$ is such that $(Ax, x) = 0$ for all $x \in \R^n$, then $A = 0$.
  \item If $A \in \M_n(\C)$ is such that $(Ax, x) = 0$ for all $x \in \C^n$, then $A = 0$.
  \item If $A \in \M_n(\C)$ is such that $(Ax, x) \ge 0$ for all $x \in \C^n$, then $A = A^*$.
\end{enumerate}

\qitem{1.9} Let $A = \begin{bmatrix} 1 & 1 & 1 \\ 0 & 2 & 1 \\ 0 & 0 & 2 \end{bmatrix}$. Let $W = \{x \in \R^3 \mid Ax = 2x\}$. Construct a linear functional $\varphi$ on $\R^3$ such that $\varphi(x_0) = 1$, where $x_0^T = (1, 2, 3)$, and $\varphi(x) = 0$ for all $x \in W$.

\qitem{1.10} Let $V$ be a finite dimensional vector space over $\R$ and let $f$ and $g$ be two non-zero linear functionals on $V$ such that whenever $f(x) \ge 0$, we also have that $g(x) \ge 0$. Which of the following statements are true?\\
a.\ $\ker(f) \subset \ker(g)$. \qquad b.\ $\ker(f) = \ker(g)$. \qquad c.\ $f = \alpha g$ for some $\alpha > 0$.

\examsec{Section 2: Analysis}

\qitem{2.1} Let $\{a_n\}_{n=1}^{\infty}$ be a sequence of non-negative real numbers. Which of the following statements are true?
\begin{enumerate}[a.]
  \item If $\sum_{n=1}^{\infty} a_n < +\infty$, then $\sum_{n=1}^{\infty} a_n^5 < +\infty$.
  \item If $\sum_{n=1}^{\infty} a_n^5 < +\infty$, then $\sum_{n=1}^{\infty} a_n < +\infty$.
  \item If $\sum_{n=1}^{\infty} a_n^{\frac{3}{2}} < +\infty$, then $\sum_{n=1}^{\infty} \frac{a_n}{n} < +\infty$.
\end{enumerate}

\qitem{2.2} Which of the following functions are uniformly continuous on $\R$?
\begin{enumerate}[a.]
  \item $f(x) = |\sin^3 x|$.
  \item $f(x) = \tan^{-1} x$.
  \item $f(x) = \sum_{n=1}^{\infty} f_n(x)$, where $f_n(x) = \begin{cases} n(x - n + \frac{1}{n}), & \text{if } x \in [n - \frac{1}{n}, n], \\ n(n + \frac{1}{n} - x), & \text{if } x \in [n, n + \frac{1}{n}], \\ 0, & \text{otherwise.} \end{cases}$
\end{enumerate}

\qitem{2.3} Let $f$ and $g$ be defined on $\R$ by
\[
  f(x) = \begin{cases} 1, & \text{if } x \text{ is rational,} \\ 0, & \text{if } x \text{ is irrational,} \end{cases}
  \quad\text{and}\quad
  g(x) = \begin{cases} 1, & \text{if } x \ge 0, \\ 0, & \text{if } x < 0. \end{cases}
\]
Which of the following statements are true?
\begin{enumerate}[a.]
  \item The function $f$ is continuous almost everywhere.
  \item The function $f$ is equal to a continuous function almost everywhere.
  \item The function $g$ is equal to a continuous function almost everywhere.
\end{enumerate}

\qitem{2.4} Let $\{f_n\}_{n=1}^{\infty}$ be a sequence of continuous functions defined on $[0, 1]$. Assume that $f_n(x) \to f(x)$ for every $x \in [0, 1]$. Which of the following conditions imply that this convergence is uniform?
\begin{enumerate}[a.]
  \item The function $f$ is continuous.
  \item $f_n(x) \downarrow f(x)$ for every $x \in [0, 1]$.
  \item The function $f$ is continuous and $f_n(x) \downarrow f(x)$ for every $x \in [0, 1]$.
\end{enumerate}

\qitem{2.5} Let $\{f_n\}_{n=1}^{\infty}$ be a sequence of non-negative continuous functions defined on $[0, 1]$. Assume that $f_n(x) \to f(x)$ for every $x \in [0, 1]$. Which of the following conditions imply that $\lim_{n \to \infty} \int_0^1 f_n(x)\,dx = \int_0^1 f(x)\,dx$?\\
a.\ $f_n(x) \uparrow f(x)$ for every $x \in [0, 1]$. \qquad b.\ $f_n(x) \le f(x)$ for every $x \in [0, 1]$. \qquad c.\ $f$ is continuous.

\qitem{2.6} Let $\{f_n\}_{n=1}^{\infty}$ and $f$ be integrable functions on $[0, 1]$ such that $\lim_{n \to \infty} \int_0^1 |f_n(x) - f(x)|\,dx = 0$. Which of the following statements are true?
\begin{enumerate}[a.]
  \item $f_n(x) \to f(x)$, as $n \to \infty$, for almost every $x \in [0, 1]$.
  \item $\int_0^1 f_n(x)\,dx \to \int_0^1 f(x)\,dx$, as $n \to \infty$.
  \item If $\{g_n\}_{n=1}^{\infty}$ is a uniformly bounded sequence of continuous functions converging pointwise to a function $g$, then $\int_0^1 |f_n(x)g_n(x) - f(x)g(x)|\,dx \to 0$ as $n \to \infty$.
\end{enumerate}

\qitem{2.7} Let $f : [-\pi, \pi] \to \R$ be a continuous $2\pi$-periodic function whose Fourier series is given by $\frac{a_0}{2} + \sum_{k=1}^{\infty} (a_k\cos kt + b_k\sin kt)$. Let, for each $n \in \N$, $f_n(t) = \frac{a_0}{2} + \sum_{k=1}^{n} (a_k\cos kt + b_k\sin kt)$, and let $f_0$ denote the constant function $a_0/2$. Which of the following statements are true?
\begin{enumerate}[a.]
  \item $f_n \to f$ uniformly on $[-\pi, \pi]$.
  \item If $\sigma_n = (f_0 + f_1 + \dots + f_n)/(n + 1)$, then $\sigma_n \to f$ uniformly on $[-\pi, \pi]$.
  \item $\int_{-\pi}^{\pi} |f_n(x) - f(x)|^2\,dx \to 0$, as $n \to \infty$.
\end{enumerate}

\qitem{2.8} Let $f \in \mathcal{C}^1[-\pi, \pi]$. Define, for $n \in \N$, $b_n = \int_{-\pi}^{\pi} f(t)\sin nt\,dt$. Which of the following statements are true?\\
a.\ $b_n \to 0$, as $n \to \infty$. \qquad b.\ $nb_n \to 0$, as $n \to \infty$. \qquad c.\ The series $\sum_{n=1}^{\infty} n^3 b_n^3$ is absolutely convergent.

\qitem{2.9} Let $\Gamma$ denote the circle in the complex plane, centred at zero and of radius $2$, described in the counter-clockwise sense. Evaluate: $\displaystyle\int_\Gamma \frac{e^{-z}}{(z - 1)^2}\,dz$.

\qitem{2.10} Which of the following statements are true?
\begin{enumerate}[a.]
  \item There exists an entire function defined on $\C$ such that $f(0) = 1$ and $|f(z)| \le |z|^{-2}$ for all $z \in \C$ such that $|z| \ge 10$.
  \item If $f : \C \to \C$ is a non-constant entire function, then its image is dense in $\C$.
  \item Let $\gamma : [0, 1] \to \{z \in \C \mid |z| \le 1\}$ be a non-constant continuous mapping such that $\gamma(0) = 0$. Let $f$ be an analytic function in the disc $\{z \in \C \mid |z| < 2\}$, such that $f(0) = 0$ and $f(1) = 1$. Then, there exists $\tau$ such that $0 < \tau < 1$ and such that for all $0 < t < \tau$, we have that $f(\gamma(t)) \neq 0$.
\end{enumerate}

\examsec{Section 3: Topology}

\qitem{3.1} Let $(X, d)$ be a metric space and let $\{x_n\}_{n=1}^{\infty}$ be a sequence in $X$. Let $x \in X$. Define a sequence $\{y_n\}_{n=1}^{\infty}$ by $y_{2n-1} = x_n$ and $y_{2n} = x$, $n \in \N$. Which of the following statements are true?
\begin{enumerate}[a.]
  \item If $x_n \to x$ as $n \to \infty$, then the sequence $\{y_n\}$ is Cauchy.
  \item If the sequence $\{y_n\}$ is Cauchy, then $x_n \to x$ as $n \to \infty$.
  \item Let $f : X \to X$ be a mapping that maps Cauchy sequences to Cauchy sequences. Then $f$ is continuous.
\end{enumerate}

\qitem{3.2} Let $\{V_n\}_{n=1}^{\infty}$ be a sequence of open and dense subsets of $\R^N$. Set $V = \bigcap_{n=1}^{\infty} V_n$. Which of the following statements are true?\\
a.\ $V \neq \emptyset$. \qquad b.\ $V$ is an open set. \qquad c.\ $V$ is dense in $\R^N$.

\qitem{3.3} Let $X$ be a topological space and let $U \subset X$. In which of the following cases is $U$ open?
\begin{enumerate}[a.]
  \item Let $U$ be the set of invertible upper triangular matrices in $\M_n(\R)$, where $n \ge 2$, and $X = \M_n(\R)$.
  \item Let $U$ be the set of all $2 \times 2$ matrices with real entries such that all their eigenvalues belong to $\C \setminus \R$, and $X = \M_2(\R)$.
  \item Let $U$ be the set of all complex numbers $\lambda$ such that $A - \lambda I$ is invertible, where $A$ is a given $3 \times 3$ matrix with complex entries, and $X = \C$.
\end{enumerate}

\qitem{3.4} Let $X$ be an infinite set. Define a topology $\tau$ on $X$ as follows: $\tau = \{X, \emptyset\} \cup \{U \mid X \setminus U \text{ is a non-empty finite set}\}$. Which of the following statements are true?\\
a.\ $(X, \tau)$ is Hausdorff. \qquad b.\ $(X, \tau)$ is compact. \qquad c.\ $(X, \tau)$ is connected.

\qitem{3.5} Which of the following sets are connected?
\begin{enumerate}[a.]
  \item The set of orthogonal matrices in $\M_n(\R)$, where $n \ge 2$.
  \item The set $S = \left\{f \in \mathcal{C}[0, 1] \;\middle|\; \int_0^{\frac{1}{2}} f(t)\,dt - \int_{\frac{1}{2}}^{1} f(t)\,dt = 1\right\}$ in $\mathcal{C}[0, 1]$.
  \item The set of all points in $\R^2$ with at least one coordinate being a transcendental number.
\end{enumerate}

\qitem{3.6} Which of the following sets are nowhere dense?
\begin{enumerate}[a.]
  \item $S = \{A \in \M_n(\R) \mid \tr(A) = 0\}$ in $\M_n(\R)$, where $n \ge 2$.
  \item $S = \{x \in \ell_2 \mid x = (x_n),\ x_n = 0 \text{ for infinitely many } n\}$ in $\ell_2$.
  \item The Cantor set in $[0, 1]$.
\end{enumerate}

\qitem{3.7} For $x = (x_n) \in \ell_2$, define $T(x) = (0, x_1, x_2, \dots)$ and $S(x) = (x_2, x_3, \dots)$. Which of the following statements are true?
\begin{enumerate}[a.]
  \item $\|T\| = \|S\| = 1$.
  \item If $A : \ell_2 \to \ell_2$ is a continuous linear operator such that $\|A - T\| < 1$, then $SA$ is invertible.
  \item If $A$ is as above, then $A$ is not invertible.
\end{enumerate}

\qitem{3.8} Let $V$ and $W$ be normed linear spaces and let $T : V \to W$ be a continuous linear operator. Let $B$ be the closed unit ball in $V$. In which of the following cases is $\overline{T(B)}$ compact?
\begin{enumerate}[a.]
  \item $V = \mathcal{C}^1[0, 1]$, $W = \mathcal{C}[0, 1]$ and $T(f) = f$.
  \item $V = W = \ell_2$ and $T(x) = (0, x_1, x_2, \dots)$, where $x = (x_n) \in \ell_2$.
  \item $V = W = \ell_2$ and $T(x) = (x_1, x_2, \dots, x_{10}, 0, \dots, 0, \dots)$, where $x = (x_n) \in \ell_2$.
\end{enumerate}

\qitem{3.9} Which of the following statements are true?
\begin{enumerate}[a.]
  \item The equation $x^5 + \cos^2 x = 0$ has a solution in $\R$.
  \item The equation $2x - \cos^2 x = 0$ has a solution in $[0, 1]$.
  \item The equation $x^3 - \cos^2 x = 0$ has a solution in $[-1, 0]$.
\end{enumerate}

\qitem{3.10} Let $D$ denote the closed unit disc and let $S^1$ denote the unit circle in $\R^2$. Let $E = \left\{(x, y) \in \R^2 \;\middle|\; \frac{x^2}{4} + \frac{y^2}{9} \le 1\right\}$. Which of the following statements are true?
\begin{enumerate}[a.]
  \item If $f : E \to E$ is continuous, then there exists $x \in E$ such that $f(x) = x$.
  \item If $f : D \to S^1$ is continuous, there exists $x \in S^1$ such that $f(x) = x$.
  \item If $f : S^1 \to S^1$ is continuous, there exists $x \in S^1$ such that $f(x) = x$.
\end{enumerate}

\examsec{Section 4: Calculus \& Differential Equations}

\qitem{4.1} Let $A = [0, 1] \times [0, 1] \subset \R^2$. Evaluate: $\displaystyle\iint_A \cos(\pi\max\{x, y\})\,dx\,dy$.

\qitem{4.2} Let $S$ denote the unit sphere in $\R^3$. Evaluate: $\displaystyle\int_S (x^4 + y^4 + z^4)\,dS$.

\qitem{4.3} If $\Gamma$ denotes the usual gamma function, evaluate $\Gamma(\frac{5}{2})$, given that $\Gamma(\frac{1}{2}) = \sqrt{\pi}$.

\qitem{4.4} Find the general solution of the differential equation: $yy'' = (y')^2$.

\qitem{4.5} A particle falling freely from rest under the influence of gravity suffers air resistance proportional to the square of its velocity at each instant. If $g$ stands for the acceleration due to gravity and if $c > 0$ is the constant of proportionality for the air resistance, write down the initial value problem satisfied by $y(t)$, the distance travelled by the particle at time $t$.

\qitem{4.6} In the preceding problem, if $v(t)$ is the velocity at time $t$, evaluate: $\lim_{t \to +\infty} v(t)$.

\qitem{4.7} Write down the following ordinary differential equation as a system of first order differential equations: $y''' = y'' - x^2(y')^2$.

\qitem{4.8} Let $\Omega \subset \R^N$ be a bounded domain in $\R^N$, $N \ge 2$. Let $\Delta$ denote the Laplace operator in $\R^N$. Consider the eigenvalue problem: $-\Delta u = \lambda u$ in $\Omega$, $u = 0$ on $\partial\Omega$. If $(u_i, \lambda_i)$, $i = 1, 2$ are two solutions such that $\lambda_1 \neq \lambda_2$ and $\int_\Omega |u_1(x)|^2\,dx = \int_\Omega |u_2(x)|^2\,dx = 1$, evaluate $\int_\Omega u_1(x)u_2(x)\,dx$.

\qitem{4.9} Let $\Omega$ denote the unit ball in $\R^3$. Let $\Delta$ denote the Laplace operator in $\R^3$. Let $u$ be such that $\Delta u = c$ in $\Omega$, $\frac{\partial u}{\partial \nu} = 1$ on $\partial\Omega$, where $\frac{\partial u}{\partial \nu}$ is the outer normal derivative of $u$ on $\partial\Omega$. Given that $c$ is a constant, find its value.

\qitem{4.10} Let $u(x, t)$ be the solution of the following initial value problem: $\frac{\partial^2 u}{\partial t^2} = \frac{\partial^2 u}{\partial x^2}$, $t > 0,\ x \in \R$; $u(x, 0) = u_0(x)$, $x \in \R$; $\frac{\partial u}{\partial t}(x, 0) = 0$, $x \in \R$. If $u_0(x)$ vanishes outside the interval $[-1, 1]$, find the interval outside which $u(\cdot, t)$ vanishes, when $t > 1$.

\examsec{Section 5: Miscellaneous}

\qitem{5.1} Let $X$ be a non-empty set. Let $E$ and $F$ be subsets of $X$. Define $E\Delta F = (E \setminus F) \cup (F \setminus E)$. Simplify the following expressions: (a) $(E\Delta F)\Delta(E \cap F)$, and (b) $E\Delta(E\Delta F)$.

\qitem{5.2} Let $X$ be a non-empty set and let $\{E_n\}_{n=1}^{\infty}$ be a sequence of subsets of $X$. Let $E \subset X$ be the set of all points in $X$ which lie in infinitely many of the $E_n$. Express $E$ in terms of the $E_n$ using the set theoretic operations of union and intersection.

\qitem{5.3} Let $n \in \N$ be fixed. If $0 \le r \le n$, let $C_r = \binom{n}{r}$. Evaluate the sum up to $n$ terms of the series: $3C_1 + 7C_2 + 11C_3 + \cdots$

\qitem{5.4} Eight different dolls are to be packed in eight different boxes. If two of the boxes are too small to hold five of the dolls, in how many ways can the dolls be packed?

\qitem{5.5} Given six consonants and three vowels, five-letter words are formed. What is the probability that a randomly chosen word contains three consonants and two vowels?

\qitem{5.6} Find the lengths of the semi-axes of the ellipse: $5x^2 - 8xy + 5y^2 = 1$.

\qitem{5.7} Find the sum of the infinite series: $\dfrac{1}{5} + \dfrac{1}{3}\cdot\dfrac{1}{5^3} + \dfrac{1}{5}\cdot\dfrac{1}{5^5} + \cdots$

\qitem{5.8} How many zero's are there at the end of $61!$?

\qitem{5.9} Which of the following statements are true?
\begin{enumerate}[a.]
  \item The product of $r$ consecutive positive integers is always divisible by $r!$.
  \item If $n$ is a prime number and if $0 < r < n$, then $n$ divides $\binom{n}{r}$.
  \item If $n \in \N$, then $n(n + 1)(2n + 1)$ is divisible by $6$.
\end{enumerate}

\qitem{5.10} Let $a, b, c \in \R$. Find the maximum value of $ax + by + cz$ over the set $\{(x, y, z) \in \R^3 \mid |x|^3 + |y|^3 + |z|^3 = 1\}$.

\examsec{Answer Key --- NBHM 2018}
\begin{multicols}{2}
\textbf{Section 1: Algebra}\\
\textbf{1.1} $6!$\\
\textbf{1.2} $2^7$\\
\textbf{1.3} $N(H) = H$\\
\textbf{1.4} b, c\\
\textbf{1.5} c\\
\textbf{1.6} Any two linearly independent elements in $W$. Example: $\begin{bmatrix} 0 & 1 & -1 \\ -1 & 0 & 0 \\ 1 & 0 & 0 \end{bmatrix},\ \begin{bmatrix} 0 & 0 & 0 \\ 0 & 0 & 1 \\ 0 & -1 & 0 \end{bmatrix}$\\
\textbf{1.7} c\\
\textbf{1.8} b, c\\
\textbf{1.9} Any functional of the form $ax - ay + cz$, where $3c - a = 1$\\
\textbf{1.10} a, b, c

\medskip
\textbf{Section 2: Analysis}\\
\textbf{2.1} a, c\\
\textbf{2.2} a, b\\
\textbf{2.3} b\\
\textbf{2.4} c\\
\textbf{2.5} a, b\\
\textbf{2.6} b, c\\
\textbf{2.7} b, c\\
\textbf{2.8} a, b, c\\
\textbf{2.9} $-\frac{2\pi i}{e}$\\
\textbf{2.10} b, c

\medskip
\textbf{Section 3: Topology}\\
\textbf{3.1} a, b, c\\
\textbf{3.2} a, c\\
\textbf{3.3} b, c\\
\textbf{3.4} b, c\\
\textbf{3.5} b, c\\
\textbf{3.6} a, c\\
\textbf{3.7} a, b, c\\
\textbf{3.8} a, c\\
\textbf{3.9} a, b\\
\textbf{3.10} a, b

\medskip
\textbf{Section 4: Calculus \& Differential Equations}\\
\textbf{4.1} $-\frac{4}{\pi^2}$\\
\textbf{4.2} $\frac{12\pi}{5}$\\
\textbf{4.3} $\frac{3}{4}\sqrt{\pi}$\\
\textbf{4.4} $y = c_1 e^{c_2 x}$\\
\textbf{4.5} $y'' = g - c(y')^2$ for $t > 0$; $y(0) = y'(0) = 0$.\\
\textbf{4.6} $\sqrt{\frac{g}{c}}$\\
\textbf{4.7} $y' = u;\ u' = v;\ v' = v - x^2 u^2$\\
\textbf{4.8} $0$\\
\textbf{4.9} $3$\\
\textbf{4.10} $[-(t + 1), (t + 1)]$

\medskip
\textbf{Section 5: Miscellaneous}\\
\textbf{5.1} a.\ $E \cup F$; b.\ $F$\\
\textbf{5.2} $E = \bigcap_{n=1}^{\infty}\bigcup_{m=n}^{\infty} E_m$\\
\textbf{5.3} $2^n(2n - 1) + 1$\\
\textbf{5.4} $4320$\\
\textbf{5.5} $\frac{10}{21}$\\
\textbf{5.6} Semi-major axis $= 1$; semi-minor axis $= \frac{1}{3}$\\
\textbf{5.7} $\frac{1}{2}\log\frac{3}{2}$\\
\textbf{5.8} $14$\\
\textbf{5.9} a, b, c\\
\textbf{5.10} $\left(|a|^{\frac{3}{2}} + |b|^{\frac{3}{2}} + |c|^{\frac{3}{2}}\right)^{\frac{2}{3}}$

\medskip
\noindent\emph{Note: Please accept any correct equivalent form of the answers.}
\end{multicols}
% ==================== NBHM 2019 ====================
\examheader{NBHM 2019}{Research Scholarships Screening Test}{Saturday, January 19, 2019}{Time Allowed: 150 Minutes \quad|\quad Maximum Marks: 40}

\subsection*{Instructions to Candidates \& Notation}
\begin{itemize}
  \item Booklet contains 9 numbered pages; 5 sections of 10 questions each (Algebra, Analysis, Topology, Calculus \& Differential Equations, Miscellaneous); assessment based on best 4 sections (max 40 marks); no partial credit; calculators not allowed.
  \item \textbf{Notation:} Standard $\N, \Z, \Q, \R, \C, \R^n, \C^n, (\cdot,\cdot), \M_n(\R), \M_n(\C), \binom{n}{r}, i, ]a, b[, [a, b], \mathcal{C}[a, b], \mathcal{C}^1[a, b], f', f'', f''', f^{(n)}, x^T, x^*, A^T, A^*, I, V^\perp, r(A), GL_n(F), S_n, \F_p$, logarithms to base $e$.
\end{itemize}

\examsec{Section 1: Algebra}

\qitem{1.1} Which of the following statements are true?
\begin{enumerate}[a.]
  \item If $G$ is a finite group, then there exists $n \in \N$ such that $G$ is isomorphic to a subgroup of $GL_n(\R)$.
  \item There exists an infinite group $G$ such that every element, other than the identity element, is of order $2$.
  \item The group $GL_2(\R)$ contains a cyclic subgroup of order $5$.
\end{enumerate}

\qitem{1.2} Let $p$ be a prime number. Let $n \in \N$, $n > 1$. What is the order of a $p$-Sylow subgroup of $GL_n(\F_p)$?

\qitem{1.3} Give an example of a $5$-Sylow subgroup of $GL_3(\F_5)$.

\qitem{1.4} What is the number of elements of order $2$ in $S_4$?

\qitem{1.5} Let $G$ be a group of order $10$. Which of the following could be the class equation of $G$?\\
a.\ $10 = 1 + \dots + 1$ ($10$ times). \qquad b.\ $10 = 1 + 2 + 3 + 4$. \qquad c.\ $10 = 1 + 1 + 1 + 2 + 5$.

\qitem{1.6} Find the number of irreducible monic polynomials of degree $2$ in $\F_p$, where $p$ is a prime number.

\qitem{1.7} Let $A$ be an $m \times n$ matrix with real entries. Which of the following statements are true?\\
a.\ $r(A^T A) \le r(A)$. \qquad b.\ $r(A^T A) = r(A)$. \qquad c.\ $r(A^T A) > r(A)$.

\qitem{1.8} Let $V = \{(x, y, z, t) \in \R^4 \mid x + y - z = 0 \text{ and } x + y + t = 0\}$. Write down a basis for $V^\perp$.

\qitem{1.9} Let $x_0 \in \R^3$ be the column vector such that $x_0^T = (1, 1, 1)$. Let $V = \{A \in \M_3(\R) \mid Ax_0 = 0\}$. What is the dimension of $V$?

\qitem{1.10} Let $A = \begin{bmatrix} 19 & 2019 \\ 2019 & 1 \end{bmatrix}$. Which of the following statements are true?
\begin{enumerate}[a.]
  \item The matrix $A$ is diagonalizable over $\R$.
  \item There exists a basis of $\R^2$ consisting of eigenvectors $\{w_1, w_2\}$ of the matrix $A$ such that $w_1^T w_2 = 0$.
  \item There exists a matrix $B \in GL_2(\R)$ such that $B^3 = A$.
\end{enumerate}

\examsec{Section 2: Analysis}

\qitem{2.1} Let $f$ be a function that is known to be analytic in a neighbourhood of the origin in the complex plane. Furthermore it is known that for $n \in \N$,
\[
  f^{(n)}(0) = (n - 1)!(n + 1)\left(\frac{n + 1}{n}\right)^{(n+1)(n-1)}.
\]
Find the radius of the largest circle with centre at the origin inside which the Taylor series of $f$ defines an analytic function.

\qitem{2.2} Which of the following statements are true?
\begin{enumerate}[a.]
  \item The function $f(x) = \sin^2 x$ is uniformly continuous on $]0, \infty[$.
  \item If $f : ]0, \infty[ \to \R$ is uniformly continuous, then $\lim_{x \to 0} f(x)$ exists.
  \item If $f : ]0, \infty[ \to \R$ is uniformly continuous, then $\lim_{x \to \infty} f(x)$ exists.
\end{enumerate}

\qitem{2.3} Which of the following statements are true?\\
a.\ $\log x \le \frac{x}{e}$ for all $x > 0$. \qquad b.\ $\log x \ge \frac{x}{e}$ for all $x > 0$. \qquad c.\ $e^\pi > \pi^e$.

\qitem{2.4} Let $\{r_1, \dots, r_n, \dots\}$ be an enumeration of the rationals in the interval $[0, 1]$. Define, for $n \in \N$ and for each $x \in [0, 1]$, $f_n(x) = \begin{cases} 1, & \text{if } x = r_1, \dots, r_n, \\ 0, & \text{otherwise.} \end{cases}$ Which of the following statements are true?
\begin{enumerate}[a.]
  \item The function $f_n$ is Riemann integrable over the interval $[0, 1]$ for each $n \in \N$.
  \item The sequence $\{f_n\}$ is pointwise convergent and the limit function is Riemann integrable over the interval $[0, 1]$.
  \item The sequence $\{f_n\}$ is pointwise convergent but the limit function is not Riemann integrable over the interval $[0, 1]$.
\end{enumerate}

\qitem{2.5} Let $[a, b] \subset \R$ be a finite interval. A function $f : [a, b] \to \R$ is said to be of \emph{bounded variation} if there exists a real number $L > 0$ such that for every partition $\mathcal{P} = \{a = x_0 < x_1 < \dots < x_n = b\}$, where $n \in \N$, we have $\sum_{i=1}^{n} |f(x_i) - f(x_{i-1})| \le L$. Which of the following statements are true?
\begin{enumerate}[a.]
  \item If $f : [a, b] \to \R$ is monotonic, then it is of bounded variation.
  \item If $f \in \mathcal{C}^1[a, b]$, then it is of bounded variation.
  \item If $f : [a, b] \to \R$ and $g : [a, b] \to \R$ are of bounded variation, then $f + g$ is also of bounded variation.
\end{enumerate}

\qitem{2.6} Let $[a, b] \subset \R$ be a finite interval. Let $f : [a, b] \to \R$ be a bounded and Riemann integrable function. Define, for $x \in [a, b]$, $F(x) = \int_a^x f(t)\,dt$. Which of the following statements are true?\\
a.\ The function $F$ is uniformly continuous. \qquad b.\ The function $F$ is of bounded variation. \qquad c.\ The function $F$ is differentiable on $]a, b[$.

\qitem{2.7} Which of the following statements are true?
\begin{enumerate}[a.]
  \item The sequence of functions $\{f_n\}_{n=1}^{\infty}$, defined by $f_n(x) = x^n(1 - x)$, is uniformly convergent on the interval $[0, 1]$.
  \item The sequence of functions $\{f_n\}_{n=1}^{\infty}$, defined by $f_n(x) = n\log\left(1 + \frac{x^2}{n}\right)$, is uniformly convergent on $\R$.
  \item The series $\displaystyle\sum_{n=1}^{\infty} 2^n \sin\frac{1}{3^n x}$ is uniformly convergent on the interval $[1, \infty[$.
\end{enumerate}

\qitem{2.8} Let $\{f_n\}_{n=1}^{\infty}$ be a sequence of functions in $\mathcal{C}^1[0, 1]$ such that $f_n(0) = 0$ for all $n \in \N$. Which of the following statements are true?
\begin{enumerate}[a.]
  \item If the sequence $\{f_n\}$ converges uniformly on the interval $[0, 1]$, then the limit function is in $\mathcal{C}^1[0, 1]$.
  \item If the sequence $\{f_n'\}$ is uniformly convergent over the interval $[0, 1]$, then the sequence $\{f_n\}$ is also uniformly convergent over the same interval.
  \item If the series $\sum_{n=1}^{\infty} f_n'$ converges uniformly over the interval $[0, 1]$ to a function $g$, then $g$ is Riemann integrable and $\int_0^1 g(t)\,dt = \sum_{n=1}^{\infty} f_n(1)$.
\end{enumerate}

\qitem{2.9} a.\ Write down the Laurent series expansion of the function $f(z) = \dfrac{-1}{(z - 1)(z - 2)}$ in the region $\{z \in \C \mid 1 < |z| < 2\}$.\\
b.\ Let the circle $\Gamma = \{z \in \C \mid |z| = \frac{3}{2}\}$ be described in the counter-clockwise sense. Evaluate: $\int_\Gamma f(z)\,dz$.

\qitem{2.10} Which of the following statements are true?
\begin{enumerate}[a.]
  \item There exists an analytic function $f : \C \to \C$ such that its real part is the function $e^x$, where $z = x + iy$.
  \item There exists an analytic function $f : \C \to \C$ such that $f(z) = z$ for all $z$ such that $|z| = 1$ and $f(z) = z^2$ for all $z$ such that $|z| = 2$.
  \item There exists an analytic function $f : \C \to \C$ such that $f(0) = 1$, $f(4i) = i$ and for all $z_j$ such that $1 < |z_j| < 3$, $j = 1, 2$, we have $|f(z_1) - f(z_2)| \le |z_1 - z_2|^{\frac{\pi}{3}}$.
\end{enumerate}

\examsec{Section 3: Topology}

\qitem{3.1} Let $X$ be a topological space. If $A \subset X$, we denote by $A^\circ$, the interior of $A$. Which of the following statements are true?
\begin{enumerate}[a.]
  \item If $A$ and $B$ are subsets of $X$, then $(A \cup B)^\circ = A^\circ \cup B^\circ$.
  \item If $A$ and $B$ are subsets of $X$, then $(A \cap B)^\circ = A^\circ \cap B^\circ$.
  \item If $A \subset X$, then $A^\circ = (\overline{A^c})^c$, where for $B \subset X$, we denote by $B^c$, its complement, \emph{i.e.}\ $B^c = X \setminus B$, and by $\overline{B}$, the closure of $B$.
\end{enumerate}

\qitem{3.2} Which of the following statements are true?
\begin{enumerate}[a.]
  \item Let $X$ be a topological space such that the set $\Delta = \{(x, x) \mid x \in X\}$ is closed in $X \times X$ (with the product topology). Then $X$ is Hausdorff.
  \item Let $X$ be a set and let $Y$ be a Hausdorff space. Let $f : X \to Y$ be a given mapping. Define $U \subset X$ to be open in $X$ if, and only if, $U = f^{-1}(V)$ for some set $V$ open in $Y$. This defines a Hausdorff topology on $X$.
  \item The weakest topology on $\R$ such that all polynomials (in a single variable) are continuous, is Hausdorff.
\end{enumerate}

\qitem{3.3} Which of the following statements are true?
\begin{enumerate}[a.]
  \item The map $f : [0, 2\pi[ \to S^1$, defined by $f(t) = e^{it}$, is a homeomorphism.
  \item The spaces $S^1$ and $S^2$, with their topologies inherited from $\R^2$ and $\R^3$, respectively, are homeomorphic.
  \item Let $X$ and $Y$ be topological spaces and let $f : X \to Y$ be a continuous map. Define $G(f) = \{(x, f(x)) \mid x \in X\} \subset X \times Y$. If $G(f)$ inherits the product topology from $X \times Y$, then $X$ is homeomorphic to $G(f)$ via the map $x \mapsto (x, f(x))$.
\end{enumerate}

\qitem{3.4} Let $(X, d)$ be a metric space. Let $J$ be an indexing set. Consider a set of the form $S = \{x_j \in X \mid j \in J\}$ with the property that $d(x_j, x_k) = 1$ for all $j \neq k$, $j, k \in J$. Which of the following statements are true?
\begin{enumerate}[a.]
  \item If such a set exists in $X$, then there exist open sets $\{U_j\}_{j \in J}$ in $X$ such that $U_j \cap U_k = \emptyset$ for all $j \neq k$, $j, k \in J$.
  \item There exists such a set $S$ in $\mathcal{C}[0, 1]$ with $J$ being uncountable.
  \item If such a set exists in $X$, and if $X$ is compact, then $J$ must be finite.
\end{enumerate}

\qitem{3.5} Let $E = \left\{(x, y) \in \R^2 \;\middle|\; \frac{x^2}{4} + \frac{y^2}{9} = 1\right\}$, with the topology inherited from $\R^2$. Which of the following statements are true?
\begin{enumerate}[a.]
  \item There exists $f : E \to \R$ that is continuous and onto.
  \item If $[\alpha, \beta]$ is any finite interval in $\R$, there exists $f : E \to [\alpha, \beta]$ that is continuous and onto.
  \item If $[\alpha, \beta]$ is any finite interval in $\R$, there exists $f : E \to [\alpha, \beta]$ that is continuous, one-one and onto.
\end{enumerate}

\qitem{3.6} Let $f : \R \to \R$ be a continuous $2\pi$-periodic function, \emph{i.e.}\ for every $t \in \R$, we have $f(t) = f(t + 2\pi)$. Which of the following statements are true?\\
a.\ There exists $t_0 \in \R$ such that $f(t_0) = f(-t_0)$.\\
b.\ There exists $t_0 \in \R$ such that $f(t_0) = f(t_0 + \frac{\pi}{2})$.\\
c.\ There exists $t_0 \in \R$ such that $f(t_0) = f(t_0 + \frac{\pi}{4})$.

\qitem{3.7} Which of the following spaces are connected?
\begin{enumerate}[a.]
  \item $GL_2(\R)$, with the topology inherited from $\M_2(\R)$.
  \item $GL_2(\C)$, with the topology inherited from $\M_2(\C)$.
  \item The space $X$ of all symmetric matrices in $GL_2(\R)$ with both the eigenvalues belonging to the interval $]0, 2[$, with the topology inherited from $\M_2(\R)$.
\end{enumerate}

\qitem{3.8} Let $X$ be a topological space. A mapping $f : X \to \R$ is said to be \emph{lower semi-continuous} if the set $\{x \in X \mid f(x) \le \alpha\}$ is closed for every $\alpha \in \R$. Which of the following statements are true?
\begin{enumerate}[a.]
  \item If $f$ is continuous, then it is lower semi-continuous.
  \item If the set $\{x \in X \mid f(x) > \alpha\}$ is open for every $\alpha \in \Q$, then $f$ is lower semi-continuous.
  \item If $\{f_n\}_{n=1}^{m}$ is a finite collection of lower semi-continuous functions defined on $X$, then $f : X \to \R$ defined by $f(x) = \min_{1 \le n \le m} f_n(x)$ is lower semi-continuous.
\end{enumerate}

\qitem{3.9} Which of the following statements are true?
\begin{enumerate}[a.]
  \item Let $\mathcal{F}$ be an infinite family of continuous real-valued functions on the interval $[0, 1]$ with the property that given any finite subfamily of functions $\mathcal{F}' \subset \mathcal{F}$, there exists at least one point $t \in [0, 1]$ (depending on the subfamily) such that $f(t) = 0$ for all $f \in \mathcal{F}'$. Then, there exists at least one point $t_0 \in [0, 1]$ such that $f(t_0) = 0$ for all $f \in \mathcal{F}$.
  \item Let $\mathcal{F} = \{f \in \mathcal{C}^1[0, 1] \mid |f(t)| \le 1 \text{ and } |f'(t)| \le 1 \text{ for all } t \in [0, 1]\}$. Given any sequence in $\mathcal{F}$, there exists a subsequence which converges uniformly on $[0, 1]$.
  \item If $f : [0, 1] \to \R$ is lower semi-continuous, then there exists $t_0 \in [0, 1]$ such that $f(t_0) = \min_{t \in [0, 1]} f(t)$.
\end{enumerate}

\qitem{3.10} Let $D$ be the closed unit disc in $\R^2$. Let $d(\cdot, \cdot)$ denote the Euclidean distance in $D$. Let $T : D \to D$ be a mapping such that $d(T(x), T(y)) = d(x, y)$ for all points $x$ and $y$ in $D$. Which of the following statements are true?\\
a.\ There exists $x_0 \in D$ such that $T(x_0) = x_0$. \qquad b.\ The image of $T$ is closed. \qquad c.\ The mapping $T$ is surjective.

\examsec{Section 4: Calculus \& Differential Equations}

\qitem{4.1} Evaluate: $\displaystyle\int_0^{\frac{\pi}{4}} \tan^3 x\,dx$.

\qitem{4.2} Evaluate: $\displaystyle\iint_{\R^2} e^{-(19x^2 + 2xy + 19y^2)}\,dx\,dy$.

\qitem{4.3} Evaluate: $\displaystyle\iint_S (x^4 + y^4 + z^4)\,dS$, where $S = \{(x, y, z) \in \R^3 \mid x^2 + y^2 + z^2 = a^2\},\ a > 0$.

\qitem{4.4} Evaluate $f'(3)$, where $f(x) = \displaystyle\int_{-x}^{x} \frac{1 - e^{-xy}}{y}\,dy,\ x > 0$.

\qitem{4.5} Find the maximum value of $3x^2 + 2xy + y^2$ over the set $\{(x, y) \in \R^2 \mid x^2 + y^2 = 1\}$.

\qitem{4.6} Solve: $(x + y)\,dx = (x - y)\,dy$.

\qitem{4.7} Find the general solution of the system of differential equations: $\frac{dx}{dt} = x + 2y$, $\frac{dy}{dt} = 3x + 2y$.

\qitem{4.8} Find a function whose Laplace transform is given by $\dfrac{1}{s^4 + s^2}$.

\qitem{4.9} Find all solutions $(u, \lambda)$ where $\lambda \in \R$ and $u \not\equiv 0$ of the boundary value problem: $-u''(x) = \lambda u(x)$ for $0 < x < 1$, $u'(0) = u(1) = 0$.

\qitem{4.10} Solve the boundary value problem: $-\Delta u = 1$ in $B(0; R) \subset \R^2$, $u = 0$ on $\partial B(0; R)$, where $\Delta$ is the usual Laplace operator in $\R^2$ and $B(0; R)$ is the open ball of radius $R (> 0)$ with centre at the origin.

\examsec{Section 5: Miscellaneous}

\qitem{5.1} Let $n \in \N$ be fixed. Let $C_r = \binom{n}{r}$ for $0 \le r \le n$. Evaluate: $C_0^2 + 3C_1^2 + \dots + (2n + 1)C_n^2$.

\qitem{5.2} Find all positive integers which leave remainders $5, 4, 3$ and $2$ when divided by $6, 5, 4$ and $3$, respectively.

\qitem{5.3} Evaluate: $\displaystyle\sum_{n=0}^{\infty} \frac{5n + 1}{(2n + 1)!}$.

\qitem{5.4} Find the value(s) of $c$ such that the straight line $y = 2x + c$ is tangent to the ellipse $\frac{x^2}{4} + \frac{y^2}{9} = 1$.

\qitem{5.5} Find the equation of the plane passing through the point $(-1, 3, 2)$ and which is perpendicular to the planes $x + 2y + 2z = 5$ and $3x + 3y + 2z = 8$.

\qitem{5.6} Let $k \in \N$ and let $N = (k + 1)^2$. Evaluate the following sum as a closed form expression in terms of $k$: $\sum_{n=1}^{N} [\sqrt{n}]$, where $[x]$ denotes the largest integer less than, or equal to, $x$.

\qitem{5.7} A triangle has perimeter of length $16$ units and a fixed base of length $6$ units. What is the maximum value of the area of the triangle?

\qitem{5.8} Let $B$ be a square of side $3$ units in length. Ten points are marked in it at random. What is the probability that there exist at least two points amongst them which are of distance less than or equal to $\sqrt{2}$ units apart?

\qitem{5.9} Let $X$ be a set and let $V$ be a (real) vector space of real-valued functions defined on $X$, with pointwise addition and scalar multiplication as the operations. Assume in addition that if $f \in V$, then $f^2 \in V$ and $|f| \in V$. Which of the following statements are true?\\
a.\ If $f$ and $g$ are in $V$, then $fg \in V$. \qquad b.\ If $f$ and $g$ are in $V$, then $\max\{f, g\} \in V$. \qquad c.\ If $f \in V$, then $f^3 \in V$.

\qitem{5.10} A portion of a wooden cube is sawed off at each vertex so that a small equilateral triangle is formed at each corner with vertices on the edges of the cube. The $24$ vertices of the new object are all connected to each other by straight lines. How many of these lines (with the exception, of course, of their end-points) lie entirely in the interior of the original cube?

\examsec{Answer Key --- NBHM 2019}
\begin{multicols}{2}
\textbf{Section 1: Algebra}\\
\textbf{1.1} a, b, c\\
\textbf{1.2} $p^{\frac{n(n-1)}{2}}$\\
\textbf{1.3} A subgroup of order $125$. Example: $3 \times 3$ upper triangular matrices with $1$ on the diagonal.\\
\textbf{1.4} $9$\\
\textbf{1.5} a\\
\textbf{1.6} $\frac{p(p-1)}{2}$\\
\textbf{1.7} a, b\\
\textbf{1.8} Any two linearly independent vectors of the form $(u, u, w, u + w)$. Example: $\{(1, 1, 0, 1), (1, 1, -1, 0)\}$.\\
\textbf{1.9} $6$\\
\textbf{1.10} a, b, c

\medskip
\textbf{Section 2: Analysis}\\
\textbf{2.1} $\frac{1}{e}$\\
\textbf{2.2} a, b\\
\textbf{2.3} a, c\\
\textbf{2.4} a, c\\
\textbf{2.5} a, b, c\\
\textbf{2.6} a, b\\
\textbf{2.7} a, c\\
\textbf{2.8} b, c\\
\textbf{2.9} a.\ $\sum_{n=0}^{\infty} \frac{1}{z^{n+1}} + \sum_{n=0}^{\infty} \frac{z^n}{2^{n+1}}$; b.\ $2\pi i$\\
\textbf{2.10} None

\medskip
\textbf{Section 3: Topology}\\
\textbf{3.1} b, c\\
\textbf{3.2} a, c\\
\textbf{3.3} c\\
\textbf{3.4} a, c\\
\textbf{3.5} b\\
\textbf{3.6} a, b, c\\
\textbf{3.7} b, c\\
\textbf{3.8} a, b, c\\
\textbf{3.9} a, b, c\\
\textbf{3.10} a, b, c

\medskip
\textbf{Section 4: Calculus \& Differential Equations}\\
\textbf{4.1} $\frac{1}{2}(1 - \log 2)$\\
\textbf{4.2} $\frac{\pi}{6\sqrt{10}}$\\
\textbf{4.3} $\frac{12\pi}{5}a^6$\\
\textbf{4.4} $\frac{2}{3}(e^9 - e^{-9})$\\
\textbf{4.5} $2 + \sqrt{2}$\\
\textbf{4.6} $\tan^{-1}\left(\frac{y}{x}\right) = \log\sqrt{x^2 + y^2} + c$\\
\textbf{4.7} $x(t) = 2c_1 e^{4t} + c_2 e^{-t},\ y(t) = 3c_1 e^{4t} - c_2 e^{-t}$\\
\textbf{4.8} $x - \sin x$\\
\textbf{4.9} $(\lambda_n, u_n)$, $n \in \N \cup \{0\}$, where $\lambda_n = (2n+1)^2\frac{\pi^2}{4}$, $u_n(x) = \cos(2n+1)\frac{\pi}{2}x$\\
\textbf{4.10} $u(x, y) = \frac{R^2 - x^2 - y^2}{4}$

\medskip
\textbf{Section 5: Miscellaneous}\\
\textbf{5.1} $(n + 1)\frac{(2n)!}{(n!)^2}$\\
\textbf{5.2} $60k + 59,\ k \ge 0$\\
\textbf{5.3} $\frac{e}{2} + \frac{2}{e}$\\
\textbf{5.4} $\pm 5$\\
\textbf{5.5} $2x - 4y + 3z + 8 = 0$\\
\textbf{5.6} $\frac{1}{6}(k + 1)(4k^2 + 5k + 6)$\\
\textbf{5.7} $12$ sq.\ units\\
\textbf{5.8} $1$\\
\textbf{5.9} a, b, c\\
\textbf{5.10} $120$

\medskip
\noindent\emph{Note: Please accept any correct equivalent form of the answers.}
\end{multicols}
% ==================== NBHM 2020 ====================
\examheader{NBHM 2020}{Doctoral Scholarship Scheme 2020 --- Written Test}{Saturday, 25th January 2020}{Time Allowed: 180 Minutes (Three Hours) \quad|\quad 36 Questions (3 Sections)}

\subsection*{Instructions}
\begin{itemize}
  \item Test booklet consists of 6 pages of questions and 1 cover page; 36 questions distributed over 3 sections; answer all of them.
  \item Each question carries 4 marks. There is \textbf{negative marking} in Section C (but not in Sections A and B).
  \item Fill in only your final answer in the box provided. Mobile phones, calculators, and smart watches are prohibited.
\end{itemize}

\examsec{Section A (Questions 1 to 18)}
\noindent\emph{18 short-answer questions, 4 marks each, no partial credit, no negative marking.}

\qitem{(1)} Find rational numbers $a, b,$ and $c$ such that $(1 + \sqrt[3]{2})^{-1} = a + b\sqrt[3]{2} + c\sqrt[3]{2}^2$.

\qitem{(2)} Let $u$ and $v$ be the real and imaginary parts respectively of the function $f(z) = 1/(z^2 - 6z + 8)$ of a complex variable $z = x + iy$. Let $C$ be the simple closed curve $|z| = 3$ oriented in the counter clockwise direction. Evaluate: $\displaystyle\oint_C u\,dy + v\,dx$.

\qitem{(3)} A point is moving along the curve $y = x^2$ with unit speed. What is the magnitude of its acceleration at the point $(1/2, 1/4)$?

\qitem{(4)} Evaluate $\displaystyle\int_{-\infty}^{\infty} (1 + 2x^4)e^{-x^2}\,dx$.

\qitem{(5)} Let $p(x)$ be the minimal polynomial of $\sqrt{2} + \sqrt{-2}$ over the field $\Q$ of rational numbers. Evaluate $p(\sqrt{2})$.

\qitem{(6)} Find the volume of the tetrahedron in $\R^3$ bounded by the coordinate planes $x = 0$, $y = 0$, $z = 0$, and the tangent plane at the point $(4, 5, 5)$ to the sphere $(x - 3)^2 + (y - 3)^2 + (z - 3)^2 = 9$.

\qitem{(7)} From the collection of all permutation matrices of size $10 \times 10$, one such matrix is randomly picked. What is the expected value of its trace?

\qitem{(8)} You are given $20$ identical balls and $5$ bins that are coloured differently (so that any two of the bins can be distinguished from each other). In how many ways can the balls be distributed into the bins in such a way that each bin has at least two balls?

\qitem{(9)} Let $G$ be the symmetric group $S_5$ of permutations of five symbols. Consider the set $\mathcal{S}$ of subgroups of $G$ that are isomorphic to the non-cyclic group of order $4$. Let us call two subgroups $H$ and $K$ belonging to $\mathcal{S}$ as \emph{equivalent} if they are conjugate (that is, there exists $g \in G$ such that $gHg^{-1} = K$). How many equivalence classes are there in $\mathcal{S}$?

\qitem{(10)} Let $M$ be a $7 \times 6$ real matrix. The entries of $M$ in the positions $(1, 3), (1, 4), (3, 3), (3, 4),$ and $(5, 4)$ are changed to obtain another $7 \times 6$ real matrix $\widetilde{M}$. Suppose that the rank of $\widetilde{M}$ is $4$. What could be the rank of $M$? List all possibilities.

\qitem{(11)} What are the maximum and minimum values of $x + y$ in the region $S = \{(x, y) : x^2 + 4y^2 \le 1\}$?

\qitem{(12)} Let $k$ be the field obtained by adjoining to the field $\Q$ of rational numbers the roots of the polynomial $x^4 - 2$. Let $k'$ be the field obtained by adjoining to $k$ the roots of the polynomial $x^4 + 2$. What is the degree of $k'$ over $k$?

\qitem{(13)} Evaluate the (absolute value of the) surface integral $\left|\int_S \overline{F}\cdot d\overline{A}\right|$ of the vector field $\overline{F}(x, y, z) := (e^y, 0, e^x)$ on the surface $S := \{(x, y, z) \mid x^2 + y^2 = 25,\ 0 \le z \le 2,\ x \ge 0,\ y \ge 0\}$.

\qitem{(14)} Let $k$ be a field with five elements. Let $V$ be the $k$-vector space of $5 \times 1$ matrices with entries in $k$. Let $S$ be a subset of $V$ such that $u^t v = 0$ for all $u$ and $v$ in $S$: here $u^t$ denotes the transpose of $u$ and $u^t v$ the usual matrix product. What is the maximum possible cardinality of $S$?

\qitem{(15)} Suppose $f : \C \to \C$ is a holomorphic function such that the real part of $f''(z)$ is strictly positive for all $z \in \C$. What is the maximum possible number of solutions of the equation $f(z) = az + b$, as $a$ and $b$ vary over complex numbers?

\qitem{(16)} What is the expected minimum number of tosses of a fair coin required to get both heads and tails each at least once?

\qitem{(17)} How many real solutions does the equation $f(x) = 0$ have, where $f(x) := \displaystyle\sum_{i=1}^{2020} \frac{i^2}{x - i}$?

\qitem{(18)} Let $SL_2(\Z)$ denote the group (under usual matrix multiplication) of $2 \times 2$ matrices with integer entries and determinant $1$. Let $H$ be the subgroup of $SL_2(\Z)$ consisting of those matrices such that: (i) the diagonal entries are all equivalent to $1 \pmod{3}$; (ii) the off-diagonal entries are all divisible by $3$. What is the index of $H$ in $SL_2(\Z)$?

\examsec{Section B (Questions 19--22)}
\noindent\emph{4 questions, 4 parts each, 1 mark per part, partial credit allowed, no negative marking.}

\qitem{(19)} For each of the following numbers $q$ in turn, consider a field $k$ of order $q$. In each case, determine the number of elements $\alpha$ in $k$ such that the smallest subfield of $k$ containing $\alpha$ is $k$ itself.\\
(a) $2^4$ \qquad (b) $3^5$ \qquad (c) $5^{10}$ \qquad (d) $7^{12}$

\qitem{(20)} Let $B_r$ denote the closed disk $\{z \in \C : |z| \le r\}$. State whether $\infty$ is a removable singularity (RS), pole (P), or essential singularity (ES) in each of the following cases. There may be more than one possibility in each case.
\begin{enumerate}[(a)]
  \item $f$ is a non-constant polynomial in $z$.
  \item $f(z) = \frac{p(z)}{q(z)}$, where $p, q$ are non-zero polynomials of the same degree.
  \item $f$ is an entire function for which $f^{-1}(B_1)$ is bounded.
  \item $f$ is an entire function for which $f^{-1}(B_r)$ is bounded for all $r > 0$.
\end{enumerate}

\qitem{(21)} Let $R := \Z/2020\Z$ be the quotient ring of the integers $\Z$ by the ideal $2020\Z$.
\begin{enumerate}[(a)]
  \item What is the number of ideals in $R$?
  \item What is the number of units in $R$?
  \item What is the number of elements $r$ in $R$ such that $r^n = 1$ for some integer $n \ge 1$?
  \item What is the number of elements $r$ in $R$ such that $r^n = 0$ for some integer $n \ge 1$?
\end{enumerate}

\qitem{(22)} Let $X$ be a three element set. For each of the following numbers $n$, determine the number of distinct homeomorphism classes of topologies on $X$ with exactly $n$ open subsets (including the empty set and the whole set).\\
(a) $3$ \qquad (b) $4$ \qquad (c) $5$ \qquad (d) $7$

\examsec{Section C (Questions 23--36)}
\noindent\emph{14 questions, 4 parts each, 1 mark per correct part, $-1$ mark penalty per incorrect part.}

\qitem{(23)} Consider $f : \R^2 \to \R$ defined by $f(x, y) = x + y$. State whether true or false:
\begin{enumerate}[(a)]
  \item Image under $f$ of any open set is open.
  \item Image under $f$ of any closed set is closed.
  \item Image under $f$ of any dense set is dense.
  \item Image under $f$ of any discrete set is discrete.
\end{enumerate}

\qitem{(24)} Listed below are four subsets of $\C^2$. Write ``Bounded'' or ``Unbounded'' ($\Re(z)$ denotes the real part of $z$):
\begin{enumerate}[(a)]
  \item $\{(z, w) \in \C^2 : z^2 + w^2 = 1\}$
  \item $\{(z, w) \in \C^2 : |\Re(z)|^2 + |\Re(w)|^2 = 1\}$
  \item $\{(z, w) \in \C^2 : |z|^2 + |w|^2 = 1\}$
  \item $\{(z, w) \in \C^2 : |z|^2 - |w|^2 = 1\}$
\end{enumerate}

\qitem{(25)} Write ``convergent'' or ``divergent'' for each series:\\
(a) $\displaystyle\sum_{n \ge 2} \frac{1}{n\log n}$ \quad
(b) $\displaystyle\sum_{n \ge 2} \frac{\log^2 n}{n^2}$ \quad
(c) $\displaystyle\sum_{n \ge 2} \frac{1}{n\log^2 n}$ \quad
(d) $\displaystyle\sum_{n \ge 2} \frac{\sqrt{n+1} - \sqrt{n}}{n}$

\qitem{(26)} Let $T$ be a nilpotent linear operator on $\R^5$. Let $d_i$ denote the dimension of the kernel of $T^i$. Which of the following can possibly occur as $(d_1, d_2, d_3)$? (Yes/No)\\
(a) $(1, 2, 3)$ \qquad (b) $(2, 3, 5)$ \qquad (c) $(2, 2, 4)$ \qquad (d) $(2, 4, 5)$

\qitem{(27)} For $n$ a positive integer, let $\Q/n\Z$ be the quotient of $(\Q, +)$ by $n\Z$. State true or false:
\begin{enumerate}[(a)]
  \item Every element of $\Q/n\Z$ is of finite order.
  \item There are only finitely many elements in $\Q/n\Z$ of any given finite order.
  \item Every proper subgroup of $\Q/n\Z$ is finite.
  \item $\Q/2\Z$ and $\Q/5\Z$ are isomorphic as groups.
\end{enumerate}

\qitem{(28)} Let $\mathcal{S} := \{f : (0, \infty) \to \R \text{ continuous} \mid f(x) = f(2x)\ \forall x \in (0, \infty)\}$. State true or false:
\begin{enumerate}[(a)]
  \item Any element $f \in \mathcal{S}$ is bounded.
  \item Any element $f \in \mathcal{S}$ is uniformly continuous.
  \item Any element $f \in \mathcal{S}$ is differentiable.
  \item Any uniformly bounded sequence in $\mathcal{S}$ has a uniformly converging subsequence.
\end{enumerate}

\qitem{(29)} Let $B_1 := \{z \in \C : |z| \le 1\}$, and let $C^0(B_1, \C)$ be the space of continuous complex-valued functions on $B_1$ with uniform convergence topology. Write ``Dense'' or ``Not dense'':
\begin{enumerate}[(a)]
  \item Restrictions to $B_1$ of polynomials in $z$.
  \item Restrictions to $B_1$ of polynomials in $z$ and $\overline{z}$.
  \item The set of smooth functions $f : B_1 \to \C$ that vanish on the boundary $\partial B_1$.
  \item The set of smooth functions $f : B_1 \to \C$ whose normal derivative vanishes along the boundary $\partial B_1$.
\end{enumerate}

\qitem{(30)} Let $p$ be a prime number, and let $S = [0, 1] \cap \{q/p^n \mid q \in \Z,\ n \in \Z_{\ge 0}\}$ with subspace topology from $[0, 1]$. State true or false:
\begin{enumerate}[(a)]
  \item Any bounded function on $S$ uniquely extends to a bounded function on $[0, 1]$.
  \item Any continuous function on $S$ uniquely extends to a continuous function on $[0, 1]$.
  \item Any uniformly continuous function on $S$ uniquely extends to a uniformly continuous function on $[0, 1]$.
  \item Any bounded continuous function on $S$ uniquely extends to a bounded continuous function on $[0, 1]$.
\end{enumerate}

\qitem{(31)} Let $X = GL_2(\R) \subset \M_2(\R) \cong \R^4$ with subspace topology. Which of the following spaces is obtained as the image under a continuous surjection from $X$? (Yes/No)
\begin{enumerate}[(a)]
  \item the real line $\R$.
  \item the subspace $\{(x, 1/x) \mid x \in \R,\ x \neq 0\}$ of $\R^2$.
  \item the complement in $\R^2$ of the set $\{(x, 1/x) \mid x \in \R,\ x \neq 0\}$.
  \item the closed disk $\{(x, y) \mid x^2 + y^2 \le 1\}$ of $\R^2$.
\end{enumerate}

\qitem{(32)} For $n$ a positive integer, let $f_n : \R \to \R$ be defined by $f_n(x) = x/(1 + nx^2)$. State true or false:
\begin{enumerate}[(a)]
  \item The sequence $\{f_n(x)\}$ of functions converges uniformly on $\R$.
  \item The sequence $\{f_n(x)\}$ of functions converges uniformly on $[1, b]$ for any $b > 1$.
  \item The sequence $\{f_n'(x)\}$ of derivatives converges uniformly on $\R$.
  \item The sequence $\{f_n'(x)\}$ of derivatives converges uniformly on $[1, b]$ for any $b > 1$.
\end{enumerate}

\qitem{(33)} For which of the following subspaces $X$ of $\R$ does every continuous surjective map $f : X \to X$ have a fixed point? (Yes/No)\\
(a) $[1, 2]$ \qquad (b) $[1, 2] \cup [3, 7]$ \qquad (c) $[3, \infty)$ \qquad (d) $[1, 2] \cup [3, \infty)$

\qitem{(34)} Let $A$ be an arbitrary real $5 \times 5$ matrix row equivalent to
\[
  R = \begin{pmatrix} 1 & 0 & 0 & -3 & -1 \\ 0 & 1 & 0 & -2 & -1 \\ 0 & 0 & 1 & -1 & -1 \\ 0 & 0 & 0 & 0 & 0 \\ 0 & 0 & 0 & 0 & 0 \end{pmatrix}.
\]
State whether each statement is true or false:
\begin{enumerate}[(a)]
  \item The first two rows of $A$ are linearly independent.
  \item The last four rows of $A$ generate a space of dimension at least $2$.
  \item The first two columns of $A$ are linearly independent.
  \item The last four columns of $A$ generate a space of dimension $3$.
\end{enumerate}

\qitem{(35)} Consider the space $X := \R^{[0,1]}$ of real valued functions on $[0, 1]$ given the product topology. Write ``Closed'' or ``Not closed'':
\begin{enumerate}[(a)]
  \item The subset consisting of all continuous functions.
  \item The subset consisting of functions that take integer values everywhere.
  \item The subset consisting of all unbounded functions.
  \item The subset consisting of all bounded functions.
\end{enumerate}

\qitem{(36)} Let $X$ be the space of all real polynomials $a_5 t^5 + a_4 t^4 + a_3 t^3 + a_2 t^2 + a_1 t + a_0$ of degree at most $5$, identified with $\R^6$. Write ``Connected'' or ``Disconnected'':
\begin{enumerate}[(a)]
  \item All polynomials in $X$ that do not vanish at $t = 2$.
  \item All polynomials in $X$ whose derivatives vanish at $t = 3$.
  \item All polynomials in $X$ that vanish at both $t = 4$ and $t = 5$.
  \item All polynomials in $X$ that are increasing (as functions from $\R$ to $\R$).
\end{enumerate}

\examsec{Answer Key --- NBHM 2020}
\begin{multicols}{2}
\textbf{Section A}\\
\textbf{(1)} $a = 1/3,\ b = -1/3,\ c = 1/3$\\
\textbf{(2)} $-\pi$\\
\textbf{(3)} $1/\sqrt{2}$ or $\sqrt{2}/2$\\
\textbf{(4)} $5\sqrt{\pi}/2$\\
\textbf{(5)} $20$\\
\textbf{(6)} $576$\\
\textbf{(7)} $1$\\
\textbf{(8)} $1001$\\
\textbf{(9)} $2$\\
\textbf{(10)} $2, 3, 4, 5, 6$\\
\textbf{(11)} maximum $= \sqrt{5}/2$, minimum $= -\sqrt{5}/2$\\
\textbf{(12)} $1$\\
\textbf{(13)} $2(e^5 - 1)$\\
\textbf{(14)} $5^2 = 25$\\
\textbf{(15)} $2$\\
\textbf{(16)} $3$\\
\textbf{(17)} $2019$\\
\textbf{(18)} $24$

\medskip
\textbf{Section B}\\
\textbf{(19)} (a) $2^4 - 2^2 = 12$; (b) $3^5 - 3^1 = 240$; (c) $5^{10} - 5^2 - 5^5 + 5^1$; (d) $7^{12} - 7^6 - 7^4 + 7^2$\\
\textbf{(20)} (a) P; (b) RS; (c) RS, P; (d) P\\
\textbf{(21)} (a) $12$; (b) $800$; (c) $800$; (d) $2$\\
\textbf{(22)} (a) $2$; (b) $2$; (c) $2$; (d) $0$

\medskip
\textbf{Section C}\\
\textbf{(23)} (a) True; (b) False; (c) True; (d) False\\
\textbf{(24)} (a) Unbounded; (b) Unbounded; (c) Bounded; (d) Unbounded\\
\textbf{(25)} (a) divergent; (b) convergent; (c) convergent; (d) convergent\\
\textbf{(26)} (a) Yes; (b) No; (c) No; (d) Yes\\
\textbf{(27)} (a) True; (b) True; (c) False; (d) True\\
\textbf{(28)} (a) True; (b) False; (c) False; (d) False\\
\textbf{(29)} (a) Not dense; (b) Dense; (c) Not dense; (d) Dense\\
\textbf{(30)} (a) False; (b) False; (c) True; (d) False\\
\textbf{(31)} (a) Yes; (b) Yes; (c) No; (d) Yes\\
\textbf{(32)} (a) True; (b) True; (c) False; (d) True\\
\textbf{(33)} (a) Yes; (b) No; (c) Yes; (d) Yes\\
\textbf{(34)} (a) False; (b) True; (c) True; (d) True\\
\textbf{(35)} (a) Not closed; (b) Closed; (c) Not closed; (d) Not closed\\
\textbf{(36)} (a) Disconnected; (b) Connected; (c) Connected; (d) Connected
\end{multicols}
% ==================== NBHM 2021 ====================
\examheader{NBHM 2021}{Doctoral Scholarship Scheme 2021 --- Written Test}{Sunday, 11th April 2021 (Held: Wednesday, 18th August 2021)}{Time Allowed: 180 Minutes (Three Hours) \quad|\quad 39 Questions (2 Sections)}

\subsection*{Instructions}
\begin{itemize}
  \item 39 questions distributed over two sections; 180 minutes; questions arranged randomly.
  \item Section A: 30 short-answer questions (2 marks each, no negative marking). Section B: 9 True/False questions comprising 20 assertions (1 mark each, $-1$ mark penalty for each incorrect response).
\end{itemize}

\examsec{Section A (Questions 1 to 30) --- Short Answer Type}

\qitem{(1)} A smooth function $y$ of a single variable satisfies the differential equation $y'' - 2y' + y = 4$ and the initial conditions $y(0) = 5$, $y'(0) = 2$. Find $y(1)$.

\qitem{(2)} Given that the following Laurent series expansion is valid in the annulus $1 < |z| < 2$, what is $a_3$?
\[
  \frac{1}{(z - 1)(z - 2)} = \sum_{n=-\infty}^{\infty} a_n z^n
\]

\qitem{(3)} Let $\R^3$ be the real vector space consisting of $1 \times 3$ real matrices. Let $A$ be the $3 \times 3$ real matrix with the property that the linear transformation $x \mapsto xA$ from $\R^3$ to itself projects every vector $x$ in $\R^3$ orthogonally onto the line in the direction of $(1, 0, 1)$. What is the sum of the entries of $A$?

\qitem{(4)} Let $G$ denote the group of rational numbers $\Q$ with respect to addition. What is the cardinality of the automorphism group of $G$?\\
(a) $1$ (that is, the only automorphism is the identity map) \quad (b) finite but not $1$ \quad (c) countable (but not finite) \quad (d) uncountable

\qitem{(5)} Given a positive real number $x$, a sequence $\{a_n(x)\}_{n \ge 1}$ is defined as follows: $a_1(x) := x$ and $a_n(x) := x^{a_{n-1}(x)}$ recursively for all $n \ge 2$. Determine the largest value of $x$ for which $\lim_{n \to \infty} a_n(x)$ exists.

\qitem{(6)} Among $3 \times 3$ invertible matrices with entries in the finite field $\Z/3\Z$ containing $3$ elements, how many are similar to $\begin{pmatrix} 2 & 0 & 0 \\ 0 & 2 & 0 \\ 0 & 0 & 1 \end{pmatrix}$?

\qitem{(7)} For any integer $n$, let $I_n$ denote the ideal $\{m \in \Z \mid m^r \in n\Z \text{ for some } r \ge 1\}$. What is the cardinality of the quotient ring $\Z/(I_{63} \cap I_{84})$?

\qitem{(8)} Let $A$ be an invertible $3 \times 3$ real matrix such that $A$ and $A^2$ have the same characteristic polynomial. What are the possible traces of such a matrix $A$?

\qitem{(9)} Let $v$ be a fixed non-zero vector of an $n$-dimensional real vector space $V$. Let $\mathcal{S}(v)$ be the subspace of the vector space of linear operators on $V$ consisting of those operators that admit $v$ as an eigenvector. What is the dimension of $\mathcal{S}(v)$ as a real vector space (in terms of $n$)?

\qitem{(10)} Find the largest positive real number $\delta$ such that, for all real numbers $x$ and $y$, we have $|\cos x - \cos y| < \sqrt{2}$ whenever $|x - y| < \delta$.

\qitem{(11)} A particle is at the origin of the Cartesian coordinate plane to begin with. At the end of every second, it jumps one unit either to the East or to the West or to the North or to the South with equal probability. What is the probability that the particle is back at the origin after six seconds?

\qitem{(12)} As $f(x)$ varies over all continuously differentiable functions from $\R$ to $\R$ with the property that $f(0) = 10$ and $f(1) = 0$, find the infimum of $\int_0^1 \sqrt{1 + f'(x)^2}\,dx$.

\qitem{(13)} A particle is moving on the $x$-axis such that $\frac{dx}{dt} = (x - 1)(x + 2)(x - 3)$. The particle is so positioned initially that it does not wander off to infinity. Which point of equilibrium will it be close to after a sufficiently long time?

\qitem{(14)} Let $G$ be the group of homeomorphisms of the real line $\R$ with its usual topology, the group operation being composition. Consider the elements $f(x) = 2x$ and $g(x) = 8x$ of $G$. Suppose $h(x)$ in $G$ is such that it conjugates $f$ to $g$, that is, $(h \circ f \circ h^{-1})(x) = g(x)$, and further satisfies $h(1) = 5$. What is $h(2)$?

\qitem{(15)} An element $a$ of a ring $R$ is said to be an \emph{idempotent} if $a^2 = a$. Note that $0$ and $1$ are idempotents. How many idempotents (including $0$ and $1$) are there in the ring $\Z/120\Z$?

\qitem{(16)} Consider the Taylor expansion of the function $\frac{1}{1+x^3}$ in powers of $x - \frac{1}{2}$: $\frac{1}{1+x^3} = \sum_{n \ge 0} a_n\left(x - \frac{1}{2}\right)^n$. What is the radius of convergence of this series?

\qitem{(17)} Let $n$ be a positive integer and $\mathcal{P}$ the real vector space of polynomials with real coefficients and degree at most $n$. Let $T$ be the linear operator on $\mathcal{P}$ defined by $Tf(x) = xf'(x) - f(x - 1)$, where $f'(x)$ is the derivative of $f(x)$. What is the trace of $T$?

\qitem{(18)} Let $D$ be the disc in the complex plane centred at the point $\pi/4$ and of radius $r$. Let $D'$ be the image of this disk under the map $z \mapsto \sin z$. Evaluate $\lim_{r \to 0} \frac{\operatorname{Area}(D')}{\operatorname{Area}(D)}$.

\qitem{(19)} A prime $p$ is such that $1/p$ when represented in octal (base $8$) notation equals $0.\overline{00331}$ (meaning that the digits $0, 0, 3, 3, 1$ get repeated ad infinitum). What is the order of $2$ in the group of units of the ring of integers modulo $p$?

\qitem{(20)} Let $M_5(\C)$ be the set of $5 \times 5$ complex matrices and let $A$ be a fixed element of $M_5(\C)$. Consider the linear transformation $T_A : M_5(\C) \to M_5(\C)$ given by $X \mapsto AX$. If $A$ has rank $2$, what is the rank of $T_A$?

\qitem{(21)} Let $\Z$ denote the set of integers. For $c$ and $r$ in $\Z$, define $B(c, r) := \{c + kr \mid k \in \Z\}$. As $c$ varies over all integers and $r$ over all positive integers, the sets $B(c, r)$ form a basis for a topology on $\Z$. Does the limit $\lim_{n \to \infty} (n! - 2)^2$ exist with respect to this topology? If so, write the value of the limit; if not, write ``No''.

\qitem{(22)} Let $S := \R[x]$ denote the polynomial ring in one variable over the field $\R$ of real numbers. Find a monic polynomial of least degree in $S$ that is a square root of $-4$ modulo the ideal generated by $(x^2 + 1)^2$.

\qitem{(23)} Let $M_n(\R)$ be the real vector space of $n \times n$ matrices with entries in $\R$. Consider the subset
\[
  M := \left\{A = (a_{ij})_{1 \le i, j \le n} \in M_n(\R) \;\middle|\; \sum_{j=1}^{n} a_{ij} = 0 \text{ for every } i \text{ and } \sum_{i=1}^{n} a_{ij} = 0 \text{ for every } j\right\}.
\]
What is the dimension of $M$ (as a real vector space)?

\qitem{(24)} Consider the entire function $f(z) = z(z - i)$. Put $S := \left\{\frac{1}{|f(z)|} \;\middle|\; |z| \ge 2\right\}$. At what value(s) of $z$ is the maximum of the set $S$ attained?

\qitem{(25)} Consider the following system of linear equations over the field $\Z/5\Z$. How many solutions does it have?
\[
  \begin{pmatrix} 1 & 1 & 0 & 0 & 1 \\ 1 & 1 & 1 & 1 & 3 \\ 1 & 1 & 4 & 4 & 4 \end{pmatrix}
  \begin{pmatrix} x_1 \\ x_2 \\ x_3 \\ x_4 \\ x_5 \end{pmatrix} =
  \begin{pmatrix} 1 \\ 0 \\ 2 \end{pmatrix}
\]

\qitem{(26)} Let $f(x)$ be an irreducible polynomial over $\Q$ of degree $2021$ and let $K$ be the field obtained by adjoining to $\Q$ a root $\alpha$ of $f(x)$. Suppose that the Galois group over $\Q$ of the splitting field of $f(x)$ is abelian. Then how many fields $F$ are there such that $\Q \subseteq F \subseteq K$? (Note: $43 \times 47$ is the prime factorisation of $2021$.)\\
(a) $2$ \qquad (b) $3$ \qquad (c) $4$ \qquad (d) Cannot be determined from the given data.

\qitem{(27)} Find the smallest positive real number $k$ such that, given any finite set $z_1, \dots, z_n$ of complex numbers, all with strictly positive real and imaginary parts, the inequality $|z_1 + \dots + z_n| \ge \frac{1}{k}(|z_1| + \dots + |z_n|)$ holds.

\qitem{(28)} A ``Fibonacci like'' sequence $\{a_n\}_{n \ge 0}$ is defined as $a_0 = 1$, $a_1 = 1$, and $a_n = 2a_{n-2} + a_{n-1}$ for all $n \ge 2$. Find a closed form formula for $a_n$ (that works for all $n \ge 0$).

\qitem{(29)} Find the real number $a$ such that $\displaystyle\oint_{|z-i|=1} \frac{dz}{z^2 - z + a} = \pi$.

\qitem{(30)} Let $\mathcal{L}$ be the set of all $k \times k$ real matrices $A$ such that $\lim_{n \to \infty} A^n$ exists as a $k \times k$ real matrix. Let $\mathcal{P}$ be the subset of $\mathcal{L}$ consisting of all matrices in $\mathcal{L}$ that are symmetric and positive definite. What is the image of the map $\varphi : \mathcal{P} \to \R$ defined by $\varphi(A) = \operatorname{trace}(A)$?

\examsec{Section B (Questions 31--39) --- True or False}

\qitem{(31)} Let $A$ be a real $n \times n$ symmetric matrix with distinct eigenvalues. Let $B$ be a real $n \times n$ matrix that commutes with $A$. Then:\\
(a) The eigenvalues of $B$ are real. \qquad (b) Every eigenvector of $A$ is also an eigenvector for $B$. \qquad (c) $B$ is diagonalizable.

\qitem{(32)} For $f : (0, 1) \to \R$ a real valued function on the open interval $(0, 1)$:
\begin{enumerate}[(a)]
  \item Suppose that $f$ is uniformly continuous. Then the image under $f$ of a Cauchy sequence is a Cauchy sequence.
  \item Suppose that, for every Cauchy sequence in $(0, 1)$, its image under $f$ is a Cauchy sequence. Then $f$ is uniformly continuous.
\end{enumerate}

\qitem{(33)} (a) Every bounded continuous real valued function on $\Q^2$ can be extended to a bounded continuous function on $\R^2$.\\
(b) There exists a non-finite subset $X$ of $\Q^2$ such that every real valued continuous function on $X$ is bounded.

\qitem{(34)} Let $C(\R)$ and $C[0, 1]$ be the rings of continuous real valued functions on $\R$ and $[0, 1]$ respectively (with pointwise operations).
\begin{enumerate}[(a)]
  \item Let $\pi : C(\R) \to C[0, 1]$ be the ring homomorphism given by restriction: $\pi(f) := f|_{[0,1]}$ for $f \in C(\R)$. The kernel of $\pi$ is a prime ideal in $C(\R)$.
  \item The ideal $I := \{f \in C(\R) \mid \text{there exists } N \in \R \text{ such that } f(x) = 0 \text{ for all } x > N\}$ is a maximal ideal in $C(\R)$.
\end{enumerate}

\qitem{(35)} Let $f : \R \to \R$ be an arbitrary continuous function such that the inverse image of any bounded set is bounded. Then:\\
(a) The image under $f$ of a closed set is closed. \qquad (b) The image of $f$ is all of $\R$ (i.e., $f$ is surjective).

\qitem{(36)} Let $a, b,$ and $c$ be three non-collinear points in the complex plane. Let $f(z) := |z - a|\cdot|z - b|\cdot|z - c|$. Let $M$ be the maximum value of $f(z)$ on the (closed) triangle $T$ with vertices $a, b,$ and $c$. Determine whether it holds always:\\
(a) $f(\alpha) \ge M$ where $\alpha$ is the centroid of $T$. \qquad (b) $f(\alpha) \ge M$ where $\alpha$ is the orthocentre of $T$.

\qitem{(37)} Let $f : \R \to \R$ be an arbitrary continuous function. Define $S := \{y \in \R \mid \text{there exists a sequence } x_n \text{ in } \R \text{ with } \lim_{n \to \infty} x_n = \infty \text{ and } \lim_{n \to \infty} f(x_n) = y\}$. Then:\\
(a) $S$ is closed. \qquad (b) $S$ is connected.

\qitem{(38)} (a) Every infinite group has infinitely many subgroups.\\
(b) Every countable group has only countably many subgroups.\\
(c) Every uncountable group has uncountably many subgroups.

\qitem{(39)} (a) $\sum_{n \ge 1} \sqrt{\frac{a_n}{n}}$ converges for every sequence $\{a_n\}_{n \ge 1}$ of positive real numbers such that $\sum_{n \ge 1} a_n$ converges.\\
(b) $\sum_{n \ge 1} \frac{\sqrt{a_n}}{n}$ converges for every sequence $\{a_n\}_{n \ge 1}$ of positive real numbers such that $\sum_{n \ge 1} a_n$ converges.

\examsec{Answer Key --- NBHM 2021}
\begin{multicols}{2}
\textbf{Section A}\\
\textbf{(1)} $2e + 4$\\
\textbf{(2)} $-1/16$\\
\textbf{(3)} $2$\\
\textbf{(4)} (c) / Countable\\
\textbf{(5)} $e^{1/e}$\\
\textbf{(6)} $117$\\
\textbf{(7)} $42$\\
\textbf{(8)} $0, 3$\\
\textbf{(9)} $n^2 - n + 1$\\
\textbf{(10)} $\pi/2$\\
\textbf{(11)} $25/256$\\
\textbf{(12)} $\sqrt{101}$\\
\textbf{(13)} $1$\\
\textbf{(14)} $40$\\
\textbf{(15)} $8$\\
\textbf{(16)} $\sqrt{3}/2$\\
\textbf{(17)} $(n + 1)(n - 2)/2$\\
\textbf{(18)} $1/2$\\
\textbf{(19)} $15$\\
\textbf{(20)} $10$\\
\textbf{(21)} $4$\\
\textbf{(22)} $x^3 + 3x$\\
\textbf{(23)} $n^2 - 2n + 1$\\
\textbf{(24)} $2i$\\
\textbf{(25)} $125$\\
\textbf{(26)} (c) ($4$ is the answer)\\
\textbf{(27)} $\sqrt{2}$\\
\textbf{(28)} $\dfrac{2^{n+1} + (-1)^n}{3}$\\
\textbf{(29)} $5/4$\\
\textbf{(30)} $(0, k]$

\medskip
\textbf{Section B}\\
\textbf{(31)} (a) True; (b) True; (c) True\\
\textbf{(32)} (a) True; (b) True\\
\textbf{(33)} (a) False; (b) True\\
\textbf{(34)} (a) False; (b) False\\
\textbf{(35)} (a) True; (b) False\\
\textbf{(36)} (a) False; (b) False\\
\textbf{(37)} (a) True; (b) True\\
\textbf{(38)} (a) True; (b) False; (c) True\\
\textbf{(39)} (a) False; (b) True
\end{multicols}
% ==================== NBHM 2022 ====================
\examheader{NBHM 2022}{Doctoral Scholarship Scheme 2022 --- Written Test}{Saturday, 18 June 2022, 10:30 A.M. to 1:30 P.M.}{Time Allowed: 180 Minutes \quad|\quad Maximum Score: 150 \quad|\quad 34 Questions}

\subsection*{Instructions}
\begin{itemize}
  \item 34 questions distributed over two parts (Part A and Part B); maximum possible score is 150.
  \item Part A: 25 short-answer questions (4 marks each, no partial credit, no negative marking).
  \item Part B: 9 True/False questions comprising 25 assertions (2 marks per correct response, $-1$ mark penalty per incorrect response).
\end{itemize}

\examsec{Part A (Questions 1 to 25) --- Short Answer Type}

\qitem{(1)} Let $n$ be the smallest positive integer that has $2022$ as a factor and has exactly $2022$ factors (counting both $1$ and $n$ as factors of $n$). Let $d$ be the largest positive integer such that $2^d$ divides $n$. What is $d$?

\qitem{(2)} A box contains $b$ blue balls and $y$ yellow balls. One ball is drawn at random from the box, but is then put back into the box along with $a$ additional balls of the same colour. Another ball is now drawn at random from the box. What is the probability that the first ball drawn was blue, given that the second ball drawn is yellow (in terms of $b, y,$ and $a$)?

\qitem{(3)} The inner surface of a radially symmetric coffee mug is described as follows: with the $x$-, $y$-, and $z$-axes being oriented as usual, consider the portion of the curve $z = x^4$ in the $x$-$z$ plane that is below the horizontal line $z = 1$; then rotate it about the $z$-axis to obtain the surface. Find the volume of the coffee mug.

\qitem{(4)} Consider the surface $S$ in $\R^3$ defined by the equation $z^2 - xy = 0$. What is the number of connected components of its complement (in $\R^3$)?

\qitem{(5)} The complex numbers $a_n$, $n \ge 0$, are defined by the condition that
\[
  \sum_{n \ge 0} a_n(z - 2i)^n = 1 + 2z + 3z^2 + \dots = \sum_{n \ge 0} (n + 1)z^n
\]
holds for each $z$ in $\C$ such that $|z| < 1$ and $\Im(z) > 0$. What is $\limsup \sqrt[n]{|a_n|}$?

\qitem{(6)} Find the largest positive real number $r$ such that $|\Im\sin(z)| \le 1$ for all $z \in \C$ with $|z| \le r$. If there is no such largest $r$, write ``infinity''.

\qitem{(7)} Given two subsets $C$ and $S$ of a vector space $V$, their sum is $C + S := \{c + s \mid c \in C,\ s \in S\}$. Take $V = \R^3$, $C := \{(x, y, z) \in \R^3 \mid |x| \le 1,\ |y| \le 1,\ |z| \le 1\}$ and $S := \{(x, y, z) \in \R^3 \mid x^2 + y^2 + z^2 \le 1\}$. Find the volume of their sum $C + S$.

\qitem{(8)} Compute $\displaystyle\iint_A (3x^2 y - y^3)\,dxdy$ where $A$ is the region $\{(x, y) \mid x^2 + (y - 1)^2 \le 1\}$.

\qitem{(9)} An element of the alternating group $A_5$ is picked at random (each element is equally likely to be picked). What is the expected value of its order?

\qitem{(10)} On the set $\mathcal{M}$ of all $m \times n$ real matrices, two matrices $A$ and $B$ are equivalent if there exist a real invertible $m \times m$ matrix $P$ and a real invertible $n \times n$ matrix $Q$ such that $A = PBQ$. Into how many disjoint equivalence classes is $\mathcal{M}$ partitioned by this equivalence relation?

\qitem{(11)} What is the area of the region $\left\{z + \frac{z^2}{2} \;\middle|\; z \in \C,\ |z| \le 1\right\}$?

\qitem{(12)} Let $\Gamma := C[-1, 1]$ with the supremum norm. Let $K$ be the kernel of the linear functional $I : f \mapsto \int_{-1}^{1} f(x)\,dx$ on $\Gamma$. Consider the function $f(x) = x^2$ in $\Gamma$. What is its distance from the subspace $K$?

\qitem{(13)} Let $X = \{1, 2, \dots, n\}$ ($n \ge 3$). Let $\mathcal{S}$ be the group of bijections of $X$, acting on $X \times X \times X$ by $\varphi\cdot(i, j, k) := (\varphi(i), \varphi(j), \varphi(k))$. How many orbits does this induced action of $\mathcal{S}$ on $X \times X \times X$ have?

\qitem{(14)} Let $k$ be a field with $8$ elements and $K$ a field with $64$ elements. How many ring homomorphisms are there from $k$ to $K$?

\qitem{(15)} Let $\mathcal{C}$ be the subset of real valued continuous functions $f(x)$ on $[-1, 1]$ such that $f(x) \ge 0$ (for all $x \in [-1, 1]$) and $\limsup_{n \to \infty} \int_{-1}^{1} f(x)^n\,dx < \infty$. What is the image of $\mathcal{C}$ in $\R$ under the map that takes $f(x)$ in $\mathcal{C}$ to $\int_{-1}^{1} f(x)\,dx$?

\qitem{(16)} Let $\mathcal{M}$ be the real vector space of all $n \times n$ real matrices with inner product $(A, B) := \operatorname{trace}(AB^t)$. Let $\mathcal{S}$ be the subspace of symmetric $n \times n$ matrices, and $P : \mathcal{M} \to \mathcal{S}$ the orthogonal projection to $\mathcal{S}$. For $X \in \mathcal{M}$, write a formula (in terms of $X$) for $P(X)$.

\qitem{(17)} Let $\mathcal{H} := \{a + bi \in \C \mid b > 0\}$ be the open upper half plane. Let $\varphi$ be an analytic function on $\mathcal{H}$ such that $\varphi(\mathcal{H}) \subseteq \mathcal{H}$, $\varphi(i) = i$, and $\varphi(2i) = i/2$. What is $|\varphi(1 + i)|$? (Write ``indeterminable'' if not enough information.)

\qitem{(18)} Let $A$ be the abelian group $\frac{\Z}{7\Z} \oplus \frac{\Z}{49\Z}$. How many subgroups of order $7$ does $A$ have?

\qitem{(19)} Let $\mathcal{M}$ be the algebra of all $3 \times 3$ complex matrices. For $A \in \mathcal{M}$, let $\Lambda_A : \mathcal{M} \to \mathcal{M}$ be defined by $X \mapsto AX - XA$. What are the possible values of the rank of $\Lambda_A$, as $A$ varies over all of $\mathcal{M}$?

\qitem{(20)} Consider the topology on $\C$ where a subset $Z$ of $\C$ is closed iff $Z = \{z \in \C \mid f(z) = 0\}$ for some polynomial $f(x) \in \Q[x]$. Write respectively the cardinalities of the closures of the three singleton sets: $\{\sqrt{2} + \sqrt{3}\}$, $\{e^{2\pi i/15}\}$, $\{\pi\}$.

\qitem{(21)} What is the number of conjugacy classes of cyclic subgroups of order $15$ in the symmetric group $S_{15}$?

\qitem{(22)} Let $f(x) := \sin(\pi x/2)$ on $\R$. For $y \in \R$, define $x_0(y) := y$ and $x_{n+1}(y) = f(x_n(y))$ for $n \ge 1$, and $g(y) := \lim_{n \to \infty} x_n(y)$. Find $\int_0^3 g(y)\,dy$.

\qitem{(23)} Let $y(t)$ be a smooth real function on an open interval containing $[0, 1]$ satisfying $y''(t) + w(t)y(t) = \lambda y(t)$, with $y(t) > 0$ on $[0, 1]$ and $y'(0) = y'(1) = 0$. Given that $\int_0^1 \left(\frac{y'}{y}\right)^2 dt = 10$ and $\int_0^1 w(t)\,dt = 20$, what is $\lambda$?

\qitem{(24)} Let $u(x)$ be a real valued continuous function on $[0, 1]$ such that $u(\frac{1}{2}) = 1$, $u(x) \ge 0$ and $u(x) \le \int_0^x u(t)\,dt$ for all $x \in [0, 1]$. What is $u(1)$? (Write ``indeterminable'' or ``no such $u$'' if applicable.)

\qitem{(25)} A smooth positive function $y = f(x)$ on $x > 0$ has the property that for any point $(x_0, y_0)$ on its graph $C$, $(x_0, y_0)$ is the mid-point of the segment $L(x_0, y_0)$ of the tangent line to $C$ at $(x_0, y_0)$ that lies in the first quadrant. Given that $f(1) = 2$, what is $f(2022)$?

\examsec{Part B (Questions 26--34) --- True or False Type}

\qitem{(26)} Suppose $A$ is an $n \times n$ complex matrix ($n \ge 2$) such that every matrix that commutes with $A$ is a polynomial in $A$.\\
(a) Such a matrix $A$ is necessarily diagonalizable.\\
(b) Such a matrix $A$ cannot be nilpotent.\\
(c) The characteristic and minimal polynomials coincide for every such $A$.

\qitem{(27)} Let $V = \{v_1, \dots, v_7\}$ be a set with $7$ elements. Which of the following tuples can occur as valency tuples $(p_1, \dots, p_7)$ of a simple graph structure on $V$? (True/False)\\
(a) $(2, 3, 4, 4, 4, 3, 3)$ \qquad (b) $(2, 2, 4, 6, 4, 2, 2)$ \qquad (c) $(2, 6, 4, 6, 2, 6, 4)$

\qitem{(28)} Let $f(z)$ be analytic on an open set containing the closed unit disk $D := \{z \in \C \mid |z| \le 1\}$. Let $m = \min_{z \in D} |f(z)|$ and $M = \min_{z \in C} |f(z)|$, where $C = \{|z| = 1\}$. Assume $m < M$. Then:\\
(a) $f(z)$ admits a zero on $D$. \qquad (b) $f(z)$ attains on $D$ every complex value $w$ such that $|w| < M$.

\qitem{(29)} Let $S = \R^2 \setminus \{(0, y) \mid y \neq 0\}$ with subspace topology $\tau_1$. Let $\pi : \R^2 \to S$ be defined by $(x, y) \mapsto (x, y)$ for $x \neq 0$ and $(0, y) \mapsto (0, 0)$ for all $y$, and let $\tau_2$ be the quotient topology on $S$. Then:
\begin{enumerate}[(a)]
  \item The topologies $\tau_1$ and $\tau_2$ are not comparable.
  \item $\tau_2$ is regular.
  \item $\tau_2$ admits a countable basis.
  \item $\tau_2$ is metrizable.
  \item Any subset of $S$ that is compact with respect to $\tau_2$ is also compact with respect to $\tau_1$.
\end{enumerate}

\qitem{(30)} Let $f : \R \to \R$ be a smooth, bounded function such that $\lim_{x \to \infty} f(x)$ and $\lim_{x \to \infty} f(-x)$ exist. Then:\\
(a) The function $f$ is uniformly continuous. \qquad (b) The derivative $f'$ is uniformly continuous.

\qitem{(31)} Let $A$ be an arbitrary special orthogonal $3 \times 3$ real matrix ($AA^t = A^t A = I$, $\det A = 1$). Then:
\begin{enumerate}[(a)]
  \item $A$ admits $1$ as an eigenvalue.
  \item If all the eigenvalues of $A$ are real, then $A$ is the identity $3 \times 3$ matrix.
  \item There is a complex invertible matrix $P$ such that $PAP^{-1}$ is diagonal.
\end{enumerate}

\qitem{(32)} Consider a sequence of i.i.d.\ random variables $\{X_n\}_{n \ge 1}$ taking values in $\{1, -1\}$ and satisfying $\lim_{n \to \infty} X_n = X$ almost surely, for some random variable $X$. Then:\\
(a) $X$ is a constant, almost surely. \qquad (b) $X_1$ is a constant, almost surely. \qquad (c) $X$ is distributed as $X_1$.

\qitem{(33)} Let $\R'$ be $\R$ with the topology having $\{[a, \infty) \mid a \in \R\}$ as a basis. Consider $\R' \times \R'$ with product topology, $S = \{(x, \sin(x)) \mid x \in \R\}$ with subspace topology, and $\varphi : \R \to S$ defined by $\varphi(x) := (x, \sin x)$. Then:\\
(a) The image of $(\pi/2, \pi)$ under $\varphi$ is closed in $S$. \qquad (b) The image of $(0, \pi/2)$ under $\varphi$ is closed in $S$.

\qitem{(34)} Let $d, e$ be positive integers and $T : \R^d \to \R^e$ a linear transformation. Then:
\begin{enumerate}[(a)]
  \item $T$ is an open map if and only if $T$ is surjective.
  \item $T$ is a closed map if and only if either $T = 0$ or $T$ is injective.
\end{enumerate}

\examsec{Answer Key --- NBHM 2022}
\begin{multicols}{2}
\textbf{Part A}\\
\textbf{(1)} $336$\\
\textbf{(2)} $b/(b + y + a)$\\
\textbf{(3)} $2\pi/3$\\
\textbf{(4)} $3$\\
\textbf{(5)} $1/\sqrt{5}$\\
\textbf{(6)} $\log(1 + \sqrt{2})$ or $\ln(1 + \sqrt{2})$\\
\textbf{(7)} $32 + \frac{22\pi}{3}$\\
\textbf{(8)} $-\pi$\\
\textbf{(9)} $211/60$\\
\textbf{(10)} $\min(m, n) + 1$\\
\textbf{(11)} $3\pi/2$\\
\textbf{(12)} $1/3$\\
\textbf{(13)} $5$\\
\textbf{(14)} $3$\\
\textbf{(15)} $[0, 2]$\\
\textbf{(16)} $(X + X^t)/2$\\
\textbf{(17)} $1/\sqrt{2}$\\
\textbf{(18)} $8$\\
\textbf{(19)} $0, 4, 6$\\
\textbf{(20)} $4$;\quad $8$;\quad $\infty$ or infinity\\
\textbf{(21)} $5$\\
\textbf{(22)} $1$\\
\textbf{(23)} $30$\\
\textbf{(24)} no such $u$\\
\textbf{(25)} $1/1011$

\medskip
\textbf{Part B}\\
\textbf{(26)} (a) False; (b) False; (c) True\\
\textbf{(27)} (a) False; (b) True; (c) False\\
\textbf{(28)} (a) True; (b) True\\
\textbf{(29)} (a) False; (b) True; (c) False; (d) False; (e) False\\
\textbf{(30)} (a) True; (b) False\\
\textbf{(31)} (a) True; (b) False; (c) True\\
\textbf{(32)} (a) True; (b) True; (c) True\\
\textbf{(33)} (a) False; (b) False\\
\textbf{(34)} (a) True; (b) True
\end{multicols}

\end{document}
